% Complete English EGA III2 editable text.
%
\documentclass[10pt,oneside]{book}

\usepackage[T1]{fontenc}
\usepackage[utf8]{inputenc}
\usepackage[english]{babel}
\usepackage{amsmath,amssymb,mathrsfs}
\usepackage{aurical}
\usepackage[all]{xy}
\usepackage[a4paper,margin=22mm]{geometry}
\usepackage[colorlinks=true,linkcolor=blue,hypertexnames=false,unicode=true,pdfencoding=auto]{hyperref}
\hypersetup{
  pdftitle={Elements of Algebraic Geometry III, Part Two -- English Edition},
  pdfauthor={Alexandre Grothendieck and Jean Dieudonne},
  pdfsubject={Standalone English edition rechecked against the diplomatic French NUMDAM source}
}

\renewcommand{\thesection}{\arabic{section}}
\renewcommand{\thesubsection}{\thesection.\arabic{subsection}}
\newcommand{\sh}[1]{\mathscr{#1}}
\newcommand{\shTor}{\sh{T}\textup{\kern-2.2pt{\Fontauri\slshape or}}}
\newcommand{\RR}{\mathrm R}
\newcommand{\HH}{\mathrm H}
\newcommand{\bb}[1]{\mathbf{#1}}
\newcommand{\cat}[1]{\mathtt{#1}}
\newcommand{\kres}{\kappa}
\newcommand{\vphi}{\phi}
\newcommand{\supertilde}{{\,\widetilde{\,}\,}}
\newcommand{\Spec}{\operatorname{Spec}}
\newcommand{\Proj}{\operatorname{Proj}}
\newcommand{\shProj}{\mathscr{P}\!roj}
\newcommand{\Hom}{\operatorname{Hom}}
\newcommand{\Tor}{\operatorname{Tor}}
\newcommand{\Ker}{\operatorname{Ker}}
\newcommand{\Coker}{\operatorname{Coker}}
\newcommand{\isoto}{\xrightarrow{\sim}}
\newcommand{\simto}{\xrightarrow{\sim}}
\newcommand{\shHom}{\mathscr{H}\!om}
\newcommand{\sref}[2][]{(#2)}
\newcommand{\oldpage}[2][]{\setcounter{footnote}{0}\marginpar{\scriptsize\textbf{#1}~|~#2}\ignorespaces}
\newenvironment{env}[1][]{\par\medskip\noindent\textbf{(#1)}\ }{\par}
\newenvironment{theorem}[1][]{\par\medskip\noindent\itshape Theorem (#1). ---\ }{\par}
\newenvironment{proposition}[1][]{\par\medskip\noindent\itshape Proposition (#1). ---\ }{\par}
\newenvironment{lemma}[1][]{\par\medskip\noindent\itshape Lemma (#1). ---\ }{\par}
\newenvironment{corollary}[1][]{\par\medskip\noindent\itshape Corollary (#1). ---\ }{\par}
\newenvironment{remark}[1][]{\par\medskip\noindent\itshape Remark (#1). ---\ }{\par}
\newenvironment{remarks}[1][]{\par\medskip\noindent\itshape Remarks (#1). ---\ }{\par}
\newenvironment{scholium}[1][]{\par\medskip\noindent\textit{Scholium (#1). ---\ }}{\par}
\newenvironment{definition}[1][]{\par\medskip\noindent\itshape Definition (#1). ---\ }{\par}
\newenvironment{proof}{\par\medskip\noindent\textit{Proof. ---\ }}{\par}

\begin{document}
\begin{titlepage}
  \centering
  {\Large A. Grothendieck and J. Dieudonne\par}
  \vfill
  {\Huge Elements of Algebraic Geometry III\par}
  \vspace{1.5ex}
  {\LARGE Part Two\par}
  \vfill
  {\large Standalone English edition\par}
  \vspace{1ex}
  {\small Rechecked against the diplomatic French NUMDAM source\par}
\end{titlepage}

\setcounter{section}{5}
\setcounter{subsection}{0}
\setcounter{equation}{0}
\setcounter{footnote}{0}
% BEGIN NATIVE INPUT ega3/ega3-6.tex
\section{Local and global Tor functors; K\"unneth formula}
\label{section:III.6}


\subsection{Introduction}
\label{subsection:III.6.1}

\oldpage[III]{137}

\begin{env}[6.1.1]
\label{III.6.1.1}
Let $f:X\to Y$ be a morphism of preschemes, and let $\sh{F}$ be
a quasi-coherent $\sh{O}_X$-module.  In studying the higher direct
images $\RR^n f_*(\sh{F})$, one is led to consider the following
general problem.  Given a base-change morphism $g:Y'\to Y$, put
\[
  X'=X_{(Y')}=X\times_Y Y',\qquad
  \sh{F}'=\sh{F}\otimes_{\sh{O}_Y}\sh{O}_{Y'},\qquad
  f'=f_{(Y')}:X'\to Y'.
\]
We seek information about the higher direct images
$\RR^n f'_*(\sh{F}')$, assuming that the
$\RR^n f_*(\sh{F})$ are known.  One sees readily
(cf.\ for example \sref{III.7.7.2}) that this reduces to studying
the variation of
\[
  \RR^n f_*\bigl(\sh{F}\otimes_{\sh{O}_Y}\sh{G}\bigr)
\]
as $\sh{G}$ ranges over the quasi-coherent $\sh{O}_Y$-modules;
in other words, to studying the functor
\[
  \sh{G}\longmapsto
  \RR^n f_*\bigl(\sh{F}\otimes_{\sh{O}_Y}\sh{G}\bigr).
\]
If $\sh{F}$ is flat over $Y$, the functor
$\sh{G}\mapsto\sh{F}\otimes_{\sh{O}_Y}\sh{G}$ is exact
\sref[0]{0.6.7.4}, and consequently the composite functors
\[
  \sh{G}\longmapsto
  \RR^n f_*\bigl(\sh{F}\otimes_{\sh{O}_Y}\sh{G}\bigr)
\]
again form a cohomological functor.  This is no longer true in
general.  In order to apply cohomological methods, one is therefore
led to replace
$\sh{G}\mapsto
\RR^n f_*(\sh{F}\otimes_{\sh{O}_Y}\sh{G})$
by other functors which are always cohomological.

These functors, which generalize the Tor functors of module theory,
are defined in Sections~\hyperref[subsection:III.6.3]{6.3}--\hyperref[subsection:III.6.7]{6.7}.  There are in fact two such
generalizations, one local and the other global, related by spectral
sequences discussed in Section~\hyperref[subsection:III.6.7]{6.7}.  As an application of those
spectral sequences one obtains, under suitable hypotheses, a
K\"unneth formula expressing
\[
  \RR^n(f_1\times f_2)_*
  \bigl(\sh{F}_1\otimes_Y\sh{F}_2\bigr)
\]
in terms of the higher direct images
$\RR^p(f_1)_*(\sh{F}_1)$ and
$\RR^q(f_2)_*(\sh{F}_2)$.  Other spectral sequences
\sref{III.6.8} generalize the associativity spectral sequences for
the Tor functor of modules; finally, the base-change problem also
leads to spectral sequences \sref{III.6.9}.
\end{env}

\begin{env}[6.1.2]
\label{III.6.1.2}
Moreover, one observes in \sref{III.6.10} that the cohomological
functors
\[
  \sh{G}\longmapsto\sh{T}_\bullet(\sh{G})
\]
so defined are locally on $Y$ of the form
\[
  \sh{G}\longmapsto
  \sh{H}_\bullet\bigl(\sh{L}_\bullet\otimes_Y\sh{G}\bigr),
\]
where $\sh{L}_\bullet$ is a complex of locally free
$\sh{O}_Y$-modules, defined up to homotopy, and
$\sh{H}_\bullet$ denotes homology.  It is therefore useful to set
aside the particular situation which produced
$\sh{T}_\bullet$ and to study, in general,
\oldpage[III]{138}
functors of the preceding form, imposing when necessary appropriate
finiteness hypotheses on the $\sh{L}_i$ or on the
$\sh{H}_i(\sh{L}_\bullet)$.  This is the subject of \hyperref[section:III.7]{Section~7},
which can for the most part be read independently of \hyperref[section:III.6]{Section~6}.

The most important properties of these functors concern exactness
of a component $\sh{T}_i$ of $\sh{T}_\bullet$.  We shall give
various criteria for establishing such properties.  As an
application, we shall obtain conditions under which, with the
notation of \sref{III.6.1.1}, the functor
\[
  \sh{G}\longmapsto
  \RR^n f_*\bigl(\sh{F}\otimes_{\sh{O}_Y}\sh{G}\bigr)
\]
is exact; we shall express this by saying that $\sh{F}$ is
\emph{cohomologically flat over $Y$ in dimension $n$}.
Another important property of the components $\sh{T}_i$ of
\[
  \sh{T}_\bullet(\sh{G})
  =
  \sh{H}_\bullet\bigl(\sh{L}_\bullet\otimes_Y\sh{G}\bigr)
\]
is semicontinuity of the function
\[
  y\longmapsto
  \dim_{\kres(y)}\bigl(T_i(\kres(y))\bigr).
\]
When $\sh{T}_i$ is exact, semicontinuity is replaced by continuity;
by a theorem of Grauert the converse also holds when $Y$ is reduced
\sref{III.7.8.4}.
\end{env}

\begin{env}[6.1.3]
\label{III.6.1.3}
In \hyperref[section:III.6]{Sections~6} and \hyperref[section:III.7]{7} we systematically use hypercohomology, taking
complexes of sheaves rather than sheaves as arguments throughout,
although the necessity of this point of view will become apparent
only in later chapters.  The cohomological formalism developed here
will moreover become more transparent in the chapter of this
Treatise devoted to an algebra of cohomological functors of coherent
sheaves, including the formalism of duality.  This, however, will
require developments beyond the scope of the present chapter.
\end{env}

\begin{env}[6.1.4]
\label{III.6.1.4}
For brevity, let $K^\bullet$ and $K'^\bullet$ be two complexes in
an abelian category $\mathsf{C}$.  We shall call a morphism of
complexes
\[
  f:K^\bullet\longrightarrow K'^\bullet
\]
a \emph{homotopism} if there is a morphism
$g:K'^\bullet\to K^\bullet$ such that both composites
$f\circ g$ and $g\circ f$ are homotopic to the identity.  By abuse
of language, when such a homotopism exists we shall also say that
$K^\bullet$ and $K'^\bullet$ are homotopic.

Suppose that the hypercohomology of a covariant additive functor
$T:\mathsf{C}\to\mathsf{C}'$, with values in an abelian category
$\mathsf{C}'$, can be defined relative to a complex of
$\mathsf{C}$ \sref[0]{0.11.4.3}.  It is then immediate that a
homotopism $K^\bullet\to K'^\bullet$ canonically defines an
isomorphism
\[
  \mathbf{R}^\bullet T(K^\bullet)
  \xrightarrow{\ \sim\ }
  \mathbf{R}^\bullet T(K'^\bullet)
\]
on hypercohomology \emph{(loc.\ cit.)}.
\end{env}


\subsection{Hypercohomology of complexes of modules on a prescheme}
\label{subsection:III.6.2}

\begin{env}[6.2.1]
\label{III.6.2.1}
Let $X$ be a prescheme, and let
$\sh{K}^\bullet=(\sh{K}^i)_{i\in\bb{Z}}$ be a complex of
$\sh{O}_X$-modules whose differential has degree $+1$.
Recall that, for every morphism $f:X\to Y$ of preschemes, the
$\sh{O}_Y$-modules of hypercohomology
$\mathcal{H}^n(f,\sh{K}^\bullet)$, also denoted
$\mathcal{H}_f^n(\sh{K}^\bullet)$ or
$\RR^n f_*(\sh{K}^\bullet)$, have been defined for every
$n\in\bb{Z}$ \sref[0]{0.12.4.1}.  The hypercohomology
$\mathcal{H}^\bullet(f,\sh{K}^\bullet)$ is the common abutment of
the two spectral functors
$'\sh{E}(f,\sh{K}^\bullet)$ and
$''\sh{E}(f,\sh{K}^\bullet)$ whose $E_2$-terms are
\[
  '\sh{E}_2^{pq}
  =
  \sh{H}^p\bigl(\sh{H}^q(f,\sh{K}^\bullet)\bigr),
  \tag{6.2.1.1}\label{III.6.2.1.1}
\]
and
\[
  ''\sh{E}_2^{pq}
  =
  \sh{H}^p\bigl(f,\sh{H}^q(\sh{K}^\bullet)\bigr)
  =
  \RR^p f_*\bigl(\sh{H}^q(\sh{K}^\bullet)\bigr).
  \tag{6.2.1.2}\label{III.6.2.1.2}
\]
\oldpage[III]{139}
Here $\sh{H}^q(f,\sh{K}^\bullet)$ is the complex whose component
of degree $i$ is
\[
  \sh{H}^q(f,\sh{K}^i)=\RR^q f_*(\sh{K}^i)
\]
\emph{(loc.\ cit.)}.  Recall also that, when $Y$ is reduced to a
point, the corresponding hypercohomology is denoted
$\mathbf{H}^\bullet(X,\sh{K}^\bullet)$; it consists of modules
over $\Gamma(X,\sh{O}_X)$ and is independent of the one-point
prescheme $Y$ chosen.  When $Y=X$ and $f=1_X$,
\[
  \mathcal{H}^n(f,\sh{K}^\bullet)
  =
  \sh{H}^n(\sh{K}^\bullet),
\]
the cohomology of the complex $\sh{K}^\bullet$.  When
$\sh{K}^i=0$ except for $i=i_0$, one has
\[
  \mathcal{H}^n(f,\sh{K}^\bullet)
  =
  \RR^{n-i_0}f_*(\sh{K}^{i_0}).
\]

The spectral sequence $'\sh{E}(f,\sh{K}^\bullet)$ is always
regular; both spectral sequences are biregular when
$\sh{K}^\bullet$ is bounded below \sref[0]{0.12.4.1}.
Every homotopism
$h:\sh{K}^\bullet\to\sh{K}'^\bullet$ of complexes of
$\sh{O}_X$-modules \sref{III.6.1.4} induces an isomorphism
\[
  \mathcal{H}^\bullet(f,\sh{K}^\bullet)
  \xrightarrow{\ \sim\ }
  \mathcal{H}^\bullet(f,\sh{K}'^\bullet)
\]
on hypercohomology.  The same conclusion holds if one assumes only
that
\[
  \sh{H}^\bullet(h):
  \sh{H}^\bullet(\sh{K}^\bullet)
  \xrightarrow{\ \sim\ }
  \sh{H}^\bullet(\sh{K}'^\bullet)
\]
is an isomorphism and that both complexes are bounded below.  This
follows immediately from \sref[0]{0.11.1.5}, applied to the
spectral sequence \sref{III.6.2.1.2} and its analogue for
$\sh{K}'^\bullet$.

Finally, for every open covering
$\mathfrak{U}=(U_\alpha)$ of $X$, the hypercohomology
$\mathbf{H}^\bullet(\mathfrak{U},\sh{K}^\bullet)$ has also been
defined \sref[0]{0.12.4.5} as the cohomology of the bicomplex
$\mathrm{C}^\bullet(\mathfrak{U},\sh{K}^\bullet)$, whose component
of bidegree $(i,j)$ is by definition
$\mathrm{C}^i(\mathfrak{U},\sh{K}^j)$.  The modules
$\mathbf{H}^n(\mathfrak{U},\sh{K}^\bullet)$ are again modules over
$\Gamma(X,\sh{O}_X)$.
\end{env}

\begin{proposition}[6.2.2]
\label{III.6.2.2}
Let $X$ be a scheme, and let $\mathfrak{U}=(U_\alpha)$ be a
covering of $X$ by affine open subsets.  For every complex
$\sh{K}^\bullet$ of quasi-coherent $\sh{O}_X$-modules, the
hypercohomology modules
\[
  \mathbf{H}^\bullet(X,\sh{K}^\bullet)
  \quad\text{and}\quad
  \mathbf{H}^\bullet(\mathfrak{U},\sh{K}^\bullet)
\]
are canonically isomorphic.
\end{proposition}

\begin{proof}
Every finite intersection $V$ of members of $\mathfrak{U}$ is
affine \sref[I]{I.5.5.6}; hence
$\HH^q(V,\sh{K}^i)=0$ for every $i$ and every $q>0$
\sref{III.1.3.1}.  The proposition is therefore a special case of
\sref[0]{0.12.4.7}.
\end{proof}

\begin{proposition}[6.2.3]
\label{III.6.2.3}
Let $f:X\to Y$ be a quasi-compact separated morphism of
preschemes.  For every complex $\sh{K}^\bullet$ of quasi-coherent
$\sh{O}_X$-modules, the $\sh{O}_Y$-modules
$\mathcal{H}^n(f,\sh{K}^\bullet)$ are quasi-coherent.
\end{proposition}

\begin{proof}
The modules
\[
  \sh{H}^q(f,\sh{K}^i)=\RR^qf_*(\sh{K}^i)
\]
are quasi-coherent $\sh{O}_Y$-modules \sref{III.1.4.10}.  Hence so
is $'\sh{E}_2^{pq}$: by \sref{III.6.2.1.1}, it is the quotient of
the kernel of a homomorphism of quasi-coherent modules by the image
of such a homomorphism \sref[I]{I.4.1.1}.  For the same reason, all
terms, cycles, and boundaries of the first spectral sequence are
quasi-coherent.  Regularity implies that
$\mathrm{Z}_\infty('\sh{E}_2^{pq})$ equals one of the
$\mathrm{Z}_k('\sh{E}_2^{pq})$ and is therefore quasi-coherent.
Likewise,
\[
  \mathrm{B}_\infty('\sh{E}_2^{pq})
  =
  \varinjlim_k\mathrm{B}_k('\sh{E}_2^{pq})
\]
is quasi-coherent \sref[0]{0.11.2.4}, \sref[I]{I.4.1.1}; thus the
$'\sh{E}_\infty^{pq}$ are quasi-coherent as well.

Since the spectral sequence is regular, the filtration
$F^\bullet\mathcal{H}^n(f,\sh{K}^\bullet)$ is discrete and
exhaustive.  Equivalently,
$\mathcal{H}^n(f,\sh{K}^\bullet)$ is the union of an increasing
sequence $(\sh{G}_k)_{k\geq0}$ of $\sh{O}_Y$-modules such that
$\sh{G}_0=0$ and every quotient
$\sh{G}_k/\sh{G}_{k-1}$ is one of the
$'\sh{E}_\infty^{pq}$, and hence is quasi-coherent.  It follows
inductively that the $\sh{G}_k$ are quasi-coherent
\sref{III.1.4.17}.  Since
\[
  \mathcal{H}^n(f,\sh{K}^\bullet)
  =
  \varinjlim_k\sh{G}_k,
\]
the proposition follows from \sref[I]{I.4.1.1}.
\end{proof}

\begin{corollary}[6.2.4]
\label{III.6.2.4}
\oldpage[III]{140}
Under the hypotheses of \sref{III.6.2.3}, for every affine open
subset $V$ of $Y$, the canonical homomorphism
\[
  \mathbf{H}^n\bigl(f^{-1}(V),\sh{K}^\bullet\bigr)
  \longrightarrow
  \Gamma\bigl(V,\mathcal{H}^n(f,\sh{K}^\bullet)\bigr)
  \tag{6.2.4.1}\label{III.6.2.4.1}
\]
is bijective for every $n\in\bb{Z}$.
\end{corollary}

\begin{proof}
The proof is the same as that of \sref{III.1.4.11}, using
\sref{III.6.2.2}: one replaces the sheaf there by
$\sh{K}^\bullet$, uses the Cech bicomplex
$f_*\mathrm{C}^\bullet(\mathfrak{U},\sh{K}^\bullet)$, and replaces
ordinary cohomology by
$\mathcal{H}^\bullet(f,\sh{K}^\bullet)$.  The latter is a
quasi-coherent $\sh{O}_Y$-module by \sref{III.6.2.3}.
\end{proof}

\begin{proposition}[6.2.5]
\label{III.6.2.5}
Let $Y$ be a locally Noetherian prescheme, let $f:X\to Y$ be
proper, and let $\sh{K}^\bullet$ be a complex of
$\sh{O}_X$-modules such that the $\sh{O}_X$-modules
$\sh{H}^q(\sh{K}^\bullet)$ are coherent.  Then the
$\sh{O}_Y$-modules $\mathcal{H}^n(f,\sh{K}^\bullet)$ are
coherent.
\end{proposition}

\begin{proof}
The question is local on $Y$, so we may suppose that $Y$ is
Noetherian and affine.  By \sref{III.6.2.4}, it is enough to prove
that the $\mathbf{H}^n(X,\sh{K}^\bullet)$ are finite-type modules
over $\Gamma(Y,\sh{O}_Y)$ \sref{III.3.2.1}.  By
\sref{III.6.2.2},
\[
  \mathbf{H}^\bullet(X,\sh{K}^\bullet)
  =
  \mathbf{H}^\bullet(\mathfrak{U},\sh{K}^\bullet),
\]
where $\mathfrak{U}$ may be taken finite, since $X$ is
quasi-compact.  The cochains in each complex
$\mathrm{C}^\bullet(\mathfrak{U},\sh{K}^j)$ are alternating by
definition.  Hence there is an integer $r>0$ such that
\[
  \mathrm{C}^i(\mathfrak{U},\sh{K}^j)=0
  \qquad\text{for }i<0\text{ and }i>r.
\]
It follows from \sref[0]{0.11.3.3} that both spectral sequences of
the bicomplex
$\mathrm{C}^\bullet(\mathfrak{U},\sh{K}^\bullet)$ are biregular.

Intersections of members of $\mathfrak{U}$ are affine open subsets
\sref[I]{I.5.5.6}, so each functor
\[
  \sh{F}\longmapsto\mathrm{C}^i(\mathfrak{U},\sh{F})
\]
is exact in the category of quasi-coherent $\sh{O}_X$-modules.
Consequently the $E_2$-terms of the second spectral sequence are
\[
  ''E_2^{pq}
  =
  \HH^p\bigl(
    \mathrm{C}^\bullet(
      \mathfrak{U},\sh{H}^q(\sh{K}^\bullet)
    )
  \bigr)
  =
  \mathbf{H}^p\bigl(
    \mathfrak{U},\sh{H}^q(\sh{K}^\bullet)
  \bigr)
  =
  \HH^p\bigl(X,\sh{H}^q(\sh{K}^\bullet)\bigr)
\]
\sref[0]{0.11.3.2}, \sref{III.1.4.1}.  Since $f$ is proper, these
are finite-type $\Gamma(Y,\sh{O}_Y)$-modules
\sref{III.3.2.1}.  Biregularity now implies that the
$\mathbf{H}^n(X,\sh{K}^\bullet)$ are finite-type
$\Gamma(Y,\sh{O}_Y)$-modules \sref[0]{0.11.1.8}.

\emph{A fortiori}, if $\sh{K}^\bullet$ is a complex of coherent
$\sh{O}_X$-modules, then the
$\mathcal{H}^n(f,\sh{K}^\bullet)$ are coherent under the
hypotheses on $Y$ and $f$ in \sref{III.6.2.5}
\sref[0]{0.5.3.4}.
\end{proof}

\begin{env}[6.2.6]
\label{III.6.2.6}
The hypercohomology
$\mathcal{H}^\bullet(f,\sh{K}^\bullet)$ is a cohomological functor
on the category of bounded-below complexes of $\sh{O}_X$-modules
\sref[0]{0.12.4.4}.  It is a cohomological functor on the category
of all complexes of $\sh{O}_X$-modules when $f$ is quasi-compact
and the underlying space of $X$ is locally Noetherian.  Indeed,
$f_*$ then commutes with inductive limits
\emph{(G, II, 3.10.1)}---the question being local on $Y$---and
\sref[0]{0.11.5.2} applies.

Finally, if $f$ is separated,
$\mathcal{H}^\bullet(f,\sh{K}^\bullet)$ is a cohomological functor
on the category of complexes of quasi-coherent
$\sh{O}_X$-modules.  This is immediate when $Y$ is affine, for
then $X$ is a scheme and the canonical isomorphism
\sref{III.6.2.2} reduces the assertion to the fact that
\[
  \sh{K}^\bullet\longmapsto
  \mathbf{H}^\bullet(\mathfrak{U},\sh{K}^\bullet)
\]
is cohomological.  This in turn is immediate because
\[
  \sh{K}^\bullet\longmapsto
  \mathrm{C}^\bullet(\mathfrak{U},\sh{K}^\bullet)
\]
is exact in this category \sref[I]{I.1.3.7}.

In general, for every affine open subset $V$ of $Y$,
\oldpage[III]{141}
$f^{-1}(V)$ is a scheme.  To apply the preceding argument, it is
enough to verify that, for an exact sequence
\[
  0\longrightarrow\sh{K}'^\bullet
  \longrightarrow\sh{K}^\bullet
  \longrightarrow\sh{K}''^\bullet
  \longrightarrow0
\]
of complexes of quasi-coherent $\sh{O}_X$-modules, the connecting
homomorphism
\[
  \partial:
  \mathcal{H}^n(f,\sh{K}''^\bullet)|V
  \longrightarrow
  \mathcal{H}^{n+1}(f,\sh{K}'^\bullet)|V
\]
does not depend on the affine open covering $\mathfrak{U}$ of
$f^{-1}(V)$ used to define it.  If $\mathfrak{U}'$ is an affine
open refinement of $\mathfrak{U}$, this follows from commutativity
of the triangle of canonical isomorphisms
\[
  \xymatrix@C=34pt@R=11pt{
    \mathbf{H}^\bullet(\mathfrak{U},\sh{K}^\bullet)
      \ar[rrd]^-{\sim}
      \ar[dd]^{\iota}
    && \\
    &&
    \mathbf{H}^\bullet(f^{-1}(V),\sh{K}^\bullet)
    \\
    \mathbf{H}^\bullet(\mathfrak{U}',\sh{K}^\bullet)
      \ar[rru]_-{\sim}
    &&
  }
\]
and of the diagram
\[
  \xymatrix@C=46pt@R=25pt{
    \mathbf{H}^n(\mathfrak{U},\sh{K}''^\bullet)
      \ar[r]^-{\partial}
      \ar[d]^{\iota}
    &
    \mathbf{H}^{n+1}(\mathfrak{U},\sh{K}'^\bullet)
      \ar[d]^{\iota}
    \\
    \mathbf{H}^n(\mathfrak{U}',\sh{K}''^\bullet)
      \ar[r]_-{\partial}
    &
    \mathbf{H}^{n+1}(\mathfrak{U}',\sh{K}'^\bullet).
  }
\]

Whenever one of the preceding conditions holds and
$\sh{K}^\bullet\to\sh{K}'^\bullet$ is a homotopism
\sref{III.6.1.4}, the resulting isomorphism
\[
  \mathcal{H}^\bullet(f,\sh{K}^\bullet)
  \xrightarrow{\ \sim\ }
  \mathcal{H}^\bullet(f,\sh{K}'^\bullet)
\]
is an isomorphism of $\partial$-functors
\sref[0]{0.11.4.4}.
\end{env}

\begin{env}[6.2.7]
\label{III.6.2.7}
Everything above applies naturally, with only a change of notation,
to a complex $\sh{K}_\bullet$ of quasi-coherent
$\sh{O}_X$-modules whose differential has degree $-1$: it is
enough to consider the complex
$\sh{K}^\bullet=(\sh{K}^i)$ defined by
\[
  \sh{K}^i=\sh{K}_{-i}
  \qquad(i\in\bb{Z}).
\]
\end{env}


\subsection{Hypertor of two complexes of modules}
\label{subsection:III.6.3}

\begin{env}[6.3.1]
\label{III.6.3.1}
Let $A$ be a commutative ring, and let $P_\bullet$ and
$Q_\bullet$ be two complexes of $A$-modules whose differentials
have degree $-1$.  Let $L_{\bullet\bullet}$ and
$M_{\bullet\bullet}$ be projective Cartan--Eilenberg resolutions
of $P_\bullet$ and $Q_\bullet$, respectively
\sref[0]{0.11.6.1}.  Taking the sum of the first degrees and the
sum of the second degrees,
\[
  L_{\bullet\bullet}\otimes_A M_{\bullet\bullet}
\]
is a bicomplex whose differentials have degree $-1$.  Its homology
\[
  \HH_\bullet(
    L_{\bullet\bullet}\otimes_A M_{\bullet\bullet}
  )
\]
does not depend on the chosen Cartan--Eilenberg resolutions and is,
by definition, the hyperhomology of the bifunctor
$P_\bullet\otimes_AQ_\bullet$ in $P_\bullet$ and $Q_\bullet$
\sref[0]{0.11.6.5}.  We put
\[
  \mathbf{Tor}_n^A(P_\bullet,Q_\bullet)
  =
  \HH_n(L_{\bullet\bullet}\otimes_A M_{\bullet\bullet})
  \tag{6.3.1.1}\label{III.6.3.1.1}
\]
\oldpage[III]{142}
and call this $A$-module the \emph{hypertor of index $n$} of the
two complexes $P_\bullet$ and $Q_\bullet$.  In the category of
bounded-below complexes of $A$-modules, the modules
$\mathbf{Tor}_n^A(P_\bullet,Q_\bullet)$ form a homological
bifunctor in $P_\bullet$ and $Q_\bullet$
\sref[0]{0.11.6.5}.  Moreover:
\end{env}

\begin{proposition}[6.3.2]
\label{III.6.3.2}
The bifunctor $\mathbf{Tor}_\bullet^A(P_\bullet,Q_\bullet)$ is the
common abutment of two bifunctorial spectral sequences
$'E(P_\bullet,Q_\bullet)$ and $''E(P_\bullet,Q_\bullet)$ whose
$E_2$-terms are
\[
  'E_{pq}^2
  =
  \HH_p\bigl(\Tor_q^A(P_\bullet,Q_\bullet)\bigr),
  \tag{6.3.2.1}\label{III.6.3.2.1}
\]
and
\[
  ''E_{pq}^2
  =
  \bigoplus_{q'+q''=q}
  \Tor_p^A\bigl(
    \HH_{q'}(P_\bullet),\HH_{q''}(Q_\bullet)
  \bigr).
  \tag{6.3.2.2}\label{III.6.3.2.2}
\]
In \sref{III.6.3.2.1},
$\Tor_q^A(P_\bullet,Q_\bullet)$ denotes the bicomplex formed by
the $A$-modules $\Tor_q^A(P_i,Q_j)$.  The spectral sequence
\sref{III.6.3.2.2} is always regular.  If $P_\bullet$ and
$Q_\bullet$ are bounded below, or if $A$ has finite cohomological
dimension, then both spectral sequences
\sref{III.6.3.2.1} and \sref{III.6.3.2.2} are biregular.
\end{proposition}

\begin{proof}
This follows from \sref[0]{0.11.6.5}.  Indeed, if $A$ has finite
cohomological dimension $n$, every $A$-module has a projective
resolution of length $n$ \emph{(M, VI, 2.1)}.
\end{proof}

\begin{corollary}[6.3.3]
\label{III.6.3.3}
Let $P'_\bullet$ and $Q'_\bullet$ be complexes of $A$-modules, and
let
\[
  u:P_\bullet\to P'_\bullet,
  \qquad
  v:Q_\bullet\to Q'_\bullet
\]
be homomorphisms of complexes.  If the induced homomorphisms
\[
  \HH_\bullet(u):
  \HH_\bullet(P_\bullet)\longrightarrow\HH_\bullet(P'_\bullet),
  \qquad
  \HH_\bullet(v):
  \HH_\bullet(Q_\bullet)\longrightarrow\HH_\bullet(Q'_\bullet)
\]
are bijective, then the homomorphism
\[
  \mathbf{Tor}_n^A(P_\bullet,Q_\bullet)
  \longrightarrow
  \mathbf{Tor}_n^A(P'_\bullet,Q'_\bullet)
\]
induced by $u$ and $v$ is bijective.
\end{corollary}

\begin{proof}
The induced homomorphism
\[
  ''E(P_\bullet,Q_\bullet)
  \longrightarrow
  ''E(P'_\bullet,Q'_\bullet)
\]
is an isomorphism on the $E_2$-terms.  The conclusion follows from
regularity of these spectral sequences \sref{III.6.3.2} and
\sref[0]{0.11.1.5}.
\end{proof}

\begin{proposition}[6.3.4]
\label{III.6.3.4}
Let $P_\bullet$ and $Q_\bullet$ be bounded-below complexes of
$A$-modules.  Let $L_{\bullet\bullet}$ and
$M_{\bullet\bullet}$ be bicomplexes of flat $A$-modules such that,
for every $i$, $L_{i,\bullet}$ is a resolution of $P_i$ and
$M_{i,\bullet}$ is a resolution of $Q_i$.  Then there are canonical
isomorphisms
\[
  \begin{split}
  \mathbf{Tor}_n^A(P_\bullet,Q_\bullet)
  &\xrightarrow{\ \sim\ }
  \HH_n(L_{\bullet\bullet}\otimes_AQ_\bullet)\\
  &\xrightarrow{\ \sim\ }
  \HH_n(P_\bullet\otimes_AM_{\bullet\bullet})\\
  &\xrightarrow{\ \sim\ }
  \HH_n(L_{\bullet\bullet}\otimes_AM_{\bullet\bullet}).
  \end{split}
  \tag{6.3.4.1}\label{III.6.3.4.1}
\]
\end{proposition}

\begin{proof}
This follows from \sref[0]{0.11.6.5}, parts (ii) and (iii), and
from the definition of flat $A$-modules.
\end{proof}

\begin{remarks}[6.3.5]
\label{III.6.3.5}
\emph{(i)} With the notation of \sref{III.6.3.1}, the bicomplexes
$L_{\bullet\bullet}\otimes_AM_{\bullet\bullet}$ and
$M_{\bullet\bullet}\otimes_AL_{\bullet\bullet}$ are canonically
isomorphic.  Hence there is a canonical isomorphism
\[
  \mathbf{Tor}_n^A(P_\bullet,Q_\bullet)
  \xrightarrow{\ \sim\ }
  \mathbf{Tor}_n^A(Q_\bullet,P_\bullet).
\]

\emph{(ii)} Let $F$ and $G$ be $A$-modules, and let $P_\bullet$
and $Q_\bullet$ be the complexes concentrated in degree $0$ with
values $F$ and $G$, respectively.  Projective resolutions
$L_\bullet$ of $F$ and $M_\bullet$ of $G$ may be regarded, after
adjoining zeros, as Cartan--Eilenberg resolutions of $P_\bullet$
and $Q_\bullet$.  Thus in this case
\[
  \mathbf{Tor}_n^A(P_\bullet,Q_\bullet)=\Tor_n^A(F,G).
\]
\end{remarks}

\begin{proposition}[6.3.6]
\label{III.6.3.6}
Let $(P_\bullet^\lambda)$ and $(Q_\bullet^\mu)$ be two filtered
inductive systems of complexes of $A$-modules.  There is a
canonical isomorphism
\[
  \varinjlim_{\lambda,\mu}
  \mathbf{Tor}_n^A(P_\bullet^\lambda,Q_\bullet^\mu)
  \xrightarrow{\ \sim\ }
  \mathbf{Tor}_n^A\left(
    \varinjlim_\lambda P_\bullet^\lambda,
    \varinjlim_\mu Q_\bullet^\mu
  \right).
  \tag{6.3.6.1}\label{III.6.3.6.1}
\]
\end{proposition}

\begin{proof}
Put
\[
  P_\bullet=\varinjlim_\lambda P_\bullet^\lambda,
  \qquad
  Q_\bullet=\varinjlim_\mu Q_\bullet^\mu.
\]
By functoriality, the modules
$\mathbf{Tor}_n^A(P_\bullet^\lambda,Q_\bullet^\mu)$ form an
inductive system, and
\oldpage[III]{143}
the canonical maps to
$\mathbf{Tor}_n^A(P_\bullet,Q_\bullet)$ give the homomorphism
\sref{III.6.3.6.1}.  More generally, they give a canonical
homomorphism of spectral sequences
\[
  \varinjlim_{\lambda,\mu}
  ''E(P_\bullet^\lambda,Q_\bullet^\mu)
  \longrightarrow
  ''E(P_\bullet,Q_\bullet),
\]
whose homomorphism on abutments is \sref{III.6.3.6.1}.

The spectral sequence $''E(P_\bullet,Q_\bullet)$ is regular by
\sref{III.6.3.2}.  So is
$\varinjlim_{\lambda,\mu}''E(P_\bullet^\lambda,Q_\bullet^\mu)$,
as follows from the definitions \sref[0]{0.11.1.7} and the proof
of \sref[0]{0.11.3.3}.  By \sref[0]{0.11.1.5}, it is therefore
enough to prove that
\[
  \varinjlim_{\lambda,\mu}
  ''E(P_\bullet^\lambda,Q_\bullet^\mu)
  \longrightarrow
  ''E(P_\bullet,Q_\bullet)
  \tag{6.3.6.2}\label{III.6.3.6.2}
\]
is bijective on the $E^2$-terms.

Since homology commutes with inductive limits of complexes of
modules, it remains only to prove that, for filtered inductive
systems $(F^\lambda)$ and $(G^\mu)$ of $A$-modules, the canonical
homomorphism
\[
  \varinjlim_{\lambda,\mu}\Tor_n^A(F^\lambda,G^\mu)
  \longrightarrow
  \Tor_n^A\left(
    \varinjlim_\lambda F^\lambda,
    \varinjlim_\mu G^\mu
  \right)
\]
is bijective.

For each $F^\lambda$, consider its canonical free resolution
\[
  L_\bullet^\lambda:\quad
  \cdots\longrightarrow L_{i+1}^\lambda
  \longrightarrow L_i^\lambda
  \longrightarrow\cdots\longrightarrow
  L_1^\lambda\longrightarrow L_0^\lambda\longrightarrow0,
\]
where $L_0^\lambda$ is the $A$-module of formal linear
combinations of elements of $F^\lambda$, and
$L_{i+1}^\lambda$ is the $A$-module of formal linear combinations
of elements of
$\Ker(L_i^\lambda\to L_{i-1}^\lambda)$.  The
$L_\bullet^\lambda$ form an inductive system of complexes.  Put
\[
  F=\varinjlim_\lambda F^\lambda,
  \qquad
  L_i=\varinjlim_\lambda L_i^\lambda.
\]
Since filtered inductive limits are exact, the $L_i$ form a
resolution $L_\bullet$ of $F$; moreover, each $L_i$ is flat, being
an inductive limit of free $A$-modules \sref[0]{0.6.1.2}.

Likewise, let $M_\bullet^\mu$ be the canonical free resolution of
$G^\mu$, and put
\[
  M_\bullet=\varinjlim_\mu M_\bullet^\mu,
  \qquad
  G=\varinjlim_\mu G^\mu.
\]
Then $M_\bullet$ is a flat resolution of $G$.  By
\sref{III.6.3.5} and \sref{III.6.3.4},
\[
  \Tor_n^A(F,G)=\HH_n(L_\bullet\otimes_AM_\bullet).
\]
Since homology commutes with inductive limits of complexes of
modules,
\[
  \HH_n(L_\bullet\otimes_AM_\bullet)
  =
  \varinjlim_{\lambda,\mu}
  \HH_n(L_\bullet^\lambda\otimes_AM_\bullet^\mu).
\]
The last modules are $\Tor_n^A(F^\lambda,G^\mu)$, which proves
\sref{III.6.3.6.1}.

If, in addition, there is an $i_0$ such that
$P_i^\lambda=Q_i^\mu=0$ for $i<i_0$, for all $\lambda$ and
$\mu$, the same argument proves that the canonical homomorphism
\[
  \varinjlim_{\lambda,\mu}
  'E(P_\bullet^\lambda,Q_\bullet^\mu)
  \longrightarrow
  'E(P_\bullet,Q_\bullet)
  \tag{6.3.6.3}\label{III.6.3.6.3}
\]
is bijective.
\end{proof}

\begin{proposition}[6.3.7]
\label{III.6.3.7}
Suppose that $P_\bullet$ and $Q_\bullet$ are bounded below.  If
$P_\bullet$ consists of flat $A$-modules, there is a canonical
$A$-linear isomorphism of $\partial$-functors in $Q_\bullet$
\[
  \mathbf{Tor}_n^A(P_\bullet,Q_\bullet)
  \xrightarrow{\ \sim\ }
  \HH_n(P_\bullet\otimes_AQ_\bullet).
  \tag{6.3.7.1}\label{III.6.3.7.1}
\]
\end{proposition}

\begin{proof}
The spectral sequence \sref{III.6.3.2.1} is biregular and
degenerate, so existence of \sref{III.6.3.7.1} follows from
\sref[0]{0.11.1.6}.  Computing hypertor from a projective
Cartan--Eilenberg resolution of $P_\bullet$ as in
\sref{III.6.3.4} shows immediately that this isomorphism is an
isomorphism of $\partial$-functors in $Q_\bullet$.
\end{proof}

\begin{env}[6.3.8]
\label{III.6.3.8}
\oldpage[III]{144}
Let $\rho:A\longrightarrow A'$ be a homomorphism of rings.  We
shall define a functorial $A$-linear homomorphism of degree zero,
canonically associated with $\rho$,
\[
  \rho_{P_\bullet,Q_\bullet}:
  \mathbf{Tor}_\bullet^A(P_\bullet,Q_\bullet)
  \longrightarrow
  \mathbf{Tor}_\bullet^{A'}
  (P_\bullet\otimes_AA',Q_\bullet\otimes_AA').
  \tag{6.3.8.1}\label{III.6.3.8.1}
\]

For this purpose, let $L_{\bullet\bullet}$ be a projective
Cartan--Eilenberg resolution of $P_\bullet$, and let
$L'_{\bullet\bullet}$ be a projective Cartan--Eilenberg resolution
of $P_\bullet\otimes_AA'$.  We first show that there is an
$A'$-linear homomorphism of complexes
\[
  L_{\bullet\bullet}\otimes_AA'
  \longrightarrow L'_{\bullet\bullet},
\]
determined up to homotopy.  Indeed, the construction of
$L_{\bullet\bullet}$ is completely determined once one chooses,
arbitrarily for each $i$, a projective resolution
$(X_{ij}^{\mathrm B})_{j\geq0}$ of $\mathrm B_i(P_\bullet)$ and a
projective resolution $(X_{ij}^{\mathrm H})_{j\geq0}$ of
$\HH_i(P_\bullet)$.  These are respectively
$\mathrm B_i^{\mathrm I}(L_{\bullet\bullet})$ and
$\HH_i^{\mathrm I}(L_{\bullet\bullet})$.  One then obtains
successively
\[
  \mathrm Z_i^{\mathrm I}(L_{\bullet\bullet})
  =
  \HH_i^{\mathrm I}(L_{\bullet\bullet})
  \oplus
  \mathrm B_i^{\mathrm I}(L_{\bullet\bullet}),
  \qquad
  L_{i,\bullet}
  =
  \mathrm Z_i^{\mathrm I}(L_{\bullet\bullet})
  \oplus
  \mathrm B_{i-1}^{\mathrm I}(L_{\bullet\bullet}).
\]
Now $X_{i,\bullet}^{\mathrm B}\otimes_AA'$ is not in general a
resolution of $\mathrm B_i(P_\bullet)\otimes_AA'$, but it is still
a complex of projective $A'$-modules.  Hence there is an
$A'$-linear homomorphism
\[
  X_{i,\bullet}^{\mathrm B}\otimes_AA'
  \longrightarrow
  \mathrm B_i^{\mathrm I}(L'_{\bullet\bullet})
\]
compatible with the augmentations and determined up to homotopy
(\emph{M}, V, 1.1).  Likewise there is an $A'$-linear homomorphism
\[
  X_{i,\bullet}^{\mathrm H}\otimes_AA'
  \longrightarrow
  \HH_i^{\mathrm I}(L'_{\bullet\bullet})
\]
determined up to homotopy.  By the construction just recalled,
these maps give, for every $i$, an $A'$-linear homomorphism
\[
  L_{i,\bullet}\otimes_AA'\longrightarrow L'_{i,\bullet}.
\]
For $i\in\mathbb Z$ these homomorphisms are compatible with the
differentials $L_{i,\bullet}\longrightarrow L_{i-1,\bullet}$ and
their analogues for $L'_{\bullet\bullet}$; together they therefore
give the required homomorphism
$L_{\bullet\bullet}\otimes_AA'\longrightarrow L'_{\bullet\bullet}$.

To define \sref{III.6.3.8.1}, take in the same way projective
Cartan--Eilenberg resolutions $M_{\bullet\bullet}$ and
$M'_{\bullet\bullet}$ of $Q_\bullet$ and
$Q_\bullet\otimes_AA'$, respectively, and an $A'$-linear
homomorphism
$M_{\bullet\bullet}\otimes_AA'\longrightarrow M'_{\bullet\bullet}$.
The resulting homomorphism
\[
  (L_{\bullet\bullet}\otimes_AA')
  \otimes_{A'}
  (M_{\bullet\bullet}\otimes_AA')
  \longrightarrow
  L'_{\bullet\bullet}\otimes_{A'}M'_{\bullet\bullet}
\]
gives, by composition, an $A$-linear homomorphism of bicomplexes
\[
  L_{\bullet\bullet}\otimes_AM_{\bullet\bullet}
  \longrightarrow
  L'_{\bullet\bullet}\otimes_{A'}M'_{\bullet\bullet}.
\]
Passing to homology yields \sref{III.6.3.8.1}; it is well defined
because it is induced by a homomorphism of complexes determined up
to homotopy.

If $\rho':A'\longrightarrow A''$ is a second homomorphism of rings
and $\rho''=\rho'\circ\rho:A\longrightarrow A''$, then
\[
  \rho''_{P_\bullet,Q_\bullet}
  =
  \rho'_{P'_\bullet,Q'_\bullet}
  \circ
  \rho_{P_\bullet,Q_\bullet},
  \qquad
  P'_\bullet=P_\bullet\otimes_AA',
  \quad
  Q'_\bullet=Q_\bullet\otimes_AA'.
\]
The preceding homomorphism of bicomplexes also defines functorial
homomorphisms of spectral sequences
\[
  'E_{pq}^r(P_\bullet,Q_\bullet)
  \longrightarrow
  'E_{pq}^r(P_\bullet\otimes_AA',Q_\bullet\otimes_AA')
\]
and
\[
  ''E_{pq}^r(P_\bullet,Q_\bullet)
  \longrightarrow
  ''E_{pq}^r(P_\bullet\otimes_AA',Q_\bullet\otimes_AA').
\]
They are independent of the Cartan--Eilenberg resolutions chosen
and satisfy the same transitivity property.
\end{env}

\begin{proposition}[6.3.9]
\label{III.6.3.9}
Let $\rho:A\longrightarrow A'$ be a homomorphism of rings such that
$A'$ is flat as an $A$-module.  There are canonical functorial
isomorphisms
\[
  \mathbf{Tor}_\bullet^{A'}
  (P_\bullet\otimes_AA',Q_\bullet\otimes_AA')
  \xrightarrow{\ \sim\ }
  \mathbf{Tor}_\bullet^A(P_\bullet,Q_\bullet)\otimes_AA',
  \tag{6.3.9.1}\label{III.6.3.9.1}
\]
\[
  \left\{
  \begin{aligned}
  'E(P_\bullet\otimes_AA',Q_\bullet\otimes_AA')
    &\xrightarrow{\ \sim\ }
      'E(P_\bullet,Q_\bullet)\otimes_AA',\\
  ''E(P_\bullet\otimes_AA',Q_\bullet\otimes_AA')
    &\xrightarrow{\ \sim\ }
      ''E(P_\bullet,Q_\bullet)\otimes_AA'.
  \end{aligned}
  \right.
  \tag{6.3.9.2}\label{III.6.3.9.2}
\]
\end{proposition}

\begin{proof}
\oldpage[III]{145}
Since the functor $M\longmapsto M\otimes_AA'$ is exact,
$L_{\bullet\bullet}\otimes_AA'$ and
$M_{\bullet\bullet}\otimes_AA'$ are projective
Cartan--Eilenberg resolutions of $P_\bullet\otimes_AA'$ and
$Q_\bullet\otimes_AA'$, respectively.  The assertion follows.
\end{proof}

\begin{env}[6.3.10]
\label{III.6.3.10}
Let $\rho:A\longrightarrow A'$ be a homomorphism of rings.  For
every complex $P'_\bullet$ of $A'$-modules, let
$P'_{\bullet[\rho]}$ denote the same complex regarded as a complex
of $A$-modules.  The identity map
$P'_{\bullet[\rho]}\longrightarrow P'_\bullet$ is the composite of
the canonical maps
\[
  P'_{\bullet[\rho]}
  \longrightarrow
  P'_{\bullet[\rho]}\otimes_AA'
  \xrightarrow{\ \mu\ }P'_\bullet,
\]
where $\mu$ is the $A'$-linear homomorphism
$\mu(x\otimes a')=a'x$.  If $Q'_\bullet$ is a second complex of
$A'$-modules, there are therefore canonical functorial
homomorphisms of degree zero
\[
  \begin{split}
  \mathbf{Tor}_\bullet^A(P'_{\bullet[\rho]},Q'_{\bullet[\rho]})
  &\longrightarrow
  \mathbf{Tor}_\bullet^{A'}
  (P'_{\bullet[\rho]}\otimes_AA',
   Q'_{\bullet[\rho]}\otimes_AA')\\
  &\longrightarrow
  \mathbf{Tor}_\bullet^{A'}(P'_\bullet,Q'_\bullet).
  \end{split}
  \tag{6.3.10.1}\label{III.6.3.10.1}
\]
The first arrow is the $A$-linear homomorphism defined in
\sref{III.6.3.8}; the second follows by functoriality from the
$A'$-linear homomorphisms
\[
  P'_{\bullet[\rho]}\otimes_AA'\longrightarrow P'_\bullet,
  \qquad
  Q'_{\bullet[\rho]}\otimes_AA'\longrightarrow Q'_\bullet.
\]
There are analogous homomorphisms for the spectral sequences of
\sref{III.6.3.2}, together with the evident transitivity
properties, whose formulation is left to the reader.
\end{env}

\begin{proposition}[6.3.11]
\label{III.6.3.11}
Let $\rho:A\longrightarrow A'$ be a homomorphism of rings such that
$A'$ is flat as an $A$-module.  For every complex $P'_\bullet$ of
$A'$-modules and every bounded-below complex $Q_\bullet$ of
$A$-modules, there is a canonical functorial isomorphism
\[
  \mathbf{Tor}_\bullet^A(P'_{\bullet[\rho]},Q_\bullet)
  \xrightarrow{\ \sim\ }
  \mathbf{Tor}_\bullet^{A'}(P'_\bullet,Q_\bullet\otimes_AA').
  \tag{6.3.11.1}\label{III.6.3.11.1}
\]
\end{proposition}

\begin{proof}
If $M_{\bullet\bullet}$ is a projective Cartan--Eilenberg
resolution of $Q_\bullet$, then
$M_{\bullet\bullet}\otimes_AA'$ is a projective
Cartan--Eilenberg resolution of $Q_\bullet\otimes_AA'$.  Moreover,
up to canonical isomorphism,
\[
  P'_{\bullet[\rho]}\otimes_AM_{\bullet\bullet}
  =
  P'_\bullet\otimes_{A'}
  (M_{\bullet\bullet}\otimes_AA').
\]
The assertion follows from \sref{III.6.3.4}.
\end{proof}

\begin{remark}[6.3.12]
\label{III.6.3.12}
Let $(A^\lambda)$ be a filtered inductive system of rings, and let
$(P_\bullet^\lambda)$ and $(Q_\bullet^\lambda)$ be two inductive
systems of complexes of $A^\lambda$-modules.  There is a canonical
isomorphism generalizing \sref{III.6.3.6.1},
\[
  \varinjlim_\lambda
  \mathbf{Tor}_\bullet^{A^\lambda}
  (P_\bullet^\lambda,Q_\bullet^\lambda)
  \xrightarrow{\ \sim\ }
  \mathbf{Tor}_\bullet^A(P_\bullet,Q_\bullet),
  \tag{6.3.12.1}\label{III.6.3.12.1}
\]
where
\[
  A=\varinjlim_\lambda A^\lambda,
  \qquad
  P_\bullet=\varinjlim_\lambda P_\bullet^\lambda,
  \qquad
  Q_\bullet=\varinjlim_\lambda Q_\bullet^\lambda.
\]
Once the homomorphisms
\[
  \mathbf{Tor}_\bullet^{A^\lambda}
  (P_\bullet^\lambda,Q_\bullet^\lambda)
  \longrightarrow
  \mathbf{Tor}_\bullet^{A^\mu}
  (P_\bullet^\mu,Q_\bullet^\mu)
  \qquad(\lambda\leq\mu)
\]
have been defined by \sref{III.6.3.10}, the proof is the same as
that of \sref{III.6.3.6}.
\end{remark}

\begin{proposition}[6.3.13]
\label{III.6.3.13}
Let $S$ be a multiplicative subset of $A$, and let $P_\bullet$ and
$Q_\bullet$ be complexes of $A$-modules on which multiplication by
every element of $S$ is bijective.  Thus, if $A'=S^{-1}A$, both
complexes consist of $A'$-modules.  Then there is a canonical
isomorphism
\[
  \mathbf{Tor}_\bullet^A(P_\bullet,Q_\bullet)
  \xrightarrow{\ \sim\ }
  \mathbf{Tor}_\bullet^{A'}(P_\bullet,Q_\bullet).
\]
\end{proposition}

\begin{proof}
The hypothesis implies that the canonical homomorphisms
\[
  P_\bullet\longrightarrow P_\bullet\otimes_AA',
  \qquad
  Q_\bullet\longrightarrow Q_\bullet\otimes_AA'
\]
are bijective.  Functoriality of hypertor also shows that each
$s\in S$ acts bijectively on
$\mathbf{Tor}_\bullet^A(P_\bullet,Q_\bullet)$; hence
\[
  \mathbf{Tor}_\bullet^A(P_\bullet,Q_\bullet)
  \longrightarrow
  \mathbf{Tor}_\bullet^A(P_\bullet,Q_\bullet)\otimes_AA'
\]
\oldpage[III]{146}
is bijective.  Since $A'$ is flat over $A$, the assertion follows
from \sref{III.6.3.9}.  The same argument gives canonical
isomorphisms of the associated spectral sequences.
\end{proof}


\subsection{Local hypertor functors of complexes of quasi-coherent modules: the affine scheme case}
\label{subsection:III.6.4}
\label{III.6.4}

\begin{env}[6.4.1]
\label{III.6.4.1}
Let $S$ be an affine scheme with ring $A$, and let $X$ and $Y$ be
affine $S$-schemes with rings $B$ and $C$, respectively; thus $B$
and $C$ are $A$-algebras.  Every complex $\sh P_\bullet$ of
quasi-coherent $\sh O_X$-modules is of the form
$\widetilde{P_\bullet}$ for a complex $P_\bullet$ of $B$-modules,
and similarly every complex $\sh Q_\bullet$ of quasi-coherent
$\sh O_Y$-modules is of the form $\widetilde{Q_\bullet}$ for a
complex $Q_\bullet$ of $C$-modules
(\sref[I]{I.1.3.7} and \sref[I]{I.1.3.8}).

Regarded as complexes of $A$-modules, $P_\bullet$ and $Q_\bullet$
have hypertor modules
$\mathbf{Tor}_n^A(P_\bullet,Q_\bullet)$.  By bifunctoriality, the
$A$-algebras $B$ and $C$ act on this $A$-module; these actions make
it a $(B,C)$-bimodule, equivalently a module over
\[
  B\otimes_AC=A(X\times_SY).
\]
We have thus defined a quasi-coherent
$\sh O_{X\times_SY}$-module
\[
  \shTor_n^{\sh O_S}(\sh P_\bullet,\sh Q_\bullet)
  =
  \bigl(\mathbf{Tor}_n^A(P_\bullet,Q_\bullet)\bigr)\supertilde.
  \tag{6.4.1.1}\label{III.6.4.1.1}
\]
It is called the \emph{local hypertor of index $n$} of
$\sh P_\bullet$ and $\sh Q_\bullet$, and is also denoted
$\shTor_n^S(\sh P_\bullet,\sh Q_\bullet)$.
\end{env}

\begin{lemma}[6.4.2]
\label{III.6.4.2}
With the notation of \sref{III.6.4.1}, suppose that
$A=R^{-1}A'$, where $A'$ is a ring and $R$ is a multiplicative
subset of $A'$.  Let $S'=\Spec(A')$.  Then $X$ and $Y$ may be
regarded as $S'$-schemes and
$X\times_{S'}Y=X\times_SY$
(\sref[I]{I.1.6.2} and \sref[I]{I.3.2.4}); moreover,
\[
  \shTor_\bullet^{S'}(\sh P_\bullet,\sh Q_\bullet)
  =
  \shTor_\bullet^S(\sh P_\bullet,\sh Q_\bullet).
\]
\end{lemma}

\begin{proof}
This follows from \sref{III.6.4.1.1} and \sref{III.6.3.13}.
\end{proof}

\begin{env}[6.4.3]
\label{III.6.4.3}
With the notation and hypotheses of \sref{III.6.4.1}, let
$\sh F=\widetilde F$ be a quasi-coherent $\sh O_X$-module and
$\sh G=\widetilde G$ a quasi-coherent $\sh O_Y$-module.  Regarding
$\sh F$ and $\sh G$ as complexes concentrated in degree zero, we
write $\shTor_n^{\sh O_S}(\sh F,\sh G)$, or
$\shTor_n^S(\sh F,\sh G)$, for their local hypertor of index $n$.
By \sref{III.6.3.5}(ii),
\[
  \shTor_n^{\sh O_S}(\sh F,\sh G)
  =
  \bigl(\Tor_n^A(F,G)\bigr)\supertilde.
  \tag{6.4.3.1}\label{III.6.4.3.1}
\]

Return now to two arbitrary complexes of quasi-coherent modules
$\sh P_\bullet$ and $\sh Q_\bullet$.  Formulas
\sref{III.6.4.1.1} and \sref{III.6.4.3.1}, together with
\sref{III.6.3.2}, show that
$\shTor_\bullet^S(\sh P_\bullet,\sh Q_\bullet)$ is the common abutment
of two spectral sequences
$'\sh E(\sh P_\bullet,\sh Q_\bullet)$ and
$''\sh E(\sh P_\bullet,\sh Q_\bullet)$ whose $E_2$-terms are
\[
  '\sh E_{pq}^2
  =
  \sh H_p\bigl(\shTor_q^S(\sh P_\bullet,\sh Q_\bullet)\bigr)
  \tag{6.4.3.2}\label{III.6.4.3.2}
\]
and
\[
  ''\sh E_{pq}^2
  =
  \bigoplus_{q'+q''=q}
  \shTor_p^S
  \bigl(\sh H_{q'}(\sh P_\bullet),
        \sh H_{q''}(\sh Q_\bullet)\bigr).
  \tag{6.4.3.3}\label{III.6.4.3.3}
\]
Here $\shTor_q^S(\sh P_\bullet,\sh Q_\bullet)$ in
\sref{III.6.4.3.2} denotes the bicomplex of quasi-coherent
$\sh O_{X\times_SY}$-modules
$\shTor_q^S(\sh P_i,\sh Q_j)$.
\end{env}

\begin{env}[6.4.4]
\label{III.6.4.4}
Let
\[
  X^{(1)}=\Spec(B^{(1)}),
  \qquad
  Y^{(1)}=\Spec(C^{(1)})
\]
be two further affine schemes, where $B^{(1)}$ and $C^{(1)}$ are
$A$-algebras, and let
\oldpage[III]{147}
\[
  u:X^{(1)}\longrightarrow X,
  \qquad
  v:Y^{(1)}\longrightarrow Y
\]
be $S$-morphisms corresponding to the $A$-homomorphisms
$\vphi:B\longrightarrow B^{(1)}$ and
$\psi:C\longrightarrow C^{(1)}$.  Consider the complexes
\[
  u^*(\sh P_\bullet)
  =
  (P_\bullet\otimes_BB^{(1)})\supertilde,
  \qquad
  v^*(\sh Q_\bullet)
  =
  (Q_\bullet\otimes_CC^{(1)})\supertilde.
\]
The canonical $A$-linear homomorphisms
\[
  P_\bullet\longrightarrow P_\bullet\otimes_BB^{(1)},
  \qquad
  Q_\bullet\longrightarrow Q_\bullet\otimes_CC^{(1)}
\]
give, by functoriality, an $A$-linear homomorphism
\[
  \mathbf{Tor}_\bullet^A(P_\bullet,Q_\bullet)
  \longrightarrow
  \mathbf{Tor}_\bullet^A
  (P_\bullet\otimes_BB^{(1)},
   Q_\bullet\otimes_CC^{(1)}).
\]
Again by functoriality, this is in fact a homomorphism of
$(B\otimes_AC)$-modules.  It therefore defines a
$(u\times_Sv)$-morphism
\[
  \theta:
  \shTor_\bullet^S(\sh P_\bullet,\sh Q_\bullet)
  \longrightarrow
  \shTor_\bullet^S(u^*\sh P_\bullet,v^*\sh Q_\bullet),
  \tag{6.4.4.1}\label{III.6.4.4.1}
\]
and hence a homomorphism of
$\sh O_{X^{(1)}\times_SY^{(1)}}$-modules
\[
  \theta^\sharp:
  (u\times_Sv)^*
  \shTor_\bullet^S(\sh P_\bullet,\sh Q_\bullet)
  \longrightarrow
  \shTor_\bullet^S(u^*\sh P_\bullet,v^*\sh Q_\bullet).
  \tag{6.4.4.2}\label{III.6.4.4.2}
\]
This is plainly a morphism of bi-$\partial$-functors on the
categories of bounded-below quasi-coherent modules.  The
homomorphism \sref{III.6.4.4.2} need not be bijective; nevertheless:
\end{env}

\begin{lemma}[6.4.5]
\label{III.6.4.5}
With the notation of \sref{III.6.4.4}, if $u$ and $v$ are open
immersions, then \sref{III.6.4.4.2} is bijective.
\end{lemma}

\begin{proof}
Identify $X^{(1)}$ and $Y^{(1)}$ with open subsets of $X$ and $Y$.
They are unions of open subsets of the form $\mathrm D(f)$ and
$\mathrm D(g)$, where $f\in B$ and $g\in C$; the schemes induced
on $\mathrm D(\vphi(f))$ and $\mathrm D(f)$, and likewise on
$\mathrm D(\psi(g))$ and $\mathrm D(g)$, are isomorphic.

It is enough to prove the lemma when
$X^{(1)}=\mathrm D(f)$ and $Y^{(1)}=\mathrm D(g)$.  Indeed, once
this is known, in the general case it suffices to show that the
restriction of $\theta^\sharp$ to every
$\mathrm D(f)\times_S\mathrm D(g)$ is an isomorphism.  If
\[
  u_1:\mathrm D(f)\longrightarrow X^{(1)},
  \qquad
  v_1:\mathrm D(g)\longrightarrow Y^{(1)}
\]
are the canonical inclusions, this restriction is
$(u_1\times_Sv_1)^*(\theta^\sharp)$.  From
\sref{III.6.4.4} and \sref[0]{0.4.4.8}, it is immediate that
composing it with the canonical homomorphism
\[
  (u_1\times_Sv_1)^*
  \shTor_\bullet^S(u^*\sh P_\bullet,v^*\sh Q_\bullet)
  \longrightarrow
  \shTor_\bullet^S((u')^*\sh P_\bullet,(v')^*\sh Q_\bullet),
  \tag{6.4.5.1}\label{III.6.4.5.1}
\]
where $u'=u\circ u_1$ and $v'=v\circ v_1$, gives the canonical
homomorphism
\[
  (u'\times_Sv')^*
  \shTor_\bullet^S(\sh P_\bullet,\sh Q_\bullet)
  \longrightarrow
  \shTor_\bullet^S((u')^*\sh P_\bullet,(v')^*\sh Q_\bullet).
  \tag{6.4.5.2}\label{III.6.4.5.2}
\]
If \sref{III.6.4.5.1} and \sref{III.6.4.5.2} are isomorphisms,
then so is $(u_1\times_Sv_1)^*(\theta^\sharp)$.

Suppose therefore that $X^{(1)}=\mathrm D(f)$ and
$Y^{(1)}=\mathrm D(g)$, so that $B^{(1)}=B_f$ and
$C^{(1)}=C_g$.  Then $u^*(\sh P_\bullet)$ and
$v^*(\sh Q_\bullet)$ identify with
$((P_\bullet)_f)\supertilde$ and
$((Q_\bullet)_g)\supertilde$, respectively.  Moreover,
$X^{(1)}\times_SY^{(1)}$ identifies with the open subscheme
$\mathrm D(f\otimes g)$ of
\[
  X\times_SY=\Spec(B\otimes_AC)
\]
(\sref[II]{II.4.3.2.4}).  It remains to prove that the
homomorphism
\[
  \bigl(\mathbf{Tor}_\bullet^A(P_\bullet,Q_\bullet)\bigr)_{f\otimes g}
  \longrightarrow
  \mathbf{Tor}_\bullet^A((P_\bullet)_f,(Q_\bullet)_g),
  \tag{6.4.5.3}\label{III.6.4.5.3}
\]
\oldpage[III]{148}
induced by the canonical homomorphisms
$P_\bullet\longrightarrow(P_\bullet)_f$ and
$Q_\bullet\longrightarrow(Q_\bullet)_g$, is bijective.

By \sref[0]{0.1.6.1},
\[
  (P_\bullet)_f=\varinjlim_m P_\bullet^{(m)},
\]
where every $P_\bullet^{(m)}$ is the same complex of $B$-modules
as $P_\bullet$, and for $m\leq n$ the map
$P_\bullet^{(m)}\longrightarrow P_\bullet^{(n)}$ is
multiplication by $f^{\,n-m}$.  The analogous statement holds for
$Q_\bullet$, with $g$ in place of $f$.  The homomorphism
\[
  \mathbf{Tor}_\bullet^A(P_\bullet^{(m)},Q_\bullet^{(m)})
  \longrightarrow
  \mathbf{Tor}_\bullet^A(P_\bullet^{(n)},Q_\bullet^{(n)})
\]
corresponding to these transition maps is, by definition,
multiplication by $(f\otimes g)^{n-m}$.  The conclusion follows
from \sref[0]{0.1.6.1}, applied to the left-hand side of
\sref{III.6.4.5.3}, and from \sref{III.6.3.6}.
\end{proof}

\begin{env}[6.4.6]
\label{III.6.4.6}
With the notation of \sref{III.6.4.4}, one defines in the same way
canonical homomorphisms of spectral functors
\[
  \begin{aligned}
  (u\times_Sv)^*\,{}'\sh E(\sh P_\bullet,\sh Q_\bullet)
  &\longrightarrow
  {}'\sh E(u^*\sh P_\bullet,v^*\sh Q_\bullet),\\
  (u\times_Sv)^*\,{}''\sh E(\sh P_\bullet,\sh Q_\bullet)
  &\longrightarrow
  {}''\sh E(u^*\sh P_\bullet,v^*\sh Q_\bullet).
  \end{aligned}
  \tag{6.4.6.1}\label{III.6.4.6.1}
\]
The argument of \sref{III.6.4.5} shows that when $u$ and $v$ are
open immersions, these homomorphisms are bijective.  Taking
\sref{III.6.3.6.2} and \sref{III.6.3.6.3} into account, the
argument proves that they are isomorphisms on the $E^2$-terms,
while \sref{III.6.4.5} proves that they are isomorphisms on the
abutments.  The assertion therefore follows from
\sref[0]{0.11.1.2} and \sref[0]{0.11.2.4}.
\end{env}


\subsection{Local hypertor functors of complexes of quasi-coherent modules: the general case}
\label{subsection:III.6.5}

\begin{env}[6.5.1]
\label{III.6.5.1}
Let $S$ now be an arbitrary prescheme and let $X$ and $Y$ be two
$S$-preschemes.  Let $\sh P_\bullet$ be a complex of quasi-coherent
$\sh O_X$-modules and $\sh Q_\bullet$ a complex of quasi-coherent
$\sh O_Y$-modules.  Set $Z=X\times_SY$.  We shall define
quasi-coherent $\sh O_Z$-modules
\[
  \shTor_n^S(\sh P_\bullet,\sh Q_\bullet),
\]
called the \emph{local hypertors} of $\sh P_\bullet$ and
$\sh Q_\bullet$, which reduce to those of \sref{III.6.4} when
$S$, $X$, and $Y$ are affine.

When $\sh P_\bullet$ and $\sh Q_\bullet$ are concentrated in
degree zero, with respective terms $\sh F$ and $\sh G$, we write
$\shTor_n^S(\sh F,\sh G)$ instead of
$\shTor_n^S(\sh P_\bullet,\sh Q_\bullet)$.
\end{env}

\begin{env}[6.5.2]
\label{III.6.5.2}
Suppose first that $S$ is affine, and let $(X_\lambda)$ and
$(Y_\mu)$ be coverings of $X$ and $Y$ by affine open subsets.
Then
\[
  Z_{\lambda\mu}=X_\lambda\times_SY_\mu
\]
form an affine open covering of $Z$.  Put
\[
  \sh P_{\lambda,\bullet}=\sh P_\bullet|X_\lambda,
  \qquad
  \sh Q_{\mu,\bullet}=\sh Q_\bullet|Y_\mu.
\]
For each pair $(\lambda,\mu)$ we thus have a quasi-coherent
$\sh O_{Z_{\lambda\mu}}$-module
\[
  \sh F_{\lambda\mu}
  =
  \shTor_n^S(\sh P_{\lambda,\bullet},\sh Q_{\mu,\bullet}).
\]
We must show that the $\sh F_{\lambda\mu}$ satisfy the gluing
condition \sref[0]{0.3.3.1}.  It is enough to verify that, for
every affine open
\[
  U\subset X_\lambda\cap X_{\lambda'},
  \qquad
  V\subset Y_\mu\cap Y_{\mu'},
\]
the restrictions of $\sh F_{\lambda\mu}$ and
$\sh F_{\lambda'\mu'}$ to $U\times_SV$ are canonically
isomorphic.  This follows at once from their canonical
isomorphisms with
\[
  \shTor_n^S(\sh P_\bullet|U,\sh Q_\bullet|V)
\]
given by \sref{III.6.4.5}.

The definition and \sref{III.6.4.5} also show that the resulting
$\sh O_Z$-module is independent, up to a unique isomorphism
compatible with the local identifications, of the chosen open
coverings.  We denote it by
\oldpage[III]{149}
\[
  \shTor_n^S(\sh P_\bullet,\sh Q_\bullet).
\]
Finally, for every open subset $U$ of $X$ and every open subset
$V$ of $Y$, its restriction to $U\times_SV$ is canonically
isomorphic to
\[
  \shTor_n^S(\sh P_\bullet|U,\sh Q_\bullet|V).
\]
\end{env}

\begin{env}[6.5.3]
\label{III.6.5.3}
Pass now to the general case, where $S$ is arbitrary.  Let
$(S_\alpha)$ be a covering of $S$ by affine open subsets, and let
$X_\alpha$ and $Y_\alpha$ be the inverse images of $S_\alpha$ in
$X$ and $Y$.  We must show that the sheaves
\[
  \sh G_\alpha
  =
  \shTor_n^{S_\alpha}
  (\sh P_\bullet|X_\alpha,\sh Q_\bullet|Y_\alpha)
\]
satisfy the gluing condition.

For every affine open subset
$T\subset S_\alpha\cap S_\beta$, let $U$ and $V$ be its inverse
images in $X$ and $Y$.  It is enough to define canonical
isomorphisms between the restrictions of $\sh G_\alpha$ and
$\sh G_\beta$ to $U\times_SV$ and
\[
  \shTor_n^T(\sh P_\bullet|U,\sh Q_\bullet|V).
\]
We may further restrict to the case in which
$T=\mathrm D(f_\alpha)=\mathrm D(f_\beta)$, where $f_\alpha$ and
$f_\beta$ are sections of $\sh O_S$ over $S_\alpha$ and
$S_\beta$, respectively.  By \sref{III.6.4.2},
\[
  \shTor_n^T(\sh P_\bullet|U,\sh Q_\bullet|V)
\]
is canonically isomorphic both to
$\shTor_n^{S_\alpha}(\sh P_\bullet|U,\sh Q_\bullet|V)$ and to
$\shTor_n^{S_\beta}(\sh P_\bullet|U,\sh Q_\bullet|V)$.
Combining these isomorphisms with the local identifications of
\sref{III.6.5.2} completes the definition of the
$\sh O_Z$-module
\[
  \shTor_n^S(\sh P_\bullet,\sh Q_\bullet).
\]
For every open subset $U$ of $X$ and every open subset $V$ of
$Y$, the restriction of this module to $U\times_SV$ is
canonically isomorphic to
\[
  \shTor_n^S(\sh P_\bullet|U,\sh Q_\bullet|V).
\]

This construction plainly defines, on the categories of
bounded-below complexes of quasi-coherent modules, a
bi-$\partial$-functor
\[
  \shTor_\bullet^S(\sh P_\bullet,\sh Q_\bullet)
\]
with values in $\sh O_Z$-modules.  Indeed the assertion is local
on $X$, $Y$, and $S$, by \sref{III.6.4.5} and by the fact that
\sref{III.6.4.4.2} is a morphism of bi-$\partial$-functors.

If $\sh P_\bullet$ and $\sh Q_\bullet$ are concentrated in degree
zero, with respective terms $\sh F$ and $\sh G$, then
\[
  \shTor_0^S(\sh F,\sh G)
  =
  \sh F\otimes_S\sh G,
\]
the external tensor product defined in \sref[I]{I.9.1.2}; this
follows from \sref{III.6.4.1.1} and \sref[I]{I.9.1.3}.
\end{env}

\begin{env}[6.5.4]
\label{III.6.5.4}
The preceding construction and \sref{III.6.4.6} show that
$\shTor_\bullet^S(\sh P_\bullet,\sh Q_\bullet)$ is the common abutment
of two spectral functors
$'\sh E(\sh P_\bullet,\sh Q_\bullet)$ and
$''\sh E(\sh P_\bullet,\sh Q_\bullet)$, whose $E_2$-terms are
\[
  '\sh E_{pq}^2
  =
  \sh H_p\bigl(\shTor_q^S(\sh P_\bullet,\sh Q_\bullet)\bigr)
  \tag{6.5.4.1}\label{III.6.5.4.1}
\]
and
\[
  ''\sh E_{pq}^2
  =
  \bigoplus_{q'+q''=q}
  \shTor_p^S
  \bigl(\sh H_{q'}(\sh P_\bullet),
        \sh H_{q''}(\sh Q_\bullet)\bigr).
  \tag{6.5.4.2}\label{III.6.5.4.2}
\]
The spectral sequence \sref{III.6.5.4.2} is always regular.  Both
spectral sequences are biregular when $\sh P_\bullet$ and
$\sh Q_\bullet$ are bounded below.  A further case in which they
are biregular is described next.
\end{env}

\begin{env}[6.5.5]
\label{III.6.5.5}
Let $T$ be a topological space.  A sheaf of rings $\sh A$ on $T$
is said to have \emph{cohomological dimension at most $n$} if the
ring $\sh A_t$ has cohomological dimension at most $n$ for every
$t\in T$.  In this case the ringed space $(T,\sh A)$ is also said
to have cohomological dimension at most $n$.  A sheaf of rings, or
a ringed space, has \emph{finite cohomological dimension} if this
holds for some integer $n$.

If the $\sh A_t$ are Noetherian commutative local rings, saying
that they have cohomological dimension at most $n$ means that
they are regular and have Krull dimension at most $n$
\oldpage[III]{150}
(\sref[0]{0.17.3.1}).  In the terminology of dimension theory to
be introduced in Chapter IV, a locally Noetherian prescheme $T$ has
cohomological dimension at most $n$ exactly when it is regular
(\sref[0]{0.4.1.4}) and has dimension at most $n$.  Equivalently,
for every affine open subset $U$ of $T$, the ring
$\Gamma(U,\sh O_T)$ has cohomological dimension at most $n$
(\sref[0]{0.17.2.6}).

Consequently, if $S$ is locally Noetherian and has finite
cohomological dimension, then \sref{III.6.3.2} shows that both
$'\sh E(\sh P_\bullet,\sh Q_\bullet)$ and
$''\sh E(\sh P_\bullet,\sh Q_\bullet)$ are biregular.

Finally, under the canonical isomorphism
$X\times_SY\xrightarrow{\sim}Y\times_SX$,
$\shTor_\bullet^S(\sh P_\bullet,\sh Q_\bullet)$ corresponds
canonically to
$\shTor_\bullet^S(\sh Q_\bullet,\sh P_\bullet)$.
\end{env}

\begin{proposition}[6.5.6]
\label{III.6.5.6}
Let $(\sh P_{\alpha\bullet})$ be a filtered direct system of
complexes of quasi-coherent $\sh O_X$-modules.  There is a
canonical isomorphism
\[
  \varinjlim_\alpha
  \shTor_\bullet^S(\sh P_{\alpha\bullet},\sh Q_\bullet)
  \xrightarrow{\sim}
  \shTor_\bullet^S
  \left(\varinjlim_\alpha\sh P_{\alpha\bullet},\sh Q_\bullet\right).
  \tag{6.5.6.1}\label{III.6.5.6.1}
\]
\end{proposition}

\begin{proof}
The question is local on $S$, $X$, and $Y$.  We may therefore
suppose that $S$, $X$, and $Y$ are affine, in which case the
proposition reduces to \sref{III.6.3.6}.
\end{proof}

\begin{env}[6.5.7]
\label{III.6.5.7}
\emph{Remarks.}

\textnormal{(i)} Consider in particular the case $S=X=Y$, with
$\sh P_\bullet$ and $\sh Q_\bullet$ two complexes of
quasi-coherent $\sh O_S$-modules.  Then the
$\shTor_n^S(\sh P_\bullet,\sh Q_\bullet)$ are quasi-coherent
$\sh O_S$-modules.  Moreover, for every point $z\in S$,
\sref{III.6.5.6} gives a canonical isomorphism
\[
  \bigl(\shTor_n^S(\sh P_\bullet,\sh Q_\bullet)\bigr)_z
  \xrightarrow{\sim}
  \mathbf{Tor}_n^{\sh O_z}
  \bigl((\sh P_\bullet)_z,(\sh Q_\bullet)_z\bigr),
  \tag{6.5.7.1}\label{III.6.5.7.1}
\]
because the question is local and, by \sref{III.6.4.1.1},
reduces to the case of modules.

\textnormal{(ii)} The definition of hypertor can be generalized
to two complexes $\sh P_\bullet,\sh Q_\bullet$ of
$\sh O_X$-modules on the same ringed space $(X,\sh O_X)$.  For
every open subset $U$ of $X$, put
\[
  A(U)=\Gamma(U,\sh O_X),\qquad
  P_\bullet(U)=\Gamma(U,\sh P_\bullet),\qquad
  Q_\bullet(U)=\Gamma(U,\sh Q_\bullet).
\]
The $A(U)$-modules
\[
  \mathbf{Tor}_n^{A(U)}\bigl(P_\bullet(U),Q_\bullet(U)\bigr)
\]
then form a presheaf on $X$; denote the associated
$\sh O_X$-module by
$\shTor_n^X(\sh P_\bullet,\sh Q_\bullet)$.  When $X$ is a
prescheme, \sref{III.6.3.12} shows that this module is
canonically isomorphic to the hypertor defined above.  We shall
not pursue this generalization further.
\end{env}

\begin{proposition}[6.5.8]
\label{III.6.5.8}
Let $X$ and $Y$ be two $S$-preschemes, and let $\sh F$
(respectively $\sh G$) be a quasi-coherent $\sh O_X$-module
(respectively $\sh O_Y$-module).  If $\sh F$ or $\sh G$ is
$S$-flat, then
\[
  \shTor_n^S(\sh F,\sh G)=0\qquad(n\ne0).
\]
\end{proposition}

\begin{proof}
The question is local on $X$ and $Y$, so we may suppose that
$X$, $Y$, and $S$ are affine, with respective rings $B$, $C$,
and $A$, and write $\sh F=\widetilde M$ and
$\sh G=\widetilde N$, where $M$ is a $B$-module and $N$ a
$C$-module.  Suppose, for example, that $\sh F$ is $S$-flat.
This means that, for every $s\in S$, $M_s$ is a flat
$A_s$-module (\sref[0]{0.6.7.1}); hence $M$ is a flat
$A$-module (\sref[0]{0.6.3.3}).  Thus
$\Tor_n^A(M,N)=0$ for $n>0$ and every $C$-module $N$
(\sref[0]{0.6.1.1}), and the conclusion follows from
\sref{III.6.4.1.1}.
\end{proof}

\begin{corollary}[6.5.9]
\label{III.6.5.9}
Let $X$ and $Y$ be two $S$-preschemes, and let
$\sh P_\bullet$ (respectively $\sh Q_\bullet$) be a
\oldpage[III]{151}
bounded-below complex of quasi-coherent $\sh O_X$-modules
(respectively $\sh O_Y$-modules).  Suppose that every
$\sh P_i$ is $S$-flat.  Then there is a canonical isomorphism of
$\partial$-functors in $\sh Q_\bullet$
\[
  \shTor_\bullet^S(\sh P_\bullet,\sh Q_\bullet)
  \xrightarrow{\sim}
  \sh H_\bullet(\sh P_\bullet\otimes_S\sh Q_\bullet).
  \tag{6.5.9.1}\label{III.6.5.9.1}
\]
\end{corollary}

\begin{proof}
When $S$, $X$, and $Y$ are affine, this is precisely
\sref{III.6.3.7}; the general case follows by the arguments of
\sref{III.6.5.2} and \sref{III.6.5.3}.
\end{proof}

\begin{env}[6.5.10]
\label{III.6.5.10}
\emph{Corollary.}
Suppose that $X$ is flat over $S$ (\sref[0]{0.6.7.1}), that
$\sh P_\bullet$ and $\sh Q_\bullet$ are bounded below, and that
every $\sh P_i$ is a locally free $\sh O_X$-module, not
necessarily of finite type.  Then the homomorphism
\sref{III.6.5.9.1} is bijective.

Indeed, the hypothesis of \sref{III.6.5.9} is satisfied:
flatness is, by definition, a pointwise property on $X$, and
every direct sum of flat modules is flat (\sref[0]{0.6.1.2}).
\end{env}

\begin{proposition}[6.5.11]
\label{III.6.5.11}
Let $X'$ and $Y'$ be two $S$-preschemes, and let
$f:X\to X'$ and $g:Y\to Y'$ be two affine $S$-morphisms.  For a
complex $\sh P_\bullet$ of quasi-coherent $\sh O_X$-modules and
a complex $\sh Q_\bullet$ of quasi-coherent $\sh O_Y$-modules,
there is a canonical functorial isomorphism
\[
  (f\times_Sg)_*
  \bigl(\shTor_\bullet^S(\sh P_\bullet,\sh Q_\bullet)\bigr)
  \xrightarrow{\sim}
  \shTor_\bullet^S(f_*\sh P_\bullet,g_*\sh Q_\bullet).
  \tag{6.5.11.1}\label{III.6.5.11.1}
\]
\end{proposition}

\begin{proof}
Since $f$ and $g$ are affine, $f_*\sh P_\bullet$ and
$g_*\sh Q_\bullet$ are complexes of quasi-coherent modules
(\sref[II]{II.1.2.6}).  If $Z'=X'\times_SY'$, both sides of
\sref{III.6.5.11.1} are quasi-coherent $\sh O_{Z'}$-modules by
\sref{III.6.5.1}.  We readily reduce to the case in which $S$,
$X'$, and $Y'$ are affine; then $X$ and $Y$ are affine as well,
and the assertion follows immediately from
\sref{III.6.4.1.1} and \sref[I]{I.1.6.3}.
\end{proof}

\begin{env}[6.5.12]
\label{III.6.5.12}
\emph{Remark.}
Let $X'$ and $Y'$ be two $S$-preschemes, with $X'$ flat over
$S$.  Let $X$ be a closed subprescheme of $X'$, and let
$i:X\to X'$ be the canonical immersion.  Let $\sh P_\bullet$
(respectively $\sh Q_\bullet$) be a bounded-below complex of
quasi-coherent $\sh O_X$-modules (respectively
$\sh O_{Y'}$-modules).  Finally, let $\sh L'_{\bullet\bullet}$
be a resolution of $i_*\sh P_\bullet$ consisting of locally free
$\sh O_{X'}$-modules, such that every point of $X'$ has an
affine open neighborhood $U$ for which
$\sh L'_{j\bullet}|U$ is a free resolution of
$i_*\sh P_j|U$, for every $j$.  There is then a canonical
isomorphism
\[
  (i\times_S1)_*
  \bigl(\shTor_\bullet^S(\sh P_\bullet,\sh Q_\bullet)\bigr)
  \xrightarrow{\sim}
  \sh H_\bullet(\sh L'_{\bullet\bullet}\otimes_S\sh Q_\bullet).
  \tag{6.5.12.1}\label{III.6.5.12.1}
\]

If $S$, $X'$, and $Y'$ are affine, and if
$\sh L'_{j\bullet}$ is a free resolution of $i_*\sh P_j$ for
every $j$, then \sref{III.6.5.11} reduces us to the case
$X'=X$, with $X$ flat over $S$, and it is enough to apply
\sref{III.6.3.4}.  In general, \sref{III.6.5.12.1} is defined
locally, and it remains to verify that these local definitions
give a global isomorphism.

For this, recall the definition of the first isomorphism
\sref{III.6.3.4.1}.  It comes from an isomorphism of spectral
sequences (\sref[0]{0.11.6.5} and \sref[0]{0.11.5.3}), itself
obtained from a morphism of bicomplexes
\[
  L_{\bullet\bullet}\longrightarrow L''_{\bullet\bullet},
\]
where $L_{\bullet\bullet}$ is the given resolution of
$P_\bullet$, and
\oldpage[III]{152}
$L''_{\bullet\bullet}$ is a projective resolution of
$P_\bullet$ in the category of bounded-below complexes of
$A$-modules (\sref[0]{0.11.5.2.2}).  The assertion follows from
the fact that \sref{III.6.3.4.1} is independent of the chosen
projective resolution $L''_{\bullet\bullet}$, since any two such
resolutions are related by a homotopism (\emph{M}, V, 1.2).
\end{env}

\begin{proposition}[6.5.13]
\label{III.6.5.13}
Let $X$ and $Y$ be two $S$-preschemes, and suppose that one of
the following conditions holds:
\begin{enumerate}
  \item[(i)] $X$ and $Z=X\times_SY$ are locally Noetherian, and
    $X$ is flat over $S$;
  \item[(ii)] $S$ and $X$ are locally Noetherian, and $Y$ is of
    finite type over $S$.
\end{enumerate}
Let $\sh P_\bullet$ (respectively $\sh Q_\bullet$) be a
bounded-below complex of quasi-coherent $\sh O_X$-modules
(respectively $\sh O_Y$-modules).  Suppose in addition that, for
every $n$, $\sh H_n(\sh P_\bullet)$ (respectively
$\sh H_n(\sh Q_\bullet)$) is an $\sh O_X$-module
(respectively $\sh O_Y$-module) of finite type.  Then
\[
  \shTor_n^S(\sh P_\bullet,\sh Q_\bullet)
\]
is a coherent $\sh O_Z$-module for every $n$.
\end{proposition}

\begin{proof}
Since $\sh P_\bullet$ and $\sh Q_\bullet$ are bounded below, the
spectral sequence
$''\sh E(\sh P_\bullet,\sh Q_\bullet)$ is biregular by
\sref{III.6.5.4}.  In either case (i) or (ii), $Z$ is locally
Noetherian, so \sref[0]{0.11.1.8} reduces the assertion to the
coherence of the terms $''\sh E_{pq}^2$.  The finite-type
hypothesis on the homology modules and formula
\sref{III.6.5.4.2} reduce the proposition to the special case in
which $\sh P_\bullet$ and $\sh Q_\bullet$ are concentrated in
degree zero, namely the following corollary.
\end{proof}

\begin{corollary}[6.5.14]
\label{III.6.5.14}
Suppose that one of conditions \emph{(i)} and \emph{(ii)} of
\sref{III.6.5.13} holds, and let $\sh F$ (respectively $\sh G$)
be a quasi-coherent $\sh O_X$-module (respectively
$\sh O_Y$-module) of finite type.  Then
$\shTor_n^S(\sh F,\sh G)$ is a coherent $\sh O_Z$-module for
every $n$.
\end{corollary}

\begin{proof}
The question is local on $X$ and $Y$, so we may suppose that
$S$, $X$, and $Y$ are affine.

\textnormal{(i)} Under the hypotheses of (i), $X$ and $Z$
are Noetherian.  There is therefore a locally free resolution
$\sh L_\bullet$ of $\sh F$ consisting of $\sh O_X$-modules of
finite type (\sref[I]{I.1.3.7}).  Since $X$ is flat over $S$,
\sref{III.6.3.4} gives
\[
  \shTor_n^S(\sh F,\sh G)
  =
  \sh H_n(\sh L_\bullet\otimes_S\sh G).
\]
The modules $\sh L_i\otimes_S\sh G$ are quasi-coherent
$\sh O_Z$-modules of finite type (\sref[I]{I.9.1.1}), hence are
coherent.  It follows that
$\sh H_n(\sh L_\bullet\otimes_S\sh G)$ is coherent
(\sref[0]{0.5.3.4}).

\textnormal{(ii)} Now suppose that condition (ii) holds.  Since
$A(Y)$ is a quotient of a polynomial $A(S)$-algebra in finitely
many indeterminates (\sref[I]{I.6.3.3}), $Y$ is a closed
subprescheme of an affine $S$-scheme $Y'$ that is flat and of
finite type over $S$.  The scheme $Y'$ is Noetherian
(\sref[I]{I.6.3.7}), so there exists a locally free resolution
$\sh M_\bullet$ of $j_*\sh G$ by $\sh O_{Y'}$-modules of finite
type, where $j:Y\to Y'$ is the canonical immersion.  By
\sref{III.6.5.12}, $\shTor_n^S(\sh F,\sh G)$ is the inverse image
under $1\times j$ of the $\sh O_{Z'}$-module
\[
  \sh H_n(\sh F\otimes_S\sh M_\bullet),
  \qquad Z'=X\times_SY'.
\]
As in (i), the $\sh F\otimes_S\sh M_i$ are coherent
$\sh O_{Z'}$-modules, and \sref[0]{0.5.3.4} again gives the
conclusion.
\end{proof}

\begin{env}[6.5.15]
\label{III.6.5.15}
The theory developed above for two complexes
$\sh P_\bullet,\sh Q_\bullet$ of quasi-coherent sheaves on two
$S$-preschemes $X$ and $Y$ extends without difficulty to an
arbitrary finite family of $S$-preschemes $X^{(i)}$
$(1\leq i\leq m)$, each equipped with a complex
$\sh P_\bullet^{(i)}$ of quasi-coherent
$\sh O_{X^{(i)}}$-modules.  If
\[
  Z=X^{(1)}\times_SX^{(2)}\times_S\cdots\times_SX^{(m)},
\]
one thereby defines a quasi-coherent $\sh O_Z$-module
\[
  \shTor_q^S
  \bigl(\sh P_\bullet^{(1)},\sh P_\bullet^{(2)},\ldots,
        \sh P_\bullet^{(m)}\bigr).
\]
We leave the development of the theory in this general case to
the reader, and record only, for later reference, the $E_2$-term
of the second, regular spectral sequence whose abutment is this
hypertor:

\oldpage[III]{153}

\[
  ''\sh E_{pq}^2
  =
  \bigoplus_{q_1+q_2+\cdots+q_m=q}
  \shTor_p^S
  \bigl(
    \sh H_{q_1}(\sh P_\bullet^{(1)}),\ldots,
    \sh H_{q_m}(\sh P_\bullet^{(m)})
  \bigr).
  \tag{6.5.15.1}\label{III.6.5.15.1}
\]
In \sref{III.6.8} we shall study the associativity spectral
sequences to which these hypertor functors of an arbitrary
finite number of complexes give rise.
\end{env}

\subsection{Global hypertor functors of complexes of
quasi-coherent modules and K\"unneth spectral sequences: the
affine-base case}
\label{subsection:III.6.6}

\begin{env}[6.6.1]
\label{III.6.6.1}
Let $S=\Spec(A)$ be an affine scheme, and let $X^{(i)}$
$(i=1,2)$ be two quasi-compact $S$-schemes.  For $i=1,2$, let
$\sh P_\bullet^{(i)}$ be a bounded-below complex of
quasi-coherent $\sh O_{X^{(i)}}$-modules whose differential has
degree $-1$.  Choose a finite covering
\[
  \mathfrak U^{(i)}=(U_\alpha^{(i)})
\]
of $X^{(i)}$ by affine open subsets.  Put
\[
  X=X^{(1)}\times_SX^{(2)}.
\]
This is a quasi-compact $S$-scheme
(\sref[I]{I.5.5.1} and \sref[I]{I.6.6.4}), and
$\mathfrak U=\mathfrak U^{(1)}\times_S\mathfrak U^{(2)}$ is the
covering of $X$ by the affine open subsets
$U_\alpha^{(1)}\times_SU_\beta^{(2)}$.

For $p\leq0$ and $q\in\bb Z$, the group of $(-p)$-alternating
cochains
\[
  C^{-p}(\mathfrak U^{(i)},\sh P_q^{(i)})
\]
of the covering $\mathfrak U^{(i)}$ with coefficients in
$\sh P_q^{(i)}$ (\emph{G}, II, 5.1) is an $A$-module.  For
$p>0$, put
$C^p(\mathfrak U^{(i)},\sh P_q^{(i)})=0$.  This defines a
bicomplex of $A$-modules
\[
  C^\bullet(\mathfrak U^{(i)},\sh P_\bullet^{(i)})
\]
whose two differentials have degree $-1$.

It follows from the definitions (\sref[0]{0.12.1.2}) that the
homology module
\[
  H_n\bigl(C^\bullet
    (\mathfrak U^{(i)},\sh P_\bullet^{(i)})\bigr)
\]
of this bicomplex, viewed in the usual way as a simple complex
for the total degree, is the hypercohomology module
\[
  \HH^{-n}
  \bigl(\mathfrak U^{(i)},\sh P^{(i)\bullet}\bigr),
\]
where $\sh P^{(i)\bullet}$ is the complex with differential of
degree $+1$ whose component of degree $q$ is
$\sh P_{-q}^{(i)}$.  By abuse of notation we write this as
\[
  \HH^{-n}
  \bigl(\mathfrak U^{(i)},\sh P_\bullet^{(i)}\bigr).
\]
By \sref{III.6.2.2}, this module is canonically isomorphic to
\[
  \HH^{-n}\bigl(X^{(i)},\sh P^{(i)\bullet}\bigr),
\]
also written
$\HH^{-n}(X^{(i)},\sh P_\bullet^{(i)})$.  It is therefore
independent of the chosen finite covering
$\mathfrak U^{(i)}$.
\end{env}

\begin{env}[6.6.2]
\label{III.6.6.2}
Apply the general theory of hyperhomology of functors with
respect to bicomplexes (\sref[0]{0.11.7.4}) to the two
bicomplexes of $A$-modules
\[
  L_{\bullet\bullet}^{(i)}
  =
  C^\bullet(\mathfrak U^{(i)},\sh P_\bullet^{(i)})
  \qquad(i=1,2)
\]
and to the covariant bifunctor
$L_{\bullet\bullet}^{(1)}\otimes_A
 L_{\bullet\bullet}^{(2)}$.
Because the cochains are alternating and the coverings
$\mathfrak U^{(i)}$ are finite, the modules
\[
  C^{-p}(\mathfrak U^{(i)},\sh P_q^{(i)})
\]
are nonzero for only finitely many values of $p$, independently
of $q$.  In particular, both degrees of each
$L_{\bullet\bullet}^{(i)}$ are bounded below.

Denote by
\[
  \mathbf{Tor}_n^A
  (L_{\bullet\bullet}^{(1)},L_{\bullet\bullet}^{(2)})
\]
or by
\[
  \mathbf{Tor}_n^S
  (\mathfrak U^{(1)},\mathfrak U^{(2)};
   \sh P_\bullet^{(1)},\sh P_\bullet^{(2)})
\]
the $n$-th hyperhomology module of
$L_{\bullet\bullet}^{(1)}\otimes_A
 L_{\bullet\bullet}^{(2)}$.  We call it the hypertor of index
$n$ of $\sh P_\bullet^{(1)}$ and $\sh P_\bullet^{(2)}$ relative
to the coverings $\mathfrak U^{(1)}$ and
$\mathfrak U^{(2)}$.  When these complexes are concentrated in
degree zero, with terms $\sh F^{(1)}$ and $\sh F^{(2)}$, write
\[
  \mathbf{Tor}_n^S
  (\mathfrak U^{(1)},\mathfrak U^{(2)};
   \sh F^{(1)},\sh F^{(2)})
\]
for their hypertor.  Following the general conventions, the
same notation
\[
  \mathbf{Tor}_n^S
  (\mathfrak U^{(1)},\mathfrak U^{(2)};
   \sh P_\bullet^{(1)},\sh P_\bullet^{(2)})
\]
also denotes the bicomplex whose component of bidegree $(j,k)$
is
\[
  \mathbf{Tor}_n^S
  (\mathfrak U^{(1)},\mathfrak U^{(2)};
   \sh P_j^{(1)},\sh P_k^{(2)}).
\]

The construction $L_{\bullet\bullet}^{(i)}$ is exact in
$\sh P_\bullet^{(i)}$, since all intersections of members of
\oldpage[III]{154}
$\mathfrak U^{(i)}$ are affine
(\sref[I]{I.5.5.6} and \sref[I]{I.1.3.11}).  Consequently
\[
  \mathbf{Tor}_\bullet^S
  (\mathfrak U^{(1)},\mathfrak U^{(2)};
   \sh P_\bullet^{(1)},\sh P_\bullet^{(2)})
\]
is a covariant bi-$\partial$-functor in
$\sh P_\bullet^{(1)}$ and $\sh P_\bullet^{(2)}$, with values in
the category of $A$-modules (\sref[0]{0.11.7.3}).  By
\sref[0]{0.11.7.2}, it is the common abutment of six biregular
spectral functors, denoted
\[
  {}^{(t)}E^S
  (\mathfrak U^{(1)},\mathfrak U^{(2)};
   \sh P_\bullet^{(1)},\sh P_\bullet^{(2)}),
\]
or simply
\[
  {}^{(t)}E
  (\mathfrak U^{(1)},\mathfrak U^{(2)};
   \sh P_\bullet^{(1)},\sh P_\bullet^{(2)}),
\]
where $t$ is one of $a,b,a',b',c,d$.  Their $E_2$-terms are
\[
\begin{aligned}
  {}^{(a)}E_{pq}^2
  &=
  \bigoplus_{q_1+q_2=q}
  \Tor_p^A
  \bigl(H_{q_1}(L_{\bullet\bullet}^{(1)}),
        H_{q_2}(L_{\bullet\bullet}^{(2)})\bigr),\\
  {}^{(b)}E_{pq}^2
  &=
  H_p\bigl(
    \mathbf{Tor}_q^{A,\mathrm{II}}
    (L_{\bullet\bullet}^{(1)},L_{\bullet\bullet}^{(2)})
  \bigr),\\
  {}^{(a')}E_{pq}^2
  &=
  \bigoplus_{q_1+q_2=q}
  \Tor_p^A
  \bigl(H_{q_1}^{\mathrm I}(L_{\bullet\bullet}^{(1)}),
        H_{q_2}^{\mathrm I}(L_{\bullet\bullet}^{(2)})\bigr),\\
  {}^{(b')}E_{pq}^2
  &=
  H_p\bigl(
    \mathbf{Tor}_q^A
    (L_{\bullet\bullet}^{(1)},L_{\bullet\bullet}^{(2)})
  \bigr),\\
  {}^{(c)}E_{pq}^2
  &=
  \bigoplus_{q_1+q_2=q}
  \Tor_p^A
  \bigl(H_{q_1}^{\mathrm{II}}(L_{\bullet\bullet}^{(1)}),
        H_{q_2}^{\mathrm{II}}(L_{\bullet\bullet}^{(2)})\bigr),\\
  {}^{(d)}E_{pq}^2
  &=
  H_p\bigl(
    \mathbf{Tor}_q^{A,\mathrm I}
    (L_{\bullet\bullet}^{(1)},L_{\bullet\bullet}^{(2)})
  \bigr).
\end{aligned}
\]
These notations agree with those of the general theory of
hyperhomology.  We now make the initial terms more explicit.
\end{env}

\begin{env}[6.6.3]
\label{III.6.6.3}
\emph{Spectral sequences \textnormal{(a)} and
\textnormal{(a$'$)}.}
By \sref{III.6.6.1},
\[
  H_n(L_{\bullet\bullet}^{(i)})
  =
  \HH^{-n}(X^{(i)},\sh P_\bullet^{(i)}),
\]
and hence
\[
  {}^{(a)}E_{pq}^2
  =
  \bigoplus_{q_1+q_2=q}
  \Tor_p^A
  \bigl(
    \HH^{-q_1}(X^{(1)},\sh P_\bullet^{(1)}),
    \HH^{-q_2}(X^{(2)},\sh P_\bullet^{(2)})
  \bigr).
\]

By definition, the component of degree $k$ of the complex
$H_n^{\mathrm I}(L_{\bullet\bullet}^{(i)})$ is
\[
  H_n\bigl(
    C^\bullet(\mathfrak U^{(i)},\sh P_k^{(i)})
  \bigr)
  =
  \HH^{-n}(\mathfrak U^{(i)},\sh P_k^{(i)}).
\]
By \sref{III.1.4.1}, this module is canonically isomorphic to
$\HH^{-n}(X^{(i)},\sh P_k^{(i)})$.  Therefore
\[
  {}^{(a')}E_{pq}^2
  =
  \bigoplus_{q_1+q_2=q}
  \Tor_p^A
  \bigl(
    \HH^{-q_1}(X^{(1)},\sh P_\bullet^{(1)}),
    \HH^{-q_2}(X^{(2)},\sh P_\bullet^{(2)})
  \bigr),
\]
where each $\HH^{-q_i}(X^{(i)},\sh P_\bullet^{(i)})$ in this
last formula denotes the complex obtained by taking cohomology
term by term.
\end{env}

\begin{env}[6.6.4]
\label{III.6.6.4}
\emph{Spectral sequences \textnormal{(b)} and
\textnormal{(b$'$)}.}
By definition,
\[
  \mathbf{Tor}_q^{A,\mathrm{II}}
  (L_{\bullet\bullet}^{(1)},L_{\bullet\bullet}^{(2)})
\]
is a bicomplex whose component of bidegree $(h,k)$ is
\[
  \Tor_q^A
  \bigl(
    C^{-h}(\mathfrak U^{(1)},\sh P_\bullet^{(1)}),
    C^{-k}(\mathfrak U^{(2)},\sh P_\bullet^{(2)})
  \bigr).
\]
Let $\Phi^{(i)}$ be the index set of $\mathfrak U^{(i)}$.  For
$r\geq0$, the complex
$C^r(\mathfrak U^{(i)},\sh P_\bullet^{(i)})$ is the direct sum
of the complexes
\[
  \Gamma(U_\rho^{(i)},\sh P_\bullet^{(i)}),
\]
where $U_\rho^{(i)}$ is the intersection of the
$U_\xi^{(i)}$ for $\xi\in\rho$, and $\rho$ ranges over the
appropriate elements of $\mathfrak P(\Phi^{(i)})$.
Consequently
\[
  \Tor_q^A
  \bigl(
    C^{-h}(\mathfrak U^{(1)},\sh P_\bullet^{(1)}),
    C^{-k}(\mathfrak U^{(2)},\sh P_\bullet^{(2)})
  \bigr)
\]
is the direct sum of the $A$-modules
\[
  \Tor_q^A
  \bigl(
    \Gamma(U_\sigma^{(1)},\sh P_\bullet^{(1)}),
    \Gamma(U_\tau^{(2)},\sh P_\bullet^{(2)})
  \bigr),
\]
where $\sigma$ and $\tau$ range over
$\mathfrak P(\Phi^{(1)})$ and $\mathfrak P(\Phi^{(2)})$ with
\[
  \operatorname{Card}(\sigma)=1-h,
  \qquad
  \operatorname{Card}(\tau)=1-k,
\]
respectively.  Since $X^{(1)}$ and $X^{(2)}$ are schemes, the
$U_\rho^{(i)}$ are affine.  Thus \sref{III.6.4.1.1} gives
\[
  \Tor_q^A
  \bigl(
    \Gamma(U_\sigma^{(1)},\sh P_\bullet^{(1)}),
    \Gamma(U_\tau^{(2)},\sh P_\bullet^{(2)})
  \bigr)
  =
  \Gamma\bigl(
    U_\sigma^{(1)}\times_SU_\tau^{(2)},
    \shTor_q^S(\sh P_\bullet^{(1)},\sh P_\bullet^{(2)})
  \bigr).
\]

\oldpage[III]{155}

It follows that ${}^{(b)}E_{pq}^2$ is the $(-p)$-th cohomology
module of the complex
\[
  L^\bullet(\Phi^{(1)},\Phi^{(2)};\mathcal S)
\]
of bialternating cochains on $\Phi^{(1)}$ and $\Phi^{(2)}$ with
values in the coefficient system
\[
  \mathcal S:(\sigma,\tau)\longmapsto
  \Gamma\bigl(
    U_\sigma^{(1)}\times_SU_\tau^{(2)},
    \shTor_q^S(\sh P_\bullet^{(1)},\sh P_\bullet^{(2)})
  \bigr)
  \tag{6.6.4.1}\label{III.6.6.4.1}
\]
(\sref[0]{0.11.8.4}).  By \sref[0]{0.11.8.5} and
\sref[0]{0.11.8.6}, this complex has the same cohomology as the
complex of all cochains
\[
  C^\bullet(\Phi^{(1)},\Phi^{(2)};\mathcal S)
\]
and as the complex
\[
  P^\bullet(\Phi^{(1)},\Phi^{(2)};\mathcal S),
\]
whose elements are linear combinations of
\[
  \lambda(\sigma,\tau)\in
  \Gamma\bigl(
    U_\sigma^{(1)}\times_SU_\tau^{(2)},
    \shTor_q^S(\sh P_\bullet^{(1)},\sh P_\bullet^{(2)})
  \bigr),
\]
where
\[
  \sigma=(\alpha_0,\ldots,\alpha_h),
  \qquad
  \tau=(\beta_0,\ldots,\beta_h)
\]
are sequences with the same number of terms.  Now
\[
\begin{split}
  U_\sigma^{(1)}\times_SU_\tau^{(2)}
  &=
  (U_{\alpha_0}^{(1)}\times_SU_{\beta_0}^{(2)})
  \cap\cdots\cap
  (U_{\alpha_h}^{(1)}\times_SU_{\beta_h}^{(2)})
\end{split}
\]
by \sref[I]{I.3.2.7}.  Let $\mathfrak U$ be the covering of
$Z=X^{(1)}\times_SX^{(2)}$ by the affine open subsets
$U_\alpha^{(1)}\times_SU_\beta^{(2)}$.  Since $Z$ is a scheme,
\sref{III.1.3.1} yields
\[
  {}^{(b)}E_{pq}^2
  =
  \HH^{-p}
  \bigl(
    Z,\shTor_q^S
    (\sh P_\bullet^{(1)},\sh P_\bullet^{(2)})
  \bigr).
  \tag{6.6.4.2}\label{III.6.6.4.2}
\]

Next,
\[
  \mathbf{Tor}_q^A
  (L_{\bullet\bullet}^{(1)},L_{\bullet\bullet}^{(2)})
\]
is a bicomplex whose component of bidegree $(h,k)$ is the direct
sum of the $A$-modules
\[
  \Tor_q^A
  \bigl(
    C^{-h_1}(\mathfrak U^{(1)},\sh P_{k_1}^{(1)}),
    C^{-h_2}(\mathfrak U^{(2)},\sh P_{k_2}^{(2)})
  \bigr),
\]
where $h_1+h_2=h$ and $k_1+k_2=k$.  Expanding the cochain
modules as above shows that this component is also the direct
sum of
\[
  \Gamma\bigl(
    U_\sigma^{(1)}\times_SU_\tau^{(2)},
    \shTor_q^S(\sh P_{k_1}^{(1)},\sh P_{k_2}^{(2)})
  \bigr),
\]
where $k_1+k_2=k$ and
\[
  \operatorname{Card}(\sigma)
  +\operatorname{Card}(\tau)=2-h.
\]

The term ${}^{(b')}E_{pq}^2$ is the $(-p)$-th cohomology module
of a bicomplex $N^{\bullet\bullet}=(N^{hk})$.  For fixed $k$,
the simple complex $N^{\bullet k}$ consists of bialternating
cochains on $\Phi^{(1)}$ and $\Phi^{(2)}$ with values in the
coefficient system
\[
  \mathcal S_k:(\sigma,\tau)\longmapsto
  \Gamma\left(
    U_\sigma^{(1)}\times_SU_\tau^{(2)},
    \bigoplus_{k_1+k_2=k}
    \shTor_q^S(\sh P_{-k_1}^{(1)},\sh P_{-k_2}^{(2)})
  \right).
\]
These coefficient systems form a complex
$\mathcal S^\bullet$, whose differential is induced by that of
the simple complex associated to the bicomplex
\[
  \shTor_q^S
  (\sh P^{(1)\bullet},\sh P^{(2)\bullet}).
\]
By \sref[0]{0.11.8.9}, $N^{\bullet\bullet}$ has the same
cohomology as
\[
  C^\bullet
  (\Phi^{(1)},\Phi^{(2)};\mathcal S^\bullet),
\]
and also the same cohomology as
\[
  P^\bullet
  (\Phi^{(1)},\Phi^{(2)};\mathcal S^\bullet),
\]
whose elements of bidegree $(h,k)$ are linear combinations of
\[
  \lambda(\sigma,\tau)\in
  \Gamma\left(
    U_\sigma^{(1)}\times_SU_\tau^{(2)},
    \bigoplus_{k_1+k_2=k}
    \shTor_q^S(\sh P_{-k_1}^{(1)},\sh P_{-k_2}^{(2)})
  \right),
\]
with $\sigma=(\alpha_0,\ldots,\alpha_h)$ and
$\tau=(\beta_0,\ldots,\beta_h)$ having the same number of terms
(\sref[0]{0.11.8.10}).  As above, ${}^{(b')}E_{pq}^2$ is
therefore the $(-p)$-th cohomology module of
\[
  C^\bullet(\mathfrak U,\mathcal Q^\bullet),
\]
where $\mathcal Q^\bullet$ is the simple complex associated to
the bicomplex
\[
  \shTor_q^S
  (\sh P^{(1)\bullet},\sh P^{(2)\bullet})
\]
of $\sh O_Z$-modules.  With the conventions of
\sref{III.6.6.1}, we obtain
\[
  {}^{(b')}E_{pq}^2
  =
  \HH^{-p}
  \bigl(
    Z,\shTor_q^S
    (\sh P_\bullet^{(1)},\sh P_\bullet^{(2)})
  \bigr).
  \tag{6.6.4.3}\label{III.6.6.4.3}
\]
\end{env}

\oldpage[III]{156}

\begin{env}[6.6.5]
\label{III.6.6.5}
\emph{Spectral sequences \textnormal{(c)} and
\textnormal{(d)}.}
By definition, the component of degree $h$ of the complex
$H_n^{\mathrm{II}}(L_{\bullet\bullet}^{(i)})$ is the
$A$-module
\[
  C^{-h}
  \bigl(
    \mathfrak U^{(i)},
    \sh H_n(\sh P_\bullet^{(i)})
  \bigr),
\]
because the functor $C^{-h}$ is exact.  Hence, by the definition
of the hypertor of two modules relative to two coverings
(\sref{III.6.6.2}),
\[
  {}^{(c)}E_{pq}^2
  =
  \bigoplus_{q_1+q_2=q}
  \mathbf{Tor}_p^S
  \bigl(
    \mathfrak U^{(1)},\mathfrak U^{(2)};
    \sh H_{q_1}(\sh P_\bullet^{(1)}),
    \sh H_{q_2}(\sh P_\bullet^{(2)})
  \bigr).
\]

Finally, by definition,
\[
  \mathbf{Tor}_q^{A,\mathrm I}
  (L_{\bullet\bullet}^{(1)},L_{\bullet\bullet}^{(2)})
\]
is a bicomplex whose component of bidegree $(h,k)$ is the
$A$-module
\[
  \mathbf{Tor}_q^S
  \bigl(
    \mathfrak U^{(1)},\mathfrak U^{(2)};
    \sh P_h^{(1)},\sh P_k^{(2)}
  \bigr).
\]
Consequently
\[
  {}^{(d)}E_{pq}^2
  =
  H_p\left(
    \mathbf{Tor}_q^S
    \bigl(
      \mathfrak U^{(1)},\mathfrak U^{(2)};
      \sh P_\bullet^{(1)},\sh P_\bullet^{(2)}
    \bigr)
  \right).
\]
\end{env}

\begin{env}[6.6.6]
\label{III.6.6.6}
The theory of hyperhomology for functors of bicomplexes
(\sref[0]{0.11.7.3}) shows, as in \sref{III.6.3.4}, that for
every \emph{flat} Cartan--Eilenberg resolution
$M_{\bullet\bullet}^{(i)}$ of $L_{\bullet\bullet}^{(i)}$ in the
category of complexes of modules bounded below $(i=1,2)$, there
are canonical isomorphisms of bi-$\partial$-functors
\begin{env}[6.6.6.1]
\label{III.6.6.6.1}
\[
\begin{split}
  \mathbf{Tor}_\bullet^S
  \bigl(
    \mathfrak U^{(1)},\mathfrak U^{(2)};
    \sh P_\bullet^{(1)},\sh P_\bullet^{(2)}
  \bigr)
  &\xrightarrow{\sim}
  H_\bullet
  \bigl(
    M_{\bullet\bullet}^{(1)}
    \otimes_A M_{\bullet\bullet}^{(2)}
  \bigr)\\
  &\xrightarrow{\sim}
  H_\bullet
  \bigl(
    M_{\bullet\bullet}^{(1)}
    \otimes_A L_{\bullet\bullet}^{(2)}
  \bigr)\\
  &\xrightarrow{\sim}
  H_\bullet
  \bigl(
    L_{\bullet\bullet}^{(1)}
    \otimes_A M_{\bullet\bullet}^{(2)}
  \bigr).
\end{split}
\tag{6.6.6.1}
\]
\end{env}
\end{env}

\begin{env}[6.6.7]
\label{III.6.6.7}
We shall now show that the global hypertor defined in
\sref{III.6.6.2}, and the six corresponding spectral sequences,
do not depend, up to canonical isomorphisms, on the finite affine
open coverings $\mathfrak U^{(i)}$ used to define them.  It is
enough to show that if $\mathfrak V^{(i)}$ are two other
coverings of the same kind, with $\mathfrak V^{(i)}$ finer than
$\mathfrak U^{(i)}$ for $i=1,2$, then there are canonical
isomorphisms of spectral functors
\begin{env}[6.6.7, $t$]
\label{III.6.6.7.t}
\[
  {}^{(t)}E
  \bigl(
    \mathfrak U^{(1)},\mathfrak U^{(2)};
    \sh P_\bullet^{(1)},\sh P_\bullet^{(2)}
  \bigr)
  \xrightarrow{\sim}
  {}^{(t)}E
  \bigl(
    \mathfrak V^{(1)},\mathfrak V^{(2)};
    \sh P_\bullet^{(1)},\sh P_\bullet^{(2)}
  \bigr),
  \tag{6.6.7, $t$}\label{III.6.6.7.t.eq}
\]
\end{env}
where $t$ is one of $a,b,a',b',c,d$.

For $i=1,2$ there are homomorphisms of bicomplexes
\[
  C^\bullet
  (\mathfrak U^{(i)},\sh P_\bullet^{(i)})
  \longrightarrow
  C^\bullet
  (\mathfrak V^{(i)},\sh P_\bullet^{(i)})
\]
that are well defined up to homotopy
(\sref[G]{G.II.5.7.1}).  It follows at once that the
homomorphisms \eqref{III.6.6.7.t.eq} are canonically defined and
compatible with the boundary operators in their abutments
(\sref[0]{0.11.3.2}).  Moreover, the calculation of the
$E_2$-terms of the spectral sequences
\textnormal{(a)}, \textnormal{(b)}, \textnormal{(a$'$)}, and
\textnormal{(b$'$)} shows that
\eqref{III.6.6.7.t.eq} is an isomorphism on their $E_2$-terms.
Since these spectral sequences are biregular, it follows that
\eqref{III.6.6.7.t.eq} is an isomorphism for these spectral
functors, and hence an isomorphism of bi-$\partial$-functors on
their common abutment (\sref[0]{0.11.1.5}).

In particular, if $\sh F^{(i)}$ is a quasi-coherent
$\sh O_{X^{(i)}}$-module $(i=1,2)$, the canonical homomorphism
\[
  \mathbf{Tor}_\bullet^S
  \bigl(
    \mathfrak U^{(1)},\mathfrak U^{(2)};
    \sh F^{(1)},\sh F^{(2)}
  \bigr)
  \longrightarrow
  \mathbf{Tor}_\bullet^S
  \bigl(
    \mathfrak V^{(1)},\mathfrak V^{(2)};
    \sh F^{(1)},\sh F^{(2)}
  \bigr)
\]
is bijective.  In view of the calculation in
\sref{III.6.6.5}, the homomorphism
\eqref{III.6.6.7.t.eq} is therefore also an isomorphism on
$E_2$-terms for $t=c$ and $t=d$.  As above, we conclude that it
is an isomorphism of spectral sequences for these two values of
$t$.

\oldpage[III]{157}

The isomorphisms \eqref{III.6.6.7.t.eq} may be regarded as
defining inductive systems of spectral functors on the filtering
set of pairs
$(\mathfrak U^{(1)},\mathfrak U^{(2)})$ of finite affine open
coverings of $X^{(1)}$ and $X^{(2)}$.  We denote the inductive
limit by
\[
  {}^{(t)}E
  \bigl(
    X^{(1)},X^{(2)};
    \sh P_\bullet^{(1)},\sh P_\bullet^{(2)}
  \bigr)
  \quad\text{or}\quad
  {}^{(t)}E^S
  \bigl(
    X^{(1)},X^{(2)};
    \sh P_\bullet^{(1)},\sh P_\bullet^{(2)}
  \bigr).
\]
Its abutment is denoted
\[
  \mathbf{Tor}_\bullet^S
  \bigl(
    X^{(1)},X^{(2)};
    \sh P_\bullet^{(1)},\sh P_\bullet^{(2)}
  \bigr)
\]
and is called the \emph{global hypertor} of the two complexes
$\sh P_\bullet^{(1)}$ and $\sh P_\bullet^{(2)}$.  If the two
complexes are concentrated in degree zero, with terms
$\sh F^{(1)}$ and $\sh F^{(2)}$, we write
\[
  \mathbf{Tor}_\bullet^S
  \bigl(
    X^{(1)},X^{(2)};
    \sh F^{(1)},\sh F^{(2)}
  \bigr).
\]
In accordance with the general conventions,
\[
  \mathbf{Tor}_q^S
  \bigl(
    X^{(1)},X^{(2)};
    \sh P_\bullet^{(1)},\sh P_\bullet^{(2)}
  \bigr)
\]
will denote the bicomplex whose component of bidegree $(h,k)$ is
\[
  \mathbf{Tor}_q^S
  \bigl(
    X^{(1)},X^{(2)};
    \sh P_h^{(1)},\sh P_k^{(2)}
  \bigr).
\]
\end{env}

\begin{env}[6.6.8]
\label{III.6.6.8}
Under the hypotheses of \sref{III.6.6.1}, consider two
$S$-morphisms
\[
  f_i:X^{(i)}\longrightarrow Y^{(i)},
\]
where
$Y^{(i)}=\Spec(B_i)$ is an affine $S$-scheme, so that $B_i$ is
an $A$-algebra $(i=1,2)$.  This defines an $A$-homomorphism
\[
  B_i\longrightarrow
  \Gamma(X^{(i)},\sh O_{X^{(i)}})
\]
(\sref[I]{I.2.2.4}), and consequently each
$L_{\bullet\bullet}^{(i)}$ defined in \sref{III.6.6.2} is a
bicomplex of $B_i$-modules.  Thus
\[
  L_{\bullet\bullet}^{(1)}
  \otimes_A L_{\bullet\bullet}^{(2)}
\]
is a quadricomplex of $(B_1\otimes_AB_2)$-modules, and its six
hyperhomology spectral functors may be regarded as taking values
in the category of spectral sequences of
$(B_1\otimes_AB_2)$-modules.

Put
\[
  Y=Y^{(1)}\times_SY^{(2)}
   =\Spec(B_1\otimes_AB_2).
\]
We may then consider the quasi-coherent $\sh O_Y$-modules
associated with these modules (\sref[I]{I.1.3.4}).  For
$t=a,b,a',b',c,d$, write
\[
  {}^{(t)}\sh E
  \bigl(
    f_1,f_2;
    \sh P_\bullet^{(1)},\sh P_\bullet^{(2)}
  \bigr)
\]
for the six spectral sequences of $\sh O_Y$-modules
\[
  \left(
    {}^{(t)}E
    \bigl(
      X^{(1)},X^{(2)};
      \sh P_\bullet^{(1)},\sh P_\bullet^{(2)}
    \bigr)
  \right)\supertilde,
\]
and write
\[
  \mathfrak{Tor}_\bullet^S
  \bigl(
    f_1,f_2;
    \sh P_\bullet^{(1)},\sh P_\bullet^{(2)}
  \bigr)
\]
for their common abutment
\[
  \left(
    \mathbf{Tor}_\bullet^S
    \bigl(
      X^{(1)},X^{(2)};
      \sh P_\bullet^{(1)},\sh P_\bullet^{(2)}
    \bigr)
  \right)\supertilde.
\]
When $\sh P_\bullet^{(i)}$ is concentrated in degree zero, with
term $\sh F^{(i)}$ $(i=1,2)$, we write this abutment as
\[
  \mathfrak{Tor}_\bullet^S
  \bigl(
    f_1,f_2;
    \sh F^{(1)},\sh F^{(2)}
  \bigr).
\]
\end{env}


\subsection{Global hypertor functors of complexes of quasi-coherent modules and K\"unneth spectral sequences: the general case}
\label{subsection:III.6.7}

\begin{env}[6.7.1]
\label{III.6.7.1}
We now generalize the definitions of \sref{III.6.6.8} to the
case where $S$ is an arbitrary prescheme, $X^{(i)}$ and
$Y^{(i)}$ are $S$-preschemes, and
\[
  f_i:X^{(i)}\longrightarrow Y^{(i)}
\]
are separated and quasi-compact $S$-morphisms.  For every pair
of bounded-below complexes $\sh P_\bullet^{(i)}$ of
quasi-coherent $\sh O_{X^{(i)}}$-modules $(i=1,2)$, the aim is to
define a quasi-coherent $\sh O_Y$-module
\[
  \mathfrak{Tor}_n^S
  \bigl(
    f_1,f_2;
    \sh P_\bullet^{(1)},\sh P_\bullet^{(2)}
  \bigr)
\]
for every $n$, together with six spectral functors, in such a way
that these definitions reduce to those of \sref{III.6.6.8} when
$S$, $Y^{(1)}$, and $Y^{(2)}$ are affine; here
$Y=Y^{(1)}\times_SY^{(2)}$.

First suppose that $S=\Spec(A)$ is affine, while $Y^{(1)}$ and
$Y^{(2)}$ are arbitrary.  Let $W^{(i)}$ be an affine open subset
of $Y^{(i)}$.  Then $f_i^{-1}(W^{(i)})$ is a quasi-compact
$S$-scheme, and
\[
  W=W^{(1)}\times_SW^{(2)}
\]
is an affine open subset of $Y$.  Let
\[
  f_i':f_i^{-1}(W^{(i)})\longrightarrow W^{(i)}
\]
be the restriction of $f_i$, and let
\[
  \sh P_\bullet^{\prime(i)}
  =
  \sh P_\bullet^{(i)}
  \bigm|_{f_i^{-1}(W^{(i)})}
  \qquad(i=1,2).
\]
By \sref{III.6.6.8} we then have the spectral sequences of
quasi-coherent $(\sh O_Y|W)$-modules
\[
  {}^{(t)}\sh E
  \bigl(
    f_1',f_2';
    \sh P_\bullet^{\prime(1)},\sh P_\bullet^{\prime(2)}
  \bigr),
\]
and it remains to verify that they satisfy the gluing conditions
of \sref[0]{0.3.3.1}.

We are immediately reduced to the case in which
\[
  Y^{(i)}=\Spec(B_i),
  \qquad
  W^{(i)}=D(g_i),
  \qquad g_i\in B_i,
\]
so that
\[
  W=D(g_1\otimes g_2)
\]
in $Y=\Spec(B_1\otimes_AB_2)$
(\sref[II]{II.4.3.2.1}).  Put
$X^{\prime(i)}=f_i^{-1}(W^{(i)})$.  We must establish an
isomorphism of spectral functors
\begin{env}[6.7.1.1]
\label{III.6.7.1.1}
\[
\begin{split}
  {}^{(t)}E
  \bigl(
    X^{\prime(1)},X^{\prime(2)};
    \sh P_\bullet^{\prime(1)},\sh P_\bullet^{\prime(2)}
  \bigr)
  \xrightarrow{\sim}
  {}^{(t)}E
  \bigl(
    X^{(1)},X^{(2)};
    \sh P_\bullet^{(1)},\sh P_\bullet^{(2)}
  \bigr)
  \otimes_BB_g ,
\end{split}
\tag{6.7.1.1}\label{III.6.7.1.1.eq}
\]
\end{env}
where
\[
  B=B_1\otimes_AB_2,
  \qquad
  g=g_1\otimes g_2.
\]

\oldpage[III]{158}

Choose finite affine open coverings $\mathfrak U^{(i)}$ of
$X^{(i)}$, and let $\mathfrak U^{\prime(i)}$ be their traces on
$X^{\prime(i)}$.  The latter are again coverings by affine open
subsets (\sref[I]{I.5.5.10}), and more precisely
\[
  C^\bullet
  \bigl(
    \mathfrak U^{\prime(i)},\sh P_\bullet^{\prime(i)}
  \bigr)
  =
  C^\bullet
  \bigl(
    \mathfrak U^{(i)},\sh P_\bullet^{(i)}
  \bigr)
  \otimes_{B_i}(B_i)_{g_i}.
\]
Put
\[
  L_{\bullet\bullet}^{(i)}
  =
  C^\bullet
  \bigl(
    \mathfrak U^{(i)},\sh P_\bullet^{(i)}
  \bigr),
  \qquad
  L_{\bullet\bullet}^{\prime(i)}
  =
  C^\bullet
  \bigl(
    \mathfrak U^{\prime(i)},\sh P_\bullet^{\prime(i)}
  \bigr).
\]
Then
\[
\begin{split}
  L_{\bullet\bullet}^{\prime(1)}
  \otimes_A L_{\bullet\bullet}^{\prime(2)}
  &=
  \bigl(
    L_{\bullet\bullet}^{(1)}
    \otimes_{B_1}(B_1)_{g_1}
  \bigr)
  \otimes_A
  \bigl(
    L_{\bullet\bullet}^{(2)}
    \otimes_{B_2}(B_2)_{g_2}
  \bigr).
\end{split}
\]
Since
\[
  B_g=(B_1)_{g_1}\otimes_A(B_2)_{g_2}
\]
up to canonical isomorphism, we likewise have, up to canonical
isomorphism,
\[
  L_{\bullet\bullet}^{\prime(1)}
  \otimes_A L_{\bullet\bullet}^{\prime(2)}
  =
  \bigl(
    L_{\bullet\bullet}^{(1)}
    \otimes_A L_{\bullet\bullet}^{(2)}
  \bigr)
  \otimes_BB_g.
\]

Let $M_{\bullet\bullet\bullet}^{(i)}$ be a projective
Cartan--Eilenberg resolution of $L_{\bullet\bullet}^{(i)}$,
chosen to consist of $B_i$-modules.  Because $(B_i)_{g_i}$ is
flat over $B_i$,
\[
  M_{\bullet\bullet\bullet}^{\prime(i)}
  =
  M_{\bullet\bullet\bullet}^{(i)}
  \otimes_{B_i}(B_i)_{g_i}
\]
is a projective Cartan--Eilenberg resolution of
$L_{\bullet\bullet}^{\prime(i)}$.  Moreover,
\[
  M_{\bullet\bullet\bullet}^{\prime(1)}
  \otimes_A M_{\bullet\bullet\bullet}^{\prime(2)}
  =
  \bigl(
    M_{\bullet\bullet\bullet}^{(1)}
    \otimes_A M_{\bullet\bullet\bullet}^{(2)}
  \bigr)
  \otimes_BB_g.
\]
The desired isomorphism \eqref{III.6.7.1.1.eq} now follows
immediately from the definitions of the hyperhomology of a
bicomplex and from the exactness of
$G\longmapsto G\otimes_BB_g$ on $B$-modules.
\end{env}

\begin{env}[6.7.2]
\label{III.6.7.2}
Now suppose that $S$ is arbitrary, and let
\[
  u_i:Y^{(i)}\longrightarrow S
  \qquad(i=1,2)
\]
be the structural morphisms.  Choose an affine open covering
$(S_\alpha)$ of $S$, put
\[
  Y_\alpha^{(i)}=u_i^{-1}(S_\alpha),
  \qquad
  X_\alpha^{(i)}=f_i^{-1}(Y_\alpha^{(i)}),
\]
and let
\[
  f_{i\alpha}:X_\alpha^{(i)}\longrightarrow Y_\alpha^{(i)}
\]
be the restriction of $f_i$, which is separated and
quasi-compact.  The schemes
\[
  Y_\alpha
  =
  Y_\alpha^{(1)}
  \times_{S_\alpha}Y_\alpha^{(2)}
\]
form an open covering of $Y$, and on every $Y_\alpha$,
\sref{III.6.7.1} defines the spectral functors
\[
  {}^{(t)}\sh E^{S_\alpha}
  \left(
    f_{1\alpha},f_{2\alpha};
    \sh P_\bullet^{(1)}|X_\alpha^{(1)},
    \sh P_\bullet^{(2)}|X_\alpha^{(2)}
  \right).
\]
It remains again to show that these functors satisfy the gluing
conditions.

We reduce immediately to the following situation:
$S=\Spec(A)$ is affine,
$S'=D(h)$ with $h\in A$, and
$u_i(Y^{(i)})\subset S'$.  We may moreover assume that
$Y^{(i)}=\Spec(B_i)$ is affine.  It is then enough to define
canonical isomorphisms
\begin{env}[6.7.2.1]
\label{III.6.7.2.1}
\[
  {}^{(t)}E^S
  \bigl(
    X^{(1)},X^{(2)};
    \sh P_\bullet^{(1)},\sh P_\bullet^{(2)}
  \bigr)
  \xrightarrow{\sim}
  {}^{(t)}E^{S'}
  \bigl(
    X^{(1)},X^{(2)};
    \sh P_\bullet^{(1)},\sh P_\bullet^{(2)}
  \bigr).
  \tag{6.7.2.1}\label{III.6.7.2.1.eq}
\]
\end{env}
With the notation of \sref{III.6.6.2}, the
$L_{\bullet\bullet}^{(i)}$ consist of $A_h$-modules, and hence,
up to canonical isomorphism,
\[
  L_{\bullet\bullet}^{(1)}
  \otimes_{A_h}L_{\bullet\bullet}^{(2)}
  =
  L_{\bullet\bullet}^{(1)}
  \otimes_A L_{\bullet\bullet}^{(2)}.
\]
Since a projective Cartan--Eilenberg resolution
$M_{\bullet\bullet\bullet}^{(i)}$ of $L_{\bullet\bullet}^{(i)}$ may be
chosen to consist of $A_h$-modules, the desired canonical
isomorphism follows immediately.
\end{env}

\begin{theorem}[6.7.3]
\label{III.6.7.3}
Let $S$ be a prescheme, let
$f_i:X^{(i)}\to Y^{(i)}$ be a separated and quasi-compact
$S$-morphism of $S$-preschemes, and let
$\sh P_\bullet^{(i)}$ be a bounded-below complex of
quasi-coherent $\sh O_{X^{(i)}}$-modules $(i=1,2)$.  Write
$u_i:Y^{(i)}\to S$ for the structural morphisms, and put
\[
  Y=Y^{(1)}\times_SY^{(2)}.
\]
There exists a bi-$\partial$-functor
\[
  \mathfrak{Tor}_\bullet^S
  \bigl(
    f_1,f_2;
    \sh P_\bullet^{(1)},\sh P_\bullet^{(2)}
  \bigr)
\]
with values in the category of quasi-coherent
$\sh O_Y$-modules, such that if $S_\alpha$ is an affine open
subset of $S$, $V^{(i)}$ is an affine open subset of
$Y_\alpha^{(i)}=u_i^{-1}(S_\alpha)$ $(i=1,2)$, and
$V=V^{(1)}\times_{S_\alpha}V^{(2)}$, then
\[
\begin{split}
  &\mathfrak{Tor}_\bullet^S
  \bigl(
    f_1,f_2;
    \sh P_\bullet^{(1)},\sh P_\bullet^{(2)}
  \bigr)\bigm|_V\\
  &\quad=
  \left(
    \mathbf{Tor}_\bullet^{S_\alpha}
    \left(
      f_1^{-1}(V^{(1)}),f_2^{-1}(V^{(2)});
      \sh P_\bullet^{(1)}|f_1^{-1}(V^{(1)}),
      \sh P_\bullet^{(2)}|f_2^{-1}(V^{(2)})
    \right)
  \right)\supertilde .
\end{split}
\]

\oldpage[III]{159}

This bifunctor is the common abutment of six biregular spectral
functors
\[
  {}^{(t)}\sh E
  \bigl(
    f_1,f_2;
    \sh P_\bullet^{(1)},\sh P_\bullet^{(2)}
  \bigr)
  \qquad(t=a,b,a',b',c,d),
\]
whose $E_2$-terms are
\[
\begin{aligned}
  {}^{(a)}\sh E_{pq}^2
  &=
  \bigoplus_{q_1+q_2=q}
  \shTor_p^S
  \bigl(
    \sh H^{-q_1}(f_1,\sh P_\bullet^{(1)}),
    \sh H^{-q_2}(f_2,\sh P_\bullet^{(2)})
  \bigr),\\
  {}^{(b)}\sh E_{pq}^2
  &=
  \sh H^{-p}
  \bigl(
    f_1\times_Sf_2,
    \mathfrak{Tor}_q^S
    (\sh P_\bullet^{(1)},\sh P_\bullet^{(2)})
  \bigr),\\
  {}^{(a')}\sh E_{pq}^2
  &=
  \bigoplus_{q_1+q_2=q}
  \mathfrak{Tor}_p^S
  \bigl(
    \sh H^{-q_1}(f_1,\sh P_\bullet^{(1)}),
    \sh H^{-q_2}(f_2,\sh P_\bullet^{(2)})
  \bigr),\\
  {}^{(b')}\sh E_{pq}^2
  &=
  \sh H^{-p}
  \bigl(
    f_1\times_Sf_2,
    \shTor_q^S
    (\sh P_\bullet^{(1)},\sh P_\bullet^{(2)})
  \bigr),\\
  {}^{(c)}\sh E_{pq}^2
  &=
  \bigoplus_{q_1+q_2=q}
  \mathfrak{Tor}_p^S
  \bigl(
    f_1,f_2;
    \sh H_{q_1}(\sh P_\bullet^{(1)}),
    \sh H_{q_2}(\sh P_\bullet^{(2)})
  \bigr),\\
  {}^{(d)}\sh E_{pq}^2
  &=
  \sh H_p
  \left(
    \mathfrak{Tor}_q^S
    \bigl(
      f_1,f_2;
      \sh P_\bullet^{(1)},\sh P_\bullet^{(2)}
    \bigr)
  \right).
\end{aligned}
\]
The spectral sequences \textnormal{(a)} and \textnormal{(b)} are
called the \emph{K\"unneth spectral sequences}.  Notice that
\textnormal{(a)} and \textnormal{(a$'$)}, respectively
\textnormal{(b)} and \textnormal{(b$'$)}, are identical when
$\sh P_\bullet^{(1)}$ and $\sh P_\bullet^{(2)}$ are concentrated
in degree zero.  In that case, \textnormal{(c)} and
\textnormal{(d)} are degenerate and are therefore of no
interest.
\end{theorem}

\begin{remark}[6.7.4]
\label{III.6.7.4}
The global hypertors just defined include as special cases both
the hypercohomology modules defined in \sref{III.6.2.1} and the
local hypertors defined in \sref{III.6.5.3}.  Indeed, for every
quasi-compact and separated morphism $f:X\to Y$ and every
bounded-below complex $\sh P_\bullet$ of quasi-coherent
$\sh O_X$-modules, there is a canonical isomorphism of
$\partial$-functors in $\sh P_\bullet$:
\begin{env}[6.7.4.1]
\label{III.6.7.4.1}
\[
  \mathfrak{Tor}_n^Y
  \bigl(
    f,1_Y;
    \sh P_\bullet,\sh O_Y
  \bigr)
  \xrightarrow{\sim}
  \sh H^{-n}(f,\sh P_\bullet)
  \qquad(n\in\mathbf Z).
  \tag{6.7.4.1}\label{III.6.7.4.1.eq}
\]
\end{env}
The gluing methods of \sref{III.6.7.2} immediately reduce the
proof to the case in which $Y$ is affine.  By
\sref{III.6.2.2}, both sides of
\eqref{III.6.7.4.1.eq} may then be computed using the same
finite covering $\mathfrak U$ of $X$ by affine open subsets and,
for the first side, the covering of $Y$ consisting of $Y$
itself.  With the notation of \sref{III.6.6.2}, the bicomplex
$L_{\bullet\bullet}^{(2)}$ is then concentrated in bidegree
$(0,0)$, where it is equal to $A$, and the conclusion follows
from \sref[0]{0.11.7.5}.

For a generalization, see \sref{III.6.7.7}.  Notice, however,
that if $\sh O_Y$ on the left side of
\eqref{III.6.7.4.1.eq} is replaced by an arbitrary
quasi-coherent $\sh O_Y$-module $\sh F$, there is in general no
longer an isomorphism with
\[
  \sh H^{-n}
  \bigl(
    f,\sh P_\bullet\otimes_{\sh O_Y}\sh F
  \bigr),
\]
although the bicomplex
$L_{\bullet\bullet}^{(1)}\otimes_A
 L_{\bullet\bullet}^{(2)}$ occurring in the preceding
calculation still identifies with
\[
  C^\bullet
  \bigl(
    \mathfrak U,
    \sh P_\bullet\otimes_Y\sh F
  \bigr).
\]

On the other hand, there is a canonical isomorphism of
bi-$\partial$-functors
\begin{env}[6.7.4.2]
\label{III.6.7.4.2}
\[
  \mathfrak{Tor}_n^S
  \bigl(
    1_{X^{(1)}},1_{X^{(2)}};
    \sh P_\bullet^{(1)},\sh P_\bullet^{(2)}
  \bigr)
  \xrightarrow{\sim}
  \shTor_n^S
  \bigl(
    \sh P_\bullet^{(1)},\sh P_\bullet^{(2)}
  \bigr).
  \tag{6.7.4.2}\label{III.6.7.4.2.eq}
\]
\end{env}
Indeed, \sref{III.6.7.1} and \sref{III.6.7.2} reduce us again to
the case in which $S$ and the $X^{(i)}$ are affine.  In
computing the first side of \eqref{III.6.7.4.2.eq}, we may then
take $\mathfrak U^{(i)}$ to consist of the single member
$X^{(i)}$.

\oldpage[III]{160}

With the notation of \sref{III.6.6.2},
$L_{\bullet\bullet}^{(i)}$ then reduces to
\[
  \Gamma(X^{(i)},\sh P_\bullet^{(i)}),
\]
regarded as a bicomplex whose components of first degree
different from zero vanish.  The equality of the two sides of
\eqref{III.6.7.4.2.eq} follows from \sref{III.6.4.1.1} and
\sref{III.6.3.1}.
\end{remark}

\begin{proposition}[6.7.5]
\label{III.6.7.5}
Let
\[
  u:\sh P_\bullet^{(1)}
  \longrightarrow
  \sh Q_\bullet^{(1)}
\]
be a homomorphism of bounded-below complexes of quasi-coherent
$\sh O_{X^{(1)}}$-modules such that the induced homomorphism
\[
  \sh H_\bullet(u):
  \sh H_\bullet(\sh P_\bullet^{(1)})
  \longrightarrow
  \sh H_\bullet(\sh Q_\bullet^{(1)})
\]
is an isomorphism.  Then the homomorphisms
\[
  {}^{(t)}\sh E
  \bigl(
    f_1,f_2;
    \sh P_\bullet^{(1)},\sh P_\bullet^{(2)}
  \bigr)
  \longrightarrow
  {}^{(t)}\sh E
  \bigl(
    f_1,f_2;
    \sh Q_\bullet^{(1)},\sh P_\bullet^{(2)}
  \bigr)
\]
induced by $u$ are isomorphisms for $t=a,b,c$.
\end{proposition}

\begin{proof}
For the spectral sequence \textnormal{(c)}, the assertion follows
from biregularity and from the fact that, by hypothesis, the
homomorphism under consideration is an isomorphism on
$E_2$-terms (\sref[0]{0.11.1.5}).  This already shows that
\[
\begin{split}
  \mathfrak{Tor}_n^S
  \bigl(
    f_1,f_2;
    \sh P_\bullet^{(1)},\sh P_\bullet^{(2)}
  \bigr)
  \longrightarrow
  \mathfrak{Tor}_n^S
  \bigl(
    f_1,f_2;
    \sh Q_\bullet^{(1)},\sh P_\bullet^{(2)}
  \bigr)
\end{split}
\]
is an isomorphism.  Applying \eqref{III.6.7.4.1.eq} and
\eqref{III.6.7.4.2.eq}, we first see that the homomorphisms
\[
  \sh H^{-n}(f_1,\sh P_\bullet^{(1)})
  \longrightarrow
  \sh H^{-n}(f_1,\sh Q_\bullet^{(1)})
\]
and
\[
  \shTor_n^S
  \bigl(
    \sh P_\bullet^{(1)},\sh P_\bullet^{(2)}
  \bigr)
  \longrightarrow
  \shTor_n^S
  \bigl(
    \sh Q_\bullet^{(1)},\sh P_\bullet^{(2)}
  \bigr)
\]
induced by $u$ are isomorphisms.  The assertion for
\textnormal{(a)} and \textnormal{(b)} now follows from the
biregularity of these spectral sequences
(\sref{III.6.7.3}) and the bijectivity of the induced
homomorphisms on $E_2$-terms (\sref[0]{0.11.1.5}).

Notice also that if
$u:\sh P_\bullet^{(1)}\to\sh Q_\bullet^{(1)}$ is a homotopy
equivalence, then it induces canonical isomorphisms
\[
  {}^{(t)}\sh E
  \bigl(
    f_1,f_2;
    \sh P_\bullet^{(1)},\sh P_\bullet^{(2)}
  \bigr)
  \xrightarrow{\sim}
  {}^{(t)}\sh E
  \bigl(
    f_1,f_2;
    \sh Q_\bullet^{(1)},\sh P_\bullet^{(2)}
  \bigr)
\]
for all six spectral sequences.  Indeed, when $S$ and the
$Y^{(i)}$ are affine, $u$ induces a homotopy equivalence of
bicomplexes
\[
  C^\bullet
  \bigl(
    \mathfrak U^{(1)},\sh P_\bullet^{(1)}
  \bigr)
  \longrightarrow
  C^\bullet
  \bigl(
    \mathfrak U^{(1)},\sh Q_\bullet^{(1)}
  \bigr),
\]
and the assertion follows from the general theory of
hyperhomology (\sref[0]{0.11.3.2}).  The passage to the general
case is made by gluing, using the fact that a homotopy
equivalence of complexes induces one between projective
Cartan--Eilenberg resolutions of those complexes
(\textnormal{M}, XVII, 1.2).
\end{proof}

\begin{proposition}[6.7.6]
\label{III.6.7.6}
Suppose that either $\sh P_\bullet^{(1)}$ or
$\sh P_\bullet^{(2)}$ consists of $S$-flat modules, both
complexes being bounded below.  There is then a canonical
isomorphism of bi-$\partial$-functors
\begin{env}[6.7.6.1]
\label{III.6.7.6.1}
\[
  \mathfrak{Tor}_n^S
  \bigl(
    f_1,f_2;
    \sh P_\bullet^{(1)},\sh P_\bullet^{(2)}
  \bigr)
  \xrightarrow{\sim}
  \sh H^{-n}
  \bigl(
    f_1\times_Sf_2,
    \sh P_\bullet^{(1)}
    \otimes_S\sh P_\bullet^{(2)}
  \bigr).
  \tag{6.7.6.1}\label{III.6.7.6.1.eq}
\]
\end{env}
\end{proposition}

\begin{proof}
First suppose that $S$, $Y^{(1)}$, and $Y^{(2)}$ are affine, so
that we are in the situation of \sref{III.6.6.2}, whose notation
we retain.  Suppose, for example, that
$\sh P_\bullet^{(1)}$ consists of $S$-flat modules.  Using
\sref{III.6.6.6}, the hypertor is the homology of
\[
  L_{\bullet\bullet}^{(1)}
  \otimes_A M_{\bullet\bullet\bullet}^{(2)},
\]
where $M_{\bullet\bullet\bullet}^{(2)}$ is a projective
Cartan--Eilenberg resolution of $L_{\bullet\bullet}^{(2)}$ in
the sense of \sref[0]{0.11.7.1}.  On the other hand, the modules
$L_{ij}^{(1)}$ are flat over $A$ by the hypothesis
(\sref[I]{I.1.4.15.1}).  It follows from
\sref[0]{0.11.7.5} that there is a canonical isomorphism
\begin{env}[6.7.6.2]
\label{III.6.7.6.2}
\[
  \mathbf{Tor}_\bullet^S
  \bigl(
    \mathfrak U^{(1)},\mathfrak U^{(2)};
    \sh P_\bullet^{(1)},\sh P_\bullet^{(2)}
  \bigr)
  \xrightarrow{\sim}
  H_\bullet
  \bigl(
    L_{\bullet\bullet}^{(1)}
    \otimes_A L_{\bullet\bullet}^{(2)}
  \bigr).
  \tag{6.7.6.2}\label{III.6.7.6.2.eq}
\]
\end{env}

There is also a natural homomorphism of bicomplexes
\[
  L_{\bullet\bullet}^{(1)}
  \otimes_A L_{\bullet\bullet}^{(2)}
  \longrightarrow
  C^\bullet(\mathfrak U,\sh Q_\bullet),
\]
where $\mathfrak U$ is the covering of
\[
  Z=X^{(1)}\times_SX^{(2)}
\]
by the affine open subsets
$U_\alpha^{(1)}\times_SU_\beta^{(2)}$, and
\[
  \sh Q_\bullet
  =
  \sh P_\bullet^{(1)}
  \otimes_S\sh P_\bullet^{(2)}
\]
is regarded as the simple complex for the total degree.

\oldpage[III]{161}

The substance of this homomorphism was already given in the
calculation of spectral sequence \textnormal{(b$'$)} in
\sref{III.6.6.4}, with $q=0$.  Retaining the notation used
there, one need only observe that there is, first, a natural
homomorphism from $N^{\bullet k}$ to
\[
  C^\bullet
  \bigl(
    \Phi^{(1)},\Phi^{(2)};\mathcal S_k
  \bigr)
\]
(\sref[0]{0.11.8.5}); second, a natural homomorphism from the
latter complex to
\[
  P^\bullet
  \bigl(
    \Phi^{(1)},\Phi^{(2)};\mathcal S_k
  \bigr)
\]
(\sref[0]{0.11.8.6}); and finally, a natural homomorphism from
this last cochain complex to its subcomplex of alternating
cochains (\sref[0]{0.11.8.7}).  The resulting homomorphism of
bicomplexes induces an isomorphism in homology, as was seen in
\sref{III.6.6.4}.  Composing with
\eqref{III.6.7.6.2.eq} therefore gives an isomorphism
\begin{env}[6.7.6.3]
\label{III.6.7.6.3}
\[
  \mathbf{Tor}_n^S
  \bigl(
    \mathfrak U^{(1)},\mathfrak U^{(2)};
    \sh P_\bullet^{(1)},\sh P_\bullet^{(2)}
  \bigr)
  \xrightarrow{\sim}
  \HH^{-n}
  \bigl(
    \mathfrak U,
    \sh P_\bullet^{(1)}
    \otimes_S\sh P_\bullet^{(2)}
  \bigr).
  \tag{6.7.6.3}\label{III.6.7.6.3.eq}
\]
\end{env}

It remains to prove that the isomorphism just defined is
independent of the chosen open coverings.  The right side of
\eqref{III.6.7.6.3.eq} is canonically isomorphic, by
\sref{III.6.2.2}, to
\[
  \HH^{-n}
  \bigl(
    X^{(1)}\times_SX^{(2)},
    \sh P_\bullet^{(1)}
    \otimes_S\sh P_\bullet^{(2)}
  \bigr).
\]
For the left side, use \sref{III.6.6.7} and note, with the
notation used there, that the following diagram commutes up to
homotopy:
\[
\xymatrix@C=18pt@R=28pt{
  C^\bullet(\mathfrak U^{(1)},\sh P_\bullet^{(1)})
  \otimes_A
  C^\bullet(\mathfrak U^{(2)},\sh P_\bullet^{(2)})
  \ar[r]\ar[d]
  &
  C^\bullet(\mathfrak U,\sh Q_\bullet)
  \ar[d]
  \\
  C^\bullet(\mathfrak V^{(1)},\sh P_\bullet^{(1)})
  \otimes_A
  C^\bullet(\mathfrak V^{(2)},\sh P_\bullet^{(2)})
  \ar[r]
  &
  C^\bullet(\mathfrak V,\sh Q_\bullet).
}
\]
The horizontal arrows are those defined above.  Finally, one
passes to the general case by gluing, exactly as in
\sref{III.6.7.1} and \sref{III.6.7.2}; the details are left to
the reader.
\end{proof}

\begin{proposition}[6.7.7]
\label{III.6.7.7}
Suppose that $\sh P_\bullet^{(1)}$ and
$\sh P_\bullet^{(2)}$ are bounded below, and that either all the
modules
\[
  \sh H^{-n}(f_1,\sh P_\bullet^{(1)})
\]
or all the modules
\[
  \sh H^{-n}(f_2,\sh P_\bullet^{(2)})
\]
are $S$-flat.  There is then a canonical isomorphism of
bi-$\partial$-functors, for $n\in\mathbf Z$,
\begin{env}[6.7.7.1]
\label{III.6.7.7.1}
\[
  \mathfrak{Tor}_n^S
  \bigl(
    f_1,f_2;
    \sh P_\bullet^{(1)},\sh P_\bullet^{(2)}
  \bigr)
  \xrightarrow{\sim}
  \bigoplus_{q_1+q_2=n}
  \sh H^{-q_1}(f_1,\sh P_\bullet^{(1)})
  \otimes_S
  \sh H^{-q_2}(f_2,\sh P_\bullet^{(2)}).
  \tag{6.7.7.1}\label{III.6.7.7.1.eq}
\]
\end{env}
\end{proposition}

\begin{proof}
By \sref{III.6.5.8}, spectral sequence \textnormal{(a)} of
\sref{III.6.7.3} degenerates.  The assertion therefore follows
immediately from \sref[0]{0.11.1.6}, since that spectral sequence
is biregular (\sref{III.6.7.3}).
\end{proof}

\begin{theorem}[6.7.8]
\label{III.6.7.8}
Suppose that the following conditions hold:
\begin{enumerate}
  \item[{$1^\circ$}]
    The complexes $\sh P_\bullet^{(1)}$ and
    $\sh P_\bullet^{(2)}$ are bounded below.
  \item[{$2^\circ$}]
    Either $\sh P_\bullet^{(1)}$ or
    $\sh P_\bullet^{(2)}$ consists of $S$-flat modules.

    \oldpage[III]{162}

  \item[{$3^\circ$}]
    Either all the modules
    $\sh H^{-n}(f_1,\sh P_\bullet^{(1)})$ or all the modules
    $\sh H^{-n}(f_2,\sh P_\bullet^{(2)})$ are $S$-flat.
\end{enumerate}
There is then a canonical isomorphism of bi-$\partial$-functors,
for $n\in\mathbf Z$,
\begin{env}[6.7.8.1]
\label{III.6.7.8.1}
\[
  \sh H^n
  \bigl(
    f_1\times_Sf_2,
    \sh P_\bullet^{(1)}\otimes_S\sh P_\bullet^{(2)}
  \bigr)
  \xrightarrow{\sim}
  \bigoplus_{n_1+n_2=n}
  \sh H^{n_1}(f_1,\sh P_\bullet^{(1)})
  \otimes_S
  \sh H^{n_2}(f_2,\sh P_\bullet^{(2)}).
  \tag{6.7.8.1}\label{III.6.7.8.1.eq}
\]
\end{env}
This is the \emph{K\"unneth formula}.
\end{theorem}

\begin{proof}
The formula follows from \sref{III.6.7.6} and
\sref{III.6.7.7}.

When $S$, $Y^{(1)}$, and $Y^{(2)}$ are affine, the inverse of
\eqref{III.6.7.8.1.eq} is obtained, with the notation of
\sref{III.6.7.6}, from the homomorphism of bicomplexes
\[
  C^\bullet(\mathfrak U^{(1)},\sh P_\bullet^{(1)})
  \otimes_A
  C^\bullet(\mathfrak U^{(2)},\sh P_\bullet^{(2)})
  \longrightarrow
  C^\bullet(\mathfrak U,\sh Q_\bullet)
\]
by the procedure defined in \textnormal{(G, I, 2.7)}, as follows
from \textnormal{(G, I, 5.5)}.
\end{proof}

\begin{proposition}[6.7.9]
\label{III.6.7.9}
Suppose that the following three conditions hold:
\begin{enumerate}
  \item[{$1^\circ$}]
    The schemes $S$, $Y^{(1)}$, and $Y^{(2)}$ are locally
    noetherian; $f_1$ and $f_2$ are proper; and either
    $Y^{(1)}$ or $Y^{(2)}$ is of finite type over $S$.
  \item[{$2^\circ$}]
    The complexes $\sh P_\bullet^{(1)}$ and
    $\sh P_\bullet^{(2)}$ are bounded below.
  \item[{$3^\circ$}]
    For every $n\in\mathbf Z$, the module
    $\sh H_n(\sh P_\bullet^{(i)})$ is coherent, for $i=1,2$.
\end{enumerate}
Then
\[
  \mathfrak{Tor}_n^S
  \bigl(
    f_1,f_2;
    \sh P_\bullet^{(1)},\sh P_\bullet^{(2)}
  \bigr)
\]
is a coherent $\sh O_Y$-module, where
$Y=Y^{(1)}\times_SY^{(2)}$.
\end{proposition}

\begin{proof}
It follows from \sref{III.6.5.13} that the local hypertor modules
\[
  \shTor_n^S
  \bigl(
    \sh P_\bullet^{(1)},\sh P_\bullet^{(2)}
  \bigr)
\]
are coherent $\sh O_X$-modules, where
$X=X^{(1)}\times_SX^{(2)}$.  Indeed, each $X^{(i)}$ is locally
noetherian, and one of them is, by hypothesis, of finite type over
    $S$ (\sref[I]{I.6.3.4} and \sref[I]{I.6.3.8}); hence $X$ is
locally noetherian.

Since $Y$ is locally noetherian and
$f_1\times_Sf_2$ is proper (\sref[II]{II.5.4.2}), it follows from
\sref{III.6.2.5} that the terms ${}^{(b)}\sh E_{pq}^2$ of
\sref{III.6.7.3} are coherent $\sh O_Y$-modules.  All the spectral
sequences of \sref{III.6.7.3} are biregular by hypothesis
$2^\circ$, so the conclusion follows from
\sref[0]{0.11.1.8}.
\end{proof}

\begin{env}[6.7.10]
\label{III.6.7.10}
Let $Y^{\prime(i)}$ be two $S$-preschemes and let
\[
  v_i:Y^{\prime(i)}\longrightarrow Y^{(i)}
  \qquad (i=1,2)
\]
be two $S$-morphisms.  Write
\[
  v=v_1\times_Sv_2:Y'\longrightarrow Y,
  \qquad
  Y'=Y^{\prime(1)}\times_SY^{\prime(2)}.
\]
For $i=1,2$, let $X^{\prime(i)}$ be an $S$-prescheme, and let
\[
  u_i:X^{\prime(i)}\longrightarrow X^{(i)},
  \qquad
  f_i':X^{\prime(i)}\longrightarrow Y^{\prime(i)}
\]
be $S$-morphisms such that the diagrams
\begin{env}[6.7.10.1]
\label{III.6.7.10.1}
\[
\xymatrix@C=36pt@R=28pt{
  X^{\prime(i)} \ar[r]^{u_i}\ar[d]_{f_i'}
  & X^{(i)}\ar[d]^{f_i}
  \\
  Y^{\prime(i)} \ar[r]_{v_i}
  & Y^{(i)}
}
\tag{6.7.10.1}\label{III.6.7.10.1.eq}
\]
\end{env}
commute, the morphisms $f_i'$ being separated and quasi-compact.
There are then canonical $\sh O_{Y'}$-linear homomorphisms of
spectral functors
\begin{env}[6.7.10.2]
\label{III.6.7.10.2}
\[
  v^*
  {}^{(t)}\sh E
  \bigl(
    f_1,f_2;
    \sh P_\bullet^{(1)},\sh P_\bullet^{(2)}
  \bigr)
  \longrightarrow
  {}^{(t)}\sh E
  \bigl(
    f_1',f_2';
    u_1^*(\sh P_\bullet^{(1)}),
    u_2^*(\sh P_\bullet^{(2)})
  \bigr)
  \tag{6.7.10.2}\label{III.6.7.10.2.eq}
\]
\end{env}
for $t=a,a',b,b',c,d$.

To define them, first suppose that
\[
  S=\Spec(A),\qquad
  Y^{(i)}=\Spec(B_i),\qquad
  Y^{\prime(i)}=\Spec(B_i')
\]
are affine.  The schemes $X^{(i)}$ and $X^{\prime(i)}$ are then
quasi-compact.  To compute
\[
  {}^{(t)}\sh E
  \bigl(
    f_1,f_2;
    \sh P_\bullet^{(1)},\sh P_\bullet^{(2)}
  \bigr),
\]
choose, as in \sref{III.6.6.1}, finite affine open coverings
$\mathfrak U^{(i)}$ of $X^{(i)}$.  To compute
\[
  {}^{(t)}\sh E
  \bigl(
    f_1',f_2';
    u_1^*(\sh P_\bullet^{(1)}),
    u_2^*(\sh P_\bullet^{(2)})
  \bigr),
\]
choose finite affine open coverings $\mathfrak U^{\prime(i)}$ of
$X^{\prime(i)}$, respectively finer than the inverse-image
coverings $u_i^{-1}(\mathfrak U^{(i)})$.

\oldpage[III]{163}

There is a canonical homomorphism from the bicomplex
\[
  C^\bullet
  \bigl(
    \mathfrak U^{(i)},\sh P_\bullet^{(i)}
  \bigr)
  =
  L_{\bullet\bullet}^{(i)}
\]
to
\[
  C^\bullet
  \bigl(
    u_i^{-1}(\mathfrak U^{(i)}),
    u_i^*(\sh P_\bullet^{(i)})
  \bigr)
\]
(\sref[0]{0.4.4.3.2}).  A simplicial map
\textnormal{(G, II, 5.7)} from $\mathfrak U^{\prime(i)}$ to
$u_i^{-1}(\mathfrak U^{(i)})$ defines a homomorphism of
bicomplexes
\[
  C^\bullet
  \bigl(
    u_i^{-1}(\mathfrak U^{(i)}),
    u_i^*(\sh P_\bullet^{(i)})
  \bigr)
  \longrightarrow
  C^\bullet
  \bigl(
    \mathfrak U^{\prime(i)},
    u_i^*(\sh P_\bullet^{(i)})
  \bigr).
\]
By composition one obtains a homomorphism
\[
  L_{\bullet\bullet}^{(i)}
  \longrightarrow
  L_{\bullet\bullet}^{\prime(i)}
  =
  C^\bullet
  \bigl(
    \mathfrak U^{\prime(i)},
    u_i^*(\sh P_\bullet^{(i)})
  \bigr).
\]
Changing the simplicial map replaces this homomorphism by a
homotopic one \textnormal{(G, II, 5.7.1)}.  Thus one obtains a
well-defined homomorphism of spectral functors
\begin{env}[6.7.10.3]
\label{III.6.7.10.3}
\[
  {}^{(t)}\sh E
  \bigl(
    \mathfrak U^{(1)},\mathfrak U^{(2)};
    \sh P_\bullet^{(1)},\sh P_\bullet^{(2)}
  \bigr)
  \longrightarrow
  {}^{(t)}\sh E
  \bigl(
    \mathfrak U^{\prime(1)},\mathfrak U^{\prime(2)};
    u_1^*(\sh P_\bullet^{(1)}),
    u_2^*(\sh P_\bullet^{(2)})
  \bigr).
  \tag{6.7.10.3}\label{III.6.7.10.3.eq}
\]
\end{env}

Suppose that $\mathfrak V^{(i)}$ is a finite affine open covering
of $X^{(i)}$ finer than $\mathfrak U^{(i)}$, and that
$\mathfrak V^{\prime(i)}$ is a finite affine open covering of
$X^{\prime(i)}$ finer than both $\mathfrak U^{\prime(i)}$ and
$u_i^{-1}(\mathfrak V^{(i)})$.  One verifies immediately that the
diagram
\[
\xymatrix@C=34pt@R=30pt{
  C^\bullet(\mathfrak U^{(i)},\sh P_\bullet^{(i)})
  \ar[r]\ar[d]
  &
  C^\bullet(\mathfrak V^{(i)},\sh P_\bullet^{(i)})
  \ar[d]
  \\
  C^\bullet
  \bigl(
    \mathfrak U^{\prime(i)},u_i^*(\sh P_\bullet^{(i)})
  \bigr)
  \ar[r]
  &
  C^\bullet
  \bigl(
    \mathfrak V^{\prime(i)},u_i^*(\sh P_\bullet^{(i)})
  \bigr)
}
\]
commutes.  Consequently, \eqref{III.6.7.10.3.eq} is essentially
independent of the coverings used.  We have therefore defined a
homomorphism of $A$-modules
\begin{env}[6.7.10.4]
\label{III.6.7.10.4}
\[
  {}^{(t)}E
  \bigl(
    X^{(1)},X^{(2)};
    \sh P_\bullet^{(1)},\sh P_\bullet^{(2)}
  \bigr)
  \longrightarrow
  {}^{(t)}E
  \bigl(
    X^{\prime(1)},X^{\prime(2)};
    u_1^*(\sh P_\bullet^{(1)}),
    u_2^*(\sh P_\bullet^{(2)})
  \bigr).
  \tag{6.7.10.4}\label{III.6.7.10.4.eq}
\]
\end{env}
By the definition of the $u_i^*(\sh P_\bullet^{(i)})$ and the
commutativity of \eqref{III.6.7.10.1.eq}, this is also a
homomorphism of $(B_1\otimes_AB_2)$-modules.  Since the right-hand
side of \eqref{III.6.7.10.4.eq} consists of
$(B_1'\otimes_AB_2')$-modules, it canonically induces a
homomorphism of $(B_1'\otimes_AB_2')$-modules
\begin{env}[6.7.10.5]
\label{III.6.7.10.5}
\[
\begin{split}
  {}^{(t)}E
  \bigl(
    X^{(1)},X^{(2)};
    \sh P_\bullet^{(1)},\sh P_\bullet^{(2)}
  \bigr)
  \otimes_{B_1\otimes_AB_2}
  (B_1'\otimes_AB_2')
  \quad\longrightarrow\quad{}&
  {}^{(t)}E
  \bigl(
    X^{\prime(1)},X^{\prime(2)};
\\[-2pt]
    &u_1^*(\sh P_\bullet^{(1)}),
    u_2^*(\sh P_\bullet^{(2)})
  \bigr).
\end{split}
\tag{6.7.10.5}\label{III.6.7.10.5.eq}
\]
\end{env}
In view of \sref[I]{I.1.6.5}, this is precisely
\eqref{III.6.7.10.2.eq} in the affine case.

It remains to pass to the general case by the gluing procedures of
\sref{III.6.7.1} and \sref{III.6.7.2}.  The second passage is
immediate.  For the first, take, as in \sref{III.6.7.1}, elements
$g_i\in B_i$ and their images $g_i'\in B_i'$, and set
\[
  g=g_1\otimes g_2\in B=B_1\otimes_AB_2,
  \qquad
  g'=g_1'\otimes g_2'\in B'=B_1'\otimes_AB_2'.
\]

\oldpage[III]{164}

Everything then reduces to the canonical isomorphism
\[
  (M\otimes_BB')_{g'}
  \xrightarrow{\sim}
  M_g\otimes_{B_g}B_{g'}'
\]
(\sref[0]{0.1.5.4}).  The details are left to the reader.
\end{env}

\begin{env}[6.7.11]
\label{III.6.7.11}
The theory of global hypertor developed above for two
$S$-morphisms $X^{(i)}\to Y^{(i)}$ and two bounded-below
complexes $\sh P_\bullet^{(i)}$ of quasi-coherent modules extends
immediately to the following general situation.  Let $S$ be a
prescheme, let $(Y^{(i)})_{i\in I}$ be a finite family of
$S$-preschemes, let
\[
  f_i:X^{(i)}\longrightarrow Y^{(i)}
  \qquad (i\in I)
\]
be a finite family of separated and quasi-compact
$S$-morphisms, and for each $i$ let $\sh P_\bullet^{(i)}$ be a
bounded-below complex of quasi-coherent
$\sh O_{X^{(i)}}$-modules.  If $Y$ is the product of the
$S$-preschemes $Y^{(i)}$, then for every $n\in\mathbf Z$ one
defines a quasi-coherent $\sh O_Y$-module
\[
  \mathfrak{Tor}_n^S
  \bigl(
    (f_i)_{i\in I};
    (\sh P_\bullet^{(i)})_{i\in I}
  \bigr).
\]
These modules form a covariant $\partial$-functor in each complex
$\sh P_\bullet^{(i)}$; moreover, this functor is the common
abutment of six spectral functors
\[
  {}^{(t)}\sh E
  \bigl(
    (f_i)_{i\in I};
    (\sh P_\bullet^{(i)})_{i\in I}
  \bigr).
\]
We leave to the reader the repetition, in this general case, of
the definitions and arguments given above for $I=\{1,2\}$.

When $I$ has a single element, this construction recovers the
hypercohomology $\sh H^\bullet(f,\sh P_\bullet)$ defined in
\sref{III.6.2.7}, as already noted in \sref{III.6.7.4}.  When
$I$ is the interval $1\leq i\leq m$ in $\mathbf N$, we write
\phantomsection\label{III.6.7.11.2}
\[
  \mathfrak{Tor}_n^S
  \bigl(
    f_1,\ldots,f_m;
    \sh P_\bullet^{(1)},\ldots,\sh P_\bullet^{(m)}
  \bigr)
\]
for
\[
  \mathfrak{Tor}_n^S
  \bigl(
    (f_i)_{i\in I};
    (\sh P_\bullet^{(i)})_{i\in I}
  \bigr).
\]
\end{env}

\begin{proposition}[6.7.12]
\label{III.6.7.12}
With the notation of \sref{III.6.7.11}, let $J$ be a subset of
$I$ such that, for $i\in I-J$,
\[
  X^{(i)}=Y^{(i)}=S,
\]
the morphism $f_i$ is the identity, and
$\sh P_\bullet^{(i)}$ is the complex concentrated in degree $0$
with value $\sh O_S$.  There is then a canonical isomorphism of
$\partial$-functors
\begin{env}[6.7.12.1]
\label{III.6.7.12.1}
\[
  \mathfrak{Tor}_n^S
  \bigl(
    (f_i)_{i\in I};
    (\sh P_\bullet^{(i)})_{i\in I}
  \bigr)
  \xrightarrow{\sim}
  \mathfrak{Tor}_n^S
  \bigl(
    (f_i)_{i\in J};
    (\sh P_\bullet^{(i)})_{i\in J}
  \bigr).
  \tag{6.7.12.1}\label{III.6.7.12.1.eq}
\]
\end{env}
\end{proposition}

\begin{proof}
It is enough to define this isomorphism when $S$ and the
$Y^{(i)}$ are affine, since the gluing is carried out as usual.
For $i\in I-J$, take $\mathfrak U^{(i)}$ to be the covering
consisting of the single set $S$.  Then
\[
  L_{\bullet\bullet}^{(i)}
  =
  C^\bullet
  \bigl(
    \mathfrak U^{(i)},\sh P_\bullet^{(i)}
  \bigr)
\]
has only one nonzero term, in bidegree $(0,0)$, equal to
\[
  \Gamma(S,\sh O_S)=A(S).
\]
The isomorphism \eqref{III.6.7.12.1.eq} is therefore immediate.
\end{proof}

\begin{remark}[6.7.13]
\label{III.6.7.13}
With the notation of \sref{III.6.7.3}, consider the canonical
$S$-isomorphism
\[
  Y^{(1)}\times_SY^{(2)}
  \longrightarrow
  Y^{(2)}\times_SY^{(1)}
\]
of \sref[I]{I.3.3.5}.  Under this isomorphism, the image of
\[
  \mathfrak{Tor}_n^S
  \bigl(
    f_1,f_2;
    \sh P_\bullet^{(1)},\sh P_\bullet^{(2)}
  \bigr)
\]
is
\[
  \mathfrak{Tor}_n^S
  \bigl(
    f_2,f_1;
    \sh P_\bullet^{(2)},\sh P_\bullet^{(1)}
  \bigr).
\]
This is a local assertion, so it reduces to the case considered in
\sref{III.6.6.2}.  If
$M_{\bullet\bullet\bullet}^{(i)}$ is a projective
Cartan--Eilenberg resolution of $L_{\bullet\bullet}^{(i)}$ for
$i=1,2$, the isomorphism in question interchanges
\[
  M_{\bullet\bullet\bullet}^{(1)}
  \otimes_A
  M_{\bullet\bullet\bullet}^{(2)}
  \qquad\hbox{and}\qquad
  M_{\bullet\bullet\bullet}^{(2)}
  \otimes_A
  M_{\bullet\bullet\bullet}^{(1)}.
\]
The assertion follows by taking the homology of the simple
complexes associated with these tricomplexes.
\end{remark}

\subsection{The associativity spectral sequences of global hypertor}
\label{subsection:III.6.8}
\label{III.6.8}

\begin{env}[6.8.1]
\label{III.6.8.1}
Retain the hypotheses and notation of \sref{III.6.7.11}; in
particular, the complexes $\sh P_\bullet^{(i)}$ are bounded
below.  Let $(I_j)_{j\in J}$ be a partition of the index set $I$.
We seek an ``associativity'' relation between
\[
  \mathfrak{Tor}_n^S
  \bigl(
    (f_i)_{i\in I};
    (\sh P_\bullet^{(i)})_{i\in I}
  \bigr)
\]
and the ``partial'' hypertors
\[
  \mathcal T_{\bullet j}
  =
  \mathfrak{Tor}_\bullet^S
  \bigl(
    (f_i)_{i\in I_j};
    (\sh P_\bullet^{(i)})_{i\in I_j}
  \bigr).
\]
\oldpage[III]{165}

To simplify the notation, we shall restrict attention to the case
where $I$ is the interval $1\leq i\leq m$ and the partition
consists of the two intervals
\[
  \{1,2,\ldots,r\}
  \qquad\hbox{and}\qquad
  \{r+1,\ldots,m\}.
\]
\end{env}

\begin{proposition}[6.8.2]
\label{III.6.8.2}
There is a canonical biregular spectral functor, called the
\emph{associativity spectral functor}, denoted
\[
  {}^{(e)}\sh E^S
  \bigl(
    f_1,\ldots,f_m;
    \sh P_\bullet^{(1)},\ldots,\sh P_\bullet^{(m)}
  \bigr),
\]
or simply ${}^{(e)}\sh E$ with the same arguments, whose abutment
is
\[
  \mathfrak{Tor}_\bullet^S
  \bigl(
    f_1,\ldots,f_m;
    \sh P_\bullet^{(1)},\ldots,\sh P_\bullet^{(m)}
  \bigr)
\]
and whose $E_2$-term is
\[
\begin{split}
  {}^{(e)}\sh E_{pq}^2
  =
  \bigoplus_{q_1+q_2=q}
  \shTor_p^S
  \Bigl(
    &
    \mathfrak{Tor}_{q_1}^S
    \bigl(
      f_1,\ldots,f_r;
      \sh P_\bullet^{(1)},\ldots,\sh P_\bullet^{(r)}
    \bigr),
    \\
    &
    \mathfrak{Tor}_{q_2}^S
    \bigl(
      f_{r+1},\ldots,f_m;
      \sh P_\bullet^{(r+1)},\ldots,\sh P_\bullet^{(m)}
    \bigr)
  \Bigr).
\end{split}
\]
\end{proposition}

\begin{proof}
Here $Y$ is canonically identified with
$Z^{(1)}\times_SZ^{(2)}$, where
\[
  Z^{(1)}
  =
  Y^{(1)}\times_S\cdots\times_SY^{(r)}
  \quad\hbox{and}\quad
  Z^{(2)}
  =
  Y^{(r+1)}\times_S\cdots\times_SY^{(m)}.
\]
It is enough to treat the case in which $S$ and all the
$Y^{(i)}$ are affine.  The passage from this case to the general
one uses the methods of \sref{III.6.7.1} and
\sref{III.6.7.2}; the details are left to the reader.  Thus it is
enough to prove the following affine statement.
\end{proof}

\begin{corollary}[6.8.3]
\label{III.6.8.3}
Let $A$ be a ring, let $S=\Spec(A)$, and let $X^{(i)}$,
$1\leq i\leq m$, be quasi-compact $S$-schemes.  For each $i$, let
$\sh P_\bullet^{(i)}$ be a bounded-below complex of
quasi-coherent $\sh O_{X^{(i)}}$-modules.  There is a canonical
biregular spectral functor with abutment
\[
  \mathbf{Tor}_\bullet^S
  \bigl(
    X^{(1)},\ldots,X^{(m)};
    \sh P_\bullet^{(1)},\ldots,\sh P_\bullet^{(m)}
  \bigr)
\]
and $E_2$-term
\[
\begin{split}
  {}^{(e)}E_{pq}^2
  =
  \bigoplus_{q_1+q_2=q}
  \Tor_p^A
  \Bigl(
    &
    \mathbf{Tor}_{q_1}^S
    \bigl(
      X^{(1)},\ldots,X^{(r)};
      \sh P_\bullet^{(1)},\ldots,\sh P_\bullet^{(r)}
    \bigr),
    \\
    &
    \mathbf{Tor}_{q_2}^S
    \bigl(
      X^{(r+1)},\ldots,X^{(m)};
      \sh P_\bullet^{(r+1)},\ldots,\sh P_\bullet^{(m)}
    \bigr)
  \Bigr).
\end{split}
\]
\end{corollary}

\begin{proof}
Following the definition of \sref{III.6.6.2}, choose for each $i$
a finite affine open covering $\mathfrak U^{(i)}$ of $X^{(i)}$,
form the bicomplex
\[
  L_{\bullet\bullet}^{(i)}
  =
  C^\bullet
  \bigl(
    \mathfrak U^{(i)},\sh P_\bullet^{(i)}
  \bigr),
\]
and choose a projective Cartan--Eilenberg resolution
$M_{\bullet\bullet\bullet}^{(i)}$ of this bicomplex in the sense
of \sref[0]{0.11.7.1}.  The hypertor in question is the homology
of the tensor product
\[
  M_{\bullet\bullet\bullet}
  =
  \bigotimes_{i=1}^{m}
  M_{\bullet\bullet\bullet}^{(i)}.
\]
Regard $M_{\bullet\bullet\bullet}$ as a simple complex
$N_\bullet$, the tensor product of the two simple complexes
\[
  N_\bullet'
  =
  \bigotimes_{i=1}^{r}
  M_{\bullet\bullet\bullet}^{(i)},
  \qquad
  N_\bullet''
  =
  \bigotimes_{i=r+1}^{m}
  M_{\bullet\bullet\bullet}^{(i)},
\]
where $N_\bullet'$ and $N_\bullet''$ are graded by the sums of the
total degrees of the $M_{\bullet\bullet\bullet}^{(i)}$.

The $A$-modules forming $N_\bullet'$ and $N_\bullet''$ are
projective.  Hence \sref{III.6.5.9} gives
\[
  H_\bullet(M_{\bullet\bullet\bullet})
  =
  \mathbf{Tor}_\bullet^A(N_\bullet',N_\bullet'').
\]
The required spectral sequence is therefore precisely
\sref{III.6.3.2.2} applied to $N_\bullet'$ and $N_\bullet''$,
using the preceding interpretation of their homology modules for
the two partial products
\[
  X^{(1)}\times_S\cdots\times_SX^{(r)}
  \quad\hbox{and}\quad
  X^{(r+1)}\times_S\cdots\times_SX^{(m)}.
\]
Finally, the regularity assertions follow from
\sref{III.6.3.2}.  Indeed, since the
$\sh P_\bullet^{(i)}$ are bounded below, so are the
$M_{\bullet\bullet\bullet}^{(i)}$, and therefore so are
$N_\bullet'$ and $N_\bullet''$.
\end{proof}

\oldpage[III]{166}
\subsection{The base change spectral sequences of global hypertor}
\label{subsection:III.6.9}
\label{III.6.9}

\begin{env}[6.9.1]
\label{III.6.9.1}
Retain the hypotheses and notation of \sref{III.6.7.11}; in
particular, the complexes $\sh P_\bullet^{(i)}$ are bounded below.
Let $g:S'\to S$ be a morphism of preschemes, and set
\[
  Y^{\prime(i)}=Y_{(S')}^{(i)},
  \qquad
  X^{\prime(i)}=X_{(S')}^{(i)},
  \qquad
  \sh P_\bullet^{\prime(i)}
  =
  \sh P_\bullet^{(i)}
  \otimes_{\sh O_S}\sh O_{S'}.
\]
Thus $\sh P_\bullet^{\prime(i)}$ is a complex of quasi-coherent
$\sh O_{X^{\prime(i)}}$-modules.  Let
\[
  f_i'=(f_i)_{(S')}:
  X^{\prime(i)}\longrightarrow Y^{\prime(i)};
\]
this morphism is separated and quasi-compact
(\sref[I]{I.5.5.1} and \sref{III.6.6.4}).

We shall study the relation between the quasi-coherent
$\sh O_{Y'}$-modules
\[
  \mathfrak{Tor}_n^{S'}
  \bigl(
    (f_i')_{i\in I};
    (\sh P_\bullet^{\prime(i)})_{i\in I}
  \bigr)
\]
and
\[
  \mathfrak{Tor}_n^S
  \bigl(
    (f_i)_{i\in I};
    (\sh P_\bullet^{(i)})_{i\in I}
  \bigr)
  \otimes_{\sh O_S}\sh O_{S'},
\]
where
\[
  Y'=Y\times_SS'=Y_{(S')}.
\]
The following particularly simple case reduces to
\sref[I]{I.1.4.15} when $I$ has a single element and
$\sh P_\bullet$ consists of a single module.
\end{env}

\begin{proposition}[6.9.2]
\label{III.6.9.2}
If $g:S'\to S$ is flat, there is a canonical isomorphism of
$\partial$-functors in the complexes $\sh P_\bullet^{(i)}$:
\begin{env}[6.9.2.1]
\label{III.6.9.2.1}
\[
  \mathfrak{Tor}_n^S
  \bigl(
    (f_i)_{i\in I};
    (\sh P_\bullet^{(i)})_{i\in I}
  \bigr)
  \otimes_{\sh O_S}\sh O_{S'}
  \xrightarrow{\sim}
  \mathfrak{Tor}_n^{S'}
  \bigl(
    (f_i')_{i\in I};
    (\sh P_\bullet^{\prime(i)})_{i\in I}
  \bigr).
  \tag{6.9.2.1}\label{III.6.9.2.1.eq}
\]
\end{env}
\end{proposition}

\begin{proof}
It is enough to consider the case where $S$, $S'$, and all the
$Y^{(i)}$ are affine; the gluing is carried out by the methods of
\sref{III.6.7.1} and \sref{III.6.7.2}.  Write
\[
  S=\Spec(A),
  \qquad
  S'=\Spec(A'),
\]
and choose for each $i$ a finite affine open covering
$\mathfrak U^{(i)}$ of $X^{(i)}$.  If
$u_i:X^{\prime(i)}\to X^{(i)}$ is the canonical projection, then
\[
  \mathfrak U^{\prime(i)}
  =
  u_i^{-1}(\mathfrak U^{(i)})
\]
is a finite affine open covering of $X^{\prime(i)}$
(\sref[II]{II.1.5.5}), and
\[
  C^\bullet
  \bigl(
    \mathfrak U^{\prime(i)},\sh P_\bullet^{\prime(i)}
  \bigr)
  =
  C^\bullet
  \bigl(
    \mathfrak U^{(i)},\sh P_\bullet^{(i)}
  \bigr)
  \otimes_AA'.
\]

Let $M_{\bullet\bullet\bullet}^{(i)}$ be a projective
Cartan--Eilenberg resolution, in the sense of
\sref[0]{0.11.7.1}, of
\[
  L_{\bullet\bullet}^{(i)}
  =
  C^\bullet
  \bigl(
    \mathfrak U^{(i)},\sh P_\bullet^{(i)}
  \bigr)
\]
by $A$-modules.  Since $A'$ is flat over $A$,
\[
  M_{\bullet\bullet\bullet}^{(i)}\otimes_AA'
\]
is a projective Cartan--Eilenberg resolution, over $A'$, of
$L_{\bullet\bullet}^{(i)}\otimes_AA'$.  The same flatness gives
\[
\begin{split}
  H_\bullet
  \left(
    \bigotimes_{i=1}^{m}
    (M_{\bullet\bullet\bullet}^{(i)}\otimes_AA')
  \right)
  =
  H_\bullet
  \left(
    \bigotimes_{i=1}^{m}
    M_{\bullet\bullet\bullet}^{(i)}
  \right)
  \otimes_AA',
\end{split}
\]
which proves \eqref{III.6.9.2.1.eq}.
\end{proof}

\begin{env}[6.9.2.2]
\label{III.6.9.2.2}
When $I$ consists of the single element $1$,
\eqref{III.6.9.2.1.eq} also follows directly from
\sref{III.6.7.7}.  Indeed, take
\[
  Y_2=X_2=S',
  \qquad
  f_2=1_{S'},
\]
and let $\sh P_\bullet^{(2)}$ be concentrated in degree $0$ with
value $\sh O_{S'}$.  Then
\[
  \sh H^n(1_{S'},\sh O_{S'})
\]
vanishes for $n\neq0$ and equals $\sh O_{S'}$ for $n=0$
(\sref{III.6.2.1}).
\end{env}

In the general case, replace $\sh O_{S'}$ by a bounded-below
complex $\sh Q_\bullet'$ of quasi-coherent
$\sh O_{S'}$-modules.  Taking, for simplicity,
$I=\{1,2,\ldots,m\}$, one can then consider the
$\partial$-functor
\[
  \mathfrak{Tor}_\bullet^S
  \bigl(
    f_1,\ldots,f_m,1_{S'};
    \sh P_\bullet^{(1)},\ldots,\sh P_\bullet^{(m)},
    \sh Q_\bullet'
  \bigr).
\]

\begin{proposition}[6.9.3]
\label{III.6.9.3}
There are three canonical biregular spectral functors, denoted
\[
  {}^{(t)}\sh E
  \bigl(
    f_1,\ldots,f_m;
    \sh P_\bullet^{(1)},\ldots,\sh P_\bullet^{(m)},
    \sh Q_\bullet'
  \bigr),
  \qquad
  t=e,f,f',
\]
with common abutment
\[
  \mathfrak{Tor}_\bullet^S
  \bigl(
    f_1,\ldots,f_m,1_{S'};
    \sh P_\bullet^{(1)},\ldots,\sh P_\bullet^{(m)},
    \sh Q_\bullet'
  \bigr),
\]
and whose $E_2$-terms are respectively as follows.

\oldpage[III]{167}

\begin{env}[6.9.3.1]
\label{III.6.9.3.1}
\[
\begin{split}
  {}^{(e)}\sh E_{pq}^2
  =
  \bigoplus_{q'+q''=q}
  \shTor_p^S
  \Bigl(
    &
    \mathfrak{Tor}_{q'}^S
    \bigl(
      f_1,\ldots,f_m;
      \sh P_\bullet^{(1)},\ldots,\sh P_\bullet^{(m)}
    \bigr),
    \\
    &
    \sh H_{q''}(\sh Q_\bullet')
  \Bigr).
\end{split}
  \tag{6.9.3.1}\label{III.6.9.3.1.eq}
\]
\end{env}

\begin{env}[6.9.3.2]
\label{III.6.9.3.2}
\[
\begin{split}
  {}^{(f)}\sh E_{pq}^2
  =
  \bigoplus_{q_1+\cdots+q_{m+1}=q}
  \shTor_p^{S'}
  \Bigl(
    &
    \mathfrak{Tor}_{q_1}^S
    \bigl(
      f_1,1_{S'};
      \sh P_\bullet^{(1)},\sh O_{S'}
    \bigr),
    \ldots,
    \\
    &
    \mathfrak{Tor}_{q_m}^S
    \bigl(
      f_m,1_{S'};
      \sh P_\bullet^{(m)},\sh O_{S'}
    \bigr),
    \sh H_{q_{m+1}}(\sh Q_\bullet')
  \Bigr).
\end{split}
  \tag{6.9.3.2}\label{III.6.9.3.2.eq}
\]
\end{env}

\begin{env}[6.9.3.3]
\label{III.6.9.3.3}
\[
\begin{split}
  {}^{(f')}\sh E_{pq}^2
  =
  \bigoplus_{q_1+\cdots+q_m=q}
  \mathfrak{Tor}_p^{S'}
  \Bigl(
    &
    f_1',\ldots,f_m',1_{S'};
    \\
    &
    \shTor_{q_1}^S
    \bigl(
      \sh P_\bullet^{(1)},\sh O_{S'}
    \bigr),
    \ldots,
    \shTor_{q_m}^S
    \bigl(
      \sh P_\bullet^{(m)},\sh O_{S'}
    \bigr),
    \sh Q_\bullet'
  \Bigr).
\end{split}
  \tag{6.9.3.3}\label{III.6.9.3.3.eq}
\]
\end{env}
\end{proposition}

\begin{proof}
The sequence $(e)$ is precisely the associativity sequence of
\sref{III.6.8.2}, with $r=m$.  To define the other two spectral
sequences, it is again enough to treat the case in which $S$, $S'$,
and the $Y^{(i)}$ are affine.  The passage to the general case uses
the methods of \sref{III.6.7.1} and \sref{III.6.7.2} and is left to
the reader.  Thus it remains to prove the following affine
statement.
\end{proof}

\begin{corollary}[6.9.4]
\label{III.6.9.4}
Let $A$ be a ring, let $A'$ be an $A$-algebra, and put
\[
  S=\Spec(A),
  \qquad
  S'=\Spec(A').
\]
Let $X^{(i)}$, $1\leq i\leq m$, be quasi-compact $S$-schemes, and
for each $i$ put
\[
  X^{\prime(i)}=X_{(S')}^{(i)};
\]
this is a quasi-compact $S'$-scheme.  For each $i$, let
$\sh P_\bullet^{(i)}$ be a bounded-below complex of quasi-coherent
$\sh O_{X^{(i)}}$-modules.  Finally, let $Q_\bullet'$ be a
bounded-below complex of $A'$-modules, and let
$\widetilde Q_\bullet'$ be the associated complex.

There are three biregular spectral functors in the
$\sh P_\bullet^{(i)}$ and in $Q_\bullet'$, with common abutment
\[
  \mathbf{Tor}_\bullet^S
  \bigl(
    X^{(1)},\ldots,X^{(m)},S';
    \sh P_\bullet^{(1)},\ldots,\sh P_\bullet^{(m)},
    \widetilde Q_\bullet'
  \bigr),
\]
and whose $E_2$-terms are respectively
\[
\begin{split}
  {}^{(e)}E_{pq}^2
  =
  \bigoplus_{q'+q''=q}
  \Tor_p^A
  \Bigl(
    &
    \mathbf{Tor}_{q'}^S
    \bigl(
      X^{(1)},\ldots,X^{(m)};
      \sh P_\bullet^{(1)},\ldots,\sh P_\bullet^{(m)}
    \bigr),
    \\
    &
    H_{q''}(Q_\bullet')
  \Bigr),
\end{split}
\]
\[
\begin{split}
  {}^{(f)}E_{pq}^2
  =
  \bigoplus_{q_1+\cdots+q_{m+1}=q}
  \Tor_p^{A'}
  \Bigl(
    &
    \mathbf{Tor}_{q_1}^S
    \bigl(
      X^{(1)},S';
      \sh P_\bullet^{(1)},\sh O_{S'}
    \bigr),
    \ldots,
    \\
    &
    \mathbf{Tor}_{q_m}^S
    \bigl(
      X^{(m)},S';
      \sh P_\bullet^{(m)},\sh O_{S'}
    \bigr),
    H_{q_{m+1}}(Q_\bullet')
  \Bigr),
\end{split}
\]
and
\[
\begin{split}
  {}^{(f')}E_{pq}^2
  =
  \bigoplus_{q_1+\cdots+q_m=q}
  \mathbf{Tor}_p^{S'}
  \Bigl(
    &
    X^{\prime(1)},\ldots,X^{\prime(m)},S';
    \\
    &
    \shTor_{q_1}^S
    \bigl(
      \sh P_\bullet^{(1)},\sh O_{S'}
    \bigr),
    \ldots,
    \shTor_{q_m}^S
    \bigl(
      \sh P_\bullet^{(m)},\sh O_{S'}
    \bigr),
    \widetilde Q_\bullet'
  \Bigr).
\end{split}
\]
\end{corollary}

\begin{proof}
We do not return to the first of these spectral functors, which was
treated in \sref{III.6.8.3} and is included here only for reference.
To define the other two, choose for each $i$ a finite affine open
covering $\mathfrak U^{(i)}$ of $X^{(i)}$.  If
\[
  u_i:X^{\prime(i)}\longrightarrow X^{(i)}
\]
is the canonical projection, put
\[
  \mathfrak U^{\prime(i)}
  =
  u_i^{-1}(\mathfrak U^{(i)}).
\]

By \sref{III.6.6.6}, the common abutment is obtained by choosing,
for $1\leq i\leq m$, a projective Cartan--Eilenberg resolution
$M_{\bullet\bullet\bullet}^{(i)}$ of
\[
  L_{\bullet\bullet}^{(i)}
  =
  C^\bullet
  \bigl(
    \mathfrak U^{(i)},\sh P_\bullet^{(i)}
  \bigr)
\]
in the sense of \sref[0]{0.11.7.1}, and then taking the homology of
the tricomplex
\[
  M_{\bullet\bullet\bullet}
  =
  M_{\bullet\bullet\bullet}^{(1)}
  \otimes_A\cdots\otimes_A
  M_{\bullet\bullet\bullet}^{(m)}
  \otimes_A Q_\bullet',
\]
where $Q_\bullet'$ is regarded as a tricomplex whose last two
degrees are zero.  Set
\[
  M_{\bullet\bullet\bullet}^{\prime(i)}
  =
  M_{\bullet\bullet\bullet}^{(i)}\otimes_AA'.
\]
Since $Q_\bullet'$ is a complex of $A'$-modules, we may rewrite
\[
  M_{\bullet\bullet\bullet}
  =
  M_{\bullet\bullet\bullet}^{\prime(1)}
  \otimes_{A'}\cdots\otimes_{A'}
  M_{\bullet\bullet\bullet}^{\prime(m)}
  \otimes_{A'}Q_\bullet'.
\]

Regard each $M_{\bullet\bullet\bullet}^{\prime(i)}$ as a simple
complex for its total degree.  Its terms are projective
$A'$-modules; hence \sref{III.6.3.7}, extended to any finite number
of complexes, identifies the homology of
$M_{\bullet\bullet\bullet}$ with
\[
  \Tor_\bullet^{A'}
  \bigl(
    M_{\bullet\bullet\bullet}^{\prime(1)},\ldots,
    M_{\bullet\bullet\bullet}^{\prime(m)},Q_\bullet'
  \bigr).
\]
By \sref{III.6.5.15}, this is the abutment of a spectral sequence
having the required regularity properties and $E_2$-term
\[
\begin{split}
  E_{pq}^2
  =
  \bigoplus_{q_1+\cdots+q_{m+1}=q}
  \Tor_p^{A'}
  \Bigl(
    &
    H_{q_1}(M_{\bullet\bullet\bullet}^{\prime(1)}),\ldots,
    H_{q_m}(M_{\bullet\bullet\bullet}^{\prime(m)}),
    \\
    &
    H_{q_{m+1}}(Q_\bullet')
  \Bigr).
\end{split}
\]
\oldpage[III]{168}

By the definition of global hypertors,
\[
  H_{q_i}(M_{\bullet\bullet\bullet}^{\prime(i)})
  =
  \mathbf{Tor}_{q_i}^S
  \bigl(
    X^{(i)},S';
    \sh P_\bullet^{(i)},\sh O_{S'}
  \bigr),
\]
which gives the sequence $(f)$.

On the other hand, regard
$M_{\bullet\bullet\bullet}^{\prime(i)}$ as a bicomplex whose first
degree is the sum of the first two degrees of the tricomplex and
whose second degree is its third degree.  Since its terms are
projective $A'$-modules, the general theory of hyperhomology shows
that the homology of
\[
  M_{\bullet\bullet\bullet}^{\prime(1)}
  \otimes_{A'}\cdots\otimes_{A'}
  M_{\bullet\bullet\bullet}^{\prime(m)}
  \otimes_{A'}Q_\bullet'
\]
is canonically isomorphic to its hyperhomology
(\sref[0]{0.11.6.5}).  It is therefore the abutment of a spectral
sequence with $E_2$-term
\[
\begin{split}
  E_{pq}^2
  =
  \bigoplus_{q_1+\cdots+q_{m+1}=q}
  \Tor_p^{A'}
  \Bigl(
    &
    H_{q_1}^{\mathrm{II}}
    (M_{\bullet\bullet\bullet}^{\prime(1)}),\ldots,
    H_{q_m}^{\mathrm{II}}
    (M_{\bullet\bullet\bullet}^{\prime(m)}),
    \\
    &
    H_{q_{m+1}}^{\mathrm{II}}(Q_\bullet')
  \Bigr).
\end{split}
\]
The second degree of $Q_\bullet'$ is zero, so
\[
  H_n^{\mathrm{II}}(Q_\bullet')=0
  \quad(n\neq0),
  \qquad
  H_0^{\mathrm{II}}(Q_\bullet')=Q_\bullet',
\]
and the preceding formula becomes
\[
\begin{split}
  E_{pq}^2
  =
  \bigoplus_{q_1+\cdots+q_m=q}
  \Tor_p^{A'}
  \Bigl(
    &
    H_{q_1}^{\mathrm{II}}
    (M_{\bullet\bullet\bullet}^{\prime(1)}),\ldots,
    \\
    &
    H_{q_m}^{\mathrm{II}}
    (M_{\bullet\bullet\bullet}^{\prime(m)}),
    Q_\bullet'
  \Bigr).
\end{split}
\]

Moreover, by \sref{III.6.3.4},
\[
\begin{split}
  H_{q_i}^{\mathrm{II}}
  (M_{\bullet\bullet\bullet}^{\prime(i)})
  &=
  H_{q_i}^{\mathrm{II}}
  (M_{\bullet\bullet\bullet}^{(i)}\otimes_AA')
  \\
  &=
  \Tor_{q_i}^A(L_{\bullet\bullet}^{(i)},A').
\end{split}
\]
Now $L_{-j,k}^{(i)}$ is the direct sum of the modules
$\Gamma(V,\sh P_k^{(i)})$, where $V$ runs through the affine
intersections of $j+1$ members of $\mathfrak U^{(i)}$.  If
$V'=u_i^{-1}(V)$, then $V'$ is affine in $X^{\prime(i)}$, and
\sref{III.6.4.1.1} gives
\[
  \Gamma
  \bigl(
    V',
    \shTor_{q_i}^S
    (\sh P_k^{(i)},\sh O_{S'})
  \bigr)
  =
  \Tor_{q_i}^A
  \bigl(
    \Gamma(V,\sh P_k^{(i)}),A'
  \bigr).
\]
Consequently, the bicomplex
$H_{q_i}^{\mathrm{II}}
(M_{\bullet\bullet\bullet}^{\prime(i)})$
is
\[
  C^\bullet
  \bigl(
    \mathfrak U^{\prime(i)},
    \shTor_{q_i}^S
    (\sh P_\bullet^{(i)},\sh O_{S'})
  \bigr).
\]
This gives the asserted $E_2$-term of the sequence $(f')$.  Its
biregularity under the stated hypotheses follows in the usual way,
since, when $\sh P_\bullet^{(i)}$ is bounded below, every degree of
$M_{\bullet\bullet\bullet}^{\prime(i)}$ is bounded below.
\end{proof}

\begin{remark}[6.9.5]
\label{III.6.9.5}
As in \sref{III.6.7.6}, replacing the complexes
$\sh P_\bullet^{(i)}$ and $\sh Q_\bullet'$ by respectively
homotopic complexes does not change the spectral sequences $(e)$,
$(f)$, and $(f')$, up to canonical isomorphism.

Moreover, for the sequence $(f)$, homomorphisms of complexes
\[
  \sh P_\bullet^{(i)}\longrightarrow\sh R_\bullet^{(i)},
  \qquad
  \sh Q_\bullet'\longrightarrow\sh T_\bullet'
\]
which induce isomorphisms in homology
\[
  \sh H_\bullet(\sh P_\bullet^{(i)})
  \xrightarrow{\sim}
  \sh H_\bullet(\sh R_\bullet^{(i)}),
  \qquad
  \sh H_\bullet(\sh Q_\bullet')
  \xrightarrow{\sim}
  \sh H_\bullet(\sh T_\bullet')
\]
also induce an isomorphism of spectral sequences
\[
\begin{split}
  {}^{(f)}\sh E
  \bigl(
    f_1,\ldots,f_m;
    \sh P_\bullet^{(1)},\ldots,\sh P_\bullet^{(m)},
    \sh Q_\bullet'
  \bigr)
  \xrightarrow{\sim}\;&
  {}^{(f)}\sh E
  \bigl(
    f_1,\ldots,f_m;
    \\
    &
    \sh R_\bullet^{(1)},\ldots,\sh R_\bullet^{(m)},
    \sh T_\bullet'
  \bigr).
\end{split}
\]
The proof is the same as that of \sref{III.6.7.6}, taking into
account the result of that section and the regularity of the
sequence $(f)$.
\end{remark}

\begin{corollary}[6.9.6]
\label{III.6.9.6}
Under the hypotheses of \sref{III.6.9.1}, suppose that:
\begin{enumerate}
  \item[$1^\circ$] the complexes $\sh P_\bullet^{(i)}$ consist of
  $S$-flat modules, and the $\sh O_Y$-modules
  \[
    \mathfrak{Tor}_n^S
    \bigl(
      f_1,\ldots,f_m;
      \sh P_\bullet^{(1)},\ldots,\sh P_\bullet^{(m)}
    \bigr)
  \]
  are $S$-flat;
  \item[$2^\circ$] the complexes $\sh P_\bullet^{(i)}$ and
  $\sh Q_\bullet'$ are bounded below.
\end{enumerate}

\oldpage[III]{169}

Put
\[
  \sh P_\bullet^{\prime(i)}
  =
  \sh P_\bullet^{(i)}
  \otimes_{\sh O_S}\sh O_{S'}.
\]
There are canonical functorial isomorphisms
\begin{env}[6.9.6.1]
\label{III.6.9.6.1}
\[
\begin{split}
  &
  \mathfrak{Tor}_n^{S'}
  \bigl(
    f_1',\ldots,f_m',1_{S'};
    \sh P_\bullet^{\prime(1)},\ldots,
    \sh P_\bullet^{\prime(m)},\sh Q_\bullet'
  \bigr)
  \xrightarrow{\sim}
  \\
  &\qquad
  \bigoplus_{n'+n''=n}
  \mathfrak{Tor}_{n'}^S
  \bigl(
    f_1,\ldots,f_m;
    \sh P_\bullet^{(1)},\ldots,\sh P_\bullet^{(m)}
  \bigr)
  \otimes_{\sh O_S}
  \sh H_{n''}(\sh Q_\bullet').
\end{split}
  \tag{6.9.6.1}\label{III.6.9.6.1.eq}
\]
\end{env}

In particular, when $\sh Q_\bullet'$ consists of a single term
$\sh F'$ in degree $0$, there are canonical functorial
isomorphisms
\begin{env}[6.9.6.2]
\label{III.6.9.6.2}
\[
\begin{split}
  &
  \mathfrak{Tor}_n^{S'}
  \bigl(
    f_1',\ldots,f_m',1_{S'};
    \sh P_\bullet^{\prime(1)},\ldots,
    \sh P_\bullet^{\prime(m)},\sh F'
  \bigr)
  \xrightarrow{\sim}
  \\
  &\qquad
  \mathfrak{Tor}_n^S
  \bigl(
    f_1,\ldots,f_m;
    \sh P_\bullet^{(1)},\ldots,\sh P_\bullet^{(m)}
  \bigr)
  \otimes_{\sh O_S}\sh F'.
\end{split}
  \tag{6.9.6.2}\label{III.6.9.6.2.eq}
\]
\end{env}
More particularly, for $\sh F'=\sh O_{S'}$,
\begin{env}[6.9.6.3]
\label{III.6.9.6.3}
\[
\begin{split}
  &
  \mathfrak{Tor}_n^{S'}
  \bigl(
    f_1',\ldots,f_m';
    \sh P_\bullet^{\prime(1)},\ldots,
    \sh P_\bullet^{\prime(m)}
  \bigr)
  \xrightarrow{\sim}
  \\
  &\qquad
  \mathfrak{Tor}_n^S
  \bigl(
    f_1,\ldots,f_m;
    \sh P_\bullet^{(1)},\ldots,\sh P_\bullet^{(m)}
  \bigr)
  \otimes_{\sh O_S}\sh O_{S'}.
\end{split}
  \tag{6.9.6.3}\label{III.6.9.6.3.eq}
\]
\end{env}
\end{corollary}

\begin{proof}
The flatness of the terms of $\sh P_\bullet^{(i)}$ implies that
\[
  \shTor_{q_i}^S
  \bigl(
    \sh P_\bullet^{(i)},\sh O_{S'}
  \bigr)
  =
  0
  \qquad(q_i\neq0)
\]
by \sref{III.6.5.8}.  Thus the sequence $(f')$ degenerates.
Hypothesis $2^\circ$ also makes it biregular
(\sref{III.6.9.3}), so the edge homomorphism
\[
\begin{split}
  &
  \mathfrak{Tor}_n^S
  \bigl(
    f_1,\ldots,f_m,1_{S'};
    \sh P_\bullet^{(1)},\ldots,\sh P_\bullet^{(m)},
    \sh Q_\bullet'
  \bigr)
  \longrightarrow
  {}^{(f')}\sh E_{n0}^2
  \\
  &\qquad =
  \mathfrak{Tor}_n^{S'}
  \bigl(
    f_1',\ldots,f_m',1_{S'};
    \sh P_\bullet^{\prime(1)},\ldots,
    \sh P_\bullet^{\prime(m)},\sh Q_\bullet'
  \bigr)
\end{split}
  \tag{6.9.6.4}\label{III.6.9.6.4.eq}
\]
is bijective by \sref[0]{0.11.1.6}.

The flatness of the modules
\[
  \mathfrak{Tor}_n^S
  \bigl(
    f_1,\ldots,f_m;
    \sh P_\bullet^{(1)},\ldots,\sh P_\bullet^{(m)}
  \bigr)
\]
implies, by \sref{III.6.5.8}, that
${}^{(e)}\sh E_{pq}^2=0$ for $p\neq0$.  Hence $(e)$ also
degenerates, and its biregularity makes the edge homomorphism
\[
\begin{split}
  &
  \mathfrak{Tor}_n^S
  \bigl(
    f_1,\ldots,f_m,1_{S'};
    \sh P_\bullet^{(1)},\ldots,\sh P_\bullet^{(m)},
    \sh Q_\bullet'
  \bigr)
  \longrightarrow
  {}^{(e)}\sh E_{0n}^2
  \\
  &\qquad =
  \bigoplus_{n'+n''=n}
  \mathfrak{Tor}_{n'}^S
  \bigl(
    f_1,\ldots,f_m;
    \sh P_\bullet^{(1)},\ldots,\sh P_\bullet^{(m)}
  \bigr)
  \otimes_{\sh O_S}
  \sh H_{n''}(\sh Q_\bullet')
\end{split}
  \tag{6.9.6.5}\label{III.6.9.6.5.eq}
\]
bijective.  Combining these two edge isomorphisms gives
\sref{III.6.9.6.1}.  Formula \sref{III.6.9.6.2} follows because
$\sh H_q(\sh Q_\bullet')=0$ for $q\neq0$ and
$\sh H_0(\sh Q_\bullet')=\sh F'$.  Finally, taking
$\sh F'=\sh O_{S'}$ in \sref{III.6.9.6.2} and using
\sref{III.6.7.12} gives \sref{III.6.9.6.3}.
\end{proof}

\begin{corollary}[6.9.7]
\label{III.6.9.7}
Under the hypotheses of \sref{III.6.9.1}, suppose that $S$ and
$S'$ are affine, and choose an integer $d_i$ for each
$1\leq i\leq m$.  There is an integer $N$, depending only on $S$,
the $X^{(i)}$, and the $d_i$, with the following property.  For
every integer $n_0$, the canonical isomorphisms
\sref{III.6.9.6.3} hold for $n\leq n_0$ and for every system of
complexes $\sh P_\bullet^{(i)}$ satisfying:
\begin{enumerate}
  \item[$1^\circ$] $\sh P_k^{(i)}=0$ for $k<d_i$;
  \item[$2^\circ$] $\sh P_k^{(i)}$ is $S$-flat for
  $k<n_0+N$;
  \item[$3^\circ$]
  \[
    \mathfrak{Tor}_q^S
    \bigl(
      f_1,\ldots,f_m;
      \sh P_\bullet^{(1)},\ldots,\sh P_\bullet^{(m)}
    \bigr)
  \]
  is $S$-flat for $q<n_0+N$.
\end{enumerate}
\end{corollary}

\begin{proof}
Suppose that $\sh P_k^{(i)}$ is $S$-flat for $k<r$.  Then
\[
  \shTor_{q_i}^S
  \bigl(
    \sh P_k^{(i)},\sh O_{S'}
  \bigr)
  =
  0
  \qquad
  (k<r,\ q_i\neq0).
\]
Consider
\[
\begin{split}
  \mathfrak{Tor}_p^{S'}
  \Bigl(
    &
    f_1',\ldots,f_m';
    \shTor_{q_1}^S
    (\sh P_\bullet^{(1)},\sh O_{S'}),\ldots,
    \\
    &
    \shTor_{q_m}^S
    (\sh P_\bullet^{(m)},\sh O_{S'})
  \Bigr).
\end{split}
  \tag{6.9.7.1}\label{III.6.9.7.1.eq}
\]

\oldpage[III]{170}

Compute \eqref{III.6.9.7.1.eq} by the method of
\sref{III.6.6.2}, using the inverse images of fixed finite affine
open coverings $\mathfrak U^{(i)}$ of $X^{(i)}$.  These coverings
are independent of $S'$ and of the complexes
$\sh P_\bullet^{(i)}$.  The terms $L_{jk}^{(i)}$ vanish for
$j<-N_i$, where $N_i$ depends only on $\mathfrak U^{(i)}$.

If some $q_i$ is nonzero, the simple complex whose homology in
degree $p$ is \eqref{III.6.9.7.1.eq} has zero terms in every degree
\[
  <
  r+\sum_{j\neq i}d_j-\sum_{\ell=1}^mN_\ell.
\]
Consequently \eqref{III.6.9.7.1.eq} vanishes for $p<r-N$, where
\[
  N
  =
  \max_i
  \left(
    \sum_{\ell=1}^mN_\ell-\sum_{j\neq i}d_j
  \right).
\]
It follows that
${}^{(f')}\sh E_{pq}^2=0$ for $q\neq0$ and $p<r-N$.
Since also ${}^{(f')}\sh E_{pq}^2=0$ for $q<0$, the edge
homomorphism \eqref{III.6.9.6.4.eq} is bijective for $n<r-N$
(M, XV, 5.6), when $\sh Q_\bullet'=\sh O_{S'}$.

Secondly, if
\[
  \mathfrak{Tor}_q^S
  \bigl(
    f_1,\ldots,f_m;
    \sh P_\bullet^{(1)},\ldots,\sh P_\bullet^{(m)}
  \bigr)
\]
is $S$-flat for $q<r$, then
${}^{(e)}\sh E_{pq}^2=0$ for $p\neq0$ and $q<r$; moreover,
${}^{(e)}\sh E_{pq}^2=0$ for $p<0$.  Thus the edge homomorphism
\eqref{III.6.9.6.5.eq} is bijective for $n<r$, which completes the
proof.
\end{proof}

The case of \sref{III.6.9.3} most important in applications is
$m=1$, with $\sh Q_\bullet'$ reduced to a single term $\sh F'$ in
degree $0$.  We state it again for later reference.\footnote{The
case treated in \sref{III.6.9.8}, and in particular the spectral
sequences \sref{III.6.9.8.3} and \sref{III.6.9.8.4}, was pointed
out to us by J.-P.~Serre in 1957.}

\begin{proposition}[6.9.8]
\label{III.6.9.8}
Let $S$ be a prescheme, let $g:S'\to S$ be a morphism, and let
$f:X\to Y$ be a separated and quasi-compact $S$-morphism of
$S$-preschemes.  Let $\sh P_\bullet$ be a bounded-below complex of
quasi-coherent $\sh O_X$-modules, and let $\sh F'$ be a
quasi-coherent $\sh O_{S'}$-module.  There are two biregular
spectral functors in $\sh P_\bullet$ and $\sh F'$, with values in
the category of quasi-coherent $\sh O_{Y_{(S')}}$-modules, having
the same abutment
\[
  \mathfrak{Tor}_\bullet^S
  \bigl(
    f,1_{S'};\sh P_\bullet,\sh F'
  \bigr),
\]
and whose $E_2$-terms are
\begin{env}[6.9.8.1]
\label{III.6.9.8.1}
\[
  {}'\sh E_{pq}^2
  =
  \shTor_p^S
  \bigl(
    \sh H^{-q}(f,\sh P_\bullet),\sh F'
  \bigr)
  \tag{6.9.8.1}\label{III.6.9.8.1.eq}
\]
\end{env}
and
\begin{env}[6.9.8.2]
\label{III.6.9.8.2}
\[
  {}''\sh E_{pq}^2
  =
  \sh H^{-p}
  \bigl(
    f',
    \shTor_q^S(\sh P_\bullet,\sh F')
  \bigr),
  \tag{6.9.8.2}\label{III.6.9.8.2.eq}
\]
\end{env}
where
\[
  f'=f_{(S')}:
  X_{(S')}\longrightarrow Y_{(S')}.
\]
\end{proposition}

\begin{proof}
These spectral sequences may also be obtained, without starting
from \sref{III.6.9.3}, from the sequences $(a)$ and $(b')$ of
\sref{III.6.7.3}, by taking
\[
  X^{(1)}=X,\quad
  Y^{(1)}=Y,\quad
  X^{(2)}=Y^{(2)}=S',\quad
  f_1=f,\quad
  f_2=1_{S'}.
\]
When $S=S'=Y$ and $Y$ is affine, one obtains two spectral
sequences with $E_2$-terms
\[
  {}'\sh E_{pq}^2
  =
  \shTor_p^Y
  \bigl(
    \sh H^{-q}(f,\sh P_\bullet),\sh F
  \bigr)
  \tag{6.9.8.3}\label{III.6.9.8.3}
\]
and
\[
  {}''\sh E_{pq}^2
  =
  \sh H^{-p}
  \bigl(
    f,\shTor_q^Y(\sh P_\bullet,\sh F)
  \bigr).
  \tag{6.9.8.4}\label{III.6.9.8.4}
\]
By \sref{III.6.7.6}, both abut to the hypercohomology
\[
  \sh H^\bullet
  \bigl(
    f,\sh P_\bullet\otimes_Y\sh F
  \bigr)
\]
of $f_*$ with respect to the complex
$\sh P_\bullet\otimes_Y\sh F$ of $\sh O_X$-modules, for every
quasi-coherent $Y$-flat $\sh O_Y$-module $\sh F$; if the terms of
$\sh P_\bullet$ are $Y$-flat, $\sh F$ may be any quasi-coherent
$\sh O_Y$-module.  These spectral sequences are distinct from
those of \sref{III.6.2.1}.
\end{proof}

\begin{corollary}[6.9.9]
\label{III.6.9.9}
Under the hypotheses of \sref{III.6.9.8}, suppose that
$\sh P_\bullet$ is bounded below, consists of $S$-flat modules,
and that the $\sh O_Y$-modules
$\sh H^n(f,\sh P_\bullet)$ are $S$-flat.

\oldpage[III]{171}

Put
\[
  \sh P_\bullet'
  =
  \sh P_\bullet\otimes_{\sh O_S}\sh O_{S'}.
\]
There are canonical functorial isomorphisms
\begin{env}[6.9.9.1]
\label{III.6.9.9.1}
\[
\begin{split}
  &
  \mathfrak{Tor}_n^{S'}
  \bigl(
    f',1_{S'};\sh P_\bullet',\sh F'
  \bigr)
  \xrightarrow{\sim}
  \\
  &\qquad
  \sh H^{-n}(f,\sh P_\bullet)
  \otimes_{\sh O_S}\sh F'.
\end{split}
  \tag{6.9.9.1}\label{III.6.9.9.1.eq}
\]
\end{env}
In particular, for $\sh F'=\sh O_{S'}$,
\begin{env}[6.9.9.2]
\label{III.6.9.9.2}
\[
  \sh H^n(f',\sh P_\bullet')
  \xrightarrow{\sim}
  \sh H^n(f,\sh P_\bullet)
  \otimes_{\sh O_S}\sh O_{S'}.
  \tag{6.9.9.2}\label{III.6.9.9.2.eq}
\]
\end{env}
This is the case $m=1$ of \sref{III.6.9.6}.
\end{corollary}

\begin{corollary}[6.9.10]
\label{III.6.9.10}
Let $S$ be a prescheme, let $f:X\to Y$ be a separated and
quasi-compact $S$-morphism of $S$-preschemes, and let
$\sh P_\bullet$ be a bounded-below complex of quasi-coherent
$\sh O_X$-modules that are flat over $S$.  Suppose further that
the $\sh O_Y$-modules $\sh H^n(f,\sh P_\bullet)$ are flat over
$S$.

For $s\in S$, write $X_s$ and $Y_s$ for the fibres
\[
  X\times_S\Spec(\kres(s)),
  \qquad
  Y\times_S\Spec(\kres(s)),
\]
write
\[
  f_s:X_s\longrightarrow Y_s
\]
for the morphism $f\times_S1$, and put
\[
  \sh P_\bullet^s
  =
  \sh P_\bullet\otimes_{\sh O_S}\kres(s).
\]
Then there are canonical functorial isomorphisms
\begin{env}[6.9.10.1]
\label{III.6.9.10.1}
\[
  \sh H^n(f_s,\sh P_\bullet^s)
  \xrightarrow{\sim}
  \sh H^n(f,\sh P_\bullet)
  \otimes_{\sh O_S}\kres(s).
  \tag{6.9.10.1}\label{III.6.9.10.1.eq}
\]
\end{env}
\end{corollary}

Thus, under suitable flatness hypotheses, the formation of the
derived functors $R^nf_*(\sh F)$ ``commutes with passage to the
fibres.''  We shall obtain this result again by another method in
\S\,7.

\subsection{Local structure of certain cohomological functors}
\label{subsection:III.6.10}
\label{III.6.10}

\begin{proposition}[6.10.1]
\label{III.6.10.1}
Let $S=\Spec(A)$ be an affine scheme, and let
$Y^{(i)}$ $(1\leq i\leq n)$ be a finite family of affine
$S$-schemes, flat over $S$.  For each $i$, let
\[
  f_i:X^{(i)}\longrightarrow Y^{(i)}
\]
be a separated and quasi-compact $S$-morphism, and let
$\sh P_\bullet^{(i)}$ be a bounded-below complex of
quasi-coherent $\sh O_{X^{(i)}}$-modules.  Let $Y$ be the product
over $S$ of the schemes $Y^{(i)}$.

There exists a complex $\sh R_\bullet$ of quasi-coherent
$\sh O_Y$-modules, flat over $S$, with the following property.
For every affine $S$-scheme $S'$ and every bounded-below complex
$\sh Q_\bullet'$ of quasi-coherent $\sh O_{S'}$-modules, there is
an isomorphism
\begin{env}[6.10.1.1]
\label{III.6.10.1.1}
\[
  \mathbf{Tor}_\bullet^S
  \bigl(
    f_1,\ldots,f_n,1_{S'};
    \sh P_\bullet^{(1)},\ldots,\sh P_\bullet^{(n)},
    \sh Q_\bullet'
  \bigr)
  \xrightarrow{\sim}
  \sh H_\bullet
  \bigl(
    \sh R_\bullet\otimes_{\sh O_S}\sh Q_\bullet'
  \bigr),
  \tag{6.10.1.1}\label{III.6.10.1.1.eq}
\]
\end{env}
which is an isomorphism of $\partial$-functors in
$\sh Q_\bullet'$.

Moreover, for every $S$-morphism
$u:S''\to S'$ of affine $S$-schemes, the diagram
\begin{env}[6.10.1.2]
\label{III.6.10.1.2}
\[
\xymatrix@C=12pt@R=30pt{
  {\begin{gathered}
    \mathbf{Tor}_\bullet^S
    (f_1,\ldots,f_n,1_{S'};\\
    \sh P_\bullet^{(1)},\ldots,\sh P_\bullet^{(n)},
    \sh Q_\bullet')
  \end{gathered}}
  \ar[r]^-{\sim}\ar[d]
  &
  {\sh H_\bullet
    (\sh R_\bullet\otimes_{\sh O_S}\sh Q_\bullet')}
  \ar[d]
  \\
  {\begin{gathered}
    \mathbf{Tor}_\bullet^S
    (f_1,\ldots,f_n,1_{S''};\\
    \sh P_\bullet^{(1)},\ldots,\sh P_\bullet^{(n)},
    u^*(\sh Q_\bullet'))
  \end{gathered}}
  \ar[r]^-{\sim}
  &
  {\sh H_\bullet
    (\sh R_\bullet\otimes_{\sh O_S}u^*(\sh Q_\bullet'))}
}
\tag{6.10.1.2}\label{III.6.10.1.2.eq}
\]
\end{env}
commutes.  Here the vertical arrows are the canonical
$(1_Y\times_Su)$-morphisms defined in \sref{III.6.7.10}.
\end{proposition}

\ifdefined\EGAIIIStopAfterPageOneSevenOne
  \expandafter\endinput
\fi

\begin{proof}
\oldpage[III]{172}
Compute the hypertor by the method of \sref{III.6.6.2}, taking
account of Remark \sref{III.6.6.6}.  With the notation used
there, each $L_{\bullet\bullet}^{(i)}$ is a bicomplex of
$A_i$-modules, where $A_i$ denotes the ring of $Y^{(i)}$.
It therefore has a projective Cartan--Eilenberg resolution
$M_{\bullet\bullet\bullet}^{(i)}$ (in the sense of
\sref[0]{0.11.7.1}) by projective $A_i$-modules.  By
\sref{III.6.6.6}, the left side of
\eqref{III.6.10.1.1.eq} is canonically isomorphic to
\[
  \HH_\bullet
  \bigl(
    M_{\bullet\bullet\bullet}\otimes_AQ_\bullet'
  \bigr),
\]
where
\[
  M_{\bullet\bullet\bullet}
  =
  M_{\bullet\bullet\bullet}^{(1)}
  \otimes_A
  M_{\bullet\bullet\bullet}^{(2)}
  \otimes_A\cdots\otimes_A
  M_{\bullet\bullet\bullet}^{(n)}
\]
and $\sh Q_\bullet'=\widetilde{Q_\bullet'}$.

By hypothesis, the rings $A_i$ are flat $A$-modules.  Hence the
$M_{\bullet\bullet\bullet}^{(i)}$ are tricomplexes of flat
$A$-modules \sref[0]{0.6.2.1}, and so is
$M_{\bullet\bullet\bullet}$.  Furthermore, if $B$ is the ring of
$Y$, that is, the tensor product over $A$ of the $A_i$, then
$M_{\bullet\bullet\bullet}$ is a tricomplex of $B$-modules.
Consequently the complex
\[
  \sh R_\bullet
  =
  \widetilde{M_{\bullet\bullet\bullet}},
\]
where $M_{\bullet\bullet\bullet}$ is regarded as a simple
complex, has the required properties, as follows immediately
from \sref{III.6.7.10}.
\end{proof}

\begin{corollary}[6.10.2]
\label{III.6.10.2}
In Proposition \sref{III.6.10.1}, the complex
$\sh R_\bullet$ may be taken to be bounded below.  If the
$\sh P_\bullet^{(i)}$ are bounded above and the $Y^{(i)}$ have
finite cohomological dimension, then $\sh R_\bullet$ may be
taken to be bounded above.
\end{corollary}

\begin{proof}
The first assertion follows because all three degrees of each
$M_{\bullet\bullet\bullet}^{(i)}$ are bounded below.  On the
other hand, if the rings $A_i$ have finite cohomological
dimension, the third degree of each
$M_{\bullet\bullet\bullet}^{(i)}$ assumes only finitely many
values, and by its construction the same is true of its first
degree \sref{III.6.6.2}.  Its second degree is bounded above
precisely when the degree of $\sh P_\bullet^{(i)}$ is bounded
above \sref{III.6.6.2}, which proves the second assertion.
\end{proof}

\begin{remarks}[6.10.3]
\label{III.6.10.3}
\textnormal{(i)}
With the notation of \sref{III.6.10.1},
\[
  \sh H_\bullet
  (\sh R_\bullet\otimes_S\sh Q_\bullet')
  \simeq
  \mathbf{Tor}_\bullet(\sh R_\bullet,\sh Q_\bullet'),
\]
because $\sh R_\bullet$ consists of $\sh O_Y$-modules that are
flat over $S$ \sref{III.6.5.9}.  It is therefore
\sref{III.6.5.4} the abutment of a regular spectral sequence
whose $E^2$-term is
\begin{env}[6.10.3.1]
\label{III.6.10.3.1}
\[
  E^2_{pq}
  =
  \bigoplus_{q'+q''=q}
  \shTor_p^S
  \bigl(
    \sh H_{q'}(\sh R_\bullet),
    \sh H_{q''}(\sh Q_\bullet')
  \bigr).
  \tag{6.10.3.1}\label{III.6.10.3.1.eq}
\]
\end{env}
This is precisely the spectral sequence (e) for base change in
\sref{III.6.9.3}.

\medskip
\textnormal{(ii)}
Let $\sh R_\bullet'$ be another complex of quasi-coherent
$\sh O_Y$-modules, flat over $S$, and let
\[
  g:\sh R_\bullet'\longrightarrow\sh R_\bullet
\]
be a homomorphism of complexes for which
\[
  \sh H_\bullet(g):
  \sh H_\bullet(\sh R_\bullet')
  \longrightarrow
  \sh H_\bullet(\sh R_\bullet)
\]
is an isomorphism.  By \sref{III.6.3.3} and
\sref{III.6.5.9}, $g$ induces an isomorphism of
$\partial$-functors in $\sh Q_\bullet'$,
\[
  \sh H_\bullet
  (\sh R_\bullet'\otimes_S\sh Q_\bullet')
  \xrightarrow{\sim}
  \sh H_\bullet
  (\sh R_\bullet\otimes_S\sh Q_\bullet'),
\]
such that the diagram
\[
\xymatrix@C=28pt@R=30pt{
  \sh H_\bullet
  (\sh R_\bullet'\otimes_S\sh Q_\bullet')
  \ar[r]^-{\sim}\ar[d]
  &
  \sh H_\bullet
  (\sh R_\bullet\otimes_S\sh Q_\bullet')
  \ar[d]
  \\
  \sh H_\bullet
  (\sh R_\bullet'\otimes_Su^*(\sh Q_\bullet'))
  \ar[r]^-{\sim}
  &
  \sh H_\bullet
  (\sh R_\bullet\otimes_Su^*(\sh Q_\bullet'))
}
\]
\par\medskip\noindent
commutes.  Thus the properties in \sref{III.6.10.1} do not
determine the complex $\sh R_\bullet$ completely.

\medskip
\oldpage[III]{173}
\textnormal{(iii)}
In the proof of \sref{III.6.10.1}, the
$M_{\bullet\bullet\bullet}^{(i)}$ may be chosen to consist of
free $A_i$-modules, as follows immediately by dualising the
proof of \sref[0]{0.11.5.2.1}.  The
\[
  M_{\bullet\bullet\bullet}^{(i)}\otimes_{A_i}B
\]
then consist of free $B$-modules; since
$M_{\bullet\bullet\bullet}$ is their tensor product over $B$,
one may in addition require in \sref{III.6.10.1} that
$\sh R_\bullet$ be associated with a complex of free
$B$-modules.

Moreover, by (M, XVII, 1.2), the tricomplex
$M_{\bullet\bullet\bullet}$ depends functorially on each of the
bicomplexes $L_{\bullet\bullet}^{(i)}$, hence on each
$\sh P_\bullet^{(i)}$ once a finite covering of each
$X^{(i)}$ has been fixed.  Here ``morphisms'' of tricomplexes
are understood as homomorphism classes modulo homotopy.
Replacing a covering of $X^{(i)}$ by a finer one gives, for
$L_{\bullet\bullet}^{(i)}$, homomorphisms defined up to homotopy
by \sref{III.6.6.8}; thus, with the preceding convention on
morphisms, $M_{\bullet\bullet\bullet}$ is a functor in each
$\sh P_\bullet^{(i)}$.

We shall make this functorial dependence more precise, in
particular the behaviour of $\sh R_\bullet$ with respect to
short exact sequences
\[
  0\longrightarrow\sh P_\bullet^{\prime(i)}
  \longrightarrow\sh P_\bullet^{(i)}
  \longrightarrow\sh P_\bullet^{\prime\prime(i)}
  \longrightarrow0,
\]
in the chapter devoted to a general algebra of cohomological
functors mentioned in \sref{III.6.1.3}.
\end{remarks}

\begin{scholium}[6.10.4]
\label{III.6.10.4}
The fact that $\sh R_\bullet$ consists of $\sh O_Y$-modules
flat over $S$ immediately implies that
\[
  \sh H_\bullet
  (\sh R_\bullet\otimes_S\sh Q_\bullet')
\]
is a homological functor in $\sh Q_\bullet'$; compare the
argument of \sref{III.7.7.1}.  As noted in
\sref{III.6.1.1}, this property motivates the introduction of
hypertor.

Put
\[
  X=X^{(1)}\times_SX^{(2)}\times_S\cdots\times_SX^{(n)},
  \qquad
  f=f_1\times_Sf_2\times_S\cdots\times_Sf_n,
\]
\[
  \sh P_\bullet
  =
  \sh P_\bullet^{(1)}
  \otimes_S\sh P_\bullet^{(2)}
  \otimes_S\cdots\otimes_S\sh P_\bullet^{(n)},
\]
and
\[
  X'=X\times_SS',
  \qquad
  Y'=Y\times_SS',
  \qquad
  f'=f\times_S1_{S'}.
\]
Problems of base change lead one to study the hypercohomology
\[
  \sh H^*_{f'}
  (\sh P_\bullet\otimes_S\sh N')
\]
as a functor of the quasi-coherent $\sh O_{S'}$-module
$\sh N'$, or equivalently the hypercohomology
\[
  \sh H^*_f
  (\sh P_\bullet\otimes_S\sh N)
\]
as a functor of the quasi-coherent $\sh O_S$-module $\sh N$.
When the $\sh P_\bullet^{(i)}$, and therefore
$\sh P_\bullet$, are flat over $S$, the preceding results and
\sref{III.6.7.6} show that this is indeed a cohomological
functor in $\sh N$.  This is no longer true when the flatness
of the $\sh P_\bullet^{(i)}$ is not assumed, and the study of
\[
  \sh H^*_f(\sh P_\bullet\otimes_S\sh N)
\]
can then no longer be approached by the usual methods of
homological algebra.

We shall chiefly use the case $n=1$, $Y=S$, in which
$\sh P_\bullet$ consists of $\sh O_X$-modules flat over $Y$.
In this case we have the following theorem.
\end{scholium}

\begin{theorem}[6.10.5]
\label{III.6.10.5}
Let $Y=\Spec(A)$ be an affine Noetherian scheme, let
$f:X\to Y$ be a proper morphism, and let $\sh P_\bullet$ be a
bounded-below complex of coherent $\sh O_X$-modules flat over
$Y$.  There exists a bounded-below complex $\sh L_\bullet$ of
$\sh O_Y$-modules, whose terms have the form
\[
  \sh L_i=\sh O_Y^{m_i},
\]
and an isomorphism
\begin{env}[6.10.5.1]
\label{III.6.10.5.1}
\[
  \sh H^*
  (f,\sh P_\bullet\otimes_Y\sh Q_\bullet)
  \xrightarrow{\sim}
  \sh H_\bullet
  (\sh L_\bullet\otimes_Y\sh Q_\bullet)
  \tag{6.10.5.1}\label{III.6.10.5.1.eq}
\]
\end{env}
of $\partial$-functors in the bounded-below complex
$\sh Q_\bullet$ of quasi-coherent $\sh O_Y$-modules.

Moreover, for every morphism $u:Y'\to Y$, put
\[
  X'=X_{(Y')},
  \qquad
  f'=f_{(Y')},
  \qquad
  \sh P_\bullet'
  =
  \sh P_\bullet\otimes_Y\sh O_{Y'},
  \qquad
  \sh L_\bullet'=u^*(\sh L_\bullet).
\]
\oldpage[III]{174}
Then $\sh L_\bullet'$ is a complex of locally free
$\sh O_{Y'}$-modules of finite type, and there is an isomorphism
\begin{env}[6.10.5.2]
\label{III.6.10.5.2}
\[
  \sh H^*
  (f',\sh P_\bullet'\otimes_{Y'}\sh Q_\bullet')
  \xrightarrow{\sim}
  \sh H_\bullet
  (\sh L_\bullet'\otimes_{Y'}\sh Q_\bullet')
  \tag{6.10.5.2}\label{III.6.10.5.2.eq}
\]
\end{env}
of $\partial$-functors in the bounded-below complex
$\sh Q_\bullet'$ of quasi-coherent $\sh O_{Y'}$-modules.  These
isomorphisms may be chosen so that the diagram
\begin{env}[6.10.5.3]
\label{III.6.10.5.3}
\[
\xymatrix@C=24pt@R=30pt{
  \sh H^*
  (f,\sh P_\bullet\otimes_Y\sh Q_\bullet)
  \ar[r]^-{\sim}\ar[d]
  &
  \sh H_\bullet
  (\sh L_\bullet\otimes_Y\sh Q_\bullet)
  \ar[d]
  \\
  \sh H^*
  (f',\sh P_\bullet'\otimes_{Y'}u^*(\sh Q_\bullet))
  \ar[r]^-{\sim}
  &
  \sh H_\bullet
  (\sh L_\bullet'\otimes_{Y'}u^*(\sh Q_\bullet))
}
\tag{6.10.5.3}\label{III.6.10.5.3.eq}
\]
\end{env}
commutes.
\end{theorem}

\begin{proof}
Proposition \sref{III.6.10.1} first gives a bounded-below
complex $\sh R_\bullet$ \sref{III.6.10.2} of quasi-coherent
$\sh O_Y$-modules that are free over $Y$
\sref{III.6.10.3}~(iii), together with an isomorphism
\begin{env}[6.10.5.4]
\label{III.6.10.5.4}
\[
  \sh H^*
  (f,\sh P_\bullet\otimes_Y\sh Q_\bullet)
  \xrightarrow{\sim}
  \sh H_\bullet
  (\sh R_\bullet\otimes_Y\sh Q_\bullet)
  \tag{6.10.5.4}\label{III.6.10.5.4.eq}
\]
\end{env}
of $\partial$-functors in $\sh Q_\bullet$.  The terms of
$\sh R_\bullet$ need not, \emph{a priori}, be
$\sh O_Y$-modules of finite type.

Apply \eqref{III.6.10.5.4.eq} when $\sh Q_\bullet$ is the
complex concentrated in one degree with sole term $\sh O_Y$.
It follows that $\sh H_\bullet(\sh R_\bullet)$ is isomorphic to
$\sh H^*(f,\sh P_\bullet)$ and hence consists of coherent
$\sh O_Y$-modules \sref{III.6.2.5}.  By
\sref[0]{0.11.9.2}, there is therefore a bounded-below complex
$\sh L_\bullet$ whose terms are associated with free
$A$-modules of finite type, and a homomorphism
\[
  \sh L_\bullet\longrightarrow\sh R_\bullet
\]
which induces an isomorphism
\[
  \sh H_\bullet(\sh L_\bullet)
  \xrightarrow{\sim}
  \sh H_\bullet(\sh R_\bullet).
\]
Formula \eqref{III.6.10.5.1.eq} now follows from
\sref{III.6.10.3}~(ii).

The remaining assertions of \sref{III.6.10.5} follow from
\sref{III.6.10.1} and \sref{III.6.10.3}~(ii) when $Y'$ is
affine.  In the general case, cover $Y'$ by affine open subsets
$(V_\alpha)$.  If
\eqref{III.6.10.5.2.eq} is formed over each $V_\alpha$, its
restrictions from $V_\alpha$ and $V_\beta$ to an affine open
subset
\[
  W\subset V_\alpha\cap V_\beta
\]
both coincide with the isomorphism formed over $W$.  This
follows from the commutativity of
\eqref{III.6.10.1.2.eq}, applied to the canonical inclusions
$W\to V_\alpha$ and $W\to V_\beta$, and proves that the local
isomorphisms glue.
\end{proof}

\begin{remark}[6.10.6]
\label{III.6.10.6}
In the following chapters we shall apply \sref{III.6.10.5}
chiefly when $\sh P_\bullet$ is concentrated in one degree and
its sole term is a coherent $\sh O_X$-module $\sh F$ flat over
$Y$.  Since
\[
  \sh H_n(\sh L_\bullet)
  =
  \sh H^{-n}(f,\sh F)
  =
  R^{-n}f_*(\sh F)
  \qquad\text{\sref{III.6.2.1}},
\]
the $\sh H_n(\sh L_\bullet)$ vanish for $n>0$.  We shall see
later \sref{III.7.7.12}~(i) that $\sh L_\bullet$ may then be
chosen to have terms only in degrees $\leq0$, and hence only
finitely many terms, provided the requirement that its terms
be associated with free $A$-modules of finite type is replaced
by the requirement that they be locally free of finite type.

The complex $\sh L_\bullet$ corresponding to such an
$\sh O_X$-module $\sh F$ seems to have no special property
\oldpage[III]{175}
apart from the preceding restriction on its degrees.  One may
therefore ask conversely whether, given a bounded-below complex
$\sh L_\bullet$ whose terms are $\sh O_Y$-modules associated
with projective $A$-modules of finite type and whose terms of
degree $>0$ vanish, there exist a $Y$-scheme $X$, projective and
flat over $Y$, and a locally free $\sh O_X$-module $\sh F$ such
that there is an isomorphism
\[
  \sh H^*
  (f,\sh F\otimes_Y\sh Q_\bullet)
  \xrightarrow{\sim}
  \sh H_\bullet
  (\sh L_\bullet\otimes_Y\sh Q_\bullet)
\]
functorial in $\sh Q_\bullet$.  Such a result would reduce the
cohomological theory of coherent modules flat over $Y$ on
proper $Y$-schemes entirely to the theory, ``up to homotopy,''
of complexes of projective $A$-modules of finite type over a
Noetherian ring $A$.
\end{remark}
% END NATIVE INPUT ega3/ega3-6.tex
% BEGIN NATIVE INPUT ega3/ega3-7.tex
\section{Base change for homological functors of sheaves of modules}
\label{section:III.7}


\subsection{Functors of $A$-modules}
\label{subsection:III.7.1}

\begin{env}[7.1.1]
\label{III.7.1.1}
Given a ring $A$, not necessarily commutative, denote by
$\cat{Ab}_A$ the category of left $A$-modules, and simply by
$\cat{Ab}$ the category of $\bb Z$-modules, which is the same
as the category of commutative groups.  Let
\[
  T:\cat{Ab}_A\longrightarrow\cat{Ab}
\]
be an additive covariant functor, and let $M$ be an
$(A,A)$-bimodule.  Then $T(M)$ is naturally endowed with the
structure of a right $A$-module.

Indeed, for every $a\in A$, let $h_{a,M}$, or simply $h_a$,
denote the endomorphism
\[
  x\longmapsto xa
\]
of the left $A$-module $M$.  By hypothesis, $T(h_a)$ is an
endomorphism of the $\bb Z$-module $T(M)$.  Since $T$ is
covariant and additive, for $a,b\in A$ one has
\[
  T(h_{ab})
  =
  T(h_b\circ h_a)
  =
  T(h_b)\circ T(h_a)
\]
and
\[
  T(h_{a+b})
  =
  T(h_a+h_b)
  =
  T(h_a)+T(h_b).
\]
\par\medskip\noindent Thus
\[
  (y,a)\longmapsto T(h_a)(y)
\]
defines a right $A$-module scalar multiplication on $T(M)$.
In particular, $T(A_s)$ is a right $A$-module.
\end{env}

\begin{env}[7.1.2]
\label{III.7.1.2}
Suppose that $A$ is commutative.  It follows from
\sref{III.7.1.1} that, for every $A$-module $M$, the group
$T(M)$ is naturally an $A$-module.  Moreover, if
$u:M\to N$ is an $A$-linear homomorphism, then, for every
$a\in A$,
\[
  u\circ h_{a,M}=h_{a,N}\circ u,
\]
and hence
\[
  T(u)\circ T(h_{a,M})
  =
  T(h_{a,N})\circ T(u).
\]
Consequently $T(u):T(M)\to T(N)$ is $A$-linear, and $T$ may be
regarded as an additive covariant functor from $\cat{Ab}_A$ to
itself.

More precisely, this defines a canonical equivalence between
the category of additive covariant functors
\[
  \cat{Ab}_A\longrightarrow\cat{Ab}
\]
and the category of $A$-linear covariant functors
\[
  T:\cat{Ab}_A\longrightarrow\cat{Ab}_A,
\]
that is, functors satisfying
\[
  T(h_{a,M})=h_{a,T(M)}
  \qquad(a\in A).
\]
The forgetful functor
\[
  I:\cat{Ab}_A\longrightarrow\cat{Ab},
\]
which sends an $A$-module to its underlying $\bb Z$-module, is
exact and faithful.  The two functors corresponding under the
preceding equivalence therefore have the same exactness
properties.
\end{env}

\begin{env}[7.1.3]
\label{III.7.1.3}
Continue to suppose that $A$ is commutative.  Let $B$ be an
$A$-algebra, not necessarily commutative, and let
\[
  \rho:A\longrightarrow B
\]
be the ring homomorphism defining the algebra structure.  It
defines an
\oldpage[III]{176}
additive covariant functor
\[
  \rho_*:M\longmapsto M_{[\rho]}
\]
from the category $\cat{Ab}_B$ of left $B$-modules to
$\cat{Ab}_A$.  By composition one obtains the additive
covariant functor
\[
  T_{(B)}:
  \cat{Ab}_B
  \xrightarrow{\rho_*}
  \cat{Ab}_A
  \xrightarrow{T}
  \cat{Ab}.
\]
For typographical reasons this functor will also be denoted
$T^{(B)}$ or $T\otimes_AB$, and is said to be obtained from $T$
by extension of scalars from $A$ to $B$.

If $B$ is commutative, $T_{(B)}$ may of course be regarded as a
functor from $\cat{Ab}_B$ to itself \sref{III.7.1.2}.  If $B$ is
commutative and $C$ is a $B$-algebra, then
\[
  T_{(C)}=(T_{(B)})_{(C)}.
\]
Extension of scalars is plainly functorial and additive in
$T$.  Furthermore, if $T$ commutes with direct limits or direct
sums, then so does $T_{(B)}$; and if $T$ is left exact, right
exact, or exact, respectively, the same is true of $T_{(B)}$.
Indeed, $\rho_*$ is exact and commutes with direct limits and
direct sums.
\end{env}

\begin{env}[7.1.4]
\label{III.7.1.4}
Suppose again that $A$ is commutative, and let
\[
  T:\cat{Ab}_A\longrightarrow\cat{Ab}_A
\]
be an additive covariant $A$-linear functor that commutes with
direct limits.  For every multiplicative subset $S$ of $A$ and
every $A$-module $M$, there is a canonical functorial
isomorphism of $A$-modules
\begin{env}[7.1.4.1]
\label{III.7.1.4.1}
\[
  T(S^{-1}M)
  \xrightarrow{\sim}
  S^{-1}T(M).
  \tag{7.1.4.1}\label{III.7.1.4.1.eq}
\]
\end{env}

First suppose that $S$ is the set of powers $f^n$,
$n\geq0$, of an element $f\in A$.  Then
\[
  M_f=\varinjlim M_n,
\]
where $(M_n,\vphi_{nm})$ is the direct system of $A$-modules
with $M_n=M$ and
\[
  \vphi_{nm}:z\longmapsto f^{\,n-m}z
  \qquad\text{\sref[0]{0.1.6.1}}.
\]
Formula \eqref{III.7.1.4.1.eq} follows in this case from the
hypothesis on $T$.  For arbitrary $S$, the module $S^{-1}M$ is
the direct limit of the $M_f$, for $f\in S$
\sref[0]{0.1.4.5}, and the same argument applies.

The functoriality of \eqref{III.7.1.4.1.eq} also shows that it
is an isomorphism of $S^{-1}A$-modules.  Thus, up to canonical
isomorphism, one may write
\begin{env}[7.1.4.2]
\label{III.7.1.4.2}
\[
  T_{(S^{-1}A)}(S^{-1}M)
  =
  S^{-1}T(M)
  =
  T(S^{-1}M).
  \tag{7.1.4.2}\label{III.7.1.4.2.eq}
\]
\end{env}
If $S=A-\mathfrak p$ is the complement of a prime ideal
$\mathfrak p$ of $A$, write $T_{\mathfrak p}$ instead of
$T_{(A_{\mathfrak p})}$.
\end{env}

\begin{proposition}[7.1.5]
\label{III.7.1.5}
Under the hypotheses of \sref{III.7.1.4}, suppose that
$T_{\mathfrak m}$ is left exact, respectively right exact or
exact, for every maximal ideal $\mathfrak m$ of $A$.  Then $T$
is left exact, respectively right exact or exact.
\end{proposition}

\begin{proof}
If two submodules $N$ and $P$ of an $A$-module $M$ satisfy
\[
  N_{\mathfrak m}=P_{\mathfrak m}
\]
for every maximal ideal $\mathfrak m$ of $A$, then $N=P$
(Bourbaki, \emph{Alg\`ebre commutative}, chap.~II, \S~3,
no.~3, th.~1).  Applying this fact to the kernels and images in
question proves the assertion.
\end{proof}

% BEGIN NATIVE INPUT ega3/ega3-7-2
\subsection{Characterisation of the tensor product functor}
\label{subsection:III.7.2}

\begin{env}[7.2.1]
\label{III.7.2.1}
Let $A$ be a ring, not necessarily commutative; let $M$
(respectively $N$) be a left (respectively right) $A$-module,
and let $P$ be a $\bb Z$-module.  Recall that giving a
$\bb Z$-homomorphism
\[
  v:N\otimes_A M\longrightarrow P
\]
is equivalent to giving a $\bb Z$-bilinear map
\[
  u:N\times M\longrightarrow P
\]
such that
\[
  u(ta,x)=u(t,ax)
  \qquad(a\in A,\ t\in N,\ x\in M),
\]
the two maps being related by
\[
  v(t\otimes x)=u(t,x).
\]
On the other hand, giving $u$ is equivalent to giving a
\oldpage[III]{177}
$\bb Z$-homomorphism
\[
  x\longmapsto f_x
\]
from $M$ to $\Hom_{\bb Z}(N,P)$ such that
\[
  f_{ax}(t)=f_x(ta)
  \qquad(a\in A,\ t\in N,\ x\in M),
\]
the two maps being related by $u(t,x)=f_x(t)$.
\end{env}

\begin{env}[7.2.2]
\label{III.7.2.2}
Let
\[
  T:\cat{Ab}_A\longrightarrow\cat{Ab}
\]
be an additive covariant functor.  For every left $A$-module
$M$ we shall define a canonical homomorphism of $\bb Z$-modules,
functorial in $M$,
\begin{env}[7.2.2.1]
\label{III.7.2.2.1}
\[
  t_M:T(A_s)\otimes_A M\longrightarrow T(M).
  \tag{7.2.2.1}\label{III.7.2.2.1.eq}
\]
\end{env}
By \sref{III.7.2.1}, it is enough to define a
$\bb Z$-homomorphism
\[
  x\longmapsto t'_M(x)
\]
from $M$ to $\Hom_{\bb Z}(T(A_s),T(M))$ such that
\[
  t'_M(ax)(y)=t'_M(x)(ya)
  \qquad(a\in A,\ x\in M,\ y\in T(A_s)).
\]
For this purpose, note that
$\Hom_{\bb Z}(T(A_s),T(M))$ is canonically a left $A$-module,
by virtue of the right $A$-module structure on $T(A_s)$, with
scalar multiplication defined by
\[
  (a\cdot w)(y)=w(ya)
\]
for $a\in A$, $w\in\Hom_{\bb Z}(T(A_s),T(M))$, and
$y\in T(A_s)$.  Define $t'_M$ to be the composite of the two
canonical homomorphisms
\[
  M
  \xrightarrow{\sim}
  \Hom_A(A_s,M)
  \xrightarrow{T}
  \Hom_{\bb Z}(T(A_s),T(M)).
\]

The second arrow sends $w$ to $T(w)$; the first is the canonical
isomorphism of $A$-modules sending $x$ to the map
$\theta_x:A_s\to M$ defined by
\[
  \theta_x(\xi)=\xi x
  \qquad(\xi\in A,\ x\in M).
\]
Since $\theta_{ax}=\theta_x\circ h_a$, one has
\[
  T(\theta_{ax})
  =
  T(\theta_x)\circ T(h_a),
\]
and consequently, for $y\in T(A_s)$,
\[
  T(\theta_{ax})(y)
  =
  T(\theta_x)\bigl(T(h_a)(y)\bigr)
  =
  T(\theta_x)(ya)
\]
by the definition of the right $A$-module structure on
$T(A_s)$.  This proves the existence of $t_M$.

It is immediate that this homomorphism is functorial in $M$:
for every homomorphism $w:M\to M'$ of left $A$-modules, the
diagram
\begin{env}[7.2.2.2]
\label{III.7.2.2.2}
\[
\xymatrix{
  T(A_s)\otimes_A M
    \ar[r]^-{t_M}\ar[d]_{1\otimes w}
  & T(M)\ar[d]^{T(w)}
  \\
  T(A_s)\otimes_A M'
    \ar[r]_-{t_{M'}}
  & T(M')
}
\tag{7.2.2.2}\label{III.7.2.2.2.eq}
\]
\end{env}
is commutative.  The functoriality of
\eqref{III.7.2.2.1.eq} also shows that, when $A$ is
commutative, $t_M$ is a homomorphism of $A$-modules
\sref{III.7.1.2}.
\end{env}

\begin{env}[7.2.3]
\label{III.7.2.3}
When $A$ is commutative, one can more generally define a
canonical homomorphism of $A$-modules
\begin{env}[7.2.3.1]
\label{III.7.2.3.1}
\[
  T(N)\otimes_A M
  \longrightarrow
  T(N\otimes_A M).
  \tag{7.2.3.1}\label{III.7.2.3.1.eq}
\]
\end{env}
for every $A$-module $N$: in the construction of
\sref{III.7.2.2}, replace $\theta_x$ by the homomorphism of
$A$-modules
\[
  N\longrightarrow N\otimes_A M,
  \qquad
  y\longmapsto y\otimes x.
\]
It is immediate that this homomorphism is functorial in both
$M$ and $N$.

\oldpage[III]{178}
In particular, if $B$ is an $A$-algebra, not necessarily
commutative, there is a homomorphism, functorial in $M$,
\begin{env}[7.2.3.2]
\label{III.7.2.3.2}
\[
  (T(M))_{(B)}
  =
  T(M)\otimes_A B
  \longrightarrow
  T(M\otimes_A B)
  =
  T_{(B)}(M_{(B)}).
  \tag{7.2.3.2}\label{III.7.2.3.2.eq}
\]
\end{env}
By the functoriality of \eqref{III.7.2.3.1.eq} in $M$, this is
a homomorphism of $B$-modules.

There is moreover a commutative diagram
\begin{env}[7.2.3.3]
\label{III.7.2.3.3}
\[
\xymatrix{
  T(A)\otimes_A M
    \ar[r]^-{t_M}\ar[d]
  & T(M)\ar[d]
  \\
  T_{(B)}(B_s)\otimes_B M_{(B)}
    \ar[r]_-{t_{M_{(B)}}}
  & T_{(B)}(M_{(B)})
}
\tag{7.2.3.3}\label{III.7.2.3.3.eq}
\]
\end{env}
The right vertical arrow is the composite
\[
  T(M)
  \longrightarrow
  T(M)\otimes_A B
  \longrightarrow
  T(M\otimes_A B)
  =
  T_{(B)}(M_{(B)})
\]
of the canonical homomorphism and
\eqref{III.7.2.3.2.eq}.  The left vertical arrow in
\eqref{III.7.2.3.3.eq} is the homomorphism
\[
  T(A)\otimes_A M
  \longrightarrow
  T_{(B)}(B_s)\otimes_B(B\otimes_A M)
  =
  T_{(B)}(B_s)\otimes_A M,
\]
where
\[
  T(A)\longrightarrow T_{(B)}(B_s)=T(B)
\]
is $T(\rho)$, with $\rho$ regarded as the homomorphism of
$A$-modules $A\to B_{[\rho]}$.
\end{env}

\begin{lemma}[7.2.4]
\label{III.7.2.4}
If $T$ is an additive covariant functor from $\cat{Ab}_A$ to
$\cat{Ab}$ that commutes with direct sums, then the canonical
homomorphism $t_L$ of \eqref{III.7.2.2.1.eq} is an isomorphism
for every free $A$-module $L$.
\end{lemma}

\begin{proof}
Write
\[
  L=\bigoplus_{\alpha\in I}L_\alpha,
\]
where every $L_\alpha$ is isomorphic to $A_s$.  The definition
of $t_M$ in \sref{III.7.2.2} shows that
\[
  t_L=\bigoplus_{\alpha\in I}t_{L_\alpha},
\]
because
\[
  T:
  \Hom_A(A_s,L)
  \longrightarrow
  \Hom_{\bb Z}(T(A_s),T(L))
\]
is the direct sum of the $\bb Z$-linear maps
\[
  T_\alpha:
  \Hom_A(A_s,L_\alpha)
  \longrightarrow
  \Hom_{\bb Z}(T(A_s),T(L_\alpha)),
\]
by the hypothesis on $T$.  It therefore remains only to prove
the lemma for $L=A_s$; but then $t_L$ is precisely the canonical
isomorphism
\[
  T(A_s)\otimes_A A_s
  \xrightarrow{\sim}
  T(A_s),
\]
valid for every right $A$-module.
\end{proof}

\begin{proposition}[7.2.5]
\label{III.7.2.5}
Let $T$ be an additive covariant functor from $\cat{Ab}_A$ to
$\cat{Ab}$ that commutes with direct sums.  The following
conditions are equivalent:
\begin{enumerate}
  \item[(a)] $T$ is right exact;
  \item[(b)] the canonical homomorphism $t_M$ of
    \eqref{III.7.2.2.1.eq} is an isomorphism for every left
    $A$-module $M$;
  \item[($\mathrm{b}'$)] $T$ is semi-exact and $t_M$ is
    surjective for every left $A$-module $M$;
  \item[(c)] $T$ is isomorphic to a functor of $M$ of the form
    $N\otimes_A M$, where $N$ is a right $A$-module.
\end{enumerate}
\end{proposition}

\begin{proof}
It is clear that (b) implies (c), and that (c) implies (a).
Let us prove that (a) implies (b).
\oldpage[III]{179}
Put
\[
  T'(M)=T(A_s)\otimes_A M
\]
for every left $A$-module $M$.  There is an exact sequence
\[
  L'\longrightarrow L\longrightarrow M\longrightarrow0,
\]
where $L$ and $L'$ are free left $A$-modules.  Since $T$ and
$T'$ are right exact, there is a commutative diagram
\[
\xymatrix{
  T'(L')\ar[r]\ar[d]^{t_{L'}}
  & T'(L)\ar[r]\ar[d]^{t_L}
  & T'(M)\ar[r]\ar[d]^{t_M}
  & 0
  \\
  T(L')\ar[r]
  & T(L)\ar[r]
  & T(M)\ar[r]
  & 0
}
\]
whose rows are exact.  By \sref{III.7.2.4}, $t_L$ and
$t_{L'}$ are isomorphisms; hence so is $t_M$, by the five
lemma.

It is clear that (b) implies ($\mathrm{b}'$).  To show that
($\mathrm{b}'$) implies (a), it is enough to prove the following
lemma.
\end{proof}

\begin{lemma}[7.2.5.1]
\label{III.7.2.5.1}
Let $\cat{K}$ and $\cat{K}'$ be abelian categories, let $F$ and
$G$ be additive covariant functors from $\cat{K}$ to
$\cat{K}'$, and let
\[
  f:F\longrightarrow G
\]
be a natural transformation \emph{(T, I, 1.2)} such that
\[
  f_E:F(E)\longrightarrow G(E)
\]
is an epimorphism for every object $E$ of $\cat{K}$.  If $F$ is
right exact and $G$ is semi-exact, then $G$ is right exact.
\end{lemma}

\begin{proof}
It is enough to show that, for every epimorphism
$v:E'\to E$ in $\cat{K}$,
\[
  G(v):G(E')\longrightarrow G(E)
\]
is an epimorphism.  There is a commutative diagram
\[
\xymatrix{
  F(E')\ar[r]^{F(v)}\ar[d]_{f_{E'}}
  & F(E)\ar[d]^{f_E}
  \\
  G(E')\ar[r]_{G(v)}
  & G(E)
}
\]
Here $F(v)$, $f_{E'}$, and $f_E$ are epimorphisms; consequently
$G(v)$ is also an epimorphism.
\end{proof}

\begin{remark}[7.2.6]
\label{III.7.2.6}
For every right $A$-module $N$, put
\[
  T_N(M)=N\otimes_A M
\]
\par\medskip\noindent
for every left $A$-module $M$.  Thus $T_N$ is an additive
covariant functor from $\cat{Ab}_A$ to $\cat{Ab}$, is right
exact, and commutes with direct sums.  If $T_N(A_s)$ is
canonically identified with $N$, the corresponding homomorphism
\eqref{III.7.2.2.1.eq} is immediately seen to be the identity.
It follows that the right $A$-module $N$ in
\sref{III.7.2.5}(c) is determined up to unique isomorphism and
is canonically isomorphic to $T(A_s)$.

Equivalently, the assignments
\[
  T\longmapsto T(A_s)
  \qquad\text{and}\qquad
  N\longmapsto T_N
\]
define an equivalence \emph{(T, I, 1.2)} between the category
of right $A$-modules and the category of additive covariant
functors
\[
  \cat{Ab}_A\longrightarrow\cat{Ab}
\]
that are right exact and commute with direct sums.
\end{remark}

\begin{proposition}[7.2.7]
\label{III.7.2.7}
Let $A$ be a left Artinian ring whose quotient by its radical
$\mathfrak m$ is a field $k$, and let $T$ be an additive
covariant functor from $\cat{Ab}_A$ to $\cat{Ab}$ that commutes
with direct sums.  The conditions of \sref{III.7.2.5} are then
also equivalent to the following condition:
\begin{enumerate}
  \item[(d)] $T$ is semi-exact and the homomorphism
    \[
      T(\varepsilon):T(A_s)\longrightarrow T(k)
    \]
    induced by the canonical homomorphism
    $\varepsilon:A_s\to k$ is surjective.
\end{enumerate}
\end{proposition}

\ifdefined\EGAIIIStopAfterPageOneSevenNine
  \expandafter\endinput
\fi

\begin{proof}
\oldpage[III]{180}
It is clear that condition ($\mathrm{b}'$) of
\sref{III.7.2.5} implies (d).  We prove that (d) implies
($\mathrm{b}'$).  There is an integer $n$ such that
$\mathfrak m^n=0$.  For every $A$-module $M$, put
\[
  M_h=\mathfrak m^hM.
\]
We shall prove, by descending induction on $h$, that $t_{M_h}$
is surjective.  The assertion is clear for $h=n$.  If $h<n$,
there is an exact sequence
\[
  0\longrightarrow M_{h+1}
  \longrightarrow M_h
  \longrightarrow M_h/M_{h+1}
  \longrightarrow0,
\]
and the induction hypothesis says that $t_{M_{h+1}}$ is
surjective.  Moreover, $M_h/M_{h+1}$ is annihilated by
$\mathfrak m$, and is therefore an $(A/\mathfrak m)$-module;
in other words, it is a direct sum of $A$-modules isomorphic to
$k$.  Since $T$ commutes with direct sums, it is therefore
enough, in order to prove that $t_{M_h/M_{h+1}}$ is surjective,
to prove that $t_k$ is surjective.  The diagram
\[
\xymatrix{
  T(A_s)\otimes_A A_s
    \ar[r]^-{t_{A_s}}\ar[d]_{1\otimes\varepsilon}
  & T(A_s)\ar[d]^{T(\varepsilon)}
  \\
  T(A_s)\otimes_A k
    \ar[r]_-{t_k}
  & T(k)
}
\]
is commutative.  Condition (d), together with
\sref{III.7.2.4}, therefore shows that $t_k$ is surjective.

To finish the proof, it remains to show that if
\[
  0\longrightarrow M'
  \longrightarrow M
  \xrightarrow{v}M''
  \longrightarrow0
\]
is an exact sequence of $A$-modules for which $t_{M'}$ and
$t_{M''}$ are surjective, then $t_M$ is surjective.  There is a
commutative diagram
\[
\xymatrix{
  T'(M')\ar[r]\ar[d]_{t_{M'}}
  & T'(M)\ar[r]\ar[d]_{t_M}
  & T'(M'')\ar[r]\ar[d]_{t_{M''}}
  & 0\ar[d]
  \\
  T(M')\ar[r]
  & T(M)\ar[r]_{T(v)}
  & T(M'')\ar[r]
  & \Coker(T(v))
}
\]
whose two rows are exact, by the hypothesis that $T$ is
semi-exact.  The induction hypothesis says that $t_{M'}$ and
$t_{M''}$ are epimorphisms, while the last vertical arrow is a
monomorphism.  The five lemma \emph{(M, I, 1.1)} therefore
shows that $t_M$ is an epimorphism.
\end{proof}
% END NATIVE INPUT ega3/ega3-7-2
\ifdefined\EGAIIIStopAfterPageOneSevenNine
  \expandafter\endinput
\fi
% BEGIN NATIVE INPUT ega3/ega3-7-3
\subsection{Exactness criteria for homological functors of modules}
\label{subsection:III.7.3}

\begin{proposition}[7.3.1]
\label{III.7.3.1}
Let $A$ be a ring, not necessarily commutative, and let
$T_\bullet$ be a covariant homological functor
\emph{(T, II, 2.1)} from $\cat{Ab}_A$ to $\cat{Ab}$ that
commutes with direct sums.  Let $p$ be an integer such that
$T_p$ and $T_{p-1}$ are defined.  The following conditions are
equivalent:
\begin{enumerate}
  \item[(a)] $T_p$ is right exact;
  \item[(b)] $T_{p-1}$ is left exact;

\oldpage[III]{181}

  \item[(c)] for every left $A$-module $M$, the canonical
    functorial homomorphism \eqref{III.7.2.2.1.eq}
    \[
      T_p(A_s)\otimes_A M
      \longrightarrow
      T_p(M)
      \tag{7.3.1.1}\label{III.7.3.1.1.eq}
    \]
    is an isomorphism;
  \item[(d)] for every left $A$-module $M$, the homomorphism
    \eqref{III.7.3.1.1.eq} is an epimorphism;
  \item[(e)] $T_p$ is isomorphic to a functor of the form
    \[
      M\longmapsto N\otimes_A M,
    \]
    where $N$ is a right $A$-module.
\end{enumerate}
If, in addition, the conditions of \sref{III.7.2.7} on $A$ and
$\mathfrak m$ are satisfied, the preceding conditions are also
equivalent to
\begin{enumerate}
  \item[(f)] the canonical homomorphism
    \[
      T_p(\varepsilon):T_p(A_s)\longrightarrow T_p(k)
    \]
    is an epimorphism.
\end{enumerate}
\end{proposition}

\begin{proof}
By definition of a homological functor, $T_i$ is semi-exact for
every $i$ for which it is defined; moreover, for every exact
sequence
\[
  0\longrightarrow M'
  \xrightarrow{u}M
  \xrightarrow{v}M''
  \longrightarrow0,
\]
one has
\[
  \Ker\bigl(T_{i-1}(u)\bigr)
  =
  \Coker\bigl(T_i(v)\bigr).
\]
It is therefore clear that (a) and (b) are equivalent, and the
other assertions follow immediately from \sref{III.7.2.5} and
\sref{III.7.2.7}.
\end{proof}

\begin{corollary}[7.3.2]
\label{III.7.3.2}
Let $A$ be a commutative ring.  With the notation of
\sref{III.7.3.1}, suppose that $T_p$ is right exact.  If
$f\in A$ annihilates no nonzero element of an $A$-module $M$,
then it annihilates no nonzero element of $T_{p-1}(M)$.  In
particular, if $A$ is an integral domain, the $A$-module
$T_{p-1}(A)$ is torsion-free.
\end{corollary}

\begin{proof}
If $h_f$ denotes the homothety $x\mapsto fx$ of $M$, the
hypothesis says that $h_f$ is injective.  Condition (b) of
\sref{III.7.3.1} then shows that $T_{p-1}(h_f)$ is injective
as well.
\end{proof}

\begin{proposition}[7.3.3]
\label{III.7.3.3}
Let $A$ be a ring, and let $T_\bullet$ be a covariant
homological functor from $\cat{Ab}_A$ to $\cat{Ab}$ that
commutes with direct sums.  Let $p$ be an integer such that
$T_{p-1}$, $T_p$, and $T_{p+1}$ are defined.  The following
conditions are equivalent:
\begin{enumerate}
  \item[(a)] $T_p$ is exact;
  \item[(b)] $T_{p+1}$ and $T_p$ are right exact;
  \item[(c)] $T_p$ and $T_{p-1}$ are left exact;
  \item[(d)] $T_{p+1}$ is right exact and $T_{p-1}$ is left
    exact;
  \item[(e)] for every $A$-module $M$, the canonical
    homomorphisms
    \[
      T_i(A_s)\otimes_A M
      \longrightarrow
      T_i(M)
      \tag{7.3.3.1}\label{III.7.3.3.1.eq}
    \]
    are isomorphisms for $i=p$ and $i=p+1$;
  \item[($\mathrm{e}'$)] for every $A$-module $M$, the
    homomorphisms \eqref{III.7.3.3.1.eq} are epimorphisms for
    $i=p$ and $i=p+1$;
  \item[(f)] for every $A$-module $M$, the homomorphism
    \eqref{III.7.3.3.1.eq} is an isomorphism for $i=p$, and
    $T_p(A_s)$ is flat as a right $A$-module;
  \item[($\mathrm{f}'$)] for every $A$-module $M$, the
    homomorphism \eqref{III.7.3.3.1.eq} is an epimorphism for
    $i=p$, and $T_p(A_s)$ is flat as a right $A$-module.
\end{enumerate}
\end{proposition}

\begin{proof}
The equivalence of (a), (b), (c), and (d) follows from the
equivalence of conditions (a) and (b) of \sref{III.7.3.1}.
The equivalence of (b), (e), and ($\mathrm{e}'$) follows from
the equivalence of (a), (c), and (d) in \sref{III.7.3.1}.
Finally, to say that $T_p(A_s)$ is flat means that the functor
\[
  M\longmapsto T_p(A_s)\otimes_A M
\]
is left exact; the equivalence of (a), (f), and
($\mathrm{f}'$) again follows from the equivalence of (a), (c),
and (d) in \sref{III.7.3.1}.
\end{proof}

\oldpage[III]{182}

\begin{corollary}[7.3.4]
\label{III.7.3.4}
Suppose that $A$ is commutative, that $T_p$ is exact, and that
$T_p(A)$ is an $A$-module of finite presentation.  Then the
function
\[
  x\longmapsto
  \operatorname{rank}_{\kres(x)}
  \bigl(T_p(\kres(x))\bigr)
\]
is locally constant on $X=\Spec(A)$, and is therefore constant
if $\Spec(A)$ is connected.
\end{corollary}

\begin{proof}
By condition (f) of \sref{III.7.3.3}, $T_p(A)$ is a flat
$A$-module.  It is consequently projective of finite type, so
$\widetilde{T_p(A)}$ is a locally free $\sh O_X$-module
(Bourbaki, \emph{Alg\`ebre commutative}, chap.~II, \S~5,
no.~2, th.~1).  Moreover,
\[
  T_p(\kres(x))
  =
  T_p(A)\otimes_A\kres(x)
\]
by condition (e) of \sref{III.7.3.3}, and the rank at $x$ of
the $A$-module $T_p(A)$ is locally constant
\emph{(loc.\ cit.)}.  This proves the corollary.
\end{proof}

\begin{proposition}[7.3.5]
\label{III.7.3.5}
Suppose that $A$ is a left Artinian ring whose quotient by its
radical $\mathfrak m$ is a field $k$.  The conditions of
\sref{III.7.3.3} are also equivalent to each of the following:
\begin{enumerate}
  \item[(g)] the canonical homomorphism
    \[
      T_i(\varepsilon):T_i(A_s)\longrightarrow T_i(k)
    \]
    is an epimorphism for $i=p$ and $i=p+1$;
  \item[(h)] $T_p(\varepsilon)$ is an epimorphism and
    $T_p(A_s)$ is flat as a right $A$-module, or, equivalently,
    is a free right $A$-module (Bourbaki,
    \emph{Alg\`ebre commutative}, chap.~II, \S~3, no.~2,
    cor.~2 of prop.~5).
\end{enumerate}
Suppose in addition that $A$ is commutative and that the
$A$-module $T_p(k)$ has finite length $d$.  The preceding
conditions are then also equivalent to each of the following:
\begin{enumerate}
  \item[(i)] for every $A$-module $M$ of finite length,
    \[
      \operatorname{length}\bigl(T_p(M)\bigr)
      =
      d\,\operatorname{length}(M);
      \tag{7.3.5.1}\label{III.7.3.5.1.eq}
    \]
  \item[(j)]
    \[
      \operatorname{length}\bigl(T_p(A)\bigr)
      =
      d\,\operatorname{length}(A).
      \tag{7.3.5.2}\label{III.7.3.5.2.eq}
    \]
\end{enumerate}
\end{proposition}

The equivalence of (g) and (h) with the conditions of
\sref{III.7.3.3} follows at once from \sref{III.7.2.7}.  For
the remaining assertions we shall use the following lemma.

\begin{lemma}[7.3.5.3]
\label{III.7.3.5.3}
Let $\cat{K}$ and $\cat{K}'$ be abelian categories, and let
\[
  F:\cat{K}\longrightarrow\cat{K}'
\]
be an additive covariant functor.  Suppose that $F$ is
semi-exact and that $F(S)$ has finite length in $\cat{K}'$ for
every simple object $S$ of $\cat{K}$.  Then $F(E)$ has finite
length in $\cat{K}'$ for every object $E$ of finite length in
$\cat{K}$.  For every exact sequence
\[
  0\longrightarrow E'
  \xrightarrow{u}E
  \xrightarrow{v}E''
  \longrightarrow0
\]
of objects of finite length in $\cat{K}$, one has
\[
  \operatorname{length}F(E)
  \leq
  \operatorname{length}F(E')
  +
  \operatorname{length}F(E''),
  \tag{7.3.5.4}\label{III.7.3.5.4.eq}
\]
and equality in \eqref{III.7.3.5.4.eq} holds if and only if
\[
  0\longrightarrow F(E')
  \longrightarrow F(E)
  \longrightarrow F(E'')
  \longrightarrow0
\]
is exact.
\end{lemma}

\begin{proof}
The sequence
\[
  F(E')\xrightarrow{F(u)}F(E)\xrightarrow{F(v)}F(E'')
\]
is exact by hypothesis.  If $F(E')$ and $F(E'')$ have finite
length, so do $\Im(F(u))$ and $\Im(F(v))$.  Since
$\Ker(F(v))=\Im(F(u))$, the object $F(E)$ has finite length,
and
\[
  \begin{split}
  \operatorname{length}F(E)
  &=
  \operatorname{length}\Im(F(u))
  +
  \operatorname{length}\Im(F(v))
  \\
  &\leq
  \operatorname{length}F(E')
  +
  \operatorname{length}F(E'').
  \end{split}
  \tag{7.3.5.5}\label{III.7.3.5.5.eq}
\]

\oldpage[III]{183}

Induction on the length of $E$ proves the first assertion.
Moreover, equality in \eqref{III.7.3.5.5.eq} can hold only when
\[
  \operatorname{length}\Im(F(u))
  =
  \operatorname{length}F(E')
\]
and
\[
  \operatorname{length}\Im(F(v))
  =
  \operatorname{length}F(E'').
\]
The first equality is equivalent to
$\operatorname{length}\Ker(F(u))=0$, hence to $\Ker(F(u))=0$;
the second is equivalent to
$\operatorname{length}\Coker(F(v))=0$, hence to
$\Coker(F(v))=0$.  This proves the last assertion.
\end{proof}

Now let $M$ be an $A$-module of finite length, with $A$
commutative.  The quotients in a Jordan--H\"older series of
$M$ are necessarily isomorphic to the $A$-module $k$.
Induction on the length of $M$, using
\eqref{III.7.3.5.4.eq}, therefore gives
\[
  \operatorname{length}T_p(M)
  \leq
  d\,\operatorname{length}(M).
  \tag{7.3.5.6}\label{III.7.3.5.6.eq}
\]
It also follows from \sref{III.7.3.5.3} that if $T_p$ is exact,
then equality \eqref{III.7.3.5.1.eq} holds.  Thus condition
(a) of \sref{III.7.3.3} implies (i), and (i) plainly implies
(j).  It remains to prove the following lemma.

\begin{lemma}[7.3.5.7]
\label{III.7.3.5.7}
The equality
\[
  \operatorname{length}T_p(A)
  =
  d\,\operatorname{length}A
\]
implies that $T_p(\varepsilon)$ is an epimorphism and that
$T_p(A)$ is a flat $A$-module.
\end{lemma}

\begin{proof}
Starting from the exact sequence
\[
  0\longrightarrow\mathfrak m
  \longrightarrow A
  \longrightarrow k
  \longrightarrow0,
\]
formulas \eqref{III.7.3.5.4.eq} and
\eqref{III.7.3.5.6.eq} give
\[
  \begin{split}
  \operatorname{length}T_p(A)
  &\leq
  \operatorname{length}T_p(\mathfrak m)
  +
  \operatorname{length}T_p(k)
  \\
  &\leq
  d\bigl(
    \operatorname{length}\mathfrak m
    +
    \operatorname{length}k
  \bigr)
  =
  d\,\operatorname{length}A.
  \end{split}
\]
By \sref{III.7.3.5.3}, equality can hold only if
\[
  0\longrightarrow T_p(\mathfrak m)
  \longrightarrow T_p(A)
  \longrightarrow T_p(k)
  \longrightarrow0
  \tag{7.3.5.8}\label{III.7.3.5.8.eq}
\]
is exact.  By \sref{III.7.2.7} and \sref{III.7.2.5}, $T_p$ is
isomorphic to a functor
\[
  M\longmapsto N\otimes_A M.
\]
The exactness of \eqref{III.7.3.5.8.eq}, together with the exact
sequence of $\Tor$, shows that
\[
  \Tor^A_1(N,k)=0.
\]
It follows that $N=T_p(A)$ is a flat $A$-module
\emph{(0, 10.1.3)}.
\end{proof}

\begin{lemma}[7.3.6]
\label{III.7.3.6}
Let $A$ be a ring, and let $T_\bullet$ be a covariant
homological functor from $\cat{Ab}_A$ to $\cat{Ab}$ that
commutes with direct sums.  Suppose that $T_p$ and $T_{p+1}$
are defined and that $T_p$ is left exact.  Then $T_{p+1}$ is
exact if and only if $T_{p+1}(A_s)$ is flat as a right
$A$-module.
\end{lemma}

\begin{proof}
By \sref{III.7.3.1}, the canonical homomorphism
\[
  T_{p+1}(A_s)\otimes_A M
  \longrightarrow
  T_{p+1}(M)
\]
is an isomorphism of functors.  The assertion now follows from
the definition of a flat $A$-module.
\end{proof}

\begin{proposition}[7.3.7]
\label{III.7.3.7}
Let $A$ be a ring, and let $T_\bullet$ be a covariant
homological functor from $\cat{Ab}_A$ to $\cat{Ab}$ that
commutes with direct sums.  Suppose there is an $i_0$ such that
$T_i$ is exact for $i\leq i_0$.  Then, for every integer
$p>i_0$, the following conditions are equivalent:
\begin{enumerate}
  \item[(a)] $T_q$ is exact for $q\leq p$;
  \item[(b)] $T_q(A_s)$ is flat as a right $A$-module for
    $q\leq p$;
  \item[(c)] for every $A$-module $M$, the canonical
    homomorphism
    \[
      T_q(A_s)\otimes_A M
      \longrightarrow
      T_q(M)
    \]
    is surjective for $q\leq p+1$.
\end{enumerate}
\end{proposition}

\begin{proof}
The equivalence of (a) and (b) follows by induction on $q$ from
\sref{III.7.3.6}, since $T_{i_0}$ is exact by hypothesis.  The
equivalence of (a) and (c) follows from the equivalence of
conditions (a) and ($\mathrm{e}'$) in \sref{III.7.3.3}.
\end{proof}

\oldpage[III]{184}

\begin{env}[7.3.8]
\label{III.7.3.8}
Let $A$ be a commutative ring, let $B$ be an $A$-algebra, not
necessarily commutative, and let $T_\bullet$ be a covariant
homological functor from $\cat{Ab}_A$ to $\cat{Ab}$.  It
follows from the definitions in \sref{III.7.1.3} that the
functor from $\cat{Ab}_B$ to $\cat{Ab}$ obtained by extension
of scalars from $A$ to $B$, denoted
\[
  T_\bullet^{(B)}=\bigl(T_i^{(B)}\bigr),
\]
is again a homological functor.
\end{env}

\begin{corollary}[7.3.9]
\label{III.7.3.9}
Suppose that $T_\bullet$ satisfies the general conditions of
\sref{III.7.3.7} and commutes with direct limits.  Suppose
further that $A$ is an integral domain and that every $T_n(A)$
is an $A$-module of finite presentation.  Then, for every
integer $N$, there is an $f\in A-\{0\}$ such that
\[
  T_p^{(A_f)}:\cat{Ab}_{A_f}\longrightarrow\cat{Ab}
\]
is exact for $p\leq N$.
\end{corollary}

\begin{proof}
By hypothesis, $T_i$ is exact for $i\leq i_0$, so $T_i(A)$ is
flat for these values of $i$.  By condition (b) of
\sref{III.7.3.7}, it is enough to choose $f$ such that
\[
  T_i^{(A_f)}(A_f)=T_i(A_f)
\]
is a free $A_f$-module for $i_0<i\leq N$.  Now
\[
  T_i(A_f)=\bigl(T_i(A)\bigr)_f,
\]
since $T_i$ commutes with direct limits \sref{III.7.1.4}.  If
$x$ is the generic point of $\Spec(A)$, then
$\bigl(T_i(A)\bigr)_x$ is a finite-dimensional vector space
over the fraction field of $A$.  Since each $T_i(A)$ is of
finite presentation, an $f$ with the desired property exists
(Bourbaki, \emph{Alg\`ebre commutative}, chap.~II, \S~5,
no.~1, cor.\ of prop.~2).

If only finitely many indices $i$ satisfy $T_i\neq0$, one may
choose $f\in A-\{0\}$ such that every $T_p^{(A_f)}$ is exact.
\end{proof}

\begin{corollary}[7.3.10]
\label{III.7.3.10}
Suppose that $T_\bullet$ satisfies the general conditions of
\sref{III.7.3.7} and commutes with direct limits, that $A$ is
commutative and Noetherian, and that every $T_n(A)$ is an
$A$-module of finite type.  Then, for every integer $N$, there
is a dense open subset $U$ of $\Spec(A)$ such that, for every
integer $p\leq N$, the function
\[
  x\longmapsto
  \operatorname{rank}_{\kres(x)}
  \bigl(T_p(\kres(x))\bigr)
\]
is constant on $U$.
\end{corollary}

\begin{proof}
Let $\mathfrak p$ be a minimal prime ideal of $A$.  The ring
$B=A/\mathfrak p$ is an integral domain, and $\Spec(B)$ is
identified with an irreducible component of the topological
space $\Spec(A)$.  We prove by induction on $p\leq N$ that
there is an $f_p\in B-\{0\}$ such that, if
$B'=B_{f_p}$, then $T_i^{(B')}$ is exact and $T_i(B')$ is a
$B'$-module of finite type for $i\leq p$.

The assertion holds for $p\leq i_0$ by hypothesis, with
$f_p=1$, so that $B'=B=A/\mathfrak p$.  Indeed, $T_p$ is then
exact, and $T_p(B)$ is isomorphic to
\[
  T_p(A)/T_p(\mathfrak p);
\]
it is therefore an $A$-module, and \emph{a fortiori} a
$B$-module, of finite type.

For the induction step, let $f_p$ be the canonical image in
$B$ of an element $g_p\in A$, and put $A'=A_{g_p}$.  Then
\[
  B'=A'/\mathfrak p',
\]
where $\mathfrak p'$ is the minimal prime ideal
$\mathfrak pA_{g_p}$ of $A'$.  Since
\[
  T_i(A_{g_p})=\bigl(T_i(A)\bigr)_{g_p},
\]
the $T_i(A')$ are $A'$-modules of finite type.  Thus
$T_\bullet^{(A')}$ satisfies the same hypotheses as
$T_\bullet$, with $i_0$ replaced by $p$.

We may therefore reduce to the case $A'=A$ in which $T_p$ is
exact.  The exact sequence
\[
  0\longrightarrow\mathfrak p
  \longrightarrow A
  \longrightarrow A/\mathfrak p
  \longrightarrow0
\]
then gives an exact sequence
\[
  T_{p+1}(A)
  \longrightarrow
  T_{p+1}(A/\mathfrak p)
  \xrightarrow{\partial}
  T_p(\mathfrak p)
  \longrightarrow
  T_p(A).
\]
Since $T_p$ is exact, the last arrow is injective.
Consequently $T_{p+1}(A/\mathfrak p)$ is a quotient of
$T_{p+1}(A)$, and is therefore of finite type.

The proof of \sref{III.7.3.9} used the finite-type hypothesis
on $T_i(A)$ only for $i\leq N$.  We may therefore apply that
corollary to the integral domain $B$, the functor
$T_\bullet^{(B)}$, and $N=p+1$, which completes the induction.
There is consequently an $f_N\in B-\{0\}$ such that
$\Spec(B_{f_N})$ is an open subset $V$, dense in $\Spec(B)$,
that meets no other irreducible component of $\Spec(A)$.

\oldpage[III]{185}

If the assertion has been proved for $B_{f_N}$, there is a
dense open subset $W$ of $V$ on which the functions in the
statement are constant, since
\[
  A_x=(B_{f_N})_x
  \qquad(x\in W).
\]
Repeating the same argument for every irreducible component of
$\Spec(A)$ proves the corollary.  We may therefore reduce to
the case where $A$ is an integral domain.  The argument of
\sref{III.7.3.9} then gives an $f\in A-\{0\}$ such that the
$T_p(A_f)$ are free $A_f$-modules of finite type for
$p\leq N$.  The conclusion follows from \sref{III.7.3.4}.
\end{proof}

\begin{proposition}[7.3.11]
\label{III.7.3.11}
Let $A$ be a commutative local ring with residue field $k$,
and let $T_\bullet$ be a covariant homological functor from
$\cat{Ab}_A$ to $\cat{Ab}$ that commutes with direct sums.
Suppose there is an $i_0$ such that $T_i$ is exact for
$i\leq i_0$, and that every $T_n(A)$ is an $A$-module of
finite presentation.  The equivalent conditions (a), (b), and
(c) of \sref{III.7.3.7} imply the following two conditions,
and are equivalent to them when $A$ is also reduced:
\begin{enumerate}
  \item[(d)] for every $x\in\Spec(A)$,
    \[
      \operatorname{rank}_{\kres(x)}
      T_q(\kres(x))
      =
      \operatorname{rank}_k T_q(k)
      \qquad(q\leq p);
    \]
  \item[($\mathrm{d}'$)] for every generic point $x_j$ of an
    irreducible component of $\Spec(A)$,
    \[
      \operatorname{rank}_{\kres(x_j)}
      T_q(\kres(x_j))
      =
      \operatorname{rank}_k T_q(k)
      \qquad(q\leq p).
    \]
\end{enumerate}
\end{proposition}

\begin{proof}
Since $T_q(A)$ is an $A$-module of finite presentation,
condition (b) of \sref{III.7.3.7} is equivalent to saying that
$T_q(A)$ is free for $q\leq p$ (Bourbaki,
\emph{Alg\`ebre commutative}, chap.~II, \S~3, no.~2, cor.~2
of prop.~5).  Condition (c) implies
\[
  T_q(\kres(x))
  =
  T_q(A)\otimes_A\kres(x)
  \qquad(q\leq p).
\]
Thus the equivalent conditions of \sref{III.7.3.7} imply (d),
and it is immediate that (d) implies ($\mathrm{d}'$).

It remains to show that ($\mathrm{d}'$) implies (a) when $A$ is
reduced.  We argue by induction on $q\leq p$, since $T_q$ is
exact for $q\leq i_0$.  Suppose that $T_k$ is exact for
$k\leq q<p$.  We shall show that $T_{q+1}(A)$ is a free
$A$-module.  By the induction hypothesis, the canonical map
\[
  T_{q+1}(A)\otimes_A M
  \longrightarrow
  T_{q+1}(M)
\]
is an isomorphism for every $A$-module $M$, by condition (c) of
\sref{III.7.3.7} and \sref{III.7.3.3}.  Applying this to
$M=\kres(x_j)$ and $M=k$, condition ($\mathrm{d}'$) gives
\[
  \operatorname{rank}_{\kres(x_j)}
  \bigl(T_{q+1}(A)\otimes_A\kres(x_j)\bigr)
  =
  \operatorname{rank}_k T_{q+1}(k)
\]
for every generic-point index $j$.  It follows that
$T_{q+1}(A)$ is free (Bourbaki,
\emph{Alg\`ebre commutative}, chap.~II, \S~3, no.~2, prop.~7),
which completes the proof.
\end{proof}

The preceding results will be considerably strengthened for
the homological functors of the special type to be studied in
\sref{III.7.4}: the resulting exactness criteria involve only
a single one of the $T_p$.
% END NATIVE INPUT ega3/ega3-7-3


% BEGIN NATIVE INPUT ega3/ega3-7-4
\subsection{\texorpdfstring{Exactness criteria for the functors
$H_\bullet(P_\bullet\otimes_A M)$}{Exactness criteria for the functors H.(P. tensor M)}}
\label{subsection:III.7.4}
\label{III.7.4}

\begin{env}[7.4.1]
\label{III.7.4.1}
Let $A$ be a ring, not necessarily commutative, and let
$P_\bullet$ be a complex of flat right $A$-modules.  Since the
functor
\[
  M\longmapsto P_k\otimes_A M
\]
is then exact in $\cat{Ab}_A$ for every $k$, the
$\partial$-functor
\[
  T_\bullet(M)=H_\bullet(P_\bullet\otimes_A M)
  \tag{7.4.1.1}\label{III.7.4.1.1.eq}
\]
\oldpage[III]{186}
is a \emph{homological functor} from $\cat{Ab}_A$ to
$\cat{Ab}$; it is plainly $A$-linear when $A$ is commutative
\sref{III.7.1.2}, and it commutes with inductive limits.

If $A$ is commutative, then, for every $A$-algebra $B$, the
homological functor $T_\bullet^{(B)}$ \sref{III.7.3.8} is, by
definition, given by
\[
  T_i^{(B)}(N)
  =
  H_i(P_\bullet\otimes_A N_{[\rho]}),
  \tag{7.4.1.2}\label{III.7.4.1.2.eq}
\]
where $\rho:A\to B$ is the homomorphism defining the
$A$-algebra structure on $B$.  Since one may also write
\[
  P_\bullet\otimes_A N_{[\rho]}
  =
  P_\bullet\otimes_A(B\otimes_B N)_{[\rho]}
  =
  (P_\bullet\otimes_A B)\otimes_B N,
\]
it follows that
\[
  T_i^{(B)}(N)
  =
  H_i(P'_\bullet\otimes_B N)
  \tag{7.4.1.3}\label{III.7.4.1.3.eq}
\]
for every $B$-module $N$, where $P'_\bullet$ is the complex
$P_\bullet\otimes_A B$ of flat $B$-modules
\sref[0]{0.6.2.1}.
\end{env}

\begin{proposition}[7.4.2]
\label{III.7.4.2}
Under the general hypotheses of \sref{III.7.4.1}, and for a
given integer $p\in\bb{Z}$, the following conditions are
equivalent:
\begin{enumerate}
  \item[(a)] $T_p$ is left exact, or, equivalently,
    $T_{p+1}$ is right exact;
  \item[(b)]
    \[
      Z'_p(P_\bullet)
      =
      \Coker(P_{p+1}\longrightarrow P_p)
    \]
    is a flat right $A$-module;
  \item[(c)] there exists a complex $P'_\bullet$ of flat right
    $A$-modules for which the differential
    \[
      d'_{p+1}:P'_{p+1}\longrightarrow P'_p
    \]
    is zero, together with an isomorphism of homological
    functors
    \[
      H_\bullet(P_\bullet\otimes_A M)
      \simeq
      H_\bullet(P'_\bullet\otimes_A M).
    \]
\end{enumerate}
\end{proposition}

\begin{proof}
By definition there is a sequence, exact and functorial in
$M$,
\[
  0
  \longrightarrow
  T_p(M)
  \longrightarrow
  Z'_p(P_\bullet)\otimes_A M
  \longrightarrow
  P_{p-1}\otimes_A M.
\]
Indeed, by right exactness of the tensor product,
\[
  Z'_p(P_\bullet\otimes_A M)
  =
  \Coker(P_{p+1}\otimes_A M\longrightarrow
         P_p\otimes_A M)
  =
  Z'_p(P_\bullet)\otimes_A M.
\]
For every homomorphism $f:M\to N$ one therefore has the
commutative diagram
\[
\xymatrix{
  0 \ar[r]
    & T_p(M) \ar[r] \ar[d]_{u}
    & Z'_p(P_\bullet)\otimes_A M \ar[r] \ar[d]_{v}
    & P_{p-1}\otimes_A M \ar[d]_{w} \\
  0 \ar[r]
    & T_p(N) \ar[r]
    & Z'_p(P_\bullet)\otimes_A N \ar[r]
    & P_{p-1}\otimes_A N,
}
\tag{7.4.2.1}\label{III.7.4.2.1.eq}
\]
whose rows are exact.  If $f$ is a monomorphism, then so is
$w$, because $P_{p-1}$ is flat.  If $T_p$ is left exact,
$u$ is likewise a monomorphism, and it follows that $v$ is a
monomorphism.  Thus $Z'_p(P_\bullet)$ is flat.  Conversely, if
$Z'_p(P_\bullet)$ is flat, then $v$ is a monomorphism for every
monomorphism $f:M\to N$, and diagram
\eqref{III.7.4.2.1.eq} shows that $u$ is a monomorphism.
Hence $T_p$, which is already half-exact, is left exact.  This
proves the equivalence of (a) and (b).

It is immediate that (c) implies (a).  Indeed, if
$d'_{p+1}:P'_{p+1}\to P'_p$ is zero and
\[
  0\longrightarrow M'\longrightarrow M\longrightarrow M''
  \longrightarrow 0
\]
is an exact sequence of $A$-modules, then the boundary map in
the exact sequence
\[
  H_{p+1}(P'_\bullet\otimes_A M'')
  \xrightarrow{\partial}
  H_p(P'_\bullet\otimes_A M')
  \longrightarrow
  H_p(P'_\bullet\otimes_A M)
\]
\oldpage[III]{187}
is zero by definition \emph{(M, IV, 1)}, so $T_p$ is left
exact.

Conversely, let us prove that (b) implies (c). Put
\[
 C=Z'_p(P_\bullet),\qquad
 B=B_p(P_\bullet)=\Im(d_{p+1}),\qquad
 Z=Z_{p+1}(P_\bullet)=\Ker(d_{p+1}).
\]
There are two short exact sequences
\[
 0\longrightarrow B\longrightarrow P_p\longrightarrow C\longrightarrow0,
 \qquad
 0\longrightarrow Z\longrightarrow P_{p+1}\longrightarrow B\longrightarrow0.
\]
Since $C$ and $P_p$ are flat, $B$ is flat; since $B$ and $P_{p+1}$
are flat, $Z$ is flat. Both sequences remain exact after tensoring with
any left $A$-module $M$ \sref[0]{0.6.1.2}. In particular, the canonical
maps identify $B\otimes_A M$ with the image of $d_{p+1}\otimes1_M$
and $Z\otimes_A M$ with its kernel.

Define a complex $U_\bullet$ by $U_i=P_i$ for $i\geq p+2$,
$U_{p+1}=Z$, and $U_i=0$ for $i\leq p$; its differential in degree
$p+2$ is $d_{p+2}$ with target restricted to $Z$, and its other nonzero
differentials are inherited from $P_\bullet$. Similarly, define
$L_\bullet$ by $L_i=P_i$ for $i\leq p-1$, $L_p=C$, and $L_i=0$
for $i\geq p+1$; its differential in degree $p$ is induced by $d_p$
on the quotient $C$.
Then $P'_\bullet=U_\bullet\oplus L_\bullet$ is a complex of flat
modules, with $P'_i=P_i$ for $i\ne p,p+1$, $P'_p=C$,
$P'_{p+1}=Z$, and $d'_{p+1}=0$.

The inclusion in degree $p+1$ defines a map of complexes
$a:U_\bullet\to P_\bullet$, and the quotient in degree $p$ defines
a map $b:P_\bullet\to L_\bullet$. After tensoring with $M$, the
first induces homology isomorphisms in every degree $i\geq p+1$:
in degree $p+1$ this follows from the kernel identification above;
in degree $p+2$, from the injectivity of
$Z\otimes_A M\to P_{p+1}\otimes_A M$; the higher degrees are
unchanged. The second induces homology isomorphisms in every degree
$i\leq p$: in degree $p$, one takes the quotient by the image
$B\otimes_A M$; in degree $p-1$, one uses the surjectivity of
$P_p\otimes_A M\to C\otimes_A M$; the lower degrees are unchanged.
Taking the inverses of the first isomorphisms and retaining the second
ones therefore gives natural isomorphisms, for every $i$,
\[
 H_i(P_\bullet\otimes_A M)\xrightarrow{\sim}
 H_i(P'_\bullet\otimes_A M).
\]
These commute with the boundary maps associated to every short exact
sequence $0\to M'\to M\to M''\to0$: in degrees $i\geq p+2$ and
$i\leq p$, this follows from naturality for the maps of complexes $a$
and $b$. In degree $p+1$, every cycle in $P_\bullet\otimes_A M''$
comes from $Z\otimes_A M''$ and lifts, through the surjection
$Z\otimes_A M\to Z\otimes_A M''$, to a cycle in
$P_\bullet\otimes_A M$, so the boundary map is zero. The boundary
map for $P'_\bullet$ is also zero because $d'_{p+1}=0$.
Thus these isomorphisms form the required isomorphism of
$\partial$-functors.

We also note that the conditions of \sref{III.7.4.2} imply
that $B_p(P_\bullet)$ is flat, since there is an exact sequence
\[
  0
  \longrightarrow B_p(P_\bullet)
  \longrightarrow P_p
  \longrightarrow Z'_p(P_\bullet)
  \longrightarrow 0
\]
in which $P_p$ and $Z'_p(P_\bullet)$ are flat
\sref[0]{0.6.1.2}.
\end{proof}

\begin{corollary}[7.4.3]
\label{III.7.4.3}
Suppose that $A$ is a regular Noetherian ring of dimension
$1$, in other words, a product of Dedekind rings
\sref[0]{0.17.1.3} and \sref[0]{0.17.3.7}; for example, $A$
may be a principal ideal ring.  Then $T_p$ is left exact if
and only if $T_p(A)$ is a flat $A$-module.  The functor $T_p$
is exact if and only if both $T_p(A)$ and $T_{p-1}(A)$ are
flat $A$-modules.
\end{corollary}

\begin{proof}
Recall that, for a module $M$ over a Dedekind ring, being flat
is equivalent to being torsion-free \sref[0]{0.6.3.3} and
\sref[0]{0.6.3.4}.  Under the hypotheses of
\sref{III.7.4.3}, every submodule of a flat $A$-module is
therefore flat.

The second assertion follows from the first, since saying that
$T_p$ is exact amounts to saying that $T_p$ and $T_{p-1}$ are
left exact.  To prove the first assertion, observe that there
is an exact sequence
\[
  0
  \longrightarrow H_p(P_\bullet)
  \longrightarrow Z'_p(P_\bullet)
  \longrightarrow B_{p-1}(P_\bullet)
  \longrightarrow 0,
\]
where $B_{p-1}(P_\bullet)$ is flat, being a submodule of the
flat $A$-module $P_{p-1}$.  Hence $H_p(P_\bullet)$ is flat if
and only if $Z'_p(P_\bullet)$ is flat
\sref[0]{0.6.1.2}.
\end{proof}

The most important applications of \sref{III.7.4.2} are the
following.

\begin{proposition}[7.4.4]
\label{III.7.4.4}
Let $A$ be a Noetherian ring and $P_\bullet$ a complex of flat
$A$-modules.  Suppose either that the $P_i$ are of finite
type, or that the $H_i(P_\bullet)$ are $A$-modules of finite
type and there exists an $i_0$ such that
$H_i(P_\bullet)=0$ for $i<i_0$.  Let $T_\bullet$ be the
homological functor defined by \eqref{III.7.4.1.1.eq}.  Then
the set $U$ of points $y\in\Spec(A)$ such that $(T_p)_y$
\sref{III.7.1.4} is right exact, respectively left exact or
exact, is open in $\Spec(A)$.
\end{proposition}

\begin{proof}
\oldpage[III]{188}
Under the second hypothesis on $P_\bullet$, one may replace
$P_\bullet$ by a complex $P'_\bullet$ of finite free
$A$-modules such that the functor
\[
  H_\bullet(P'_\bullet\otimes_A M)
\]
is isomorphic to $T_\bullet(M)$ as a $\partial$-functor
\sref[0]{0.11.9.3}.  Thus one may always reduce to the first
hypothesis.  In that case the modules $Z'_i(P_\bullet)$ are of
finite type.  Moreover, it is enough to prove the assertions
concerning left exactness (cf.\ \sref{III.7.4.2}, (a)).

Let $x\in U$.  Since $M\mapsto M_x$ is exact,
\[
  \bigl(Z'_p(P_\bullet)\bigr)_x
  =
  Z'_p\bigl((P_\bullet)_x\bigr).
\]
Taking \eqref{III.7.4.1.3.eq} into account, the hypothesis and
\sref{III.7.4.2}, (b), show that
$\bigl(Z'_p(P_\bullet)\bigr)_x$ is a flat $A_x$-module.  It is
therefore free, since it is of finite type and $A_x$ is a
Noetherian local ring (Bourbaki,
\emph{Alg\`ebre commutative}, chap.~II, \S~3, no.~2, cor.~2
of prop.~5).  It follows that there is an $f\in A$ for which
$\bigl(Z'_p(P_\bullet)\bigr)_f$ is free over $A_f$
(Bourbaki, \emph{Alg\`ebre commutative}, chap.~II, \S~5,
no.~1, cor.\ of prop.~2).  A fortiori,
$\bigl(Z'_p(P_\bullet)\bigr)_y$ is free over $A_y$ for every
$y\in D(f)$, and \sref{III.7.4.2}, (b), completes the proof.
\end{proof}

\begin{corollary}[7.4.5]
\label{III.7.4.5}
Under the hypotheses of \sref{III.7.4.4}, suppose in addition
that $A$ is an integral domain.  Then the set $U$ of
$x\in\Spec(A)$ such that $(T_p)_x$ is exact is open and
nonempty.
\end{corollary}

\begin{proof}
It is enough to show that $(T_p)_x$ is exact at the generic
point $x$ of $\Spec(A)$.  This is immediate, since $A_x$ is a
field and every additive functor on $\cat{Ab}_{A_x}$ is exact.
\end{proof}

\begin{proposition}[7.4.6]
\label{III.7.4.6}
Under the general hypotheses of \sref{III.7.4.4}, conditions
(a), (b), and (c) of \sref{III.7.4.2} are also equivalent to:
\begin{enumerate}
  \item[(d)] there exists an $A$-module $Q$ and a functorial
    isomorphism
    \[
      T_p(M)\simeq\Hom_A(Q,M).
      \tag{7.4.6.1}\label{III.7.4.6.1.eq}
    \]
\end{enumerate}
Moreover, this property determines the $A$-module $Q$ up to a
unique isomorphism, and $Q$ is of finite type.
\end{proposition}

\begin{proof}
The uniqueness of $Q$ is a special case of the uniqueness of a
representing object for a representable functor
\sref[0]{0.8.1.5}.  It is clear that the right-hand side of
\eqref{III.7.4.6.1.eq} is left exact.

Conversely, suppose that $T_p$ is left exact.  As in
\sref{III.7.4.4}, one may first reduce to the case where the
$P_i$ are flat and of finite type, and hence, since $A$ is
Noetherian, projective of finite type (Bourbaki,
\emph{Alg\`ebre commutative}, chap.~II, \S~5, no.~2, cor.\ of
th.~1).  The dual $\check{P}_i$ of $P_i$ is then also a
projective $A$-module of finite type, $P_i$ is canonically
isomorphic to the dual of $\check{P}_i$, and the canonical
homomorphism
\[
  P_i\otimes_A M
  \longrightarrow
  \Hom_A(\check{P}_i,M)
\]
is bijective (Bourbaki, \emph{Alg\`ebre}, chap.~II,
3rd ed., \S~4, no.~2, prop.~2).

By \sref{III.7.4.2}, (c), one may furthermore suppose that
$d_{p+1}:P_{p+1}\to P_p$ is zero.  There is then an exact
sequence
\[
  0
  \longrightarrow T_p(M)
  \xrightarrow{u}
  P_p\otimes_A M
  \xrightarrow{v}
  P_{p-1}\otimes_A M,
\]
where $v=d_p\otimes 1$.  Put $Q'=\Ker(d_p)$, so that
\[
  0
  \longrightarrow Q'
  \longrightarrow P_p
  \xrightarrow{d_p}
  P_{p-1}
\]
is exact.  Transposition gives the exact sequence
\[
  \check{P}_{p-1}
  \xrightarrow{{}^t d_p}
  \check{P}_p
  \longrightarrow
  \check{Q}'
  \longrightarrow 0.
\]
We claim that
\[
  Q=\check{Q}'=\Coker({}^t d_p)
\]
has the required property.  Indeed, there is an exact sequence
\[
  0
  \longrightarrow \Hom_A(Q,M)
  \longrightarrow \Hom_A(\check{P}_p,M)
  \xrightarrow{v'}
  \Hom_A(\check{P}_{p-1},M),
\]
\oldpage[III]{189}
where $v'=\Hom_A({}^t d_p,1)$.  Under the canonical
identifications
\[
  P_i\otimes_A M\simeq\Hom_A(\check{P}_i,M),
\]
the homomorphism $v'$ identifies with $v=d_p\otimes 1$.
Thus one obtains the required functorial isomorphism
\[
  T_p(M)\isoto\Hom_A(Q,M).
\]
Finally, $Q$ is a quotient of $\check{P}_p$, and is therefore
of finite type.
\end{proof}

\begin{proposition}[7.4.7]
\label{III.7.4.7}
Suppose that the general hypotheses of \sref{III.7.4.4} hold.
For every $A$-module $M$ of finite type:
\begin{enumerate}
  \item[(i)] the $T_i(M)$ are $A$-modules of finite type;
  \item[(ii)] for every ideal $\mathfrak m$ of $A$, the
    canonical homomorphism
    \[
      \widehat{T_i(M)}
      \longrightarrow
      \varprojlim_n
      T_i\bigl(M\otimes_A(A/\mathfrak m^{n+1})\bigr)
      \tag{7.4.7.1}\label{III.7.4.7.1.eq}
    \]
    is bijective, where the left-hand side denotes the
    separated completion of $T_i(M)$ for the
    $\mathfrak m$-adic topology.
\end{enumerate}
\end{proposition}

\begin{proof}
As in \sref{III.7.4.4}, one reduces first to the case where the
$P_i$ are of finite type.  Since $A$ is Noetherian, the
submodules of the $P_i\otimes_A M$ are of finite type, which
immediately proves (i).  Assertion (ii) follows from the more
general lemma below.
\end{proof}

\begin{lemma}[7.4.7.2]
\label{III.7.4.7.2}
Let $A$ be a Noetherian ring and $u:E\to F$ a homomorphism of
finite-type $A$-modules.  For every finite-type $A$-module
$M$, put
\[
  K(M)=\Ker(u\otimes 1_M),
  \qquad
  C(M)=\Coker(u\otimes 1_M).
\]
Then, for every ideal $\mathfrak m$ of $A$, the canonical
homomorphisms
\[
  \widehat{K(M)}
  \longrightarrow
  \varprojlim_n K(M_n),
  \qquad
  \widehat{C(M)}
  \longrightarrow
  \varprojlim_n C(M_n)
  \tag{7.4.7.3}\label{III.7.4.7.3.eq}
\]
are bijective, where
\[
  M_n
  =
  M\otimes_A(A/\mathfrak m^{n+1})
  =
  M/\mathfrak m^{n+1}M.
\]
\end{lemma}

\begin{proof}
The modules $E\otimes_A M$ and $F\otimes_A M$ are of finite
type, and the functor $M\mapsto\widehat{M}$ is exact on the
category of finite-type $A$-modules \sref[0]{0.7.3.3}.
Consequently, $\widehat{K(M)}$ and $\widehat{C(M)}$ are,
respectively, the kernel and the cokernel of
\[
  \widehat{u\otimes 1}:
  \widehat{E\otimes_A M}
  \longrightarrow
  \widehat{F\otimes_A M}.
\]
Left exactness of $\varprojlim$ therefore gives
\[
  \widehat{K(M)}
  =
  \varprojlim_n K(M_n).
\]
On the other hand, right exactness of the tensor product gives
\[
  C(M_n)
  =
  C(M)\otimes_A(A/\mathfrak m^{n+1}),
\]
and hence
\[
  \widehat{C(M)}
  =
  \varprojlim_n C(M_n)
\]
by definition.
\end{proof}

\begin{remark}[7.4.8]
\label{III.7.4.8}
Taking \sref{III.6.10.5} and \sref{III.6.10.6} into account,
one sees that, with an additional flatness hypothesis,
\sref{III.7.4.7} recovers the fact that
(4.1.7.1) is an isomorphism, which is the
essential content of the ``first comparison theorem'' for
proper morphisms.  Moreover, the statement applies not only
to a coherent $\sh{O}_X$-module but to a \emph{complex} of
such modules.  It would be interesting to obtain a statement
that contains both \sref{III.7.4.7} and
(4.1.7.1) as special cases.

Notice also that, when the $P_i$ are of finite type, the proof
of \sref{III.7.4.7} does not use their flatness.  It would be
worth examining whether the conclusion of \sref{III.7.4.7}
continues to hold when the $P_i$ are assumed neither flat nor
of finite type, but the $H_i(P_\bullet)$ are assumed to be of
finite type for every $i$ and zero for $i<i_0$.  Can one then
replace $P_\bullet$ by a complex $P'_\bullet$ of finite-type
$A$-modules in such a way that the functors
\[
  H_\bullet(P_\bullet\otimes_A M)
  \quad\hbox{and}\quad
  H_\bullet(P'_\bullet\otimes_A M),
\]
which are no longer homological, are still isomorphic?
\end{remark}
% END NATIVE INPUT ega3/ega3-7-4
\ifdefined\EGAIIIStopAfterPageOneEightNine
  \expandafter\endinput
\fi


% BEGIN NATIVE INPUT ega3/ega3-7-5
\oldpage[III]{190}
\subsection{The case of Noetherian local rings}
\label{subsection:III.7.5}

\begin{env}[7.5.1]
\label{III.7.5.1}
Let $A$ be a Noetherian local ring and $\mathfrak m$ its
maximal ideal.  For every $A$-module $M$, denote by
$\widehat M$ its separated completion for the
$\mathfrak m$-adic topology, which is isomorphic to
\[
  \varprojlim_n
  \bigl(M\otimes_A(A/\mathfrak m^{n+1})\bigr)
  =
  \varprojlim_n
  \bigl(M/\mathfrak m^{n+1}M\bigr).
\]
Let $T$ be a covariant additive functor from $\cat{Ab}_A$ to
$\cat{Ab}$.  The canonical homomorphisms
\sref{III.7.2.3.1}
\[
  T(M)\otimes_A(A/\mathfrak m^{n+1})
  \longrightarrow
  T\bigl(M\otimes_A(A/\mathfrak m^{n+1})\bigr)
\]
plainly form a projective system of $A$-homomorphisms.  On
passing to the limit they therefore give an
$\widehat A$-homomorphism, functorial in $M$,
\[
  \widehat{T(M)}
  \longrightarrow
  \varprojlim_n T(M_n),
  \tag{7.5.1.1}\label{III.7.5.1.1.eq}
\]
where
\[
  M_n=M\otimes_A(A/\mathfrak m^{n+1}),
  \qquad
  A_n=A/\mathfrak m^{n+1}.
\]
\end{env}

\begin{proposition}[7.5.2]
\label{III.7.5.2}
Let $A$ be a Noetherian local ring with maximal ideal
$\mathfrak m$ and residue field $k=A/\mathfrak m$.  Let $T$
be a covariant additive functor from $\cat{Ab}_A$ to
$\cat{Ab}$ that is half-exact and commutes with inductive
limits.  Suppose, moreover, that for every finite-type
$A$-module $M$, the module $T(M)$ is of finite type and the
canonical homomorphism \eqref{III.7.5.1.1.eq} is an
isomorphism.  Then the following conditions are equivalent:
\begin{enumerate}
  \item[(a)] $T$ is right exact;
  \item[(b)] for every $n$, the functor $N\mapsto T(N)$ is
    right exact in the category of finite-type $A_n$-modules;
    equivalently, $T$ is right exact in the category of
    finite-length $A$-modules;
  \item[(c)] the canonical homomorphism
    \[
      T(\varepsilon):T(A)\longrightarrow T(k)
    \]
    is surjective;
  \item[(d)] for every sufficiently large $n$, the canonical
    homomorphism
    \[
      T(A_n)\longrightarrow T(k)
    \]
    is surjective.
\end{enumerate}
\end{proposition}

\begin{proof}
It is clear that (a) implies (b).  Let us prove the converse.
Suppose that $u:M\to N$ is an epimorphism of $A$-modules; we
must show that $T(u)$ is an epimorphism.  Since $T$ commutes
with inductive limits and filtered inductive limits are exact
in the category of modules, it is enough to consider the case
where $M$ and $N$ are of finite type.

The modules $T(M)$ and $T(N)$ are then of finite type.  Since
$A$ is Noetherian and local, it is enough to prove that
\[
  \widehat{T(u)}:
  \widehat{T(M)}
  \longrightarrow
  \widehat{T(N)}
\]
is surjective \sref[0]{0.7.3.5} and \sref[0]{0.6.4.1}.
By hypothesis, $\widehat{T(M)}$ and $\widehat{T(N)}$ are,
respectively,
\[
  \varprojlim_n T(M_n)
  \quad\hbox{and}\quad
  \varprojlim_n T(N_n).
\]
Thus $\widehat{T(u)}$ is the inverse limit of the projective
system of homomorphisms
\[
  T(u\otimes 1_{A_n}):T(M_n)\longrightarrow T(N_n).
\]
Condition (b) says precisely that these homomorphisms are
surjective.  Moreover, $T(M_n)$ is a finite-type
$A_n$-module, and $A_n$ is Artinian.  It follows that
$\widehat{T(u)}$ is surjective \sref[0]{0.13.1.2} and
\sref[0]{0.13.2.2}.

Clearly (a) implies (c), and (c) implies (d).  Finally,
\sref{III.7.2.7} shows that (b) and (d) are equivalent,
because $T$ is half-exact in $\cat{Ab}_{A_n}$.  This completes
the proof.
\end{proof}

\begin{corollary}[7.5.3]
\label{III.7.5.3}
Under the general hypotheses of \sref{III.7.5.2}, if
$T(k)=0$, then $T(M)=0$ for every $A$-module $M$.
\end{corollary}

\begin{proof}
\oldpage[III]{191}
Since $k$ is the only simple $A$-module,
\eqref{III.7.3.5.4.eq} shows that $T(E)=0$ for every
finite-length $A$-module $E$.  If $M$ is of finite type, then
\[
  \widehat{T(M)}
  \simeq
  \varprojlim_n T(M_n).
\]
Each $M_n$ has finite length, so $\widehat{T(M)}=0$.  By
hypothesis $T(M)$ is of finite type, and it embeds in
$\widehat{T(M)}$ \sref[0]{0.7.3.5}; hence $T(M)=0$.
Finally, for an arbitrary $A$-module $M$, the module $T(M)$
is the inductive limit of the $T(N_\alpha)$, where the
$N_\alpha$ run through the finite-type submodules of $M$.
The proof is complete.
\end{proof}

\begin{proposition}[7.5.4]
\label{III.7.5.4}
Let $A$ be a Noetherian local ring with maximal ideal
$\mathfrak m$ and residue field $k=A/\mathfrak m$.  Let
$T_\bullet$ be a homological functor from $\cat{Ab}_A$ to
$\cat{Ab}$ that commutes with inductive limits.  Suppose,
moreover, that for every $i$ and every finite-type $A$-module
$M$, the module $T_i(M)$ is of finite type and the canonical
homomorphism
\[
  \widehat{T_i(M)}
  \longrightarrow
  \varprojlim_n T_i(M_n)
\]
is bijective.  For a given integer $p$, the following
conditions are equivalent:
\begin{enumerate}
  \item[(a)] $T_p$ is exact;
  \item[(b)] $T_p$ is right exact and $T_p(A)$ is a free
    $A$-module;
  \item[(c)] the canonical homomorphisms
    \[
      T_{p+1}(A)\longrightarrow T_{p+1}(k),
      \qquad
      T_p(A)\longrightarrow T_p(k)
    \]
    are surjective;
  \item[(d)] for every $n$, the canonical homomorphisms
    \[
      T_{p+1}(A_n)\longrightarrow T_{p+1}(k),
      \qquad
      T_p(A_n)\longrightarrow T_p(k)
    \]
    are surjective;
  \item[(e)] for every $n$, the functor $N\mapsto T_p(N)$ is
    exact in the category of finite-type $A_n$-modules.
\end{enumerate}
\end{proposition}

\begin{proof}
By \sref{III.7.3.3}, condition (a) is equivalent to saying
that $T_{p+1}$ and $T_p$ are right exact.  Since
$T_\bullet$ is a homological functor in $\cat{Ab}_{A_n}$, the
same argument as in \sref{III.7.3.1} shows that (e) is
equivalent to saying that $T_p$ and $T_{p+1}$ are right exact
on the category of finite-type $A_n$-modules.  It follows from
\sref{III.7.5.2} that (a) and (e) are equivalent; the same
proposition also gives the equivalence of (a), (c), and (d).

Finally, every finite-type flat module over the local ring
$A$ is free (Bourbaki, \emph{Alg\`ebre commutative},
chap.~II, \S~3, no.~2, cor.~2 of prop.~5).  The equivalence
of (a) and (b) therefore follows from \sref{III.7.3.1} and
\sref{III.7.3.3}.
\end{proof}

\begin{corollary}[7.5.5]
\label{III.7.5.5}
Suppose that the general hypotheses of \sref{III.7.5.4} hold.
\begin{enumerate}
  \item[(i)] If $T_p(k)=0$, then $T_p=0$, the functor
    $T_{p+1}$ is right exact, and $T_{p-1}$ is left exact.
  \item[(ii)] If $T_{p-1}(k)=T_{p+1}(k)=0$, then $T_p$ is
    exact, the canonical homomorphism
    \[
      T_p(A)\otimes_A M
      \longrightarrow
      T_p(M)
    \]
    is bijective, and $T_p(A)$ is a free $A$-module.
\end{enumerate}
\end{corollary}

\begin{proof}
Assertion (i) follows at once from \sref{III.7.5.3}, since
$T_p$ is half-exact; the last assertion follows from the
definition of a homological functor.  Taking
\sref{III.7.3.3} into account, (i) immediately gives the first
two assertions of (ii), and the freeness of $T_p(A)$ follows
from \sref{III.7.5.4}.
\end{proof}

\begin{corollary}[7.5.6]
\label{III.7.5.6}
Suppose that the general hypotheses of \sref{III.7.5.4} hold,
and suppose in addition that $A$ is a discrete valuation
ring.
\begin{enumerate}
  \item[(i)] The functor $T_p$ is right exact if and only if
    $T_{p-1}(A)$ is a free $A$-module.
  \item[(ii)] The functor $T_p$ is exact if and only if
    $T_p(A)$ and $T_{p-1}(A)$ are free $A$-modules.
\end{enumerate}
\end{corollary}

\oldpage[III]{192}
\begin{proof}
It is clear that (ii) implies (i) (cf.\
\sref{III.7.3.3}).  To prove (i), let $f$ be a generator of
the maximal ideal of $A$, that is, a uniformizer of $A$.  A
finite-type $A$-module $M$ is free, equivalently flat, if and
only if the homothety
\[
  h_f:M\longrightarrow M,\qquad x\longmapsto fx,
\]
is injective, since in the present situation this amounts to
saying that $M$ is torsion-free \sref[0]{0.6.3.4}.  The exact
sequence
\[
  0
  \longrightarrow A
  \xrightarrow{h_f}
  A
  \longrightarrow k
  \longrightarrow 0
\]
gives the exact homology sequence
\[
  T_p(A)
  \longrightarrow T_p(k)
  \longrightarrow T_{p-1}(A)
  \xrightarrow{h_f}
  T_{p-1}(A).
\]
Thus $T_{p-1}(A)$ is free if and only if
$T_p(A)\to T_p(k)$ is surjective, and the conclusion follows
from \sref{III.7.5.2}.
\end{proof}

\begin{remark}[7.5.7]
\label{III.7.5.7}
By \sref{III.7.4.7}, the general hypotheses of
\sref{III.7.4.4} imply that the homological functor
$T_\bullet$ defined by \eqref{III.7.4.1.1.eq} satisfies the
general hypotheses of \sref{III.7.5.4}.  In that case,
\sref{III.7.5.6} is therefore contained in
\sref{III.7.4.3}.
\end{remark}
% END NATIVE INPUT ega3/ega3-7-5

% BEGIN NATIVE INPUT ega3/ega3-7-6
\subsection{Descent of exactness properties. Semicontinuity theorem and
Grauert's exactness criterion}
\label{subsection:III.7.6}
\label{III.7.6}

\begin{proposition}[7.6.1]
\label{III.7.6.1}
Under the hypotheses of \sref{III.7.4.1}, let $B$ be a
commutative $A$-algebra.  If $T_p$ is right exact
(respectively left exact, exact), then the same is true of
$T_p^{(B)}$; the converse is true when $B$ is a faithfully
flat $A$-module.
\end{proposition}

\begin{proof}
The first assertion is a special case of a trivial assertion
of \sref{III.7.1.3}.  Conversely, suppose first that $B$ is a
flat $A$-module.  Then, for every $A$-module $M$,
\[
  H_\bullet\bigl(P_\bullet\otimes_A(M\otimes_A B)\bigr)
  =
  H_\bullet(P_\bullet\otimes_A M)\otimes_A B,
\]
which may also be written, for every $p$,
\[
  T_p(M)\otimes_A B
  \simeq
  T_p^{(B)}(M_{(B)})
  \tag{7.6.1.1}\label{III.7.6.1.1}
\]
up to canonical isomorphism.  Suppose that $T_p^{(B)}$ is
right exact (respectively left exact, exact).  Since
$M\mapsto M_{(B)}$ is an exact functor, the left-hand side of
\eqref{III.7.6.1.1} is a right exact (respectively left exact,
exact) functor of $M$.  If $B$ is now faithfully flat over
$A$, it follows that $T_p$ has the same exactness property
\sref[0]{0.6.4.1}.
\end{proof}

\long\def\egaIIISubsectionSevenSixEleven{%
\begin{remarks}[7.6.11]
\label{III.7.6.11}
One may ask whether \sref{III.7.6.9}[(i)] can be strengthened
to the inequality
\[
  \operatorname{rank}_{\kres(x)}T_p(\kres(x))
  \geq
  \operatorname{rank}_{\kres(x)}
  \bigl(T_p(A)\otimes_A\kres(x)\bigr)
  \tag{7.6.11.1}\label{III.7.6.11.1}
\]
for every $x\in\Spec(A)$.  This inequality does hold when
$A$ is a discrete valuation ring and $x$ is its maximal
ideal \sref{III.7.6.5}.  Restrict now to the case in which
$A$ is a Noetherian local ring with maximal ideal
$\mathfrak m$ and residue field $k$.  The following
conditions are equivalent:
\begin{enumerate}
  \item[a)] For every complex $P_\bullet$ of finite-type flat
    $A$-modules,
    \[
      \operatorname{rank}_k T_p(k)
      \geq
      \operatorname{rank}_k\bigl(T_p(A)\otimes_A k\bigr)
      \qquad\text{for every integer }p.
      \tag{7.6.11.2}\label{III.7.6.11.2}
    \]
  \item[b)] For every finite-type $A$-module $M$,
    \[
      \operatorname{rank}_k(M\otimes_A k)
      \geq
      \operatorname{rank}_k(\check M\otimes_A k).
      \tag{7.6.11.3}\label{III.7.6.11.3}
    \]
  \item[c)] For every finite-type $A$-module $N$,
    \[
      \operatorname{rank}_k\Tor_1^A(N,k)
      \geq
      \operatorname{rank}_k\Tor_2^A(N,k).
      \tag{7.6.11.4}\label{III.7.6.11.4}
    \]
\end{enumerate}
Equivalently, by dimension shifting (M, V, 7.2), one may
require, for every $i\geq1$,
\[
  \operatorname{rank}_k\Tor_i^A(N,k)
  \geq
  \operatorname{rank}_k\Tor_{i+1}^A(N,k).
  \tag{7.6.11.5}\label{III.7.6.11.5}
\]

Here are brief indications of the proof.  To see that
\textnormal{a)} implies \textnormal{b)}, consider an exact
sequence
\[
  L_1\xrightarrow{d}L_0\longrightarrow M\longrightarrow0,
\]
where $L_0$ and $L_1$ are finite free, and apply
\textnormal{a)} to the complex
\[
  P_1\xrightarrow{{}^t d}P_0,
  \qquad
  P_0=\check L_1,\quad P_1=\check L_0,
\]
all its other terms being zero.  Then
\[
  T_1(A)=\check M
\]
and
\[
  T_1(k)
  =
  \Hom_A(M,k)
  =
  \Hom_k(M/\mathfrak mM,k).
\]
Thus $T_1(k)$ is the dual of the vector space $M\otimes_A k$
and has the same rank.  To prove that \textnormal{b)} implies
\textnormal{c)}, we first establish the following lemma.

\begin{lemma}[7.6.11.6]
\label{III.7.6.11.6}
Given a complex of flat $A$-modules
\[
  \cdots\longrightarrow0\longrightarrow
  P_1\xrightarrow{d}P_0
  \longrightarrow0\longrightarrow\cdots,
\]
there is an exact sequence
\[
  0\longrightarrow\Tor_2^A(Z'_0,k)
  \longrightarrow T_1(A)\otimes_A k
  \longrightarrow T_1(k)
  \longrightarrow\Tor_1^A(Z'_0,k)
  \longrightarrow0.
  \tag{7.6.11.7}\label{III.7.6.11.7}
\]
\end{lemma}

\oldpage[III]{197}
\begin{proof}
From the exact sequence
\[
  0\longrightarrow Z_1\longrightarrow P_1
  \longrightarrow B_0\longrightarrow0
\]
we obtain
\[
  0\longrightarrow\Tor_1^A(B_0,k)
  \longrightarrow Z_1\otimes_A k
  \longrightarrow P_1\otimes_A k
  \xrightarrow{u}B_0\otimes_A k
  \longrightarrow0.
\]
The exact sequence
\[
  0\longrightarrow B_0\longrightarrow Z_0
  \longrightarrow Z'_0\longrightarrow0
\]
and the flatness of $Z_0=P_0$ give
\[
  \Tor_1^A(B_0,k)=\Tor_2^A(Z'_0,k).
\]
By definition, $Z_1=T_1(A)$.  Moreover,
$T_1(k)=\Ker(d\otimes1)$, and $d\otimes1$ factors as
\[
  P_1\otimes_A k
  \xrightarrow{u}B_0\otimes_A k
  \xrightarrow{v}Z_0\otimes_A k
  =
  P_0\otimes_A k.
\]
Thus $T_1(k)=u^{-1}(R)$, where $R=\Ker v$.  Since $u$ is
surjective, $R=u(T_1(k))$; finally,
$R=\Tor_1^A(Z'_0,k)$ by the definition of $v$.  This proves
\eqref{III.7.6.11.7}.
\end{proof}

To deduce \textnormal{c)} from \textnormal{b)}, take an exact
sequence
\[
  L_1\xrightarrow{d}L_0\longrightarrow N\longrightarrow0
\]
with $L_0,L_1$ finite free, and consider the functor $T$
associated with the complex formed by $L_1$ and $L_0$.
Since $L_0$ and $L_1$ are free, they identify with their
biduals.  If
\[
  M=\Coker({}^t d),
\]
then
\[
  T_1(A)=\Ker d=\check M.
\]
Furthermore,
\[
  M\otimes_A k=\Coker({}^t d\otimes1_k)
\]
has the same rank over $k$ as $\Ker(d\otimes1_k)$.
Hypothesis \textnormal{b)} therefore gives
\[
  \operatorname{rank}_k\bigl(T_1(A)\otimes_A k\bigr)
  \leq
  \operatorname{rank}_kT_1(k).
\]
Since $Z'_0=N$, inequality \eqref{III.7.6.11.4} follows from
the exact sequence \eqref{III.7.6.11.7}.

Finally, to prove that \textnormal{c)} implies
\textnormal{a)}, apply \sref{III.7.6.11.6} to the two-term
complex
\[
  P_p\longrightarrow P_{p-1}.
\]
Hypothesis \textnormal{c)}, applied to $Z'_{p-1}$, gives
\[
  \operatorname{rank}_k R\geq\operatorname{rank}_k S,
\]
where
\[
  R=\Ker(P_p\otimes_A k\longrightarrow P_{p-1}\otimes_A k),
  \qquad
  S=Z_p\otimes_A k.
\]
Factor $d_{p+1}:P_{p+1}\to P_p$ as
\[
  P_{p+1}\xrightarrow{v}Z_p\xrightarrow{j}P_p.
\]
Then
\[
  \Im(d_{p+1}\otimes1)
  =
  (j\otimes1)\bigl(\Im(v\otimes1)\bigr).
\]
Since
\[
  T_p(A)\otimes_A k
  =
  (Z_p/B_p)\otimes_A k
  =
  (Z_p\otimes_A k)/\Im(v\otimes1)
\]
and
\[
  T_p(k)=R/\Im(d_{p+1}\otimes1),
\]
inequality \eqref{III.7.6.11.2} follows.

Now suppose that the local ring $A$ is regular of dimension
$n$.  It is known [17] that the $A$-module
$\Tor_i^A(k,k)$ is isomorphic to the exterior power
\[
  \bigwedge\nolimits^i(\mathfrak m/\mathfrak m^2).
\]
Thus \eqref{III.7.6.11.4} fails for $N=k$ as soon as
$n\geq4$.  On the other hand, suppose that $A$ is a local
integral domain such that every finite-type reflexive
$A$-module is free, as happens when $A$ is a regular ring of
dimension $2$.  Then \eqref{III.7.6.11.3} holds.  Indeed, the
dual $\check M$ of a finite-type $A$-module $M$ is reflexive,
hence free, and consequently
\[
  \operatorname{rank}_k(\check M\otimes_A k)
  =
  \operatorname{rank}_K(\check M)
  =
  \operatorname{rank}_K(M),
\]
where $K$ is the field of fractions of $A$.  Every basis of
$M\otimes_A k$ over $k$ consists of the images of a system of
generators of $M$ (Bourbaki, \emph{Alg\`ebre commutative},
chap.~II, \S~3, no.~2, cor.~2 of prop.~4).  Hence
\[
  \operatorname{rank}_K(M)
  \leq
  \operatorname{rank}_k(M\otimes_A k),
\]
which proves the assertion.
\end{remarks}
}

\begin{proposition}[7.6.2]
\label{III.7.6.2}
Under the hypotheses of \sref{III.7.4.1}, suppose in addition
that $A$ is a reduced Noetherian ring and that the $P_i$ are
finite-type $A$-modules.  For $T_p$ to be right exact
(respectively left exact, exact), it is necessary and
sufficient that $T_p^{(B)}$ have the same property for every
$A$-algebra $B$ that is a discrete valuation ring.
\end{proposition}

\begin{proof}
By \sref{III.7.3.1} and \sref{III.7.3.3}, we may restrict
attention to right exactness, and only sufficiency has to be
proved.  By \sref{III.7.4.2}, it is enough to show that
$Z'_{p-1}(P_\bullet)$ is a flat $A$-module.  Since
$P_{p-1}$ is of finite type, $Z'_{p-1}(P_\bullet)$ is also of
finite type; the criterion of \sref[0]{0.10.2.8} therefore
shows that it is enough for
\[
  Z'_{p-1}(P_\bullet)\otimes_A B
\]
to be a flat $B$-module for every $A$-algebra $B$ that is a
discrete valuation ring.  Since $P_\bullet$ is a complex of
flat $A$-modules, we have
\[
  Z'_{p-1}(P_\bullet)\otimes_A B
  =
  Z'_{p-1}(P_\bullet\otimes_A B).
\]

\oldpage[III]{193}
The complex $P_\bullet\otimes_A B$ consists of flat
$B$-modules \sref[0]{0.6.2.1}, and for every $B$-module $N$
we have
\[
  H_\bullet(P_\bullet\otimes_A N)
  =
  H_\bullet\bigl((P_\bullet\otimes_A B)\otimes_B N\bigr),
\]
and hence
\[
  T_p^{(B)}(N)
  =
  H_p\bigl((P_\bullet\otimes_A B)\otimes_B N\bigr).
\]
Applying \sref{III.7.4.2} to $T_p^{(B)}$, we see that the
hypothesis that $T_p^{(B)}$ is right exact is equivalent to
$Z'_{p-1}(P_\bullet\otimes_A B)$ being a flat $B$-module.

The preceding criterion leads us to examine the case of
discrete valuation rings more closely.
\end{proof}

\begin{proposition}[7.6.3]
\label{III.7.6.3}
Under the hypotheses of \sref{III.7.4.1}, suppose that $A$ is
a regular Noetherian ring of dimension $1$ (in other words,
that $A$ is Noetherian and, for every $x\in\Spec(A)$, the ring
$A_x$ is a field or a discrete valuation ring).  Then, for
every integer $p$ and every $A$-module $M$, there is a
canonical exact sequence, functorial in $M$,
\[
  0
  \longrightarrow T_p(A)\otimes_A M
  \xrightarrow{t_M} T_p(M)
  \longrightarrow
  \Tor_1^A\bigl(T_{p-1}(A),M\bigr)
  \longrightarrow 0.
  \tag{7.6.3.1}\label{III.7.6.3.1}
\]
\end{proposition}

\begin{proof}
In what follows, to simplify the notation we omit the complex
$P_\bullet$ from the usual homological notations
$H_p(P_\bullet)$, $B_p(P_\bullet)$, $Z_p(P_\bullet)$, and
$Z'_p(P_\bullet)$.  We have the three exact sequences
\begin{gather*}
  0\longrightarrow H_p\longrightarrow Z'_p
  \longrightarrow B_{p-1}\longrightarrow 0,\\
  0\longrightarrow B_{p-1}\longrightarrow Z_{p-1}
  \longrightarrow H_{p-1}\longrightarrow 0,\\
  0\longrightarrow Z_{p-1}\longrightarrow P_{p-1}
  \longrightarrow B_{p-2}\longrightarrow 0.
\end{gather*}
Since $P_{p-1}$ and $P_{p-2}$ are flat, so are their
respective submodules $B_{p-1}$, $Z_{p-1}$, and $B_{p-2}$:
for every $x\in\Spec(A)$, flat $A_x$-modules and torsion-free
$A_x$-modules are the same.  Tensoring with $M$ therefore
gives the exact sequences
\[
  0=\Tor_1^A(B_{p-1},M)
  \longrightarrow H_p\otimes_A M
  \longrightarrow Z'_p\otimes_A M
  \xrightarrow{u}B_{p-1}\otimes_A M
  \longrightarrow 0,
  \tag{7.6.3.2}\label{III.7.6.3.2}
\]
\[
  0=\Tor_1^A(Z_{p-1},M)
  \longrightarrow \Tor_1^A(H_{p-1},M)
  \longrightarrow B_{p-1}\otimes_A M
  \xrightarrow{v} Z_{p-1}\otimes_A M,
  \tag{7.6.3.3}\label{III.7.6.3.3}
\]
and
\[
  0=\Tor_1^A(B_{p-2},M)
  \longrightarrow Z_{p-1}\otimes_A M
  \xrightarrow{w}P_{p-1}\otimes_A M.
  \tag{7.6.3.4}\label{III.7.6.3.4}
\]

By definition,
\[
  T_p(M)=
  \Ker(d_p\otimes 1)/\Im(d_{p+1}\otimes 1).
\]
It is therefore the kernel of the homomorphism
\[
  (P_p\otimes_A M)/\Im(d_{p+1}\otimes 1)
  \longrightarrow P_{p-1}\otimes_A M
\]
induced by $d_p\otimes 1$.  By the definition
$Z'_p=P_p/B_p$, this homomorphism may also be written
$Z'_p\otimes_A M\to P_{p-1}\otimes_A M$; it is the composite
\[
  Z'_p\otimes_A M
  \xrightarrow{u}B_{p-1}\otimes_A M
  \xrightarrow{v}Z_{p-1}\otimes_A M
  \xrightarrow{w}P_{p-1}\otimes_A M.
\]
Since $w$ is injective by \eqref{III.7.6.3.4}, there is an
exact sequence
\[
  0\longrightarrow\Ker u\longrightarrow T_p(M)
  \longrightarrow\Ker v\longrightarrow 0.
\]
In view of \eqref{III.7.6.3.2} and
\eqref{III.7.6.3.3}, and since $H_p=T_p(A)$ by definition,
this is precisely \eqref{III.7.6.3.1}.
\end{proof}

\begin{remarks}[7.6.4]
\label{III.7.6.4}
\textnormal{(i)}
The group $H_\bullet(P_\bullet\otimes_A M)$ is the homology of
the bicomplex $P_\bullet\otimes_A M$, where $M$ is regarded as
a complex concentrated in degree $0$.  It is consequently
\sref{III.6.3.6}\sref{III.6.3.2} the abutment of the regular
spectral sequence whose $E_2$-term is
\[
  E^2_{pq}
  =
  \Tor_p^A\bigl(H_q(P_\bullet),M\bigr)
  =
  \Tor_p^A\bigl(T_q(A),M\bigr).
\]

\oldpage[III]{194}
The hypothesis on $A$ implies that
$\Tor_p^A(E,F)=0$ for $p\geq 2$ and arbitrary $A$-modules
$E,F$ \sref[0]{0.17.2.2}.  It is known (M, XV) that this
implies the exactness of
\[
  0\longrightarrow E^2_{0,q}
  \longrightarrow H_q(P_\bullet\otimes_A M)
  \longrightarrow E^2_{1,q-1}
  \longrightarrow 0,
\]
which is precisely \eqref{III.7.6.3.1}.

\medskip
\textnormal{(ii)}
Taking \sref{III.7.3.1} into account, the exact sequence
\eqref{III.7.6.3.1} recovers \sref{III.7.4.3} as a special
case.
\end{remarks}

\begin{corollary}[7.6.5]
\label{III.7.6.5}
Under the hypotheses of \sref{III.7.4.1}, suppose that $A$ is
a discrete valuation ring with field of fractions $K$ and
residue field $k$, and that the $T_i(A)$ are finite-type
$A$-modules.  Then
\[
  \operatorname{rank}_k T_p(k)
  \geq
  \operatorname{rank}_k\bigl(T_p(A)\otimes_A k\bigr)
  \geq
  \operatorname{rank}_A T_p(A)
  =
  \operatorname{rank}_K T_p(K).
  \tag{7.6.5.1}\label{III.7.6.5.1}
\]
Moreover, equality of the two extreme terms of this
inequality holds if and only if $T_p$ is exact, or,
equivalently, if and only if $T_p(A)$ and $T_{p-1}(A)$ are
free $A$-modules.
\end{corollary}

\begin{proof}
Set $M=k$ in the exact sequence \eqref{III.7.6.3.1}.  Since
the terms are $k$-vector spaces, this gives
\[
  \operatorname{rank}_k T_p(k)
  =
  \operatorname{rank}_k\bigl(T_p(A)\otimes_A k\bigr)
  +
  \operatorname{rank}_k
  \Tor_1^A\bigl(T_{p-1}(A),k\bigr).
\]
On the other hand, since $T_p(A)$ is a finite-type module
over the local integral domain $A$, we have (Bourbaki,
\emph{Alg\`ebre commutative}, chap.~II, \S~3, no.~2,
cor.~1 of prop.~4)
\[
  \operatorname{rank}_k\bigl(T_p(A)\otimes_A k\bigr)
  \geq
  \operatorname{rank}_A T_p(A)
  =
  \operatorname{rank}_K\bigl(T_p(A)\otimes_A K\bigr),
  \tag{7.6.5.2}\label{III.7.6.5.2}
\]
and the two sides of \eqref{III.7.6.5.2} are equal if and
only if $T_p(A)$ is a free $A$-module (\emph{loc.\ cit.},
prop.~7).  Since $K$ is a flat $A$-module, we also have, by
definition,
\[
  T_p(A)\otimes_A K
  =
  H_p(P_\bullet)\otimes_A K
  =
  H_p(P_\bullet\otimes_A K)
  =
  T_p(K).
\]
This proves \eqref{III.7.6.5.1}.  It also shows that equality
can hold only if, first, $T_p(A)$ is free, and, second,
\[
  \Tor_1^A\bigl(T_{p-1}(A),k\bigr)=0,
\]
the latter condition being equivalent to $T_{p-1}(A)$ being
free \sref[0]{0.10.1.3}.  Finally, for finite-type modules
over $A$, flatness and freeness are equivalent (Bourbaki,
\emph{Alg\`ebre commutative}, chap.~II, \S~3, no.~2,
cor.~2 of prop.~5), and the result follows from
\sref{III.7.4.3}.
\end{proof}

\begin{env}[7.6.6]
\label{III.7.6.6}
The hypotheses remaining those of \sref{III.7.4.1}, for
every $x\in\Spec(A)$ set
\[
  d_p(x)=d_p^T(x)
  =
  \operatorname{rank}_{\kres(x)}T_p(\kres(x)).
  \tag{7.6.6.1}\label{III.7.6.6.1}
\]
\end{env}

\begin{lemma}[7.6.7]
\label{III.7.6.7}
Let $\varphi:A\to A'$ be a ring homomorphism, and let
\[
  f={}^a\!\varphi:\Spec(A')\longrightarrow\Spec(A)
\]
be the corresponding map \sref[I]{I.1.2.1}.  If
$T'_\bullet=T_\bullet^{(A')}$ \sref{III.7.1.3}, then
\[
  d_p^{T'}=d_p^T\circ f.
  \tag{7.6.7.1}\label{III.7.6.7.1}
\]
\end{lemma}

\oldpage[III]{195}
\begin{proof}
For $x'\in\Spec(A')$, put $x=f(x')$.  Then
\begin{align*}
  H_\bullet(P_\bullet\otimes_A\kres(x'))
  &=
  H_\bullet\bigl(
    (P_\bullet\otimes_A\kres(x))
    \otimes_{\kres(x)}\kres(x')
  \bigr)\\
  &=
  H_\bullet(P_\bullet\otimes_A\kres(x))
  \otimes_{\kres(x)}\kres(x'),
\end{align*}
because $\kres(x')$ is flat over $\kres(x)$.  This proves
\eqref{III.7.6.7.1}.
\end{proof}

\begin{lemma}[7.6.8]
\label{III.7.6.8}
If $A$ is Noetherian and $P_\bullet$ consists of finite-type
$A$-modules, the function $x\mapsto d_p^T(x)$ on $\Spec(A)$
is constructible.
\end{lemma}

\begin{proof}
For every irreducible closed subset $Y$ of
$X=\Spec(A)$, we must show that there is a nonempty open
subset $U$ of $Y$ on which $d_p$ is constant
\sref[0]{0.9.2.2}.  Write $Y=\Spec(A/\mathfrak a)$, where
$\mathfrak a$ is an ideal of $A$ such that
$A/\mathfrak a$ is reduced.  By \sref{III.7.6.7}, we may
reduce to the case $Y=X$ and $A$ a Noetherian integral
domain; the assertion then follows from \sref{III.7.4.5}.
\end{proof}

\begin{theorem}[7.6.9]
\label{III.7.6.9}
Let $A$ be a Noetherian ring, let $P_\bullet$ be a complex of
finite-type flat $A$-modules, and let
\[
  T_\bullet(M)=H_\bullet(P_\bullet\otimes_A M)
\]
be the homological functor defined by $P_\bullet$.  For
$x\in\Spec(A)$, set
\[
  d_p(x)=\operatorname{rank}_{\kres(x)}T_p(\kres(x)).
\]
Then:
\begin{enumerate}
  \item[(i)] The function $d_p$ is constructible and upper
    semicontinuous on $\Spec(A)$.
  \item[(ii)] If $T_p$ is exact, $d_p$ is continuous (and
    hence locally constant) on $\Spec(A)$.  The converse is
    true when $A$ is reduced.
\end{enumerate}
\end{theorem}

\begin{proof}
\textnormal{(i)}
Constructibility was proved in \sref{III.7.6.8}.  To prove
upper semicontinuity, it is therefore enough
\sref[0]{0.9.2.4} to show that if $x'\neq x$ is a
generalization of $x$ in $\Spec(A)$, then
$d_p(x')\leq d_p(x)$.  There is a discrete valuation ring
$B$ and a morphism
\[
  f:\Spec(B)\longrightarrow\Spec(A)
\]
such that, if $a$ is the closed point of $\Spec(B)$ and $b$
its generic point, then $f(a)=x$ and $f(b)=x'$
\sref[II]{II.7.1.9}.  By \eqref{III.7.6.7.1}, we are reduced
to proving $d_p(a)\geq d_p(b)$ on $\Spec(B)$, which is exactly
\eqref{III.7.6.5.1}.\footnote{The principle of proving
\textnormal{(i)} by reduction to the case of a discrete
valuation ring was communicated to us orally by Hironaka.}

\medskip
\textnormal{(ii)}
The first assertion was already proved in
\sref{III.7.3.4}.  For the converse, use the valuative
criterion \sref{III.7.6.2}.  In view of
\eqref{III.7.6.7.1}, we are reduced to the case in which $A$
is a discrete valuation ring.  Since $\Spec(A)$ then has only
two points, the hypothesis that $d_p$ is constant implies
that $T_p$ is exact, by \sref{III.7.6.5}.
\end{proof}

\begin{corollary}[7.6.10]
\label{III.7.6.10}
Let $A$ be a Noetherian ring, let
$\mathfrak p_i$ $(1\leq i\leq r)$ be its minimal prime
ideals, and let $k_i$ be the residue field of
$A_{\mathfrak p_i}$.
\begin{enumerate}
  \item[(i)] For every $x\in\Spec(A)$, there is an index $i$
    such that
    \[
      d_p(x)\geq
      \operatorname{rank}_{k_i}T_p(k_i).
      \tag{7.6.10.1}\label{III.7.6.10.1}
    \]
\end{enumerate}
In particular, if $A$ is an integral domain and $K$ is its
field of fractions, then
\[
  d_p(x)\geq\operatorname{rank}_K T_p(K)
  \tag{7.6.10.2}\label{III.7.6.10.2}
\]
for every $x\in\Spec(A)$.

\oldpage[III]{196}
\begin{enumerate}
  \item[(ii)] Suppose in addition that $A$ is local and
    reduced, and let $k$ be its residue field.  Then $T_p$ is
    exact if and only if
    \[
      \operatorname{rank}_k T_p(k)
      =
      \operatorname{rank}_{k_i}T_p(k_i)
      \qquad(1\leq i\leq r).
      \tag{7.6.10.3}\label{III.7.6.10.3}
    \]
\end{enumerate}
\end{corollary}

\begin{proof}
Assertion \textnormal{(i)} is immediate: every neighbourhood
of $x$ contains one of the $\mathfrak p_i$, and it is enough
to apply the definition of semicontinuity.  If $A$ is local,
the only neighbourhood of the maximal ideal $\mathfrak m$ in
$\Spec(A)$ is all of $\Spec(A)$; hence
\[
  d_p(x)\leq\operatorname{rank}_k T_p(k)
\]
for every $x\in\Spec(A)$.  Together with \textnormal{(i)},
this shows that \eqref{III.7.6.10.3} implies that $d_p$ is
constant on $\Spec(A)$, and hence that $T_p$ is exact by
\sref{III.7.6.9}[(ii)].  The converse follows at once from
the same result.
\end{proof}

\egaIIISubsectionSevenSixEleven
% END NATIVE INPUT ega3/ega3-7-6
% BEGIN NATIVE INPUT ega3/ega3-7-7
\subsection{Application to proper morphisms: I. The exchange property}
\label{subsection:III.7.7}
\label{III.7.7}

The next three paragraphs are, for the most part,
translations into the language of morphisms of preschemes
of the results of the preceding paragraphs.

\begin{env}[7.7.1]
\label{III.7.7.1}
Let $f:X\to Y$ be a quasi-compact separated morphism of
preschemes, and let $\sh{P}_\bullet$ be a complex of
quasi-coherent $\sh{O}_X$-Modules whose differential has
degree $-1$.  Suppose in addition that the $\sh{O}_X$-Modules
$\sh{P}_i$ are $Y$-flat \sref[0]{0.6.7.1}.

\oldpage[III]{198}
We shall consider the $\partial$-functor
\[
  \sh{M}\longmapsto\sh{T}_\bullet(\sh{M}),
\]
also denoted
$\sh{T}_\bullet(\sh{P}_\bullet,\sh{M})$, from the category of
quasi-coherent $\sh{O}_Y$-Modules to itself (by
\sref{III.6.2.3}), defined by
\[
  \sh{T}_n(\sh{P}_\bullet,\sh{M})
  =\sh{T}_n(\sh{M})
  =\mathscr{H}^{-n}
    \bigl(f,\sh{P}^\bullet\otimes_{\sh{O}_Y}\sh{M}\bigr)
  \qquad(n\in\bb Z),
  \tag{7.7.1.1}\label{III.7.7.1.1}
\]
where $\sh{P}^\bullet$ is the complex whose term of degree
$j$ is $\sh{P}_{-j}$, so that its differential has degree
$+1$.  The functor $\sh{T}_\bullet$ thus defined is a
\emph{homological functor} in $\sh{M}$ \sref{III.6.2.6}.
\end{env}

\long\def\egaIIISubsectionSevenSevenTheorems{%
\begin{theorem}[7.7.5]
\label{III.7.7.5}
Let $Y$ be a locally Noetherian prescheme,
$(U_\alpha)$ a covering of $Y$ by affine open subsets,
$f:X\to Y$ a proper morphism, and $\sh{P}_\bullet$ a
bounded-below complex of coherent, $Y$-flat
$\sh{O}_X$-Modules.  The homological functor
$\sh{T}_\bullet(\sh{M})$ defined by
\eqref{III.7.7.1.1} then has the following properties.

\medskip
\noindent
\textnormal{I) (The semicontinuity property).}\footnote{A
special case of this theorem already occurs in note [3] of
Chow--Igusa.  The semicontinuity property was discovered,
for analytic spaces and under rather special hypotheses, by
Kodaira--Spencer, \emph{On the variations of almost-complex
structures}, in \emph{Algebraic Geometry and Topology: A
Symposium in Honor of S. Lefschetz}, Princeton Series
no.~12, pp.~139--150, Princeton, 1957; the general version
was proved by Grauert [5].}
\emph{The function}
\[
  y\longmapsto d_p(y)
  =
  \operatorname{rank}_{\kres(y)}
    T_p^{\kres(y)}(\kres(y))
  \tag{7.7.5.1}\label{III.7.7.5.1}
\]
\emph{is upper semicontinuous.}

\oldpage[III]{200}
\medskip
\noindent
\textnormal{II) (The exchange property).}
\emph{For a fixed integer $p$, the following conditions are
equivalent:}
\begin{enumerate}
  \item[a)] $\sh{T}_p$ is right exact.

  \item[a$'$)] $\sh{T}_p(\sh{M})$ is isomorphic to a functor
    of the form
    $\sh{N}\otimes_{\sh{O}_Y}\sh{M}$, where necessarily
    \[
      \sh{N}\simeq
      \sh{T}_p(\sh{O}_Y)
      =
      \mathscr{H}^{-p}(f,\sh{P}^\bullet).
    \]

  \item[a$''$)] The canonical functorial homomorphism
    \eqref{III.7.7.3.1} is an isomorphism.

  \item[b)] $\sh{T}_{p-1}$ is left exact.

  \item[b$'$)] There is an $\sh{O}_Y$-Module $\sh{Q}$,
    necessarily coherent and determined up to unique
    isomorphism, together with an isomorphism of functors
    \[
      \sh{T}_{p-1}(\sh{M})
      \simto
      \shHom_{\sh{O}_Y}(\sh{Q},\sh{M}).
      \tag{7.7.5.2}\label{III.7.7.5.2}
    \]

  \item[c)] If $A_\alpha$ denotes the ring of the affine open
    subset $U_\alpha$, then for every $\alpha$ the functor
    $T_p^{A_\alpha}$ on $A_\alpha$-modules is right exact.

  \item[d)] For every morphism $g:Y'\to Y$, the canonical
    homomorphism
    \[
      \sh{T}_p(\sh{O}_Y)
      \otimes_{\sh{O}_Y}\sh{O}_{Y'}
      \longrightarrow
      \sh{T}^{Y'}_p(\sh{O}_{Y'})
      \tag{7.7.5.3}\label{III.7.7.5.3}
    \]
    is an isomorphism.
\end{enumerate}
\end{theorem}

\begin{proof}
The semicontinuity property is local on $Y$, and therefore
follows from \sref{III.7.7.4} and \sref{III.7.6.9}.  It is
clear that a$''$) implies a$'$), and that a$'$) implies a).
When $Y$ is affine, the equivalence of a), a$''$), b), and
b$'$) was proved in \sref{III.7.3.1} and
\sref{III.7.4.6}, in view of \sref{III.7.7.4}.

To pass to the general case, first observe that a) is
equivalent to c); this also proves that a) is local on $Y$.
The same argument proves the locality of a$''$) and b).
Condition c) plainly implies a), so it remains to prove the
converse.  It is enough to show that, for every affine open
subset $U$ of $Y$ and every exact sequence
\[
  0\longrightarrow\sh{F}'
  \longrightarrow\sh{F}
  \longrightarrow\sh{F}''
  \longrightarrow0
\]
of quasi-coherent $(\sh{O}_Y|U)$-Modules, there is an exact
sequence
\[
  0\longrightarrow\sh{G}'
  \longrightarrow\sh{G}
  \longrightarrow\sh{G}''
  \longrightarrow0
\]
of quasi-coherent $\sh{O}_Y$-Modules such that
\[
  \sh{F}'=\sh{G}'|U,
  \qquad
  \sh{F}=\sh{G}|U,
  \qquad
  \sh{F}''=\sh{G}''|U.
\]
This follows immediately from the hypothesis that $Y$ is
locally Noetherian and from \sref[I]{I.9.4.2}: extend
$\sh{F}$ to a quasi-coherent $\sh{O}_Y$-Module $\sh{G}$,
extend $\sh{F}'$ to a quasi-coherent submodule $\sh{G}'$ of
$\sh{G}$, and take $\sh{G}''=\sh{G}/\sh{G}'$.

To prove the equivalence of b) and b$'$) in the general
case, note first that, when $Y$ is affine, $\sh{Q}$ is
determined up to unique isomorphism.  If $U$ is an affine
open subset of the affine scheme $Y$, it follows that there
is a functorial isomorphism
\[
  \sh{T}^{U}_{p-1}(\sh{M}|U)
  \simeq
  \shHom_{\sh{O}_Y|U}
    (\sh{Q}|U,\sh{M}|U).
\]
For arbitrary $Y$ and each affine open subset $U$, there is
a coherent $(\sh{O}_Y|U)$-Module $\sh{Q}_U$ and a
functorial isomorphism
\[
  \sh{T}^{U}_{p-1}(\sh{M}|U)
  \simeq
  \shHom_{\sh{O}_Y|U}
    (\sh{Q}_U,\sh{M}|U).
\]
The preceding observation shows that, whenever $V\subset U$
is affine, $\sh{Q}_U|V$ is canonically identified with
$\sh{Q}_V$.  Hence there exists a unique
$\sh{O}_Y$-Module $\sh{Q}$ satisfying
\eqref{III.7.7.5.2}.

It remains to prove the equivalence of a) and d).  Condition
d) is local on $Y$, and we have just seen that the same is
true of a); moreover, d) is also local on $Y'$.  When
$Y=\Spec(A)$ and $Y'=\Spec(A')$, the functor
$T^{A'}_\bullet$ is obtained from $T^A_\bullet$ by extension
of scalars to $A'$.  It is then clear that a$'$) implies
that \eqref{III.7.7.5.3} is an isomorphism.

\oldpage[III]{201}
Conversely, continue to suppose that $Y=\Spec(A)$ is affine
and take for $A'$ the $A$-algebra $A\oplus M$, where $M$ is
an arbitrary $A$-module and multiplication is given by
\[
  (a_1,m_1)(a_2,m_2)
  =
  (a_1a_2,a_1m_2+a_2m_1).
\]
Then
\[
  T^{A'}_p(A')
  =
  T_p(A\oplus M)
  =
  T_p(A)\oplus T_p(M).
\]
If \eqref{III.7.7.5.3} is bijective, the same is therefore
true of the canonical homomorphism
\[
  T_p(A)\otimes_A M\longrightarrow T_p(M).
\]
Thus d) implies a$''$), and the proof is complete.
\end{proof}

\begin{theorem}[7.7.6]
\label{III.7.7.6}
Let $Y$ be a locally Noetherian prescheme, let
$f:X\to Y$ be proper, and let $\sh{F}$ be a coherent,
$Y$-flat $\sh{O}_X$-Module.  There is a coherent
$\sh{O}_Y$-Module $\sh{Q}$, determined up to unique
isomorphism, and an isomorphism of functors in the
quasi-coherent $\sh{O}_Y$-Module $\sh{M}$,
\[
  f_*
  \bigl(\sh{F}\otimes_{\sh{O}_Y}\sh{M}\bigr)
  \simto
  \shHom_{\sh{O}_Y}(\sh{Q},\sh{M}).
  \tag{7.7.6.1}\label{III.7.7.6.1}
\]
Consequently there is an isomorphism of functors
\[
  \Gamma
  \bigl(X,\sh{F}\otimes_{\sh{O}_Y}\sh{M}\bigr)
  \simto
  \Hom_{\sh{O}_Y}(\sh{Q},\sh{M}).
  \tag{7.7.6.2}\label{III.7.7.6.2}
\]
\end{theorem}

\begin{proof}
The functor
$\sh{M}\mapsto
 \sh{F}\otimes_{\sh{O}_Y}\sh{M}$
is exact \sref[0]{0.6.7.4}, while $f_*$ is left exact.
Hence
\[
  \sh{M}\longmapsto
  f_*
  \bigl(\sh{F}\otimes_{\sh{O}_Y}\sh{M}\bigr)
\]
is left exact.  Apply the equivalence of b) and b$'$) in
\sref{III.7.7.5} with $p=1$.
\end{proof}
}

\begin{env}[7.7.2]
\label{III.7.7.2}
Let $g:Y'\to Y$ be a morphism, and set
\[
  X'=X_{(Y')}=X\times_Y Y',
  \qquad
  f'=f_{(Y')}:X'\longrightarrow Y'.
\]
Thus $f'$ is quasi-compact and separated.  Put also
\[
  \sh{P}'_\bullet
  =\sh{P}_\bullet\otimes_{\sh{O}_Y}\sh{O}_{Y'}.
\]
This is a complex of quasi-coherent $\sh{O}_{X'}$-Modules
that are $Y'$-flat, by \sref[I]{I.9.1.12} and
\sref[0]{0.6.2.1}.  With the same conventions on degrees, set
\[
  \sh{T}^{Y'}_\bullet(\sh{M}')
  =
  \mathscr{H}^\bullet
    \bigl(f',
      \sh{P}'^\bullet\otimes_{\sh{O}_{Y'}}\sh{M}'\bigr)
  =
  \mathscr{H}^\bullet
    \bigl(f',
      \sh{P}^\bullet\otimes_{\sh{O}_Y}\sh{M}'\bigr).
  \tag{7.7.2.1}\label{III.7.7.2.1}
\]
This is a homological functor of the quasi-coherent
$\sh{O}_{Y'}$-Module $\sh{M}'$.  When $Y'$ is an affine
scheme with ring $A'$, we write $\sh{T}^{A'}_\bullet$ in
place of $\sh{T}^{Y'}_\bullet$.  For every $A'$-module $M'$,
\[
  \sh{T}^{A'}_\bullet(\widetilde{M'})
  =
  \bigl(
    \Gamma(Y',
      \sh{T}^{Y'}_\bullet(\widetilde{M'}))
  \bigr)^\sim.
\]
We set
\[
  T^{A'}_\bullet(M')
  =
  \Gamma(Y',
    \sh{T}^{Y'}_\bullet(\widetilde{M'}));
\]
this is a homological functor from $A'$-modules to
$A'$-modules.

If $Y=\Spec(A)$ is also affine, the functor
$T^{A'}_\bullet$ on $A'$-modules is obtained from the
homological functor $T^A_\bullet$ on $A$-modules by extension
of scalars from $A$ to $A'$ \sref{III.7.1.3}.  Indeed, let
$g:Y'\to Y$ correspond to the ring homomorphism $A\to A'$,
and let $g':X'\to X$ be the corresponding morphism, which is
affine \sref[II]{II.1.6.2}.  If $\mathfrak U$ is an affine
open covering of $X$, then
$\mathfrak U'=g'^{-1}(\mathfrak U)$ is an affine open
covering of $X'$.  By \sref{III.6.2.2}, it is enough to check
\[
  C^\bullet
  \bigl(\mathfrak U,
    \sh{P}_\bullet\otimes_{\sh{O}_Y}g_*(\sh{M}')\bigr)
  =
  C^\bullet
  \bigl(\mathfrak U',
    \sh{P}_\bullet\otimes_{\sh{O}_Y}\sh{M}'\bigr).
\]
Ultimately, if $U$ is any affine open subset of $X$ and
$U'=g'^{-1}(U)$, this reduces to
\[
  \Gamma
  \bigl(U,
    \sh{P}_\bullet\otimes_{\sh{O}_Y}g_*(\sh{M}')\bigr)
  =
  \Gamma
  \bigl(U',
    \sh{P}_\bullet\otimes_{\sh{O}_Y}\sh{M}'\bigr),
\]
which is immediate from \sref[I]{I.1.3} and \sref{III.3.2}.

In particular, if $U$ is an open subset of $Y$, then for
every quasi-coherent $\sh{O}_Y$-Module $\sh{M}$,
\[
  \sh{T}^U_\bullet(\sh{M}|U)
  =
  \sh{T}_\bullet(\sh{M})|U.
  \tag{7.7.2.2}\label{III.7.7.2.2}
\]
\end{env}

\begin{env}[7.7.3]
\label{III.7.7.3}
For every quasi-coherent $\sh{O}_Y$-Module $\sh{M}$ there is
a canonical homomorphism, functorial in $\sh{M}$,
\[
  \sh{T}_p(\sh{O}_Y)
  \otimes_{\sh{O}_Y}\sh{M}
  \longrightarrow
  \sh{T}_p(\sh{M}).
  \tag{7.7.3.1}\label{III.7.7.3.1}
\]
Indeed, when $Y$ is affine this homomorphism was defined in
\sref{III.7.2.2}.  The definition extends immediately to
the general case: if $U$ and $V$ are affine open subsets of
$Y$ with $V\subset U$, the diagram
\oldpage[III]{199}
\[
\xymatrix@C=18pt@R=32pt{
  {\begin{gathered}
    \bigl(\sh{T}_p(\sh{O}_Y)
      \otimes_{\sh{O}_Y}\sh{M}\bigr)|U\\
    =
    \sh{T}^U_p(\sh{O}_Y|U)
      \otimes_{\sh{O}_Y|U}(\sh{M}|U)
  \end{gathered}}
  \ar[r]\ar[d]
  &
  {\begin{gathered}
    \sh{T}^U_p(\sh{M}|U)\\
    =\sh{T}_p(\sh{M})|U
  \end{gathered}}
  \ar[d]
  \\
  {\begin{gathered}
    \bigl(\sh{T}_p(\sh{O}_Y)
      \otimes_{\sh{O}_Y}\sh{M}\bigr)|V\\
    =
    \sh{T}^V_p(\sh{O}_Y|V)
      \otimes_{\sh{O}_Y|V}(\sh{M}|V)
  \end{gathered}}
  \ar[r]
  &
  {\begin{gathered}
    \sh{T}^V_p(\sh{M}|V)\\
    =\sh{T}_p(\sh{M})|V
  \end{gathered}}
}
\]
commutes by \sref{III.7.2.3.3}.

For every morphism $g:Y'\to Y$ there is a canonical
homomorphism
\[
  \sh{T}_p(\sh{O}_Y)
  \otimes_{\sh{O}_Y}\sh{O}_{Y'}
  \longrightarrow
  \sh{T}^{Y'}_p(\sh{O}_{Y'}),
  \tag{7.7.3.2}\label{III.7.7.3.2}
\]
which is precisely the special case of
\sref{III.6.7.11.2}, for the abutments, in which
$S=Y$, $v_1=f$, $v_2=1_Y$, and
$\sh{P}^{(2)}_\bullet$ is concentrated in degree $0$ with
sole term $\sh{M}$.

When $Y=\Spec(A)$ and $Y'=\Spec(A')$ are affine,
\eqref{III.7.7.3.2} is the homomorphism of sheaves
corresponding to the canonical homomorphism of $A'$-modules
defined in \sref{III.7.2.2},
\[
  T^A_p(A)\otimes_A A'
  \longrightarrow
  T^{A'}_p(A')
  =
  T^A_p(A').
\]
This follows immediately from \sref{III.6.7.11}, since in
the situation at hand one may take
$\mathfrak U'^{(i)}=u_i^{-1}(\mathfrak U^{(i)})$ there.
\end{env}

\begin{env}[7.7.4]
\label{III.7.7.4}
Suppose that $f$ is proper, that $Y=\Spec(A)$ is an affine
Noetherian scheme, and that $\sh{P}_\bullet$ is a
bounded-below complex of coherent, $Y$-flat
$\sh{O}_X$-Modules.  We saw in \sref{III.6.10.5} that, up to
isomorphism, one may write
\[
  \sh{T}_p(\sh{M})
  =
  \mathscr{H}_p
    \bigl(\sh{L}_\bullet
      \otimes_{\sh{O}_Y}\sh{M}\bigr),
  \qquad
  \sh{L}_\bullet=\widetilde{L_\bullet},
\]
where $L_\bullet$ is a bounded-below complex of
finite-type free $A$-modules.  Thus
$\sh{T}_\bullet$ is of the kind studied in detail in
\sref{III.7.4} and \sref{III.7.6}.  We shall now translate
the results of that study.
\end{env}

\ifdefined\EGAIIIStopAfterPageOneNineNine
  \expandafter\endinput
\fi

\egaIIISubsectionSevenSevenTheorems
% BEGIN NATIVE INPUT ega3/ega3-7-7-rest
\begin{corollary}[7.7.7]
\label{III.7.7.7}
Let $Y$ be a locally Noetherian prescheme, let
$f:X\to Y$ be proper, let $\sh{F}$ and $\sh{F}'$ be
coherent, $Y$-flat $\sh{O}_X$-Modules, and let
$u:\sh{F}\to\sh{F}'$ be a homomorphism.  Consider the two
functors of the quasi-coherent $\sh{O}_Y$-Module $\sh{M}$,
\begin{align*}
  \sh{T}(\sh{M})
  &=
  \Ker\left(
    f_*(\sh{F}\otimes_{\sh{O}_Y}\sh{M})
    \longrightarrow
    f_*(\sh{F}'\otimes_{\sh{O}_Y}\sh{M})
  \right),\\
  T(\sh{M})
  &=
  \Gamma(Y,\sh{T}(\sh{M}))\\
  &=
  \Ker\left(
    \Gamma(X,\sh{F}\otimes_{\sh{O}_Y}\sh{M})
    \longrightarrow
    \Gamma(X,\sh{F}'\otimes_{\sh{O}_Y}\sh{M})
  \right).
\end{align*}
There is a coherent $\sh{O}_Y$-Module $\sh{R}$, determined
up to unique isomorphism, and isomorphisms of functors
\[
  \sh{T}(\sh{M})
  \simto
  \shHom_{\sh{O}_Y}(\sh{R},\sh{M}),
  \tag{7.7.7.1}\label{III.7.7.7.1}
\]
\[
  T(\sh{M})
  \simto
  \Hom_{\sh{O}_Y}(\sh{R},\sh{M}).
  \tag{7.7.7.2}\label{III.7.7.7.2}
\]
\end{corollary}

\begin{proof}
It is enough to prove \eqref{III.7.7.7.2}.  This proves
\eqref{III.7.7.7.1} when $Y$ is affine, and the passage to
the general case proceeds as in the proof of the equivalence
of b) and b$'$) in \sref{III.7.7.5}, using the uniqueness,
up to unique isomorphism, of a representing object of a
representable functor \sref[0]{0.8.1.8}.

By \sref{III.7.7.6}, there are coherent
$\sh{O}_Y$-Modules $\sh{Q}$ and $\sh{Q}'$ and functorial
isomorphisms
\[
  \Gamma(X,\sh{F}\otimes_{\sh{O}_Y}\sh{M})
  \simto
  \Hom_{\sh{O}_Y}(\sh{Q},\sh{M}),
\]
\[
  \Gamma(X,\sh{F}'\otimes_{\sh{O}_Y}\sh{M})
  \simto
  \Hom_{\sh{O}_Y}(\sh{Q}',\sh{M}).
\]
The homomorphism $u:\sh{F}\to\sh{F}'$ canonically defines a
morphism of functors between the two left-hand sides.

\oldpage[III]{202}
It corresponds to a unique homomorphism
$v:\sh{Q}'\to\sh{Q}$ of $\sh{O}_Y$-Modules for which
\[
\xymatrix@C=42pt@R=28pt{
  \Gamma(X,\sh{F}\otimes_{\sh{O}_Y}\sh{M})
    \ar[r]\ar[d]_-{\sim}
  &
  \Gamma(X,\sh{F}'\otimes_{\sh{O}_Y}\sh{M})
    \ar[d]^-{\sim}
  \\
  \Hom_{\sh{O}_Y}(\sh{Q},\sh{M})
    \ar[r]
  &
  \Hom_{\sh{O}_Y}(\sh{Q}',\sh{M})
}
\]
commutes \sref[0]{0.8.1.4}.  The contravariant functor
\[
  \sh{N}\longmapsto
  \Hom_{\sh{O}_Y}(\sh{N},\sh{M})
\]
is left exact in the category of $\sh{O}_Y$-Modules.
Taking $\sh{R}=\Coker(v)$ therefore gives
\eqref{III.7.7.7.2}.
\end{proof}

\begin{corollary}[7.7.8]
\label{III.7.7.8}
Under the hypotheses of \sref{III.7.7.6} on $X$, $Y$, and
$f$, let $\sh{F}$ and $\sh{G}$ be coherent
$\sh{O}_X$-Modules satisfying:
\begin{enumerate}
  \item[(i)] $\sh{F}$ is $Y$-flat;
  \item[(ii)] $\sh{G}$ is isomorphic to the cokernel of a
    homomorphism
    $\sh{E}_1\to\sh{E}_0$ between finite-type locally free
    $\sh{O}_X$-Modules.
\end{enumerate}
Consider the two functors of the quasi-coherent
$\sh{O}_Y$-Module $\sh{M}$,
\begin{align*}
  \sh{T}(\sh{M})
  &=
  f_*\shHom_{\sh{O}_X}
  \bigl(
    \sh{G},
    \sh{F}\otimes_{\sh{O}_Y}\sh{M}
  \bigr),\\
  T(\sh{M})
  &=
  \Gamma(Y,\sh{T}(\sh{M}))\\
  &=
  \Hom_{\sh{O}_X}
  \bigl(
    \sh{G},
    \sh{F}\otimes_{\sh{O}_Y}\sh{M}
  \bigr).
\end{align*}
There is a coherent $\sh{O}_Y$-Module $\sh{N}$, determined
up to unique isomorphism, and isomorphisms of functors
\[
  \sh{T}(\sh{M})
  \simto
  \shHom_{\sh{O}_Y}(\sh{N},\sh{M}),
  \tag{7.7.8.1}\label{III.7.7.8.1}
\]
\[
  T(\sh{M})
  \simto
  \Hom_{\sh{O}_Y}(\sh{N},\sh{M}).
  \tag{7.7.8.2}\label{III.7.7.8.2}
\]
\end{corollary}

\begin{proof}
The functorial isomorphism \sref[0]{0.5.4.2.1} gives,
functorially in $\sh{M}$,
\begin{align*}
  \shHom_{\sh{O}_X}
    \bigl(\sh{E}_i,
      \sh{F}\otimes_{\sh{O}_Y}\sh{M}\bigr)
  &\simeq
  \sh{E}_i^\vee
    \otimes_{\sh{O}_X}
    \bigl(\sh{F}\otimes_{\sh{O}_Y}\sh{M}\bigr)\\
  &\simeq
  \bigl(
    \sh{E}_i^\vee\otimes_{\sh{O}_X}\sh{F}
  \bigr)
  \otimes_{\sh{O}_Y}\sh{M}\\
  &\simeq
  \shHom_{\sh{O}_X}(\sh{E}_i,\sh{F})
  \otimes_{\sh{O}_Y}\sh{M}
\end{align*}
for $i=0,1$.  Put
\[
  \sh{F}_i
  =
  \shHom_{\sh{O}_X}(\sh{E}_i,\sh{F})
  \qquad(i=0,1).
\]
These are coherent \sref[0]{0.5.3.5} and $Y$-flat
\sref[0]{0.5.4.2}.  If
$v:\sh{E}_1\to\sh{E}_0$ is the given homomorphism, put
\[
  u=\shHom(v,1_{\sh{F}}):\sh{F}_0\longrightarrow\sh{F}_1.
\]
Since
\[
  \sh{H}\longmapsto
  \shHom_{\sh{O}_X}
  \bigl(\sh{H},
    \sh{F}\otimes_{\sh{O}_Y}\sh{M}\bigr)
\]
is left exact, there are functorial isomorphisms
\begin{align*}
  \shHom_{\sh{O}_X}
  \bigl(\sh{G},
    \sh{F}\otimes_{\sh{O}_Y}\sh{M}\bigr)
  &\simeq
  \Ker\left(
    \shHom_{\sh{O}_X}
      \bigl(\sh{E}_0,
        \sh{F}\otimes_{\sh{O}_Y}\sh{M}\bigr)
    \longrightarrow
    \shHom_{\sh{O}_X}
      \bigl(\sh{E}_1,
        \sh{F}\otimes_{\sh{O}_Y}\sh{M}\bigr)
  \right)\\
  &\simeq
  \Ker\left(
    \sh{F}_0\otimes_{\sh{O}_Y}\sh{M}
    \longrightarrow
    \sh{F}_1\otimes_{\sh{O}_Y}\sh{M}
  \right).
\end{align*}
As $f_*$ is left exact, it follows that
\[
  f_*\shHom_{\sh{O}_X}
  \bigl(\sh{G},
    \sh{F}\otimes_{\sh{O}_Y}\sh{M}\bigr)
  \simeq
  \Ker\left(
    f_*(\sh{F}_0\otimes_{\sh{O}_Y}\sh{M})
    \longrightarrow
    f_*(\sh{F}_1\otimes_{\sh{O}_Y}\sh{M})
  \right).
\]
The assertion now follows from \sref{III.7.7.7}.
\end{proof}

\oldpage[III]{203}
\begin{remarks}[7.7.9]
\label{III.7.7.9}
\textnormal{(i)}
In \sref{III.7.7.6}, \sref{III.7.7.7}, and
\sref{III.7.7.8}, the formation of the $\sh{O}_Y$-Modules
$\sh{Q}$, $\sh{R}$, and $\sh{N}$ commutes with base change.
For example, retaining the notation of \sref{III.7.7.2}, in
the situation of \sref{III.7.7.6} one has, for every
quasi-coherent $\sh{O}_{Y'}$-Module $\sh{M}'$,
\[
  f'_*
  \bigl(\sh{F}\otimes_{\sh{O}_Y}\sh{M}'\bigr)
  \simeq
  \shHom_{\sh{O}_{Y'}}
  \bigl(g^*(\sh{Q}),\sh{M}'\bigr).
\]
Indeed, by the observation in \sref{III.7.7.2}, it is enough
to use
\[
  \Hom_{\sh{O}_Y}
  \bigl(\sh{Q},g_*(\sh{M}')\bigr)
  =
  \Hom_{\sh{O}_{Y'}}
  \bigl(g^*(\sh{Q}),\sh{M}'\bigr),
\]
which is \sref[0]{0.4.4.3.1}.

Likewise, if in \sref{III.7.7.7} one replaces
$Y,f,\sh{M},\sh{F},\sh{F}'$ by
\[
  Y',f',\sh{M}',
  \sh{F}\otimes_{\sh{O}_X}\sh{O}_{X'},
  \sh{F}'\otimes_{\sh{O}_X}\sh{O}_{X'},
\]
then one must replace $\sh{R}$ by $g^*(\sh{R})$.
Finally, in \sref{III.7.7.8}, if
$X,Y,f,\sh{M},\sh{F}$ are replaced by
\[
  X',Y',f',\sh{M}',
  \sh{F}\otimes_{\sh{O}_X}\sh{O}_{X'},
\]
and $\sh{G}$ by
\[
  \sh{G}'
  =
  \sh{G}\otimes_{\sh{O}_X}\sh{O}_{X'},
\]
then $\sh{N}$ must be replaced by $g^*(\sh{N})$.  This
follows because there remains an exact sequence
\[
  \sh{E}'_1\longrightarrow
  \sh{E}'_0\longrightarrow
  \sh{G}'\longrightarrow0,
  \qquad
  \sh{E}'_i
  =
  \sh{E}_i\otimes_{\sh{O}_X}\sh{O}_{X'},
\]
and because, for $i=0,1$,
\[
  (\sh{E}'_i)^\vee
  \otimes_{\sh{O}_{X'}}
  \bigl(\sh{F}\otimes_{\sh{O}_Y}\sh{M}'\bigr)
  =
  \sh{E}_i^\vee
  \otimes_{\sh{O}_X}
  \bigl(\sh{F}\otimes_{\sh{O}_Y}\sh{M}'\bigr).
\]

\medskip
\textnormal{(ii)}
Condition (ii) in \sref{III.7.7.8} is satisfied by every
coherent $\sh{O}_X$-Module $\sh{G}$ whenever there exists a
$Y$-ample invertible $\sh{O}_X$-Module; this holds, for
example, when $Y$ is affine and $f:X\to Y$ is projective.
Indeed, by \sref[II]{II.5.5.1} and \sref[II]{II.2.7.10},
there is a finite-type locally free $\sh{O}_X$-Module
$\sh{E}_0$ with $\sh{G}$ isomorphic to a quotient of
$\sh{E}_0$.  Since $\sh{E}_0$ and $\sh{G}$ are coherent, so
is the kernel $\sh{G}_1$ of $\sh{E}_0\to\sh{G}$.  Applying
the same result to $\sh{G}_1$ gives an exact sequence
\[
  \sh{E}_1\longrightarrow
  \sh{E}_0\longrightarrow
  \sh{G}\longrightarrow0
\]
with $\sh{E}_0$ and $\sh{E}_1$ finite-type locally free.

\medskip
\textnormal{(iii)}
We shall prove in Chapter V that the restrictive hypothesis
(ii) of \sref{III.7.7.8} is redundant.
\end{remarks}

\begin{proposition}[7.7.10]
\label{III.7.7.10}
\textnormal{(Local criteria for the exchange property).}
Under the general hypotheses of \sref{III.7.7.5}, let
$y\in Y$ and let $p$ be an integer.  The following
conditions are equivalent:
\begin{enumerate}
  \item[a)] The functor $T_p^{\sh{O}_y}$ is right exact.

  \item[b)] The canonical homomorphism
    \[
      T_p^{\sh{O}_y}(\sh{O}_y)
      \longrightarrow
      T_p^{\sh{O}_y}(\kres(y))
    \]
    is surjective.

  \item[c)] For every integer $n$, the canonical
    homomorphism
    \[
      T_p^{\sh{O}_y}
      (\sh{O}_y/\mathfrak m_y^{n+1})
      \longrightarrow
      T_p^{\sh{O}_y}(\kres(y))
    \]
    is surjective.
\end{enumerate}
Moreover, the set of points $y\in Y$ satisfying these
conditions is the largest open subset $U$ of $Y$ for which
$\sh{T}^U_p$ is right exact.
\end{proposition}

\begin{proof}
In view of \sref{III.7.7.4}, the equivalence of a), b), and
c) follows from \sref{III.7.4.7} and \sref{III.7.5.2}.
That the locus $U$ on which $T_p^{\sh{O}_y}$ is right exact
is open follows from \sref{III.7.4.4}.  Conversely, if
$\sh{T}^V_p$ is right exact, then
$T_p^{\sh{O}_y}$ is right exact for every $y\in V$, by
condition c) of \sref{III.7.7.5} and \sref{III.7.6.1}.
\end{proof}

\begin{corollary}[7.7.11]
\label{III.7.7.11}
If $\sh{T}_p$ is right exact (respectively left exact), then
for every morphism $g:Y'\to Y$ the functor
$\sh{T}^{Y'}_p$ is right exact (respectively left exact).
The converse holds when $g$ is faithfully flat.
\end{corollary}

\begin{proof}
\oldpage[III]{204}
The first assertion follows immediately from
\sref{III.7.6.1} and from the fact that the question is
local on $Y$ and $Y'$, by conditions c) and b) of
\sref{III.7.7.5}.  For the converse, it is enough, by
\sref{III.7.7.10}, to show that $T_p^{\sh{O}_y}$ is right
exact (respectively left exact) for every $y\in Y$.  By
hypothesis there is a point $y'\in Y'$ with $g(y')=y$, and
$\sh{O}_{y'}$ is faithfully flat over $\sh{O}_y$.  The
conclusion follows from the hypothesis and
\sref{III.7.6.1}.
\end{proof}

\begin{remarks}[7.7.12]
\label{III.7.7.12}
\textnormal{(i)}
Under the hypotheses of \sref{III.7.7.4}, suppose in
addition that $\sh{P}_\bullet$ is a finite complex.  Since
one may choose a finite affine open covering $\mathfrak U$
of $X$, it follows from \sref{III.7.7.1} that the
bicomplex
\[
  C^\bullet
  \bigl(
    \mathfrak U,
    \sh{P}^\bullet\otimes_{\sh{O}_Y}\sh{M}
  \bigr)
\]
is finite as well.  More precisely, there is a finite set
$E$ of pairs, independent of $\sh{M}$, such that
\[
  C^h
  \bigl(
    \mathfrak U,
    \sh{P}^k\otimes_{\sh{O}_Y}\sh{M}
  \bigr)
  =0
  \qquad\text{for }(h,k)\notin E.
\]
Consequently there is an integer $i_1$ such that
\[
  \sh{T}_i(\sh{M})=0
  \qquad(i\geq i_1)
\]
for every quasi-coherent $\sh{O}_Y$-Module $\sh{M}$.  Thus
$\sh{T}_i$ is trivially exact for these values of $i$, and
\sref{III.7.4.1} shows that $Z'_i(L_\bullet)$ is a flat
finite-type $A$-module, hence a finite-type projective
$A$-module since $A$ is Noetherian.

Define a complex $L'_\bullet$ of $A$-modules by
\[
  L'_i=L_i\quad(i<i_1),
  \qquad
  L'_{i_1}=Z'_{i_1}(L_\bullet),
  \qquad
  L'_i=0\quad(i>i_1),
\]
and put $\sh{L}'_\bullet=\widetilde{L'_\bullet}$.  Plainly
\[
  \mathscr{H}_i
  (\sh{L}'_\bullet\otimes_{\sh{O}_Y}\sh{M})
  =
  \mathscr{H}_i
  (\sh{L}_\bullet\otimes_{\sh{O}_Y}\sh{M})
\]
for $i<i_1-1$ and for $i\geq i_1$, the two sides being zero
in the latter range.  Since
\[
  \operatorname{Im}
  (Z'_{i_1}\otimes_A M)
  =
  \operatorname{Im}
  (L_{i_1}\otimes_A M)
\]
by definition, the same equality holds for $i=i_1-1$.
Thus in \sref{III.7.7.4} one may also take
$\sh{L}_\bullet$ to be finite, provided only that its terms
are locally free $\sh{O}_Y$-Modules associated to
finite-type projective $A$-modules.

This applies in particular when $\sh{P}_\bullet$ has only
one nonzero term $\sh{F}$, in degree $0$.  In that case
\[
  \sh{T}_n(\sh{M})
  =
  R^{-n}f_*
  \bigl(\sh{F}\otimes_{\sh{O}_Y}\sh{M}\bigr),
\]
and one may take $\sh{L}_i=0$ for $i>0$.  In this case it is
preferable to use cohomological notation, writing
$\sh{T}^{-p}$ in place of $\sh{T}_p$.

\medskip
\textnormal{(ii)}
If the terms $\sh{P}_i$ in \sref{III.7.7.5} are no longer
assumed $Y$-flat, the conclusions remain valid provided one
defines
\[
  \sh{T}_p(\sh{M})
  =
  \mathbf{Tor}^Y_p
  \bigl(f,1_Y;\sh{P}_\bullet,\sh{M}\bigr).
  \tag{7.7.12.1}\label{III.7.7.12.1}
\]
Indeed,
\[
  \mathbf{Tor}^Y_n
  \bigl(f,1_Y;\sh{P}_\bullet,\sh{O}_Y\bigr)
\]
is then a coherent $\sh{O}_Y$-Module by
\sref{III.6.7.9}.  The proof of \sref{III.6.10.5} applies
without change, in view of \sref{III.6.10.1}, and also shows
that, when $Y=\Spec(A)$ is affine,
\[
  \sh{T}_p(\sh{M})
  =
  \mathscr{H}_p
  \bigl(
    \sh{L}_\bullet\otimes_{\sh{O}_Y}\sh{M}
  \bigr),
  \qquad
  \sh{L}_\bullet=\widetilde{L_\bullet},
\]
where $L_\bullet$ is a complex of finite-type free
$A$-modules.  This proves the assertion.
\end{remarks}
% END NATIVE INPUT ega3/ega3-7-7-rest
% END NATIVE INPUT ega3/ega3-7-7
\ifdefined\EGAIIIStopAfterPageOneNineNine
  \expandafter\endinput
\fi
% BEGIN NATIVE INPUT ega3/ega3-7-8
\subsection{Application to proper morphisms: II. Cohomological flatness criteria}
\label{subsection:III.7.8}
\label{III.7.8}

\begin{definition}[7.8.1]
\label{III.7.8.1}
Let $X$ and $Y$ be two preschemes, let $f:X\to Y$ be a
quasi-compact separated morphism, let $\sh{P}_\bullet$ be a
complex of quasi-coherent, $Y$-flat $\sh{O}_X$-Modules, and
let $\sh{T}_\bullet$ be the homological functor on
quasi-coherent $\sh{O}_Y$-Modules defined by
\eqref{III.7.7.1.1}.  Let $y$ be a point of $Y$.

\oldpage[III]{205}
We say that $\sh{P}_\bullet$ is \emph{homologically flat over
$Y$ at $y$, in dimension $p$} (or \emph{cohomologically flat
over $Y$ at $y$, in dimension $-p$}) if there is an open
neighbourhood $U$ of $y$ in $Y$ such that
\[
  \sh{T}_p^U=\sh{T}^{-p}_U
\]
is exact.

We say that $\sh{P}_\bullet$ is \emph{homologically flat in
dimension $p$ over $Y$} (or \emph{cohomologically flat in
dimension $-p$ over $Y$}) if it is homologically flat over
$Y$ at every point $y\in Y$, in dimension $p$.

When $\sh{P}_\bullet$ is homologically flat over $Y$
(respectively over $Y$ at $y$) in \emph{every} dimension
$p$, we say simply that $\sh{P}_\bullet$ is
\emph{homologically flat over $Y$} (respectively over $Y$ at
$y$), or \emph{cohomologically flat over $Y$} (respectively
over $Y$ at $y$).
\end{definition}

\begin{env}[7.8.2]
\label{III.7.8.2}
By definition, the notion of homological flatness over $Y$
is \emph{local on $Y$}.

If $Y$ is locally Noetherian, or if $Y$ is a scheme, then
for $\sh{P}_\bullet$ to be homologically flat over $Y$ in
dimension $p$, it is necessary and sufficient that the
functor $\sh{T}_p$ be exact.  The proof for locally
Noetherian $Y$ was given in the course of the proof of
\sref{III.7.7.5}; when $Y$ is a scheme the argument is the
same, based on \sref[I]{I.9.4.2} applied to an affine open
subset of a quasi-compact scheme.
\end{env}

\begin{proposition}[7.8.3]
\label{III.7.8.3}
With the notation and hypotheses of \sref{III.7.8.1}, the
following conditions are equivalent:
\begin{enumerate}
  \item[a)] $\sh{P}_\bullet$ is homologically flat over $Y$
  at $y$, in dimension $p$.
  \item[b)] There is an open neighbourhood $U$ of $y$ in
  $Y$ such that $\sh{T}_p^U$ and $\sh{T}_{p+1}^U$ are right
  exact.
  \item[c)] There is an open neighbourhood $U$ of $y$ in
  $Y$ such that $\sh{T}_p^U$ and $\sh{T}_{p-1}^U$ are left
  exact.
  \item[d)] There is an open neighbourhood $U$ of $y$ in
  $Y$ such that $\sh{T}_{p+1}^U$ is right exact and
  $\sh{T}_{p-1}^U$ is left exact.
\end{enumerate}
\end{proposition}

\begin{proof}
In view of the interpretation of $\sh{T}_p$ when $Y$ is
affine, this is merely a translation of part of
\sref{III.7.3.3}.
\end{proof}

\begin{proposition}[7.8.4]
\label{III.7.8.4}
Let $Y$ be a locally Noetherian prescheme, let
$f:X\to Y$ be a proper morphism, let $\sh{P}_\bullet$ be a
bounded-below complex of coherent, $Y$-flat
$\sh{O}_X$-Modules, and let $\sh{T}_\bullet$ be the functor
defined by \eqref{III.7.7.1.1}.  For every $y\in Y$, the
following conditions are equivalent:
\begin{enumerate}
  \item[a)] $\sh{P}_\bullet$ is homologically flat over $Y$
  at $y$, in dimension $p$.
  \item[b)] The functor $\sh{T}_p^{\sh{O}_y}$ is exact.
  \item[c)] There is an integer $n_0$ such that, for
  $n\geq n_0$,
  \[
    \operatorname{length}
      \sh{T}_p^{\sh{O}_y}
        \bigl(\sh{O}_y/\mathfrak m_y^{\,n+1}\bigr)
    =
    \operatorname{length}
      \sh{T}_p^{\sh{O}_y}\bigl(\kres(y)\bigr)
    \cdot
    \operatorname{length}
      \bigl(\sh{O}_y/\mathfrak m_y^{\,n+1}\bigr).
    \tag{7.8.4.1}\label{III.7.8.4.1}
  \]
  Here all lengths are lengths of $\sh{O}_y$-modules.
  \item[d)] There is an open neighbourhood $U$ of $y$ such
  that
  \[
    \mathscr{H}^{-p}(f,\sh{P}^\bullet)|U
  \]
  is isomorphic to an $(\sh{O}_Y|U)$-Module of the form
  $(\sh{O}_Y|U)^m$ and, for every quasi-coherent
  $(\sh{O}_Y|U)$-Module $\sh{M}$, the canonical
  homomorphism
  \[
    \mathscr{H}^{-p}(f,\sh{P}^\bullet)|U
      \otimes_{\sh{O}_Y|U}\sh{M}
    \longrightarrow
    \mathscr{H}^{-p}
      \bigl(f,(\sh{P}^\bullet|U)
        \otimes_{\sh{O}_Y|U}\sh{M}\bigr)
    \tag{7.8.4.2}\label{III.7.8.4.2}
  \]
  is bijective.
\end{enumerate}
When these conditions hold, one also has:
\begin{enumerate}
  \item[e)] There is a neighbourhood of $y$ on which the
  function $z\mapsto d_p(z)$, defined in
  \eqref{III.7.7.5.1}, is constant.
\end{enumerate}
Moreover, if $Y$ is reduced at $y$ \sref[0]{0.4.1.4},
condition e) is equivalent to the other conditions.
\end{proposition}

\begin{proof}
\oldpage[III]{206}
Indeed, condition b) is equivalent to saying that
$\sh{T}_p^{\sh{O}_y}$ and $\sh{T}_{p+1}^{\sh{O}_y}$ are
right exact \sref{III.7.3.3}.  The equivalence of a) and b)
then follows from \sref{III.7.7.10} and
\sref{III.7.8.3}.

Since $\sh{O}_y/\mathfrak m_y^{\,n+1}$ is Artinian, and
\[
  \sh{T}_p^{\sh{O}_y}
    \bigl(\sh{O}_y/\mathfrak m_y^{\,n+1}\bigr)
  \quad\hbox{and}\quad
  \sh{T}_p^{\sh{O}_y}\bigl(\kres(y)\bigr)
\]
are finite-type
$(\sh{O}_y/\mathfrak m_y^{\,n+1})$-modules
\sref{III.7.7.4}, hence have finite length, the equivalence
of b) and c) follows again from \sref{III.7.7.10} and
\sref{III.7.3.5.7}.

That a) implies e), and that a) and e) are equivalent when
$\sh{O}_y$ is reduced, follows from \sref{III.7.6.9}.
Finally, a) implies that \eqref{III.7.8.4.2} is bijective,
by the definition \sref{III.7.8.1} and
\sref{III.7.7.5}.  It also implies that
\[
  \bigl(\mathscr{H}^{-p}(f,\sh{P}^\bullet)\bigr)_y
\]
is a flat $\sh{O}_y$-module by condition (f) of
\sref{III.7.3.3}, hence a
free $\sh{O}_y$-module of finite type
\sref[0]{0.10.1.3}, since it is of finite type by
\sref{III.7.7.4}.  As
$\mathscr{H}^{-p}(f,\sh{P}^\bullet)$ is a coherent
$\sh{O}_Y$-Module \sref{III.7.7.4}, it is locally free in
a neighbourhood of $y$ \sref[0]{0.5.2.7}.
Conversely, d) plainly implies a), by the definition of the
functor $\sh{T}_p^U$ \eqref{III.7.7.2.2}.
\end{proof}

\begin{proposition}[7.8.5]
\label{III.7.8.5}
Under the hypotheses of \sref{III.7.8.4}, the following
conditions are equivalent:
\begin{enumerate}
  \item[a)] $\sh{P}_\bullet$ is homologically flat over $Y$
  in every dimension $i\leq p$.
  \item[b)] For $i\leq p+1$, the functors $\sh{T}_i$ are
  right exact.
  \item[c)] For $i\geq -p$, the $\sh{O}_Y$-Modules
  $\mathscr{H}^{i}(f,\sh{P}^\bullet)$ are locally free.
\end{enumerate}
\end{proposition}

\begin{proof}
The equivalence of a) and b) is immediate from
\sref{III.7.8.3}, and a) implies c) by
\sref{III.7.8.4}.

Conversely, suppose that c) holds.  Recall also that
$\sh{L}_i=0$ for $i<i_0$ \sref{III.7.7.4}, and hence
$\sh{T}_i=0$ for $i\leq i_0$.  Every point $y\in Y$
therefore has an affine neighbourhood
$U=\Spec(A)$ such that $T_i^A(A)$ is a free $A$-module for
$i\leq p$.  It follows from \sref{III.7.3.7} that
$\sh{T}_i^U=T_i^A$ is exact for $i\leq p$.
\end{proof}

We shall chiefly apply the criteria for cohomological
flatness to the case where the complex $\sh{P}_\bullet$ is
concentrated in a \emph{single} coherent
$\sh{O}_X$-Module $\sh{F}$, flat over $Y$, taken as
$\sh{P}_0$.  Recall that in this case
\[
  \sh{T}_p(\sh{M})
  =R^{-p}f_*
    \bigl(\sh{F}\otimes_{\sh{O}_Y}\sh{M}\bigr).
\]

\begin{proposition}[7.8.6]
\label{III.7.8.6}
Let $Y$ be a locally Noetherian prescheme, let
$f:X\to Y$ be a proper flat morphism, and let $y$ be a point
of $Y$.  Denote the fibre
$f^{-1}(y)=X\otimes_Y\kres(y)$ by $X_y$.  Suppose that
\[
  \Gamma(X_y,\sh{O}_{X_y})=R
\]
is a separable $\kres(y)$-algebra (Bourbaki,
\emph{Alg\`ebre}, chap.~VIII, \S~7, no.~5), that is, a
direct sum of finitely many finite-degree separable
extensions of $\kres(y)$.  Then $\sh{O}_X$ is
cohomologically flat over $Y$ at $y$, in dimension $0$.
\end{proposition}

\begin{proof}
By \sref{III.7.8.4}, we may restrict to the case where
$Y$ is the spectrum of the local ring $A=\sh{O}_y$.
The flatness of $f$ implies $\sh{T}_{-1}=0$, so that
$T_0^A$ is already left exact and it remains only to prove
that it is right exact.  By condition c) of
\sref{III.7.7.10}, we may even reduce to the case where
$A=\sh{O}_y$ is Artinian.

Let $k'$ be a finite extension of $\kres(y)$ which splits
$R$, so that
\[
  R\otimes_{\kres(y)}k'
\]
is a direct sum of finitely many fields isomorphic to $k'$.
There is a local homomorphism from $A$ to a local ring
$A'$ such that $A'$ is a finite free $A$-algebra and its
residue field is isomorphic to $k'$ \sref[0]{0.10.3.2}.
By \sref{III.7.6.1}, it is therefore enough to prove that
$T_0^{A'}$ is right exact; in other words, we may suppose
that $R$ is a direct sum of $m$ fields isomorphic to
$\kres(y)$.

We shall use the following elementary lemma.

\begin{lemma}[7.8.6.1]
\label{III.7.8.6.1}
Let $Z$ be a locally ringed space.  For $Z$ to be connected,
\oldpage[III]{207}
it is necessary and sufficient that the ring
$\Gamma(Z,\sh{O}_Z)$ not be a product of two nonzero rings.
\end{lemma}

\begin{proof}
It is clear that if $Z$ is the union of two disjoint
nonempty open subsets $Z_1$ and $Z_2$, then
$\Gamma(Z,\sh{O}_Z)$ is isomorphic to the product of the
two nonzero rings $\Gamma(Z_1,\sh{O}_Z)$ and
$\Gamma(Z_2,\sh{O}_Z)$.

Conversely, to say that $\Gamma(Z,\sh{O}_Z)$ is such a
product is equivalent to saying that it contains an
idempotent $s$ distinct from $0$ and $1$.  For every
$z\in Z$, the germ $s_z$ is then an idempotent in the local
ring $\sh{O}_z$, and hence equals $0$ or $1$.  The set of
$z$ for which $s_z=0$ is open.  On the other hand,
$s_z=1$ implies, by definition, that $s(z)\ne0$, and the
set of $z$ for which $s_z=1$ is open as well
\sref[0]{0.5.5.2}.  The conclusion follows.
\end{proof}

It follows from the lemma that $X_y$ has exactly $m$
connected components $X_i'$ and that
\[
  \Gamma(X_i',\sh{O}_{X_i'})=\kres(y)
\]
for every $i$.  Since $A$ is local and Artinian, its
spectrum consists of one point, and consequently $X$ and
$X_y$ have the same underlying space.  Thus $X$ has $m$
connected components $X_i$ such that
\[
  X_i'=X_i\otimes_Y\kres(y).
\]
We are finally reduced to the case $R=\kres(y)$.  By
condition b) of \sref{III.7.7.10}, it remains to prove that
the canonical homomorphism
\[
  \Gamma(X,\sh{O}_X)\longrightarrow
  \Gamma(X_y,\sh{O}_{X_y})
\]
is surjective.  This is immediate, since the composite
\[
  \Gamma(Y,\sh{O}_Y)=A
  \longrightarrow\Gamma(X,\sh{O}_X)
  \longrightarrow\Gamma(X_y,\sh{O}_{X_y})=\kres(y)
\]
is already surjective.
\end{proof}

\begin{corollary}[7.8.7]
\label{III.7.8.7}
Under the hypotheses of \sref{III.7.8.6}, there is an open
neighbourhood $U$ of $y$ such that:
\begin{enumerate}
  \item[(i)] $f_*(\sh{O}_X)|U$ is isomorphic to an
  $(\sh{O}_Y|U)$-Module of the form $(\sh{O}_Y|U)^m$;
  \item[(ii)] for every $z\in U$, the canonical
  homomorphism
  \[
    \bigl(f_*(\sh{O}_X)\bigr)_z
      \otimes_{\sh{O}_z}\kres(z)
    \longrightarrow
    \Gamma(X_z,\sh{O}_{X_z})
  \]
  is bijective.
\end{enumerate}
\end{corollary}

\begin{proof}
Assertion (i) follows from \sref{III.7.8.6} and
\sref{III.7.8.4}.  Assertion (ii) follows from the
exactness of $\sh{T}_0^U$, for a suitable choice of $U$,
and from \eqref{III.7.7.5.3}.
\end{proof}

\begin{corollary}[7.8.8]
\label{III.7.8.8}
Suppose that the hypotheses of \sref{III.7.8.6} hold and,
in addition, that
\[
  \Gamma(X_y,\sh{O}_{X_y})=\kres(y).
\]
Then there is an open neighbourhood $U$ of $y$ such that
the canonical homomorphism
\[
  \sh{O}_Y|U\longrightarrow f_*(\sh{O}_X)|U
\]
is bijective.
\end{corollary}

\begin{proof}
Indeed, it follows from assertion (ii) of
\sref{III.7.8.7} that the integer $m$ occurring in
assertion (i) of that corollary is necessarily equal to
$1$.
\end{proof}

\begin{corollary}[7.8.9]
\label{III.7.8.9}
Under the hypotheses of \sref{III.7.8.6}, there are an open
neighbourhood $U$ of $y$, a coherent $\sh{O}_U$-Module
$\sh{Q}$, determined up to unique isomorphism, and an
isomorphism of functors in the quasi-coherent
$\sh{O}_U$-Module $\sh{M}$,
\[
  R^1f_*\bigl(f^*(\sh{M})\bigr)
  \simeq
  \shHom_{\sh{O}_U}(\sh{Q},\sh{M}).
  \tag{7.8.9.1}\label{III.7.8.9.1}
\]
\end{corollary}

\begin{proof}
The hypothesis implies that $\sh{T}_0^U$ is exact for a
suitable $U$.  It is therefore enough to apply the
equivalence of conditions a) and b') in
\sref{III.7.7.5}, with $p=0$ and with
$\sh{P}_\bullet$ the complex concentrated in degree $0$
with term $\sh{O}_X$.
\end{proof}

\begin{remarks}[7.8.10]
\label{III.7.8.10}
\textnormal{(i)}
Under the hypotheses of \sref{III.7.8.6}, consider the
Stein factorisation of $f$ \sref{III.4.3.3},
\[
  \xymatrix{
    X \ar[r]^{f'} & Y' \ar[r]^{g} & Y .
  }
\]
\oldpage[III]{208}
Here
\[
  Y'=\Spec\bigl(f_*(\sh{O}_X)\bigr).
\]
The finite morphism $g$ is such that
\[
  g_*(\sh{O}_{Y'})=f_*(\sh{O}_X)
\]
is locally free near $y$, and its fibre over $y$ is the
spectrum of a separable algebra over $\kres(y)$
\sref[II]{II.1.5.1}.  We shall deduce in Chapter~IV that
there is an open neighbourhood $U$ of $y$ in $Y$ such that,
for the restriction $g^{-1}(U)\to U$ of $g$, \emph{every}
fibre $g^{-1}(z)$, $z\in U$, is the spectrum of a separable
algebra over $\kres(z)$.  This is what we shall call an
\emph{\'etale covering} of $U$.  It will then follow from
assertion (ii) of \sref{III.7.8.7} that the hypothesis on
the point $y$ in \sref{III.7.8.6} also holds at every point
of a neighbourhood of $y$.

\medskip
\textnormal{(ii)}
We shall see in Chapter~V that, even if $X$ is projective
over $Y$ (and even if it is also ``simple'' over $Y$, a
property to be defined in Chapter~IV), the
$\sh{O}_U$-Module $\sh{Q}$ of \sref{III.7.8.9} need not be
locally free.  In other words, under these hypotheses
$\sh{O}_X$ need not be cohomologically flat in
\emph{dimension $1$} over $Y$ at $y$.  In Chapter~V we
shall interpret $\sh{Q}$ as the sheaf of
$1$-differentials of the Picard scheme of $X$ relative to
$Y$, along the unit section.
\end{remarks}
% END NATIVE INPUT ega3/ega3-7-8
\ifdefined\EGAIIIStopAfterPageTwoZeroEight
  \expandafter\endinput
\fi
% BEGIN NATIVE INPUT ega3/ega3-7-9
\oldpage[III]{208}
\subsection{Application to proper morphisms: III. Invariance of the
Euler--Poincar\'e characteristic and the Hilbert polynomial}
\label{subsection:III.7.9}
\label{III.7.9}

\begin{env}[7.9.1]
\label{III.7.9.1}
Let $A$ be a ring and let $M$ be a projective $A$-module of
finite type.  Recall (Bourbaki, \emph{Alg.~comm.}, chap.~II,
\S5, no.~2) that this is equivalent to saying that the
$\sh{O}_X$-Module $\widetilde M$ associated with $M$ on
$X=\Spec(A)$ is locally free of finite type.  For every
$\mathfrak p\in\Spec(A)$, the \emph{rank of $M$ at
$\mathfrak p$}, denoted $\operatorname{rank}_{\mathfrak p}(M)$,
is the rank of the free $A_{\mathfrak p}$-module
$M_{\mathfrak p}$ (or, equivalently, the rank at
$\mathfrak p$ of the locally free $\sh{O}_X$-Module
$\widetilde M$).  Thus
\[
  \operatorname{rank}_{\mathfrak p}M
  =
  \operatorname{rank}_{\mathfrak p}(M_{\mathfrak p})
  =
  \operatorname{rank}_{\kres(\mathfrak p)}
    \bigl(M\otimes_A\kres(\mathfrak p)\bigr).
  \tag{7.9.1.1}\label{III.7.9.1.1}
\]
\end{env}

\begin{proposition}[7.9.2]
\label{III.7.9.2}
Let $P_\bullet$ be a finite complex of projective $A$-modules
of finite type, and, for every $A$-module $M$, set
\[
  T_\bullet(M)=H_\bullet(P_\bullet\otimes_A M).
\]
Then, for every $\mathfrak p\in\Spec(A)$,
\[
  \sum_i(-1)^i
    \operatorname{rank}_{\kres(\mathfrak p)}
      T_i\bigl(\kres(\mathfrak p)\bigr)
  =
  \sum_i(-1)^i\operatorname{rank}_{\mathfrak p}(P_i).
  \tag{7.9.2.1}\label{III.7.9.2.1}
\]
\end{proposition}

\begin{proof}
Indeed,
$T_i(\kres(\mathfrak p))
 =H_i(P_\bullet\otimes_A\kres(\mathfrak p))$ by definition.
In view of \eqref{III.7.9.1.1}, formula
\eqref{III.7.9.2.1} is simply the invariance of the
Euler--Poincar\'e characteristic of a finite complex of
finite-dimensional vector spaces upon passage to homology
\sref[0]{0.11.10.2}.
\end{proof}

\begin{corollary}[7.9.3]
\label{III.7.9.3}
The function
\[
  \mathfrak p\longmapsto
  \sum_i(-1)^i
    \operatorname{rank}_{\kres(\mathfrak p)}
      T_i\bigl(\kres(\mathfrak p)\bigr)
\]
is locally constant on $\Spec(A)$.
\end{corollary}

\begin{theorem}[7.9.4]
\label{III.7.9.4}
Let $Y$ be a locally Noetherian prescheme, let
$f:X\to Y$ be a proper morphism, and let
$\sh{P}_\bullet$ be a finite complex of coherent,
$Y$-flat $\sh{O}_X$-Modules.  If
\[
  \sh{T}_\bullet(\sh{M})
  =
  \mathscr{H}^{-\bullet}
    \bigl(f,\sh{P}^\bullet
      \otimes_{\sh{O}_Y}\sh{M}\bigr)
  \qquad\text{\rm(cf.~\eqref{III.7.7.1.1})},
\]
then the function
\oldpage[III]{209}
\[
  y\longmapsto
  \sum_i(-1)^i
    \operatorname{rank}_{\kres(y)}
      \sh{T}_i\bigl(\kres(y)\bigr)
  \tag{7.9.4.1}\label{III.7.9.4.1}
\]
is locally constant on $Y$.
\end{theorem}

\begin{proof}
We may restrict ourselves to the case where
$Y=\Spec(A)$ is affine and $A$ is Noetherian.  Since the
complex $\sh{P}_\bullet$ is finite, \sref{III.7.7.12}{(i)}
shows that
\[
  \sh{T}_p(\sh{M})
  =
  \mathscr{H}_p(\sh{L}_\bullet\otimes_{\sh{O}_Y}\sh{M}),
\]
where $\sh{L}_\bullet=\widetilde{L}_\bullet$ and
$L_\bullet$ is a finite complex of projective $A$-modules
of finite type.  The theorem now follows from
\sref{III.7.9.3}.
\end{proof}

\begin{env}[7.9.5]
\label{III.7.9.5}
Under the hypotheses of \sref{III.7.9.4}, the function
\eqref{III.7.9.4.1} is \emph{constant} when $Y$ is
\emph{connected}.  When $Y$ is connected and nonempty, its
unique (integral) value is denoted
\[
  \operatorname{EP}(f,\sh{P}_\bullet),\qquad
  \operatorname{EP}(Y,\sh{P}_\bullet),
  \qquad\text{or simply}\qquad
  \operatorname{EP}(\sh{P}_\bullet)
\]
when no confusion is possible; this integer is called the
\emph{Euler--Poincar\'e characteristic} of
$\sh{P}_\bullet$ relative to $f$ (or to $Y$).  In general,
we likewise write
\[
  \operatorname{EP}(f,\sh{P}_\bullet;y),\qquad
  \operatorname{EP}(Y,\sh{P}_\bullet;y),\qquad
  \operatorname{EP}(\sh{P}_\bullet;y)
\]
for the right-hand side of \eqref{III.7.9.4.1}.
\end{env}

\begin{env}[7.9.6]
\label{III.7.9.6}
Under the hypotheses of \sref{III.7.9.4} concerning
$X$, $Y$, and $f$, let
\[
  0\longrightarrow\sh{P}'_\bullet
    \xrightarrow{u}\sh{P}_\bullet
    \xrightarrow{v}\sh{P}''_\bullet
    \longrightarrow 0
\]
be an exact sequence of finite complexes of coherent,
$Y$-flat $\sh{O}_X$-Modules, where the homomorphisms $u$
and $v$ have even degrees $2d$ and $2d'$, respectively.
Since $\sh{T}_\bullet$ is a homological functor
\sref{III.7.7.1}, there is an exact homology sequence
\[
  \cdots\longrightarrow
  \sh{T}_i(\sh{P}'_\bullet,\kres(y))
  \longrightarrow
  \sh{T}_{i+2d}(\sh{P}_\bullet,\kres(y))
  \longrightarrow
  \sh{T}_{i+2d+2d'}(\sh{P}''_\bullet,\kres(y))
  \longrightarrow
  \sh{T}_{i-1}(\sh{P}'_\bullet,\kres(y))
  \longrightarrow\cdots ,
\]
which moreover has only finitely many terms.  On writing
that the Euler--Poincar\'e characteristic of this complex is
zero \sref[0]{0.11.10.1}, we obtain at once
\[
  \operatorname{EP}(\sh{P}_\bullet;y)
  =
  \operatorname{EP}(\sh{P}'_\bullet;y)
  +
  \operatorname{EP}(\sh{P}''_\bullet;y)
  \tag{7.9.6.1}\label{III.7.9.6.1}
\]
for every $y\in Y$.

Suppose, for example, that
$\sh{P}_\bullet=(\sh{P}_i)$ and that $\sh{P}_i=0$ for
$i<0$.  We then have the following exact sequence of
complexes:
\[
\xymatrix@C=2.4em{
  \cdots \ar[r] & 0 \ar[r]\ar[d] & 0 \ar[r]\ar[d]
    & 0 \ar[r]\ar[d] & 0 \ar[r] & \cdots \\
  \cdots \ar[r] & 0 \ar[r]\ar[d] & 0 \ar[r]\ar[d]
    & \sh{P}_1 \ar[r]\ar[d] & \sh{P}_2 \ar[r] & \cdots \\
  \cdots \ar[r] & 0 \ar[r]\ar[d] & \sh{P}_0 \ar[r]\ar[d]
    & \sh{P}_1 \ar[r]\ar[d] & \sh{P}_2 \ar[r] & \cdots \\
  \cdots \ar[r] & 0 \ar[r]\ar[d] & \sh{P}_0 \ar[r]\ar[d]
    & 0 \ar[r]\ar[d] & 0 \ar[r] & \cdots \\
  \cdots \ar[r] & 0 \ar[r] & 0 \ar[r]
    & 0 \ar[r] & 0 \ar[r] & \cdots
}
\]
the nonzero vertical arrows being identity automorphisms.
Applying \eqref{III.7.9.6.1} to this exact sequence and
arguing by induction on the length of $\sh{P}_\bullet$
gives
\[
  \operatorname{EP}(\sh{P}_\bullet;y)
  =
  \sum_i(-1)^i\operatorname{EP}(\sh{P}_i;y).
  \tag{7.9.6.2}\label{III.7.9.6.2}
\]

\oldpage[III]{210}
Here, for each coherent $\sh{O}_X$-Module $\sh{F}$ that is
flat over $Y$, the notation
$\operatorname{EP}(\sh{F};y)$ (or
$\operatorname{EP}(f,\sh{F};y)$, or
$\operatorname{EP}(Y,\sh{F};y)$) means
$\operatorname{EP}(\sh{Q}_\bullet;y)$ (with the
corresponding relative notation) for the complex
$\sh{Q}_\bullet$ whose only nonzero term is $\sh{F}$ in
degree $0$.  Thus the study of Euler--Poincar\'e
characteristics of complexes reduces to that of complexes
with a single term.
\end{env}

\begin{proposition}[7.9.7]
\label{III.7.9.7}
Under the hypotheses of \sref{III.7.9.4}, let $Y'$ be a
locally Noetherian prescheme, let $g:Y'\to Y$ be a
morphism, put
\[
  X'=X\times_Y Y',\qquad
  f'=f_{(Y')}:X'\longrightarrow Y',
\]
and let
\[
  \sh{P}'_\bullet
  =
  \sh{P}_\bullet\otimes_{\sh{O}_Y}\sh{O}_{Y'}
\]
be the resulting finite complex of $\sh{O}_{X'}$-Modules.
Then $\sh{P}'_\bullet$ consists of coherent,
$Y'$-flat $\sh{O}_{X'}$-Modules, and, for every
$y'\in Y'$,
\[
  \operatorname{EP}(\sh{P}'_\bullet;y')
  =
  \operatorname{EP}(\sh{P}_\bullet;g(y')).
  \tag{7.9.7.1}\label{III.7.9.7.1}
\]
\end{proposition}

\begin{proof}
The $\sh{O}_{X'}$-Modules $\sh{P}'_i$, being inverse images
of the $\sh{P}_i$ under the projection $X'\to X$, are
coherent.  They are $Y'$-flat by \sref[0]{0.6.2.1} and
\sref[I]{I.4.14.5}, the question being local on $X$, $Y$,
and $Y'$.  Finally, $f'$ is proper by
\sref[II]{II.5.4.2}, so the left-hand side of
\eqref{III.7.9.7.1} is defined.  Formula
\eqref{III.7.9.7.1} now follows from
\textup{(6.10.4.2)}, \sref{III.7.7.2}, and
\sref{III.7.6.7}, after reducing, as one may always do, to
the case where $Y$ and $Y'$ are affine.
\end{proof}

\begin{proposition}[7.9.8]
\label{III.7.9.8}
Suppose that the hypotheses of \sref{III.7.9.4} hold and,
in addition, that there is an integer $i_0$ such that
\[
  \sh{T}_i(\kres(y))=0
  \qquad\text{for $i\neq i_0$ and every $y\in Y$.}
\]
Then
\[
  \sh{T}_{i_0}(\sh{O}_Y)
  =
  \mathscr{H}^{-i_0}(f,\sh{P}^\bullet)
\]
is a locally free $\sh{O}_Y$-Module, and its rank at
$y\in Y$ is
\[
  (-1)^{i_0}\operatorname{EP}(f,\sh{P}_\bullet;y).
\]
\end{proposition}

\begin{proof}
First observe that the hypotheses of \sref{III.7.4.4} are
satisfied by the functors $\sh{T}^{\sh{O}_y}_i$, so that
\sref{III.7.4.7} applies to them.  The hypothesis and
\sref{III.7.5.3} imply that
$\sh{T}^{\sh{O}_y}_i$ is zero for $i\neq i_0$.
Consequently $\sh{T}_{i_0}$ is exact by
\sref{III.7.3.3}, and \sref{III.7.8.4} shows that
$\mathscr{H}^{-i_0}(f,\sh{P}^\bullet)$ is locally free.
Its rank at $y\in Y$ is
\[
  \operatorname{rank}_{\kres(y)}
    \sh{T}_{i_0}(\kres(y))
  =
  (-1)^{i_0}\operatorname{EP}(f,\sh{P}_\bullet;y),
\]
by definition, since $\sh{T}_i(\kres(y))=0$ for
$i\neq i_0$.\footnote{The printed French proof omits the
factor $(-1)^{i_0}$ in this displayed equality; it is
required both by the definition and by the statement of
the proposition.}
\end{proof}

\begin{corollary}[7.9.9]
\label{III.7.9.9}
Let $Y$ be a locally Noetherian prescheme, let
$f:X\to Y$ be a proper morphism, and let $\sh{F}$ be a
coherent, $Y$-flat $\sh{O}_X$-Module.  Suppose that there
is an integer $i_0$ such that
\[
  H^i\bigl(f^{-1}(y),
    \sh{F}\otimes_{\sh{O}_Y}\kres(y)\bigr)=0
\]
for every $i\neq i_0$ and every $y\in Y$.  Then
$R^{i_0}f_*(\sh{F})$ is a locally free
$\sh{O}_Y$-Module, and its rank at $y$ is
$(-1)^{i_0}\operatorname{EP}(f,\sh{F};y)$.
\end{corollary}

In particular:

\begin{corollary}[7.9.10]
\label{III.7.9.10}
Under the preliminary hypotheses of \sref{III.7.9.9}
concerning $X$, $Y$, and $\sh{F}$, suppose that
$R^if_*(\sh{F})=0$ for every $i>0$.  Then
$f_*(\sh{F})$ is a locally free $\sh{O}_Y$-Module, and its
rank at $y$ is $\operatorname{EP}(f,\sh{F};y)$.
\end{corollary}

\begin{proof}
In view of \sref{III.7.9.9}, it is enough to prove the
following lemma.
\end{proof}

\begin{lemma}[7.9.10.1]
\label{III.7.9.10.1}
Under the hypotheses of \sref{III.7.9.10},
\[
  H^i\bigl(f^{-1}(y),
    \sh{F}\otimes_{\sh{O}_Y}\kres(y)\bigr)=0
\]
for every $i>0$ and every $y\in Y$.
\end{lemma}

\begin{proof}
We may restrict ourselves to the case where
$Y=\Spec(A)$ is affine.  With the notation of
\sref{III.7.9.4}, and with $\sh{P}_\bullet$ reduced to its
only nonzero term $\sh{F}$ in degree $0$, the hypothesis
gives $\sh{T}_p(\sh{O}_Y)=0$ for $p<0$.  It follows from
\sref{III.7.3.7} that $\sh{T}_p$ is exact for $p<0$, and
the lemma then follows from the equivalence of conditions
a) and d) in \sref{III.7.7.5}.
\end{proof}

\oldpage[III]{211}
\begin{proposition}[7.9.11]
\label{III.7.9.11}
Under the hypotheses of \sref{III.7.9.4}, let $\sh{L}$ be
an invertible $\sh{O}_X$-Module that is very ample over
$Y$, and, for every $n\in\bb{Z}$, put
\[
  \sh{P}_\bullet(n)
  =
  \sh{P}_\bullet
    \otimes_{\sh{O}_X}\sh{L}^{\otimes n}.
\]
Then, for every $y\in Y$, the function
\[
  n\longmapsto
  \operatorname{EP}(f,\sh{P}_\bullet(n);y)
  \tag{7.9.11.1}\label{III.7.9.11.1}
\]
is a polynomial with rational coefficients, and is the
same for all points in any one connected component of
$Y$.
\end{proposition}

\begin{proof}
It is clear that $\sh{P}_\bullet(n)$ is a complex of
$Y$-flat $\sh{O}_X$-Modules.  By
\eqref{III.7.9.6.2}, we may reduce to the case where
$\sh{P}_\bullet$ has a single nonzero term
$\sh{F}\neq0$ in degree $0$.  Since the questions are
local on $Y$, we may moreover suppose that $Y$ is affine
and that $f$ is projective \sref[II]{II.5.5.3}.  Put
\[
  X_y=f^{-1}(y),\qquad
  \sh{L}_y=\sh{L}\otimes_{\sh{O}_Y}\kres(y);
\]
then $\sh{L}_y$ is a very ample $\sh{O}_{X_y}$-Module
\sref[II]{II.4.4.10}.  By \sref{III.7.7.2}, for the
functor $\sh{T}_\bullet$ associated with
$\sh{P}_\bullet(n)$ we have
\[
  \sh{T}_i(\kres(y))
  =
  H^{-i}\bigl(
    X_y,\sh{F}_y\otimes\sh{L}_y^{\otimes n}\bigr),
  \qquad
  \sh{F}_y=\sh{F}\otimes_{\sh{O}_Y}\kres(y).
\]
Thus $\operatorname{EP}(f,\sh{F}(n);y)$ is precisely the
Euler--Poincar\'e characteristic
$\chi_{\kres(y)}(\sh{F}_y(n))$ defined in
\sref{III.2.5.1}.  It follows from \sref{III.2.5.3} that
\eqref{III.7.9.11.1} is a polynomial.  For each fixed
$n$, its value is constant on each connected component of
$Y$ by \sref{III.7.9.4}; this completes the proof.
\end{proof}

\begin{env}[7.9.12]
\label{III.7.9.12}
The polynomial \eqref{III.7.9.11.1}, with rational
coefficients, will be denoted
\[
  \operatorname{PH}(f,\sh{P}_\bullet;y)
  \qquad\text{or}\qquad
  \operatorname{PH}(\sh{P}_\bullet;y),
\]
and will be called the \emph{Hilbert polynomial at $y$}
relative to $\sh{P}_\bullet$, $f$, and $\sh{L}$ (or simply
the \emph{Hilbert polynomial at $y$} of
$\sh{P}_\bullet$, or of $f$, when no confusion is
possible).  When $Y$ is connected and nonempty, the
reference to $y$ is omitted from the notation and
terminology.  This invariant will play an essential role
in Chapter~V, in the theory of ``moduli'' of coherent
quotient sheaves of a given coherent sheaf.

With the notation of \sref{III.7.9.6} and
\sref{III.7.9.11}, we have
\[
  \operatorname{PH}(\sh{P}_\bullet;y)
  =
  \operatorname{PH}(\sh{P}'_\bullet;y)
  +
  \operatorname{PH}(\sh{P}''_\bullet;y),
  \tag{7.9.12.1}\label{III.7.9.12.1}
\]
and in particular
\[
  \operatorname{PH}(\sh{P}_\bullet;y)
  =
  \sum_i(-1)^i\operatorname{PH}(\sh{P}_i;y).
  \tag{7.9.12.2}\label{III.7.9.12.2}
\]
These formulas follow immediately from
\eqref{III.7.9.6.1} and \eqref{III.7.9.6.2}.  Likewise,
with the notation and hypotheses of \sref{III.7.9.7},
\[
  \operatorname{PH}(\sh{P}'_\bullet;y')
  =
  \operatorname{PH}(\sh{P}_\bullet;g(y')).
  \tag{7.9.12.3}\label{III.7.9.12.3}
\]
\end{env}

Formula \eqref{III.7.9.12.2} reduces the study of the
Hilbert polynomials of a complex to that of the Hilbert
polynomial of a single $Y$-flat $\sh{O}_X$-Module.  Such
polynomials have the following remarkable interpretation,
independent of homological considerations.

\begin{corollary}[7.9.13]
\label{III.7.9.13}
Let $Y$ be a Noetherian prescheme, let $f:X\to Y$ be a
proper morphism, let $\sh{L}$ be an invertible
$\sh{O}_X$-Module that is very ample over $Y$, and let
$\sh{F}$ be a coherent, $Y$-flat $\sh{O}_X$-Module.  There
is an integer $n_0$ such that, for $n\geq n_0$,
$f_*(\sh{F}(n))$ is a locally free $\sh{O}_Y$-Module whose
rank at $y\in Y$ is
\[
  \operatorname{PH}(f,\sh{F};y)(n).
\]
\end{corollary}

\begin{proof}
\oldpage[III]{212}
Since $f$ is projective \sref[II]{II.5.5.3}, there is an
integer $n_0$ such that, for $n\geq n_0$,
\[
  R^if_*(\sh{F}(n))=0
  \qquad\text{for every $i>0$}
\]
by \sref{III.2.2.1}.  The conclusion follows from
\sref{III.7.9.10}.
\end{proof}

The following flatness criterion will be important in
Chapter~V, in the theory of ``moduli'' of coherent
sheaves.

\begin{proposition}[7.9.14]
\label{III.7.9.14}
Let $Y$ be a locally Noetherian prescheme, let
$f:X\to Y$ be a projective morphism, and let $\sh{L}$ be
an invertible $\sh{O}_X$-Module that is ample for $f$.
For every $\sh{O}_X$-Module $\sh{F}$ and every
$n\in\bb{Z}$, put
\[
  \sh{F}(n)=\sh{F}\otimes\sh{L}^{\otimes n}.
\]
A coherent $\sh{O}_X$-Module $\sh{F}$ is $Y$-flat if and
only if there is an integer $n_0$ such that, for every
$n\geq n_0$, $f_*(\sh{F}(n))$ is a locally free
$\sh{O}_Y$-Module.
\end{proposition}

\begin{proof}
The necessity of the condition is proved as in
\sref{III.7.9.13}; the result of \sref{III.2.2.1} applies
to an ample sheaf because $f$ is projective.

Conversely, we may reduce to the case where
$Y=\Spec(A)$ is affine.  By the hypothesis and
\sref{III.2.2.2}{(i)}, the $A$-modules
$\Gamma(X,\sh{F}(n))$ are projective and of finite type
(Bourbaki, \emph{Alg.~comm.}, chap.~II, \S5, no.~2,
th.~1).  Let
\[
  S=\bigoplus_{n\geq0}\Gamma(X,\sh{L}^{\otimes n})
\]
be the associated graded ring.  We know that $X$ is
canonically identified with $\Proj(S)$ by
\sref[II]{II.4.5.2}{(b)} and \sref{III.5.4.4}.  Let
\[
  M=\bigoplus_{n\geq n_0}\Gamma(X,\sh{F}(n)).
\]
After replacing $\sh{L}$, if necessary, by a power
$\sh{L}^{\otimes d}$, we may suppose that $S$ is
generated by finitely many elements of degree $1$
\sref{III.2.3.5.1}; it then follows from
\sref[II]{II.2.7.5} and \sref[II]{II.2.7.2} that $\sh{F}$ is
identified with $\shProj_0(M)$.

For every homogeneous element $g\in S$ of positive
degree, we therefore have
\[
  \Gamma(X_g,\sh{F})=M_{(g)}.
\]
Now $M$, being a direct sum of projective $A$-modules, is
a flat $A$-module.  Hence $M_g$ is flat by
\sref[0]{0.6.3.2}, and so is $M_{(g)}$, which is a direct
factor of $M_g$ by \sref[0]{0.6.1.2}.  We conclude from
\sref[I]{I.4.14.5} that $\sh{F}$ is $Y$-flat at every
point of $X_g$.  Since the open subsets $X_g$ cover $X$,
the proposition is proved.
\end{proof}

\bigskip
\noindent\emph{(To be continued.)}

\clearpage
\oldpage[III]{213}
\section*{INDEX OF NOTATION}

\begin{flushleft}
$\mathbf{Tor}_n^A(P_\bullet,Q_\bullet)$: 6.3.1.

\medskip
$\shTor_n^{\sh{O}_S}(\sh{P}_\bullet,\sh{Q}_\bullet)$,
$\shTor_n^S(\sh{P}_\bullet,\sh{Q}_\bullet)$: 6.4.1 and 6.5.1.

\medskip
$\shTor_n^S(\sh{F},\sh{G})$: 6.5.1.

\medskip
$\shTor_q^S(\sh{P}_\bullet^{(1)},\sh{P}_\bullet^{(2)},\ldots,
\sh{P}_\bullet^{(m)})$: 6.5.15.

\medskip
$\mathbf H^{-n}(\mathfrak U^{(i)},\sh{P}_\bullet^{(i)})$,
$\mathbf H^{-n}(X^{(i)},\sh{P}_\bullet^{(i)})$: 6.6.1.

\medskip
$\mathbf{Tor}_n^A(L_{\bullet\bullet}^{(1)},L_{\bullet\bullet}^{(2)})$,
$\mathbf{Tor}_n^S(\mathfrak U^{(1)},\mathfrak U^{(2)};
\sh{P}_\bullet^{(1)},\sh{P}_\bullet^{(2)})$: 6.6.2.

\medskip
$\shTor_n^S(\mathfrak U^{(1)},\mathfrak U^{(2)};
\sh{F}^{(1)},\sh{F}^{(2)})$,
$\shTor_n^S(\mathfrak U^{(1)},\mathfrak U^{(2)};
\sh{P}_\bullet^{(1)},\sh{P}_\bullet^{(2)})$: 6.6.2.

\medskip
$\shTor_n^S(f_1,f_2;\sh{P}_\bullet^{(1)},
\sh{P}_\bullet^{(2)})$: 6.6.8, 6.7.1, and 6.7.3.

\medskip
$\shTor_n^S((f_i)_{i\in I};(\sh{P}_\bullet^{(i)})_{i\in I})$,
$\shTor_n^S(f_1,\ldots,f_m;\sh{P}_\bullet^{(1)},\ldots,
\sh{P}_\bullet^{(m)})$: 6.7.11.

\medskip
$\mathbf{Ab}$, $\mathbf{Ab}_A$: 7.1.1.

\medskip
$T_{(B)}$, $T^{(B)}$, $T\otimes_A B$
($T$ a functor from $\mathbf{Ab}_A$ to $\mathbf{Ab}$): 7.1.3.

\medskip
$t_M$: 7.2.2.

\medskip
$T_{\mathfrak p}$: 7.1.4.

\medskip
$\sh{T}_\bullet^{Y'}(\sh{M}')$,
$\sh{T}_\bullet^U(\sh{M}|U)$: 7.7.2.

\medskip
$\operatorname{EP}(f,\sh{P}_\bullet;y)$,
$\operatorname{EP}(Y,\sh{P}_\bullet;y)$,
$\operatorname{EP}(\sh{P}_\bullet;y)$: 7.9.5.

\medskip
$\operatorname{PH}(f,\sh{P}_\bullet;y)$,
$\operatorname{PH}(\sh{P}_\bullet;y)$: 7.9.12.
\end{flushleft}

\clearpage
\oldpage[III]{214}
\section*{TERMINOLOGICAL INDEX}

\begin{flushleft}
Euler--Poincaré characteristic of a complex of Modules: 7.9.5.

\medskip
Cohomologically flat at a point $y\in Y$ in dimension $p$,
cohomologically flat over $Y$ in dimension $p$, cohomologically flat over
$Y$ at the point $y$, cohomologically flat over $Y$ (complex of
$\sh{O}_X$-Modules flat over $Y$): 7.8.1.

\medskip
Homotopic complexes: 6.1.4.

\medskip
Ringed space of cohomological dimension $\leq n$: 6.5.5.

\medskip
Extension of scalars in an additive covariant functor: 7.1.3.

\medskip
Sheaf of rings of cohomological dimension $\leq n$: 6.5.5.

\medskip
Künneth formula: 6.7.8.

\medskip
Homologically flat at a point $y\in Y$ in dimension $p$, homologically flat
over $Y$ in dimension $p$, homologically flat over $Y$ at the point $y$,
homologically flat over $Y$ (complex of $\sh{O}_X$-Modules flat over $Y$):
7.8.1.

\medskip
Homotopy: 6.1.4.

\medskip
Hypertor of two complexes of $A$-modules: 6.3.1.

\medskip
Local hypertor of two complexes of Modules: 6.4.1.

\medskip
Global hypertor of two complexes of Modules: 6.6.2, 6.6.8, and 6.7.3.

\medskip
Hilbert polynomial relative to a complex of Modules: 7.9.12.

\medskip
Künneth spectral sequences: 6.7.3.
\end{flushleft}

\clearpage
\oldpage[III]{215}
\section*{TABLE OF CONTENTS}

\begin{flushleft}
\textsc{Chapter III.} --- \textbf{Cohomological study of coherent sheaves
(continued)} \dotfill 5

\medskip
\S~6. Local and global Tor functors; Künneth formula \dotfill 5

6.1. Introduction \dotfill 5

6.2. Hypercohomology of complexes of Modules on a prescheme \dotfill 6

6.3. Hypertor of two complexes of modules \dotfill 9

6.4. Local hypertor functors of complexes of quasi-coherent Modules; the case
of affine schemes \dotfill 14

6.5. Local hypertor functors of complexes of quasi-coherent Modules: the
general case \dotfill 16

6.6. Global hypertor functors of complexes of quasi-coherent Modules and
Künneth spectral sequences: the case of an affine base \dotfill 21

6.7. Global hypertor functors of complexes of quasi-coherent Modules and
Künneth spectral sequences: the general case \dotfill 25

6.8. Associativity spectral sequences for global hypertor \dotfill 32

6.9. Base-change spectral sequences in global hypertor \dotfill 34

6.10. Local structure of certain cohomological functors \dotfill 39

\medskip
\S~7. Study of base change in covariant homological functors of Modules
\dotfill 43

7.1. Functors of $A$-modules \dotfill 43

7.2. Characterization of the tensor-product functor \dotfill 44

7.3. Exactness criteria for homological functors of modules \dotfill 48

7.4. Exactness criteria for the functors $\mathrm H_i(P\otimes_A M)$
\dotfill 53

7.5. The case of Noetherian local rings \dotfill 58

7.6. Descent of exactness properties. Semicontinuity theorem and Grauert's
exactness criterion \dotfill 60

7.7. Application to proper morphisms: I. The exchange property \dotfill 65

7.8. Application to proper morphisms: II. Cohomological flatness criteria
\dotfill 72

7.9. Application to proper morphisms: III. Invariance of the Euler--Poincaré
characteristic and the Hilbert polynomial \dotfill 76

\begin{flushright}
\emph{Received December 15, 1962.}
\end{flushright}
\end{flushleft}

% Printed page 216 is blank.
\clearpage
\thispagestyle{empty}
\null
\clearpage

\oldpage[III]{217}
\section*{ERRATA AND ADDENDA}
\begin{center}
(List 2)

\medskip
A) \emph{Typographical errors}
\end{center}

\noindent\textbf{$(0_{\mathrm I},\,3.2.1)$} On line 2 of p. 27, replace $X$
by $U$.

\noindent\textbf{$(0_{\mathrm I},\,3.2.6)$} On line 1 from the bottom of p.
27 and line 1 of p. 28, replace (3 times) $\mathfrak S_\lambda$ by
$\sh{H}_\lambda$.

\noindent\textbf{$(0_{\mathrm I},\,3.4.5)$} Add a parenthesis before 3.4.5).

\noindent\textbf{$(0_{\mathrm I},\,4.2.1)$} On line 10 of p. 39, replace
$\psi_*(A)$ by $\psi_*(\sh{A})$.

\noindent\textbf{$(0_{\mathrm I},\,5.1.3)$} On line 16 of p. 45, replace
$\sh{F}|V$ by $\sh{F}|U$.

\noindent\textbf{$(0_{\mathrm I},\,5.3.9)$} On line 12 of p. 48, before
``such that'', add: over an open neighborhood $U$ of $x$.

\noindent\textbf{$(0_{\mathrm I},\,7.6.9)$} On line 3 from the bottom of p.
73, replace $I_\lambda$ by $J_\lambda$.

\noindent\textbf{(I, 1.1.15)} On line 3 of p. 83, replace $\mathfrak a$ by
$\mathfrak a\ne A$.

\noindent\textbf{(I, 1.4.1)} On line 2 from the bottom of p. 90, replace
$\widetilde N$ by $N$.

\noindent\textbf{(I, 2.3.1)} On line 3 from the bottom of p. 101, replace (0,
4.1.6) by (0, 4.1.7).

\noindent\textbf{(I, 2.3.2)} On line 11 of p. 101, replace $B$ by $B_s$ and
$C$ by $C_t$.

\noindent\textbf{(I, 2.5.5)} On line 19 of p. 104, replace $S$-morphism by
$S$-prescheme.

\noindent\textbf{(I, 3.3.7)} On line 17 of p. 109, replace $Y_{(S)}$ by
$Y_{(S')}$.

\noindent\textbf{(I, 3.4.5)} On line 9 of p. 113, replace $\mathrm s$ by $s$.

\noindent\textbf{(I, 3.4.8)} On line 4 of p. 114, replace $p(x)$ by $f(x)$.

\noindent\textbf{(I, 3.5.1)} On line 3 of p. 115, replace $1_X\times g$ by
$1_{X'}\times g$.

\noindent\textbf{(I, 3.7.2)} On line 8 of p. 119, replace the first $X$ by
$X'$.

\noindent\textbf{(I, 4.2.2)} On line 8 of p. 123, in the left vertical arrow
of the diagram, replace $\alpha_{\psi(y)}$ by $\vphi_{\psi(y)}$.

\noindent\textbf{(I, 4.4.3)} On line 16 of p. 126, replace \emph{isomorphée}
by \emph{isomorphes}.

\noindent\textbf{(I, 4.5.5)} On line 17 from the bottom of p. 127, replace
(4.2.4) by (4.2.5).

\noindent\textbf{(I, 5.1.4)} On line 14 from the bottom of p. 128, replace
(2.1.7) by (2.1.8).

\noindent\textbf{(I, 5.3.13)} On line 20 of p. 134, replace (4.2.4) by
(4.2.5).

\noindent\textbf{(I, 6.4.2)} On line 1 from the bottom of p. 147, replace
\emph{covering} by \emph{finite covering}.

\noindent\textbf{(I, 9.1.13)} On line 15 of p. 171, replace
$p^{-1}(\sh{F})$ by $p^*(\sh{F})$.

\noindent\textbf{(I, 9.5.11)} On line 7 of p. 179, replace $Y$ by $Y'$.

\noindent\textbf{(I, 9.6.5)} On line 17 from the bottom of p. 180, replace
$(0,\,4.1.4)$ by $(0,\,4.1.3)$.

\noindent\textbf{(I, 10.12.2)} On line 14 from the bottom of p. 206, replace
$(X_n)_{(S_n)}$ by $(X_n)_{(S_m)}$.

\clearpage
\oldpage[III]{218}
\noindent\textbf{(I), Terminological index:} On line 6 from the bottom of p.
219, replace 0, 4.1.4 by 0, 4.1.5. On lines 16, 17, and 18 from the bottom of
p. 222, replace I, 2.1.7 by I, 2.1.8; on line 19 from the bottom of p. 222,
replace I, 2.1.8 by I, 2.1.7.

\noindent\textbf{(II, 1.5.2)} On line 17 of p. 12, insert an $f$ beside the
left vertical arrow of the diagram.

\noindent\textbf{(II, 4.2.7)} On line 16 from the bottom of p. 75, replace
(III, 2.1.14) by (III, 2.1.13).

\noindent\textbf{(II, 5.2.1)} On line 16 from the bottom of p. 97, replace
$X-U$ by $U$. On line 8 from the bottom of p. 97, replace $X'$ by $X_f$.

\noindent\textbf{(II, 5.3.6)} On line 16 of p. 100, replace (4.6.18) by
(4.6.17).

\noindent\textbf{(II, 6.2.7)} On line 9 from the bottom of p. 116, replace
chap. V by chap. IV.

\noindent\textbf{(II, 6.6.5)} On line 15 from the bottom of p. 132, replace
chap. V by chap. IV.

\noindent\textbf{(II, 7.4.12)} On line 1 of p. 151, replace chap. V by chap.
III.

\noindent\textbf{(II, 8.1.4)} On line 15 of p. 154, replace (III, 2.3.8) by
(III, 2.3.7).

\noindent\textbf{(II, 8.10.5)} On line 17 from the bottom of p. 187, replace
$g_{(j)x}$ by $g_{j(x)}$.

\noindent\textbf{$(0_{\mathrm{III}},\,10.2.7)$} On line 1 of p. 20, replace
$u$-adic by $\mathfrak m$-adic. On line 13 from the bottom of p. 20, replace
$k$ by $k_i$.

\noindent\textbf{$(0_{\mathrm{III}},\,11.1.1)$} On line 18 of p. 23, replace
$\operatorname{Im}(d_r^{p+r,q-r+1})$ by
$\operatorname{Im}(d_r^{p-r,q+r-1})$.

\noindent\textbf{$(0_{\mathrm{III}},\,11.8.6)$} On line 14 from the bottom of
p. 45, replace the arrows $\to$ by $\rightsquigarrow$ (4 times).

\noindent\textbf{$(0_{\mathrm{III}},\,12.3.4)$} On line 4 of p. 61, replace
$\rightsquigarrow$ by $\to$.

\noindent\textbf{$(0_{\mathrm{III}},\,13.2.4)$} On line 2 from the bottom of
p. 67, replace $\varinjlim$ by $\varprojlim$ (2 times).

\noindent\textbf{(III, 1.1.7)} On line 15 of p. 85, replace $A^r$ by
$\bigwedge(A^r)$ (2 times).

\noindent\textbf{(III, 1.2.4)} On line 7 of p. 87, replace
$\Gamma(U,\sh{F})$ by $\Gamma(U,\sh{O}_X)$ and replace $\sh{G}$ by $\sh{F}$.
On line 8 of p. 87, replace $\HH^p(\mathfrak U,\sh{G})$ by
$\HH^p(\mathfrak U,\sh{F})$.

\noindent\textbf{(III, 2.5.3.1)} On line 3 of p. 111, replace $n>0$ by
$n\geq0$. On line 8 of p. 111, replace $-r\leq n\leq0$ by $-r\leq n<0$.

\noindent\textbf{(III, 3.1.2)} On line 13 of p. 116, replace $\sh{G}$ by
$\sh{G}\in\mathbf K'$.

\noindent\textbf{(III, 4.3.6)} On line 7 of p. 133, replace
$K'=K\otimes_A\widehat A$ by $K'\supset K\otimes_A\widehat A$. On line 16 of
p. 133, replace $R(X)$ by $R(Y)$.

\noindent\textbf{(III, 4.4.12)} On line 10 of p. 138, replace chap. V by
chap. IV.

\noindent\textbf{(III, 4.6.6)} On line 19 of p. 142, replace ``chap. V'' by
``in a later section''.

\begin{center}
B) \emph{Textual modifications}\footnote{To facilitate the search for
references, modifications belonging to the errata list inserted in chapter N
will henceforth be designated by $\mathbf{Err}_{\mathrm N}$ followed by a
number.}
\end{center}

\noindent\textbf{$(\mathbf{Err}_{\mathrm{III}},\,1)$} In
$(0_{\mathrm I},\,5.1.3)$, on line 18 of p. 45, replace ``direct sum'' by
``finite direct sum''.

\noindent\textbf{$(\mathbf{Err}_{\mathrm{III}},\,2)$} In
$(0_{\mathrm I},\,5.4.3)$, on lines 8 to 4 from the bottom of p. 49, replace
the text beginning ``In the general case\ldots{}'' by the following:

In the general case where $X$ is a ringed space such that $\sh{O}_x$ is a
\emph{local} ring for
\oldpage[III]{219}
every $x\in X$, if $\sh{L}$ is an $\sh{O}_X$-Module \emph{of finite type} for
which there exists an $\sh{O}_X$-Module $\sh{F}$ such that
$\sh{L}\otimes_{\sh{O}_X}\sh{F}$ is \emph{isomorphic to} $\sh{O}_X$, the
preceding argument shows that, for every $x\in X$, $\sh{L}_x$ is isomorphic to
$\sh{O}_x$. It follows that $\sh{L}$ is then \emph{invertible}. Indeed, for
every $x\in X$, let $U$ be an open neighborhood of $x$ such that $\sh{L}|U$ is
generated by $n$ sections $s_i\ (1\leq i\leq n)$ over $U$ (5.2.1). We may
suppose, for example, that $s=s_1$ is such that $s_x\ne0$; and since
$\sh{L}_x$ is isomorphic to $\sh{O}_x$, for each $i$ there exists a section
$t_i$ of $\sh{O}_X$ over an open neighborhood $V_i\subset U$ of $x$ such that
$(t_i)_x s_x=(s_i)_x$. Consequently there is an open neighborhood
$V\subset\bigcap_{i=1}^nV_i$ of $x$ such that $t_i s=s_i$ in $V$; in other
words, $\sh{L}|V$ is generated by the \emph{single} section $s|V$. Moreover,
if $z$ is a section over an open subset $W\subset V$ of the kernel of the
homomorphism $\sh{O}_X|V\to\sh{L}|V$ defined by $s$ (5.1.1), then $z_y$
annihilates $\sh{L}_y$ for every $y\in W$, so $z_y=0$ by hypothesis and hence
$z=0$, which completes the proof of our assertion. Finally, consideration of
the tensor product $\sh{L}^{-1}\otimes\sh{L}\otimes\sh{F}$ immediately shows
that $\sh{F}$ is isomorphic to $\sh{L}^{-1}$.

\noindent\textbf{$(\mathbf{Err}_{\mathrm{III}},\,3)$} In
$(0_{\mathrm I},\,7.2.4)$, one must impose the condition that the powers
$\mathfrak J^n$ are \emph{closed} ideals in the admissible ring $A$ in order
for the proposition to be correct. The proof given is incorrect, and the
corrected statement follows from Bourbaki, \emph{Top. gén.}, chap. III, 3rd
ed., \S~3, no.~5, cor. 1 of prop. 9. The $\mathfrak J^n$ must likewise be
assumed \emph{closed} in $A$ in statements $(0_{\mathrm I},\,7.2.5)$ and
$(0_{\mathrm I},\,7.2.6)$.

\noindent\textbf{$(\mathbf{Err}_{\mathrm{III}},\,4)$} After proposition (I,
2.2.5), add: with the notation of (I, 2.2.5), if $\mathfrak J$ is an ideal of
$A$, we denote by $\widetilde{\mathfrak J\sh{F}}$ the
$\sh{O}_X$-submodule $\mathfrak J\widetilde{\sh{F}}$ defined in
$(0,\,4.3.5)$.

\noindent\textbf{$(\mathbf{Err}_{\mathrm{III}},\,5)$} In (I, 3.4.5), on line
1 from the bottom of p. 112, replace ``\emph{a field} $K$'' by ``\emph{an
algebraically closed field} $K$''. On line 1 of p. 113, replace
``\emph{extension}'' by ``\emph{algebraically closed extension}''. On lines 1
to 4 of p. 113, replace the text beginning ``$K$ will be called\ldots{}'' by
the following:

For every point of $X$ with values in a field $K$, $K$ will be called the
\emph{field of values} of the corresponding point; and if $x$ is the locality
of this point, the latter is also said to be \emph{localized at} $x$. This
defines a map $X(K)\to X$ that associates to a point with values in $K$ its
locality.

On lines 7 and 16 of p. 113, delete ``geometric''; on lines 10 and 13 of p.
113, delete ``geometric''. On lines 10 and 11 from the bottom of p. 115,
delete ``\emph{geometric}''.

\noindent\textbf{$(\mathbf{Err}_{\mathrm{III}},\,6)$} At the end of the proof
of (I, 3.6.1), on line 12 from the bottom of p. 117, add: If we set
$X'=X\times_Y\Spec(\sh{O}_y/\mathfrak a_y)$, then, for every point $x\in X'$,
identified by $p$ with a point of $X$, we have
$\sh{O}_{X',x}=\sh{O}_{X,x}/\mathfrak a_y\sh{O}_{X,x}$. Since the question is
local on $X$ and $Y$, we may suppose that $X=\Spec(B)$ and $Y=\Spec(A)$; then
$(B/\mathfrak a_yB)_x=B_x/\mathfrak a_yB_x$ by flatness $(0,\,1.3.2)$.

\noindent\textbf{$(\mathbf{Err}_{\mathrm{III}},\,7)$} In (I, 5.1.1), on line
3 from the bottom of p. 127, replace ``quasi-coherent $\sh{O}_X$-Module'' by
``\emph{quasi-coherent ideal of} $\sh{B}$''.

\noindent\textbf{$(\mathbf{Err}_{\mathrm{III}},\,8)$} After (I, 5.1.10), on
line 16 of p. 131, add:

\emph{Remark (5.1.11).} --- A slight adaptation of the argument used in
(5.1.9) shows that the conclusion remains valid \emph{without assuming a
priori that $X$ is a prescheme}, but

\oldpage[III]{220}
assuming only that $(X,\sh{O}_X)$ is a ringed space in local rings, that
$\sh{J}$ is an ideal of $\sh{O}_X$ such that $\sh{J}^n=0$, that the ringed
space $X_0=(X,\sh{O}_X/\sh{J})$ is an affine scheme, and finally that the
ideals $\sh{J}^k/\sh{J}^{k+1}$ are quasi-coherent $\sh{O}_{X_0}$-Modules. It
suffices to use (1.8.1) (cf. $(\mathbf{Err}_{\mathrm{II}})$) instead of
(2.2.4) in the argument.

\noindent\textbf{$(\mathbf{Err}_{\mathrm{III}},\,9)$} In (I, 5.3.1), on line
11 of p. 132, after ``or $\Delta_X$'', insert ``or $\Delta_\vphi$ if
$\vphi:X\to S$ is the structural morphism''.

\noindent\textbf{$(\mathbf{Err}_{\mathrm{III}},\,10)$} In (I, 5.3.9), the
proof is insufficient, because it does not prove that $\Delta_X(X)$ is
locally closed in $X\times_SX$. To obtain a correct proof, it suffices to use
(4.2.4, $a$)): for every $x\in X$ and every affine neighborhood $U$ of $x$ in
$X$, $U\times_SU$ is an affine neighborhood of $\Delta_X(x)$; taking (5.3.16)
into account (whose proof uses only definition (5.3.1.1)), we are reduced to
proving (5.3.9) when $S=\Spec(B)$ and $X=\Spec(A)$ are affine schemes. It is
then clear from (5.3.1.1) that $\Delta_X$ corresponds to the canonical
homomorphism $A\otimes_BA\to A$ that sends $x\otimes y$ to $xy$; since this
homomorphism is surjective, $\Delta_X$ is then a closed immersion (4.2.3),
which completes the proof.

\noindent\textbf{$(\mathbf{Err}_{\mathrm{III}},\,11)$} In (I, 7.1.15) and
(I, 7.1.16), on lines 7 and 15 from the bottom of p. 158, delete
``\emph{geometric}''.

\noindent\textbf{$(\mathbf{Err}_{\mathrm{III}},\,12)$} Replace (I, 7.4.7)
by: We extend (by abuse of language) the definitions of (7.4.1) to the case
where $X$ is a \emph{reduced} prescheme every point of which admits an open
neighborhood having only a \emph{finite} number of irreducible components. It
then follows from (7.3.4) and (7.4.6) that, for a quasi-coherent
$\sh{O}_X$-Module of finite type $\sh{F}$, to say that $\sh{F}$ is a
\emph{torsion sheaf} is equivalent to saying that $\operatorname{Supp}(\sh{F})$
\emph{contains no irreducible component of $X$}.

\noindent\textbf{$(\mathbf{Err}_{\mathrm{III}},\,13)$} In (I, 10.11.7), on
lines 17 to 23 of p. 205, the end of the proof beginning ``It remains to
prove\ldots{}'' is unnecessary by (10.11.6), since the sheaves in question
are coherent by $(0,\,5.3.5)$.

\noindent\textbf{$(\mathbf{Err}_{\mathrm{III}},\,14)$} In (II, 1.7.8), after
line 10 from the bottom of p. 16, add: when $\sh{E}=\sh{O}_S^n$, we also write
$\mathbf V_S^n$ in place of $\mathbf V(\sh{E})$; if moreover $S=\Spec(A)$ is
affine, we write $\mathbf V_A^n$ in place of $\mathbf V_S^n$; we have
$\mathbf V_A^n=\Spec(A[T_1,\ldots,T_n])$, where the $T_i$ are indeterminates.

\noindent\textbf{$(\mathbf{Err}_{\mathrm{III}},\,15)$} In (II, 1.7.10), on
line 10 of p. 17, delete ``\emph{geometric}''; on lines 12 and 16--17 of p.
17, delete ``geometric''.

\noindent\textbf{$(\mathbf{Err}_{\mathrm{III}},\,16)$} In (II, 1.7.12), add:
We set similarly:
\[
  \mathbf V(\sh{O}_S^n)=\mathbf V_S^n=S[T_1,\ldots,T_n].
\]

\noindent\textbf{$(\mathbf{Err}_{\mathrm{III}},\,17)$} In (II, 4.2.6), on
line 7 of p. 75, delete ``\emph{geometric}''; on line 8 of p. 75, delete
``\emph{geometric}''.

\noindent\textbf{$(\mathbf{Err}_{\mathrm{III}},\,18)$} In (II, 4.6.13), on
line 7 of p. 92, the argument must be made more precise, because with the
notation of (4.4.10) one must prove here that the immersion $\Gamma_f$ is
\emph{quasi-compact}, in order to apply (i bis). By (4.6.4), to prove that
$\sh{L}$ is ample relative to $f$, we may restrict to the case where $Y$ is
affine. Observe also that in each of the hypotheses of (v), $f$ is
quasi-compact (I, 6.6.4). On the other hand, if $g$ is
\oldpage[III]{221}
separated, $\Gamma_f$ is a closed immersion (I, 5.4.3), hence quasi-compact
(I, 6.6.4). If, on the contrary, $X$ is locally Noetherian, then since $Y$ is
affine and $f$ is quasi-compact, the underlying space of $X$ is quasi-compact,
hence Noetherian; and (I, 6.6.4) may again be applied to prove that $\Gamma_f$
is quasi-compact.

\noindent\textbf{$(\mathbf{Err}_{\mathrm{III}},\,19)$} The statement of (II,
5.2.2) is incorrect: conditions $b)$, $c)$, and $c')$ are not local on $Y$
when it is not known whether, for an open subset $U$ of $Y$, a quasi-coherent
$(\sh{O}_X|f^{-1}(U))$-Module is the restriction to $f^{-1}(U)$ of a
quasi-coherent $\sh{O}_X$-Module. Thus, on line 16 of p. 98, replace
``\emph{Let $f:X\to Y$ be a separated quasi-compact morphism}'' by:
``\emph{Let $X$ and $Y$ be two preschemes such that $X$ is a scheme or the
underlying space of $X$ is locally Noetherian, and let $f:X\to Y$ be a
quasi-compact morphism.}'' In the proof, observe that the hypotheses imply
that $f$ is \emph{separated} when $X$ is assumed to be a scheme, by (I, 5.5.5
and 5.5.8); hence (I, 9.2.2) may be applied in both cases.

\noindent\textbf{$(\mathbf{Err}_{\mathrm{III}},\,20)$} In (II, 6.2.3), after
line 18 from the bottom of p. 115, add:

\emph{A morphism $f:X\to Y$ is said to be quasi-finite at a point $x\in X$ if
there exist an affine open neighborhood $V$ of $y=f(x)$ and an affine open
neighborhood $U$ of $x$ such that $f(U)\subset V$ and the morphism $U\to V$
induced by restricting $f$ is quasi-finite. A morphism $f:X\to Y$ is said to
be locally quasi-finite if it is quasi-finite at every point of $X$.}

\noindent\textbf{$(\mathbf{Err}_{\mathrm{III}},\,21)$} In (II, 6.4.3), the
proof is incorrect, since the elements of a finite-type $A$-submodule of $K$
are not necessarily integral over $A$. Replace the last 9 lines of p. 121 by
the following text: Since we may suppose that $E$ is torsion-free, hence
faithful, $E$ is also a faithful $A[u]$-module ($A$ being embedded in the
ring of endomorphisms of $E$). Since $E$ is an $A$-module of finite type, it
follows that $u$ is \emph{integral} over $A$ (Bourbaki, \emph{Alg. comm.},
chap. V, \S~1, no.~1, lemma 1). We immediately conclude that the eigenvalues
of $u\otimes1$ (in an algebraic closure of $K$) are elements \emph{integral}
over $A$, and consequently so are the $\sigma_i(u)$.

\noindent\textbf{$(\mathbf{Err}_{\mathrm{III}},\,22)$} In
$(0_{\mathrm{III}},\,9.1.1)$, on line 18 from the bottom of p. 12, delete
``open''.

\noindent\textbf{$(\mathbf{Err}_{\mathrm{III}},\,23)$} In
$(0_{\mathrm{III}},\,11.7.3)$, on line 10 of p. 43, after $\mathsf C''$, add:
such that for every projective object $P$ of $\mathsf C'$ (resp. every
projective object $P'$ of $\mathsf C''$), the functor
$A'\rightsquigarrow T(P,A')$ (resp. $A\rightsquigarrow T(A,P')$) is exact in
$\mathsf C'$ (resp. $\mathsf C$).

\noindent\textbf{$(\mathbf{Err}_{\mathrm{III}},\,24)$} The statement of
$(0_{\mathrm{III}},\,13.7.7)$ is inaccurate and must be modified as follows:
on line 6 of p. 78, delete ``and
$(\RR^{n+1}T(A_k))_{k\in\mathbf Z}$''; on line 7 of p. 78, replace
``$\RR'^nT(A)$ is an $S$-module of finite type'' by ``$\RR'^nT(A)$ and
$\RR'^{n+1}T(A)$ are $S$-modules of finite type''; on lines 12--13 of p. 78,
delete ``and $p+q=n+1$''. In the proof, on line 23 of p. 78, delete ``and for
$n+1$''.

\noindent\textbf{$(\mathbf{Err}_{\mathrm{III}},\,25)$} In (III, 1.4.15), on
line 6 from the bottom of p. 92, replace ``\emph{of finite type}'' by
``\emph{quasi-compact}''.

\noindent\textbf{$(\mathbf{Err}_{\mathrm{III}},\,26)$} In (III, 2.2.4), on
line 1 of p. 102, replace ``\emph{the supports of $\sh{F}$ and $\sh{H}$ are
proper over $Y$}'' by ``\emph{the support of
$\operatorname{Im}(\sh{F}\to\sh{G})$ is proper over $Y$}''. In the proof,
replace the text of lines 10 to 16, beginning ``This being so\ldots{}'', by:

\oldpage[III]{222}
Set
$\sh{G}_1=\operatorname{Im}(\sh{F}\to\sh{G})
=\operatorname{Ker}(\sh{G}\to\sh{H})$, which is coherent. Since we have the
exact sequence $0\to\sh{G}_1(n)\to\sh{G}(n)\to\sh{H}(n)$ and the functor
$f_*$ is left exact, the sequence
$0\to f_*(\sh{G}_1(n))\to f_*(\sh{G}(n))\to f_*(\sh{H}(n))$ is exact for every
$n$; it therefore suffices to show that, for sufficiently large $n$, the
sequence $f_*(\sh{F}(n))\to f_*(\sh{G}_1(n))\to0$ is exact. In other words, we
may restrict to the case where $\sh{H}=0$ and the support of $\sh{G}$ is
proper over $Y$, hence \emph{closed in $Z$} (II, 5.4.10). Consequently
$\sh{G}'=i_*(\sh{G})$ is a coherent $\sh{O}_Z$-Module such that
$\sh{G}'|X=\sh{G}$. We know from (I, 9.4.3) that there exists a coherent
$\sh{O}_Z$-Module $\sh{F}'$ such that $\sh{F}'|X=\sh{F}$. Moreover, since $X$
is open in $Z$ and $\operatorname{Supp}(\sh{G})\subset X$ is closed in $Z$, it
is immediate that one defines a surjective homomorphism
$u':\sh{F}'\to\sh{G}'$ of sheaves such that $u'|X$ is the given surjective
homomorphism $u:\sh{F}\to\sh{G}$, by taking $u'|U=0$ for every open subset $U$
of $Z$ not meeting $\operatorname{Supp}(\sh{G})$ and $u'|U=u|U$ for every open
subset $U\subset X$. This being so, as at the beginning of the proof of
(2.2.2), we see that we may reduce to the case where $Y$ is affine; it is
therefore enough to show that, for sufficiently large $n$, the homomorphism
$\Gamma(u):\Gamma(X,\sh{F}(n))\to\Gamma(X,\sh{G}(n))$ is surjective. We have
the commutative diagram
\[
\xymatrix@C=35mm@R=12mm{
\Gamma(Z,\sh{F}'(n)) \ar[r]^{\Gamma(u')} \ar[d]_v &
\Gamma(Z,\sh{G}'(n)) \ar[d]^w \\
\Gamma(X,\sh{F}(n)) \ar[r]_{\Gamma(u)} & \Gamma(X,\sh{G}(n))
}
\]
where $v$ and $w$ are the restriction homomorphisms. The definition of
$\sh{G}'$ moreover shows that $w$ is \emph{bijective}; and by (2.2.3),
$\Gamma(u')$ is \emph{surjective} for sufficiently large $n$, so the same is
true of $\Gamma(u)$.

\noindent\textbf{$(\mathbf{Err}_{\mathrm{III}},\,27)$} In (III, 2.2.5),
after line 7 from the bottom of p. 102, add:

(iii) Under the hypotheses of (2.2.4) concerning $X$, $Y$, $f$, and $\sh{L}$,
it is immediate that if $\sh{H}$ is a coherent $\sh{O}_X$-Module whose support
is \emph{proper over $Y$}, an argument analogous to that of (2.2.4) shows
that there exists an integer $N$ such that, for $n\geq N$,
$\RR^if_*(\sh{H}(n))=0$ for every $i>0$. It follows from the exact
cohomology sequence that if $u:\sh{F}\to\sh{G}$ is a homomorphism of coherent
$\sh{O}_X$-Modules such that $\Ker(u)$ and $\Coker(u)$ have supports
\emph{proper over $Y$}, then there exists $N$ such that for $n\geq N$, the
corresponding homomorphism
$\RR^if_*(\sh{F}(n))\to\RR^if_*(\sh{G}(n))$ is \emph{bijective} for every
$i>0$.

\noindent\textbf{$(\mathbf{Err}_{\mathrm{III}},\,28)$} In (III, 4.4.9), on
line 17 from the bottom of p. 137, replace ``\emph{Let $Y$ be a locally
Noetherian integral prescheme}'' by ``\emph{Let $X$ and $Y$ be two locally
Noetherian integral preschemes}''; on line 16 from the bottom of p. 137,
replace ``\emph{of finite type}'' by ``\emph{locally of finite type}''. The
proof is essentially unchanged, since $f^{-1}(y)$ is locally of finite type
over $k(y)$ and is therefore still discrete.

\noindent\textbf{$(\mathbf{Err}_{\mathrm{III}},\,29)$} In (I, 9.3.4), on line
19 from the bottom of p. 173, replace ``\emph{Noetherian}'' by
``\emph{quasi-compact}''; on lines 18 and 19 from the bottom of p. 173,
replace ``\emph{coherent}'' by ``\emph{quasi-coherent of finite type}''. In
the proof, we reduce to the case where

\oldpage[III]{223}
$X=\Spec(A)$ and $\sh{F}=\widetilde M$, where $M$ is an $A$-module of finite
type, and $\sh{J}=\widetilde{\mathfrak J}$, where $\mathfrak J$ is an ideal of
finite type in $A$; the rest of the argument is then unchanged.

\noindent\textbf{$(\mathbf{Err}_{\mathrm{III}},\,30)$} In (I, 9.3.5),
replace lines 1 to 7 from the bottom of p. 173 by the following text:

\emph{Proposition (9.3.5). --- Let $X$ be a prescheme and $\sh{F}$ a
quasi-coherent $\sh{O}_X$-Module of finite type. Then there exist a closed
subprescheme $Y$ of $X$, whose underlying space is equal to
$\operatorname{Supp}(\sh{F})$, and a quasi-coherent $\sh{O}_Y$-Module of
finite type $\sh{G}$ such that, if $j:Y\to X$ is the canonical injection,
$\sh{F}$ is isomorphic to $j_*(\sh{G})$.}

It will suffice to show that the ideal $\sh{J}$ of $\sh{O}_X$,
\emph{annihilator of $\sh{F}$}, is \emph{quasi-coherent}. We then take for $Y$
the closed subprescheme of $X$ defined by $\sh{J}$ (I, 4.1.2), and since
$\sh{J}\sh{F}=0$, $\sh{F}$ is an $(\sh{O}_X/\sh{J})$-Module and we answer the
question by taking $\sh{G}=j^*(\sh{F})$. To see that $\sh{J}$ is
quasi-coherent, we may (the question being local) restrict to the case where
$X=\Spec(A)$ and $\sh{F}=\widetilde M$, where $M$ is an $A$-module generated
by finitely many elements $x_i\ (1\leq i\leq r)$. The ideal $\sh{J}$ is then
the intersection of the annihilators of the $x_i$. But the annihilator of
$x_i$ is the kernel of the homomorphism $\sh{O}_X\to\sh{F}$ corresponding to
the homomorphism $s\mapsto sx_i$ from $A$ to $M$; it is therefore indeed a
quasi-coherent ideal (I, 4.1.1), and every finite intersection of such ideals
is also a quasi-coherent ideal (I, 1.3.10).
% END NATIVE INPUT ega3/ega3-7-9


% \subsection{Descent of exactness properties. Semi-continuity theorem and Grauert’s exactness criterion}
% \label{subsection:III.7.6}


% \subsection{Application to proper morphisms: I. The exchange property}
% \label{subsection:III.7.7}


% \subsection{Application to proper morphisms: II. Cohomological flatness criteria}
% \label{subsection:III.7.8}


% \subsection{Application to proper morphisms: III. Invariance of Euler-Poincar\'e characteristic and the Hilbert polynomial}
% \label{subsection:III.7.9}
% END NATIVE INPUT ega3/ega3-7.tex
\end{document}
